% Agent cycles.  Standalone fragment; no preamble or bibliography command.
% Status words deliberately distinguish proved deductions from unresolved tasks.
\section{Part C: cycles, Lefschetz maps, and the integral questions}
\label{sec:cycles-audit}

\paragraph{Conventions and limits of the evidence.}
We write $\mathrm{CH}^r(X)$ for cycles modulo rational equivalence with
integral coefficients, $\mathrm{CH}^r(X)_{\mathrm{alg}}$ for its subgroup of
algebraically trivial classes, and
\[
 A^r_{\mathrm{alg}}(X)=
 (\mathrm{CH}^r(X)/\mathrm{CH}^r(X)_{\mathrm{alg}})\otimes\mathbb Q,
 \qquad
 J^r(X)=\mathrm{CH}^r(X)_{\mathrm{alg}}\otimes\mathbb Q.
\]
For a fixed Weil cohomology $H$, we set
$A^r_{\mathrm{hom}}(X)=\operatorname{im}(\mathrm{cl}^r)$ and
$G^r(X)=\ker(A^r_{\mathrm{alg}}(X)\longrightarrow A^r_{\mathrm{hom}}(X))$.
Thus $G^r$ is the rational Griffiths group, not the whole homologically
trivial Chow group. Unless indicated otherwise, varieties are smooth,
projective, geometrically connected over an algebraically closed field;
$L$ denotes multiplication by the class of an ample line bundle. Betti
cohomology is used over $\mathbb C$ and $\ell$-adic cohomology, with
$\ell\ne\operatorname{char}k$, in arbitrary characteristic.

We read the supplied transcriptions at 14:44--47, 52--71, 102--117 and 126,
and checked the arguments in the corresponding sweep-3 reading. The scans
are absent: none of the following is a new palaeographic certification.
The labels ``Proved'' below apply to the stated proposition and its
hypotheses, and do not assert completion of a larger task whose open gate
is recorded separately.

\subsection{C.1: the algebraic-equivalence assertions}

\paragraph{C.1a (Known input and proved consequence: the quintic obstruction).}
Let $V\subset\mathbb P^4_{\mathbb C}$ be a very general smooth quintic,
let $h=c_1(\mathcal O_V(1))$, and put $G=G^2(V)$. Then
\begin{equation}\label{eq:quintic-decomposition}
 A^1_{\mathrm{alg}}(V)=\mathbb Qh,
 \qquad A^2_{\mathrm{alg}}(V)=\mathbb Qh^2\oplus G,
 \qquad \dim_{\mathbb Q}G=\infty.
\end{equation}
In particular
\begin{equation}\label{eq:quintic-cokernel}
 \ker(L:A^1_{\mathrm{alg}}(V)\to A^2_{\mathrm{alg}}(V))=0,
 \quad
 \operatorname{coker}L\simeq G.
\end{equation}
Here ``very general'' means outside a countable union of proper closed
subsets of the parameter space; it must not be silently replaced by an
explicitly specified Fermat quintic.

\emph{Proof of the deduction.}
The input is \citet[Theorem (0.2), equation (0.1)]{cyClemens1983}:
the rational Abel--Jacobi image for a generic quintic is infinite
dimensional, whereas the image of algebraically trivial cycles is zero.
The Abel--Jacobi map therefore factors through $G$, which must be infinite
dimensional. Weak Lefschetz gives $H^2(V,\mathbb Q)=\mathbb Qh$;
Poincar\'e duality gives $H^4(V,\mathbb Q)=\mathbb Qh^2$.
The divisor cycle map identifies $A^1_{\mathrm{alg}}(V)$ with
$\mathrm{NS}(V)_{\mathbb Q}$. The cycle class of $h^2$ is nonzero since
$\deg h^3=5$. These facts give the asserted direct sum and cokernel.
We have used Clemens's theorem as an opened, numbered primary theorem;
we do not claim an independent proof of that infinite-dimensional theorem.
The complete difference-of-lines argument with a specified geometric
generic fibre is supplied separately in C.1g.

\paragraph{C.1b (Refuted: $A,A',A''$ and algebraic up to torsion $=$ numerical).}
The letter's $A$ fails by \eqref{eq:quintic-cokernel}. If
$i:Y\hookrightarrow V$ is \emph{any} smooth hyperplane section,
$A^1_{\mathrm{alg}}(Y)=\mathrm{NS}(Y)_{\mathbb Q}$ is finite dimensional,
whereas $A^2_{\mathrm{alg}}(V)$ is infinite dimensional. Thus
$i_*:A^1_{\mathrm{alg}}(Y)\to A^2_{\mathrm{alg}}(V)$ is not onto,
which refutes $A'$. For $U=V\setminus Y$ there is an exact sequence
\begin{equation}\label{eq:alg-localization}
 A^{r-1}_{\mathrm{alg}}(Y)\longrightarrow
 A^r_{\mathrm{alg}}(V)\longrightarrow
 A^r_{\mathrm{alg}}(U)\longrightarrow0.
\end{equation}
Consequently $A^2_{\mathrm{alg}}(U)$ is infinite dimensional, refuting
$A''$ in its stated range $2r>\dim U$.

To justify the equivalence relation in \eqref{eq:alg-localization}, we
first use Chow localization. If the restriction of a cycle becomes
algebraically equivalent to zero on $U$, we express that relation by
fibres of finitely many cycles on $U\times T$, with $T$ smooth connected
curves. We close these cycles in $V\times T$ and take their refined
fibres. Their differences are algebraically trivial on $V$; the residual
cycles restrict to zero on $U$, so Chow localization puts them on $Y$.
This proves the required equality of kernels. It is not an argument for
the corresponding assertion modulo homological equivalence.

Every nonzero $g\in G$ is homologically, hence numerically, zero, but is
nonzero modulo algebraic equivalence with rational coefficients.
Thus the identification with numerical equivalence also fails for
$\tau$-equivalence understood as algebraic equivalence up to torsion.

\paragraph{C.1c (Proved: the precise divisor-to-curve statement over $\mathbb C$).}
For every smooth projective complex $n$-fold, $n\ge2$, the map
\[
 L^{n-2}:\mathrm{NS}(X)_{\mathbb Q}\longrightarrow
 A^{n-1}_{\mathrm{alg}}(X)
\]
is injective and has cokernel canonically isomorphic to $G^{n-1}(X)$.
Indeed, hard Lefschetz makes its cohomological realization injective.
The inverse of the cohomological map carries a rational class of Hodge
type $(n-1,n-1)$ to one of type $(1,1)$. The Lefschetz $(1,1)$ theorem
represents the latter by a rational divisor. The image of every algebraic
curve class is therefore the image of a unique rational divisor; after
subtracting that image, its residue lies in $G^{n-1}(X)$. The intersection
of this Griffiths group with the Lefschetz image is zero. This proves
the assertion without assuming the full Hodge conjecture.

The surjectivity criterion here is exactly $G^{n-1}(X)=0$. For each
$n\ge3$, $X=V\times\mathbb P^{n-3}$ gives a counterexample: the projective
bundle formula puts $G^2(V)t^{n-3}$ inside $G^{n-1}(X)$, where $t$ is the
hyperplane class on the second factor. Thus $A$ fails universally in
the divisor case in every dimension at least three, although it holds
for an individual variety precisely under the stated vanishing condition.

\paragraph{C.1d (Proved low-dimensional cases; scope explicitly restricted).}
Over an algebraically closed field, $A$ holds for $i=0$: degree identifies
$A^n_{\mathrm{alg}}(X)$ with $\mathbb Q$, and $L^n$ is multiplication by
the positive integer $\deg L^n$. It holds when $2i=n$ because the map is
the identity. Hence all its cases for $n\le2$ hold. Divisors on surfaces
satisfy the letter's middle-dimensional $B$ by the Hodge index theorem.
In these dimensions $A'$ only asks for a degree-one zero-cycle after
rational coefficients, or is outside the cycle range; a point of the
smooth hyperplane section provides it. Equation
\eqref{eq:alg-localization} gives $A''$ for the stated hyperplane
complements in dimensions at most two.

When $G^r(X)=0$ for all $r$, the assertions on $A_{\mathrm{alg}}(X)$ are
exactly the assertions on algebraic cohomology classes. This observation
does \emph{not} prove their general homological surjectivity. Vanishing
of Griffiths groups removes one obstruction; it does not supply an
algebraic inverse to the Lefschetz map in every degree.

\paragraph{C.1e (Refuted local inference: middle positivity alone).}
There is a smooth projective fourfold for which the letter's $B(X)$
holds but algebraic and numerical equivalence differ. Take
$X=V\times\mathbb P^1$ and $L=h+t$. Since degree is the only invariant
of zero-cycles modulo algebraic equivalence on $V$, we have $hg=0$ for
$g\in G$. The projective bundle formula gives
\[
\begin{split}
A^2_{\mathrm{alg}}(X)&=\mathbb Qh^2\oplus\mathbb Qht\oplus G,\\
A^3_{\mathrm{alg}}(X)&=\mathbb Qh^3\oplus\mathbb Qh^2t\oplus Gt,\\
L(ah^2+bht+g)&=ah^3+(a+b)h^2t+gt.
\end{split}
\]
This is an isomorphism. Its primitive kernel is zero, so $B(X)$ is
vacuously satisfied. Nevertheless $G\ne0$ consists of numerically trivial
classes. This disproves $B(X)\Rightarrow D_{\mathrm{alg}}(X)$ with a
\emph{fixed-variety} quantifier. It does not disprove an implication
whose premise assumes $B$ for every smooth projective variety and all
the auxiliary products required in its proof.

\paragraph{C.1f (Proved: the complete universal surjectivity classification for $A$).}
Among the pairs $0\le i\le n/2$, the letter's $A$ holds for every
smooth projective complex $n$-fold exactly when $i=0$ or $2i=n$.
For every $1\le i<n/2$ there is a smooth projective counterexample
with finite-dimensional $A^i_{\mathrm{alg}}$ and infinite-dimensional
$A^{n-i}_{\mathrm{alg}}$. Thus it is surjectivity that fails in these
examples. We now give one construction for every pair.

Start with $X_3=V$ as in C.1a, embedded as a smooth Cartier divisor
in $R_4=\mathbb P^4$. We inductively construct $X_{2i+1}$ as a smooth
Cartier divisor in a smooth projective $R_{2i+2}$ whose Chow motive
is a sum of Tate motives. Suppose that $X_{2i-1}\subset R_{2i}$ is
already constructed, and write $D=X_{2i-1}$ for its divisor class.
Set
\[
 S=R_{2i}\times\mathbb P^1,
 \qquad Z=X_{2i-1}\times\{0\}\subset S,
 \qquad X_{2i+1}=\operatorname{Bl}_Z(S).
\]
The centre is smooth of codimension two. It is the regular zero
scheme of the section $(s_D,s_0)$ of
\[
 F=\operatorname{pr}_1^*\mathcal O_{R_{2i}}(D)
          \oplus\operatorname{pr}_2^*\mathcal O_{\mathbb P^1}(1).
\]
Let $R_{2i+2}=\mathbb P_{\mathrm{lines}}(F)$ over $S$. In this
projective line bundle, the locus where the section of $F$ lies
in the tautological line is a Cartier divisor, locally given by
$s_0u-s_Dv=0$. Since $(s_D,s_0)$ is a regular sequence, this locus
is exactly the blow-up $X_{2i+1}$. This is the elementary Rees-algebra
presentation of the blow-up of a two-generated regular ideal.
It is smooth because both $S$ and $Z$ are smooth. The ambient
$R_{2i+2}$ is a projective bundle over $S$, so its Chow motive remains
Tate, by the projective-bundle formula
\citep[Lemma 45.6.1, Tag 0FGQ]{cyStacksProject}. The construction therefore iterates without requiring an
unproved low-codimension projective-space embedding of a blow-up.

The codimension-two blow-up formula, whose maps are algebraic
correspondences and hence preserve algebraic equivalence
\citep[Theorem 2.8(3)]{cyJiang2023}, gives
\begin{equation}\label{eq:iterated-griffiths-shift}
 A^r_{\mathrm{alg}}(X_{2i+1})
 \simeq A^r_{\mathrm{alg}}(S)
           \oplus A^{r-1}_{\mathrm{alg}}(X_{2i-1}).
\end{equation}
More explicitly, if $\pi:X_{2i+1}\to S$, $j:E\hookrightarrow
X_{2i+1}$ and $q:E\to Z$ are the blow-up maps, the forward map is
$(a,b)\mapsto\pi^*a+j_*q^*b$ and the inverse is
$c\mapsto(\pi_*c,-q_*j^*c)$. We have
$\pi_*j_*q^*=0$, $q_*j^*\pi^*=0$, and
$j^*j_*q^*b=-\xi q^*b$ with $q_*\xi=1$.
These identities check both projections directly; surjectivity is
the blow-up formula. In particular all four maps commute with
cycle classes and preserve algebraic equivalence.
The analogous decomposition holds for cycle classes. Thus the only
nonzero rational Griffiths group of $X_{2i+1}$ is
\[
 G^{i+1}(X_{2i+1})\simeq G^2(V)=G,
\]
and every $A^r_{\mathrm{alg}}(X_{2i+1})$ with $r\ne i+1$ is finite
dimensional. This follows by induction from C.1a and the Tate property
of $S$, and proves failure for $n=2i+1$.
The same codimension-two decomposition for the quotient of any
two adequate equivalence relations is already stated and proved
in \citet[Proposition 13, p.~482]{cySamuel1960}; in particular it
applies directly to the Griffiths quotient used here.

For arbitrary $n\ge2i+1$, take
$X=X_{2i+1}\times\mathbb P^{n-2i-1}$. The projective bundle formula
puts its copies of $G$ exactly in codimensions
$i+1,i+2,\ldots,n-i$. Consequently $A^i_{\mathrm{alg}}(X)$ is finite
dimensional and $A^{n-i}_{\mathrm{alg}}(X)$ is infinite dimensional.
No linear map between these spaces is onto, in particular no ample
Lefschetz map is onto. The positive cases were proved in C.1d.

It is useful to record why the simpler attempted construction
$V\times\mathbb P^r$ does not prove all these pairs. Its Griffiths
string has degrees $2,\ldots,r+2$. For $2\le i<(r+3)/2$, both
$Gt^{i-2}$ and $Gt^{r+1-i}$ exist, and $L^{r+3-2i}$ maps one
isomorphically to the other. The incidence blow-up construction
shifts the first Griffiths degree, which is the necessary new step.

% Self-contained supplement to C.1.  No preamble or bibliography command.
\paragraph{C.1g (Proved: a specified generic quintic and an infinite-order difference of lines).}
We now give the missing line argument, with a specified function field.
This proof uses the local period-map formula and the elementary cyclic
projection idea of \citet[Section III]{cyCeresaCollino1983}; the polynomial
interpolation, the non-torsion specialization, and the passage to a
specified geometric generic fibre are given explicitly below. We do not
claim that every complex specialization has the asserted property.

Let
\[
 \mathcal M=\{e\in\mathbb Z_{\ge0}^{5}:|e|=5,\quad
                  e_2+e_3+e_4>0,\quad e_0+e_1+e_4>0\},
 \qquad e_*=(0,0,0,0,5).
\]
There are $114$ monomials in $\mathcal M$. For algebraically independent
indeterminates $t_e$, $e\in\mathcal M\setminus\{e_*\}$, set
\[
 K=\mathbb Q(t_e:e\in\mathcal M\setminus\{e_*\}),\qquad
 F=x_4^5+\sum_{e\ne e_*}t_e x^e,
 \qquad V=\{F=0\}\subset\mathbb P^4_K.
\]
The two specified lines are
\[
 L_0=\{x_2=x_3=x_4=0\},\qquad
 L_1=\{x_0=x_1=x_4=0\}.
\]
Then $V$ is smooth, $[L_0]-[L_1]$ is homologically trivial, and
\begin{equation}\label{eq:specified-quintic-lines}
 m([L_0]-[L_1])\ne0
 \quad\hbox{in }\mathrm{CH}^2(V_{\overline K})/
                    \mathrm{CH}^2(V_{\overline K})_{\mathrm{alg}}
 \quad\text{for every }m\ne0.
\end{equation}
The same assertion holds after every embedding $\overline K\hookrightarrow
\mathbb C$. It therefore supplies both a field and an equation for the
counterexample, without making an assumption about fields of definition
of analytic exceptional loci.

\emph{The parameter space and all its infinitesimal deformations.}
Let $T$ be the nonempty smooth-quintic locus in the affine parameter
space just displayed. The complete linear system of quintics containing
$L_0\cup L_1$ has that union as its base locus. Bertini gives smoothness
away from the lines, and its normal derivatives generate the conormal
space at every point of either line. Thus a general member is smooth
there as well. This proves that $T$ is nonempty and the generic $V$ is
smooth.

We will also need that this family supplies every abstract first-order
deformation at a general point. Here is a direct check. In the same
linear system consider the polynomial
\[
 F_* =x_2x_0^4+x_3x_0^2x_1^2+x_4x_1^4
       +x_0x_2^4+x_1x_2^2x_3^2+x_4x_3^4+x_4^5.
\]
On each of the two lines, its three normal derivatives are
$(u^4,u^2v^2,v^4)$. The resulting map
\[
 H^0(\mathbb P^1,\mathcal O(1))^{\oplus3}
       \longrightarrow H^0(\mathbb P^1,\mathcal O(5))
\]
is an isomorphism: multiplication sends the six standard basis elements
to
\[
 u^5,\quad u^4v,\quad u^3v^2,\quad u^2v^3,\quad uv^4,\quad v^5.
\]
This is an open condition on
the coefficients. We use $F_*$ only to certify that open condition,
and do not need it itself to be a smooth quintic. Intersecting this
nonempty open set with the smooth locus gives smooth quintics for which
$H^0(N_{L_j/V})=H^1(N_{L_j/V})=0$ for both $j$.

The incidence space of a quintic with two ordered skew lines is a
projective bundle over the smooth space of ordered skew pairs. At the
points just described its map to the full quintic parameter space has
invertible differential: every first-order perturbation of the equation
has a unique compatible first-order perturbation of each line, by the
six-by-six isomorphism. The incidence map is consequently dominant.
The group $\mathrm{PGL}_5$ acts transitively on ordered skew pairs,
and its directions induce only trivial abstract deformations. It follows
that the Kodaira--Spencer map from $T$ onto $H^1(V,T_V)$ is surjective
on a nonempty open subset $T^\circ$.

For clarity, all embedded deformations do give all abstract ones here.
The normal sequence and the Euler sequence give
\[
 H^0(V,\mathcal O_V(5))\longrightarrow H^1(V,T_V)\longrightarrow
 H^1(V,T_{\mathbb P^4}|_V)=0.
\]
The last vanishing follows from $H^1(V,\mathcal O_V(1))=0$ and
$H^2(V,\mathcal O_V)=0$, themselves obtained from the hypersurface
sequence and the cohomology of line bundles on $\mathbb P^4$.

\emph{The inversion assertion needed for this family.}
Adjunction gives $K_V\simeq\mathcal O_V$. A nonzero holomorphic
three-form $\Omega$ is nowhere zero, so contraction gives an isomorphism
of sheaves $T_V\simeq\Omega_V^2$ and hence
\[
 H^1(V,T_V)\xrightarrow{\sim}H^{2,1}(V).
\]
On a simply connected small ball in $T^\circ$, trivialize the integral
homology local system and choose a holomorphic generator $\Omega_t$.
For any nonzero integral class $\alpha$ put
$f_\alpha(t)=\int_{\alpha_t}\Omega_t$. This holomorphic function is
not identically zero. Indeed, if it vanished identically, its first
derivatives would vanish. The period-map differential is contraction
of the Kodaira--Spencer class with $\Omega$
\citep[Theorem (1.27), p.~816]{cyGriffiths1968}. Surjectivity just proved
would therefore make $\alpha$ annihilate $H^{3,0}\oplus H^{2,1}$.
Since $\alpha$ is real, it would also annihilate the conjugate space;
Poincar\'e duality would imply $\alpha=0$, a contradiction.

On each such ball there are countably many integral classes.
Outside their proper analytic zero sets, therefore,
\begin{equation}\label{eq:quintic-no-period-kernel}
 \int_\alpha\Omega=0,\quad \alpha\in H_3(V,\mathbb Z)
                     \quad\Longrightarrow\quad \alpha=0.
\end{equation}
We use here that $H_3(V,\mathbb Z)$ is torsion-free. The integral
Lefschetz theorem gives $H_2(V,\mathbb Z)=\mathbb Z$;
the universal coefficient theorem and Poincar\'e duality then give the
claimed torsion-freeness in degree three.

For any algebraic family of curves on a fixed such $V$, parameterized
by a smooth projective curve $D$, the cylinder map
$H_1(D,\mathbb Z)\to H_3(V,\mathbb Z)$ pairs trivially with $\Omega$.
Its swept surface has complex dimension two, so the pulled-back
holomorphic three-form vanishes. By
\eqref{eq:quintic-no-period-kernel} this cylinder map is zero.
The Abel--Jacobi map of the family is holomorphic and its map on
$H_1$ is this cylinder map; this also follows directly by integrating
over the swept chains. A map from a smooth projective curve to a
complex torus which is zero on $H_1$ is constant: its pullbacks of
holomorphic one-forms have zero periods and hence are zero, so its
differential is zero. Differences of fibres thus have Abel--Jacobi
class zero. Closing parameter curves and using refined fibres treats
every generator of algebraic equivalence. We have proved
\[
 AJ\bigl(\mathrm{CH}^2(V)_{\mathrm{alg}}\bigr)=0
\]
for this analytic very general set. The holomorphicity used here is
the construction in \citet[Theorem (2.20), pp.~826--827]{cyGriffiths1968}.

\emph{One smooth fibre where the chosen normal function is not torsion.}
We give the data, including the auxiliary surface. In the plane
$H=\{x_0=0\}\subset\mathbb P^3$ take the smooth Fermat curve
\[
 C:\ x_1^5+x_2^5+x_3^5=0.
\]
Its genus is six. Over the function field of $C\times C$, write its
two generic points in the chart $x_3=1$ as
$z_a=(0,a,b,1)$ and $z_b=(0,c,d,1)$, so
$a^5+b^5+1=c^5+d^5+1=0$ and $b\ne d$.
For every $m>0$, $m[z_a-z_b]\ne0$ in $\operatorname{Pic}^0(C)$:
the map $(z_a,z_b)\mapsto[m]([z_a]-[z_b])$ is not the zero morphism,
since the Abel map is nonconstant and multiplication by $m$ on its
Jacobian is finite. Its value at the generic point cannot be zero.

Take the two skew lines in $\mathbb P^3$
\[
 \ell_a(t,s)=(t,as,bs,s),\qquad
 \ell_b(t,s)=(t,t+cs,ds,s).
\]
Their intersections with $H$ are $z_a,z_b$; they are skew because
$b\ne d$. Define linear forms and a quartic polynomial by
\[
 S=\frac{x_2-bx_3}{d-b},\qquad T=x_1-ax_3-(c-a)S.
\]
\[
 Q(T,S)=T^4+5cT^3S+10c^2T^2S^2+10c^3TS^3+5c^4S^4.
\]
Then $T=S=0$ on $\ell_a$ and $(T,S)=(t,s)$ on $\ell_b$.
Put
\[
 K_0=x_1^5+x_2^5+x_3^5-x_0Q(T,S),\qquad
 K=K_0+x_0G,
\]
where $G$ is a generic element of
$H^0(\mathbb P^3,\mathcal I_{\ell_a\cup\ell_b}(4))$.
Thus $K$ contains both lines and $K|_H$ is exactly the fixed equation
of $C$. These identities follow by substitution, using
$tQ(t,s)=(t+cs)^5-c^5s^5$.

We justify smoothness of $\{K=0\}$, including the base locus.
The projective linear system spanned by $K_0$ and
$x_0H^0(\mathcal I_{\ell_a\cup\ell_b}(4))$ has base locus supported on
$C\cup\ell_a\cup\ell_b$. Bertini applies off it. At every point of
$C$, the derivative within $H$ of its fixed smooth equation is
nonzero. At a point of either line outside $H$, the variations $x_0G$
generate both normal derivatives: products of one linear equation of
each skew line generate the conormal of either line along that line,
and we multiply them by degree-two polynomials nonzero at the point.
Singularity therefore imposes two independent linear conditions at
a point of a one-dimensional line. The union of these bad conditions
has positive codimension in the parameter space. A generic $G$ avoids
it, and $K$ is smooth everywhere.

Now set
\[
 W:\ x_0x_4^4+K(x_0,x_1,x_2,x_3)=0\quad\subset\mathbb P^4,
 \qquad p=(0,0,0,0,1).
\]
This threefold is smooth. A hypothetical singular point has
$x_0x_4^3=0$. When $x_4=0$, smoothness of $K$ excludes it.
When $x_0=0$ and $(x_1,x_2,x_3)\ne0$, smoothness of $C$ excludes
it. The remaining point is $p$, where the $x_0$ derivative is one.
The two lines $\ell_a,\ell_b$ lie on $W$ in $x_4=0$ and avoid $p$.

Let $W'=\operatorname{Bl}_pW$ and $B=\operatorname{Bl}_C\mathbb P^3$.
Projection from $p$ lifts to a morphism $g:W'\to B$. Here is the
needed verification of this lift. In the blow-up of $\mathbb P^4$
at $p$, write its points as $[tu_0:tu_1:tu_2:tu_3:v]$ over
$[u_0:u_1:u_2:u_3]\in\mathbb P^3$. The strict transform has equation
\[
 u_0v^4+t^4K(u)=0.
\]
The ideal $(u_0,K)$ of $C$ pulls back to $(u_0)$ on $t\ne0$, and
to $(K)$ on $v\ne0$. It is invertible, so the universal property
of the blow-up gives $g$. Each proper transform of $\ell_a,\ell_b$
maps isomorphically to the corresponding proper transform in $B$.

Write $j:E\hookrightarrow B$ and $q:E\to C$ for the exceptional
divisor and its projection. The inverse blow-up formula in C.1f
has second component $-q_*j^*:\mathrm{CH}^2(B)\to\mathrm{CH}^1(C)$.
The proper transform of either line meets $E$ transversally once,
over its point $z_a$ or $z_b$. Consequently
\[
 -q_*j^*g_*([\ell_a']-[\ell_b'])=-[z_a-z_b].
\]
The blow-up of a point does not change $H^3$ or the intermediate
Jacobian, while the blow-up of $\mathbb P^3$ in $C$ identifies its
intermediate Jacobian with $\operatorname{Pic}^0(C)$. Abel--Jacobi
functoriality, equivalently integration along this correspondence,
therefore sends $AJ_W(\ell_a-\ell_b)$ to $-[z_a-z_b]$.
The latter is not torsion, so the former is not torsion.

This is a point of the specified family $T$ after an explicit change
of coordinates. Use as the five column vectors
\[
 e_0,\ z_a,\ e_0+e_1,\ z_b,\ q_0=(1,0,0,0,1),
\]
where the first four have fifth coordinate zero. Their determinant
is $b-d$ up to sign. Also $W(q_0)=1$, because $K$ vanishes at
$(1,0,0,0)$ on $\ell_a$. This matrix takes the standard ordered
pair $L_0,L_1$ to $\ell_a,\ell_b$, and the coefficient of $x_4^5$
in the transformed equation is one. All coefficients can be
embedded in $\mathbb C$, preserving the generic points and the
non-torsion conclusion.

\emph{Completion over the specified function field.}
Integral Lefschetz and duality identify $H^4(V,\mathbb Z)$ with
$\mathbb Z$ via degree of a curve, so $L_0-L_1$ is homologically
trivial on every smooth fibre of $T$. It defines a holomorphic
normal function $\nu$ in the family of intermediate Jacobians.
Its holomorphicity follows from the same locally defined chain
integrals: antiholomorphic derivatives vanish by the types of the
forms and of the algebraic relative cycle, and changing the chain
adds a period. The nonempty Zariski open $T$ in affine space is
analytically connected. For each $m>0$, the smooth fibre just
constructed has non-torsion value, so $m\nu$ cannot vanish identically
on any nonempty analytic open subset of $T$, even if that special
fibre lies outside $T^\circ$. Given any nonempty Zariski open subset
of $T$, choose one coordinate ball in its intersection with
$T^\circ$. On this ball the countably many period kernels and the
zero loci of $m\nu$, $m>0$, are all proper closed analytic subsets.
Baire's theorem shows that their joint complement is nonempty.
At each point in that complement, $m(L_0-L_1)$ is not algebraically
trivial for any $m>0$.

Suppose now that $m(L_0-L_1)$ were algebraically trivial on the
geometric generic fibre over $\overline K$. The assertion is
witnessed by finitely many curves, marked points, cycles and rational
equivalences. They descend to a finite extension of $K$. Normalize
the parameter space in this extension and shrink: the parameter
curves become smooth proper connected families with their sections,
and all finite cycle identities hold on every fibre. The image
contains a nonempty Zariski open subset of $T$. On its complex
points our cycle would be algebraically trivial, contradicting the
preceding paragraph. This proves \eqref{eq:specified-quintic-lines}.
Finally, algebraic equivalence does not acquire new relations under
an extension of algebraically closed fields: descend any proposed
finite relation to a finitely generated algebra over the smaller
field, spread its data over a nonempty open, and specialize at a
closed point, whose residue field is the smaller algebraically
closed field. This proves the assertion after every embedding into
$\mathbb C$.

In particular this proof supplies a nonzero rational Griffiths group
without importing Clemens's stronger infinite-dimensional theorem.
It also makes the counterexamples of C.1f completely specified:
start its induction with this $V$. The source has $G^i=0$, whereas
the target contains the displayed nonzero Griffiths class. Hard
Lefschetz implies that the image of the source meets the target
Griffiths group only in zero, so surjectivity fails for every ample
class. Infinite dimensionality is unnecessary for that conclusion.
The exact monomial count, both normal-derivative matrices, the
interpolation identities and the change-of-coordinates determinant
are recorded in \texttt{check\_quintic\_lines.py} and its output.

% Generic-hyperplane image, independently proved in the cycles contribution.
\paragraph{C.1h (Proved: the exact Gysin image of a geometric generic hyperplane).}
Let $X$ be smooth projective over an algebraically closed field $k_0$
of characteristic zero, and let $L$ be very ample. Put
$P=\mathbb P(H^0(X,L))$, $F=k_0(P)$ and $\Omega=\overline F$.
Let $i:Y\hookrightarrow X_\Omega$ be the geometric generic member
of $|L|$. With rational coefficients we have
\begin{equation}\label{eq:generic-gysin-image}
 \operatorname{im}\bigl(i_*:A^r_{\mathrm{alg}}(Y)
             \longrightarrow A^{r+1}_{\mathrm{alg}}(X_\Omega)\bigr)
       =L\,A^r_{\mathrm{alg}}(X_\Omega).
\end{equation}
Consequently its cokernel is
$A^{r+1}_{\mathrm{alg}}(X_\Omega)/L A^r_{\mathrm{alg}}(X_\Omega)$,
which also equals $A^{r+1}_{\mathrm{alg}}(X_\Omega\setminus Y)$
by \eqref{eq:alg-localization}. No statement of this kind is being
asserted for each special hyperplane section.

\emph{Proof.}
Let $\mathcal H\subset X\times P$ be the universal incidence,
with maps $p:\mathcal H\to X$ and $q:\mathcal H\to P$.
The evaluation sequence
\[
 0\longrightarrow E\longrightarrow H^0(X,L)\otimes\mathcal O_X
       \longrightarrow L\longrightarrow0
\]
identifies $\mathcal H$ with $\mathbb P_{\mathrm{lines}}(E)$.
Its relative class $\xi=c_1(\mathcal O_{\mathbb P(E)}(1))$
is $q^*c_1(\mathcal O_P(1))$.
Choose a cycle $\gamma$ on $Y$. It is defined over a finite
extension $F'/F$; enlarging it if necessary does not affect the
argument. Trace it to a cycle $\gamma_{\mathrm{tr}}$ on $Y_F$
and close this cycle in $\mathcal H$. By the projective-bundle
formula this closure has an expression
\[
 [\Gamma]=\sum_{j=0}^{\operatorname{rank}E-1}\xi^j p^*a_j,
       \qquad a_j\in\mathrm{CH}^{r-j}(X).
\]
On the generic fibre of $q$, the line bundle pulled back from $P$
is trivial. Hence
$[\gamma_{\mathrm{tr}}]=i_F^*a_0$ in $\mathrm{CH}^r(Y_F)$,
and its pushforward is $L a_0$.

We must relate the trace to the original class modulo algebraic
equivalence, not merely compare their cohomology. Let $T\to P$
be an integral model of $F'/F$ over $k_0$. After shrinking $T$,
the cycle and its pushforward to the constant family $X\times T$
spread as flat cycle families. Since $k_0$ is algebraically closed,
the integral variety $T$ is geometrically integral. Its base change
to $\Omega$ is therefore connected. The $[F':F]$ points of its
geometric generic fibre over $P$ give the conjugate cycles. Their
pushforwards on $X_\Omega$ are algebraically equivalent: two
points of a smooth connected parameter variety can be joined by
a chain of algebraic curves, and restricting the flat family to
these curves gives precisely the defining relations for algebraic
equivalence. We may shrink to a smooth open containing these
geometric generic points. It follows that
\[
 [F':F]\,[i_*\gamma]=[i_*\gamma_{\mathrm{tr}}]=L[a_0]
       \quad\hbox{in }A^{r+1}_{\mathrm{alg}}(X_\Omega).
\]
Division by $[F':F]$ proves one inclusion in
\eqref{eq:generic-gysin-image}. The other follows from
$i_*i^*\alpha=L\alpha$ for every cycle class $\alpha$ on
$X_\Omega$.\hfill$\square$

This already gives a further complete family of counterexamples.
For $X=X_{2r+1}$ of C.1f, we have $G^r(X)=0$ and
$G^{r+1}(X)\ne0$. Hard Lefschetz makes the cohomological map
$L:H^{2r}(X)\to H^{2r+2}(X)$ injective. Its image on cycles
therefore intersects $G^{r+1}(X)$ trivially. Equation
\eqref{eq:generic-gysin-image} shows that $A'$ fails for a
geometric generic hyperplane, and $A''$ fails on its affine
complement, in every critical pair $n=2r+1$.
We can take the algebraic closure of the specified field in C.1g
as $k_0$ and then adjoin generic hyperplane coefficients. Algebraic
equivalence is preserved and reflected under extensions of
algebraically closed fields: for reflection, spread a finite
equivalence witness and specialize it to a rational point over
the smaller algebraically closed field. Thus every embedding of
the resulting algebraically closed field into $\mathbb C$
retains these counterexamples.

% New file while the preceding checkpoint is under visual review.
\paragraph{C.1i (Proved: all non-top cases of $A'$ and $A''$ fail universally).}
Fix integers $c<n<2c$. There is a specified smooth projective
$n$-fold $X$ over the algebraic closure of a rational function
field over $\mathbb Q$, with a very ample class $L$, and a class
\[
 0\ne z\in G^c(X),\qquad
 z\notin L A^{c-1}_{\mathrm{alg}}(X).
\]
After adjoining coefficients for a geometric generic hyperplane
$Y\subset X$, the image of $z$ is outside
$i_*A^{c-1}_{\mathrm{alg}}(Y)$, and it is nonzero in
$A^c_{\mathrm{alg}}(X\setminus Y)$. These statements remain true
after every embedding of the algebraically closed defining field
into $\mathbb C$.

In particular, in the letter's range $2c>n$, with $1\le c\le n$,
the assertions $A'$ and $A''$ hold for every smooth projective
complex $n$-fold and smooth hyperplane section exactly when $c=n$.
Here the index in $A'$ is $i=c-1$; its range $2i\ge n-1$ is
exactly $2c>n$. Groups with $c>n$ vanish for dimension reasons.

\emph{Construction and proof.}
We use the split quadric
\[
 Q=\{x_0y_0+x_1y_1+\cdots+x_cy_c=0\}
            \subset\mathbb P^{2c+1}_{\mathbb Q},
       \qquad h=c_1(\mathcal O_Q(1)),
\]
and its two maximal linear spaces
\[
 \Pi_+=\{x_0=x_1=\cdots=x_c=0\},\qquad
 \Pi_-=\{y_0=x_1=\cdots=x_c=0\}.
\]
Put $Z=[\Pi_+]-[\Pi_-]\in\mathrm{CH}^c(Q)$.
The linear section $x_1=\cdots=x_c=0$ is the reduced union
$\Pi_+\cup\Pi_-$, so $h^c=[\Pi_+]+[\Pi_-]$.
Both spaces have degree one. The normal-bundle sequence
\[
 0\longrightarrow N_{\Pi_+/Q}\longrightarrow
       \mathcal O_{\mathbb P^c}(1)^{\oplus(c+1)}
       \longrightarrow\mathcal O_{\mathbb P^c}(2)
       \longrightarrow0
\]
is exact, since the displayed quadric is smooth. It gives
\[
 [\Pi_+]^2=[t^c]\frac{(1+t)^{c+1}}{1+2t}
       =\frac{1+(-1)^c}{2},\qquad
 \deg Z^2=2(-1)^c\ne0.
\]
Thus $\operatorname{cl}(Z)\ne0$. Weak Lefschetz and duality give
$H^{2c-2}(Q,\mathbb Q)=\mathbb Qh^{c-1}$ and
$\dim H^{2c+2}(Q,\mathbb Q)=1$. Equal degrees imply
$h\operatorname{cl}(Z)=0$. Consequently
\begin{equation}\label{eq:quadric-primitive-obstruction}
 \operatorname{cl}(mZ-ha)\ne0
 \quad\text{for every }m\ne0\text{ and }
                 a\in\mathrm{CH}^{c-1}(Q)_{\mathbb Q}.
\end{equation}
Indeed, the class of $ha$ is a multiple of $h^c$ and is
orthogonal to $Z$, whereas $Z^2\ne0$.

Set $s=2c-n$ and $D=4c+2$. Choose a basis of
$V=H^0(Q,\mathcal O_Q(D))$ over $\mathbb Q$, let
$S=\mathbb P(V^{\oplus s})$, and let $X$ be the geometric
generic common zero locus of its $s$ sections. This is smooth
of dimension $n$ by Bertini. To specify the field and equations
without any unnamed generality choice, use the degree-$D$
monomials not divisible by $x_0y_0$ as the basis of
$\mathbb Q[x_0,y_0,\ldots,x_c,y_c]/(\sum x_jy_j)$.
In the first section set the coefficient of $x_0^D$ equal to one;
take every other coefficient of the $s$ sections to be an
independent indeterminate. Thus the field is
$K=\mathbb Q(t_1,\ldots,t_{s\dim V-1})$ and we work over
$\Omega=\overline K$. Write $j:X\hookrightarrow Q_\Omega$,
$L=j^*h$ and $z=j^*Z$.

We make the only deep input explicit. We use the effective
connectivity theorem of \citet[Theorem 3.13, p.~23 of the opened
author version]{cyNagel2002}, with the required condition that
its connectivity parameter be at most $\dim X$, as made explicit
in \citet[Theorem 5.11]{cyNagel2004}. The regularity constant
$m_Q=1$ for a smooth quadric is recorded in
\citet[p.~3 of the opened author version, immediately before
Theorem 2]{cyNagel1998}. Put its parameter equal to $n$;
then $\mu=n$ and $r=s-1$. Since all degrees equal $D$, condition
$(C)$ has left side at least $D>2c$, and each condition $(C_i)$
has left side
\[
  \sum_{\nu=i}^{r}D+iD=sD\ge D>1+2c+n-i.
\]
Thus all the hypotheses hold. For every smooth morphism
$T\to S^\circ$, where $S^\circ$ is the smooth-fibre open,
the theorem gives
\begin{equation}\label{eq:nori-effective-connectivity}
 H^k(Q\times T,X_T;\mathbb Q)=0\qquad(k\le2n).
\end{equation}
The parameter $S=\mathbb P H^0(Q,\mathcal O_Q(D)^{\oplus s})$
is precisely the one in that theorem. No ineffective degree
threshold or unspecified exceptional field is used here.

For clarity, we reconstruct the consequence of connectivity that
we need. It is the mechanism of
\citet[Theorem 1, pp.~349--350 and Sections 5--6]{cyNori1993}.
If $W\in\mathrm{CH}^c(Q)_{\mathbb Q}$ becomes algebraically
trivial on the geometric generic $X$, clear denominators and
spread the finite equivalence witness after a finite extension
of $\mathbb Q(S)$. After shrinking, it gives a smooth base
$T$ finite \emph{\'etale} over an open of $S^\circ$, a smooth
projective family of curves $C\to T$, a divisor $\delta$ of
degree zero on every connected fibre component, and a
correspondence $\Gamma\in\mathrm{CH}^c(X_T\times_T C)$ with
\[
       j_T^*W=(p_X)_*(\Gamma\cdot p_C^*\delta)
                   \quad\text{in }\mathrm{CH}^c(X_T)_{\mathbb Q}.
\]
A disjoint union of curve families combines finitely many
relations. Rational-equivalence witnesses also spread after
shrinking, so the displayed equality is an actual Chow equality
on the total space, not merely an equality at its fibres.
Base change the finite-type data to $\mathbb C$. Both
$T\to S^\circ$ and $C\to S^\circ$ are smooth. Because $c<n$,
\eqref{eq:nori-effective-connectivity} makes
\[
 H^{2c}(Q\times C,\mathbb Q)
         \longrightarrow H^{2c}(X_T\times_T C,\mathbb Q)
\]
surjective, and makes
$H^{2c}(Q\times T,\mathbb Q)\to H^{2c}(X_T,\mathbb Q)$
injective. Lift $\operatorname{cl}(\Gamma)$ to a class $\beta$
on $Q\times C$, and push $\beta\,p_C^*\operatorname{cl}(\delta)$
along the proper curve projection $Q\times C\to Q\times T$.
The resulting class restricts to $j_T^*\operatorname{cl}(W)$,
so injectivity identifies it with the pullback of
$\operatorname{cl}(W)$. Restriction to one fibre over $T$ is
zero, because a degree-zero divisor on a projective curve has
zero cohomology class. We conclude
\begin{equation}\label{eq:nori-geometric-generic-injectivity}
 j^*W=0\text{ modulo algebraic equivalence over }\Omega
       \quad\Longrightarrow\quad\operatorname{cl}(W)=0.
\end{equation}
This proof deals directly with a putative finite witness on the
specified geometric generic fibre; it does not assume that
analytic exceptional loci are defined over $\overline{\mathbb Q}$.

The class $z$ is homologically trivial. Indeed, for every
$b\in H^{2n-2c}(X,\mathbb Q)$, weak Lefschetz gives
$b=j^*b_0$ for a class on $Q$, since $2n-2c<n$.
The projection formula gives
\[
 \int_X j^*\operatorname{cl}(Z)\,b
       =D^s\int_Q\operatorname{cl}(Z)\,h^s b_0=0,
\]
because $s\ge1$ and $h\operatorname{cl}(Z)=0$.
Poincar\'e duality proves the assertion. Equation
\eqref{eq:nori-geometric-generic-injectivity} proves $z\ne0$
modulo algebraic equivalence.

We now prove the stronger obstruction $z\notin L A^{c-1}_{\mathrm{alg}}(X)$.
Suppose otherwise and choose a finite extension $K'/K$ over
which a cycle $\alpha$ and the finite algebraic-equivalence
witness for $z=L\alpha$ are both defined. Trace to $K$ and
write $e=[K':K]$. The trace relation is
\[
       e z=L\alpha_{\mathrm{tr}}
              \quad\text{modulo algebraic equivalence over }\Omega.
\]
The universal incidence $\mathcal X\subset Q\times S$ is a
projective bundle over $Q$: it is the projectivization of the
kernel of the surjective evaluation map
\[
 V^{\oplus s}\otimes\mathcal O_Q
       \longrightarrow\mathcal O_Q(D)^{\oplus s}.
\]
Its relative hyperplane class is pulled back from $S$ and
restricts to zero on its generic fibre. Close
$\alpha_{\mathrm{tr}}$ in $\mathcal X$ and use the
projective-bundle formula. Restriction to the generic fibre
therefore writes
$\alpha_{\mathrm{tr}}=j^*a$ for
$a\in\mathrm{CH}^{c-1}(Q)_{\mathbb Q}$.
We obtain $j^*(eZ-ha)=0$ modulo algebraic equivalence.
Equation \eqref{eq:nori-geometric-generic-injectivity} contradicts
\eqref{eq:quadric-primitive-obstruction}. This proves the claim.

Adjoin a geometric generic hyperplane to $X$ and apply C.1h.
The obstruction survives this algebraically closed extension:
any putative preimage and relation over the larger field spread
over a finite-type model and specialize to a point over $\Omega$,
contradicting the obstruction just proved. The same finite-data
argument proves persistence under embeddings into $\mathbb C$.
Finally, when $c=n$, degree identifies
$A^n_{\mathrm{alg}}(X)$ with $\mathbb Q$; a point of $Y$
pushes to a class of degree one. Hence $A'$ is onto and
localization gives $A^n_{\mathrm{alg}}(X\setminus Y)=0$.
This proves all the asserted universal cases.\hfill$\square$

% Further C.1 deductions; owned by cycles. No preamble.
\paragraph{C.1j (Proved: the exact middle-positivity obstruction over $\mathbb C$).}
Let $X/\mathbb C$ be smooth projective of dimension $2m$, and put
\[
 P^m_{\mathrm{alg},L}(X)
 =\ker\bigl(L:A^m_{\mathrm{alg}}(X)\longrightarrow
                         A^{m+1}_{\mathrm{alg}}(X)\bigr).
\]
The form $q(x,y)=(-1)^m\deg(xy)$ is positive semidefinite on this
space. Its radical is exactly
\begin{equation}\label{eq:Balg-exact-gate}
 \operatorname{rad}(q)=
 \ker\bigl(L:G^m(X)\longrightarrow G^{m+1}(X)\bigr).
\end{equation}
Consequently the letter's $B(X,L)$ is equivalent to the injectivity
of this particular map on the rational Griffiths group.

\emph{Proof.}
The cycle class of $x\in P^m_{\mathrm{alg},L}(X)$ is a primitive
rational Hodge class of type $(m,m)$. The Hodge--Riemann bilinear
relations say that $(-1)^m\int_X\mathrm{cl}(x)^2$ is positive
unless $\mathrm{cl}(x)=0$; this is precisely the complex positivity
input recorded in \citet[Remark 3.10, p.~380]{cyKleiman1968}.
The intersection number agrees with that integral. Thus $q(x,x)\ge0$,
and it is zero exactly for $x\in G^m(X)$. Such an $x$ has intersection
zero with every algebraic cycle, so it is in the radical. Conversely
a radical vector has $q(x,x)=0$, hence has zero cycle class. Taking
the intersection with the defining primitive kernel proves
\eqref{eq:Balg-exact-gate}.

This proves $B$ for $m=0,1$, for every $X$ with $G^m(X)=0$, and more
generally whenever the displayed Griffiths map is injective. In all
complex dimensions the weaker nonnegative form is already a theorem.
Our counterexamples to Lefschetz \emph{surjectivity} do not supply a
nonzero vector in \eqref{eq:Balg-exact-gate}. We therefore leave the
universal strict-positivity assertion for $m\ge2$ undecided, rather
than infer its failure from the existence of a Griffiths group.

\paragraph{C.1k (Proved: exact universal ranges for three consequences).}
Fix $n\ge0$ and $0\le r\le n$. Over $\mathbb C$, each of the following
properties holds for every smooth projective connected $n$-fold
exactly when $r\in\{0,1,n\}$, with repeated or inadmissible entries
omitted:
\begin{enumerate}
\item the map $A^r_{\mathrm{alg}}(X)\to H^{2r}(X,\mathbb Q)$ is
injective;
\item algebraic equivalence up to torsion equals numerical equivalence
in codimension $r$;
\item $A^r_{\mathrm{alg}}(X)$ is finite dimensional over $\mathbb Q$.
\end{enumerate}
The corresponding assertion that cycles modulo algebraic equivalence
up to torsion form a finitely generated abelian group has the same
universal range.

For $r=0,n$, degree gives the assertions. For $r=1$, the group is
the N\'eron--Severi group modulo its finite torsion, which is finitely
generated. The divisor cycle map is injective after tensoring with
$\mathbb Q$; the Hodge index theorem identifies numerical and
homological equivalence for divisors. These are the usual
N\'eron--Severi and Lefschetz $(1,1)$ statements, also used in
C.1c--d and \citet[\S3, especially 3.4 and 3.10]{cyKleiman1968}.

Every remaining pair has $n\ge3$ and $2\le r\le n-1$. Take the
very general quintic $V$ of C.1a and
$X=V\times\mathbb P^{n-3}$. The projective-bundle formula supplies
the direct summand
\[
 G^2(V)t^{r-2}\ \subset\ G^r(X),
 \qquad 0\le r-2\le n-3.
\]
It is infinite dimensional by the opened Clemens theorem. Every
element is numerically trivial, and a nonzero one disproves (1) and
(2); the infinite-dimensional summand disproves (3) and finite
generation. This proof concerns universal assertions. For a specified
variety the first property is exactly $G^r(X)=0$, while equality with
numerical equivalence may additionally require the homological
standard conjecture $D$.

The letter's image criterion also fails in each of these target
codimensions, even for the zero correspondence: its cohomological
image contains the class of every $g\in G^r(X)$, but its cycle image
does not contain a nonzero $g$. More generally, that criterion for
all correspondences into a fixed $A^r_{\mathrm{alg}}(X)$ necessarily
implies $G^r(X)=0$. This necessary condition alone does not prove
the criterion, since an arbitrary cohomological preimage must still
be represented by an algebraic cycle.

\paragraph{C.1l (Proved: a nonzero radical in the correspondence algebra).}
If $G^r(X)\ne0$ for some $r$, the algebra
\[
 R=A^{n}_{\mathrm{alg}}(X\times X),\qquad n=\dim X,
\]
with composition of degree-zero correspondences, has nonzero
Jacobson radical. In particular it is not a semisimple algebra.
This supplies counterexamples to the letter's correspondence-ring
consequence in every dimension $n\ge3$.

Choose $0\ne g\in G^r(X)$, an ample class $h$, and
$d=\deg h^n>0$. Define
\[
 a=d^{-1}(h^{n-r}\times g)\in R.
\]
It acts by
$a_*(x)=d^{-1}\deg(xh^{n-r})g$ on $A^r_{\mathrm{alg}}(X)$.
In particular $a_*(h^r)=g$, so $a\ne0$ as a correspondence class.
For every $b\in R$, the usual composition formula for decomposable
correspondences gives
\[
 aba=d^{-1}\deg\bigl(h^{n-r}\,b_*g\bigr)a=0.
\]
Indeed $b_*g$ is homologically trivial, hence numerically trivial.
It follows that $(ba)^2=0$, so $1-ba$ is invertible, with inverse
$1+ba$, for every $b\in R$. The defining radical criterion for a
unital ring now gives $a\in\operatorname{Jac}(R)$. This proves
nonsemisimplicity; merely exhibiting a nilpotent element would not
have sufficed, since matrix algebras themselves contain nilpotents.
The product examples in C.1k finish the claimed dimension range.

There is also a direct obstruction to a semisimple category of
motives modulo algebraic equivalence. A nonzero Griffiths cycle
gives a nonzero morphism from a Tate motive to the appropriate
twist of $h(X)$, whereas composition with every morphism back is
zero by numerical triviality. In a semisimple category with
endomorphism ring $\mathbb Q$ for that Tate object, every nonzero
such morphism would split. Thus this category cannot be semisimple.
These arguments neither refute semisimplicity of numerical motives
nor claim semisimplicity for an individual variety without checking
its correspondences on the product. C.1m supplies the missing
universal surface case.

\paragraph{C.1m (Proved under an opened classical input: the exact universal
dimension range for finite semisimple correspondence algebras).}
For smooth projective connected complex $n$-folds, the assertion
that $A^n_{\mathrm{alg}}(X\times X)$ is a finite semisimple
$\mathbb Q$-algebra holds universally exactly for $n\le1$.

For $n=0$ the algebra is $\mathbb Q$. For a smooth projective
curve $C$, normalized divisor correspondences and a chosen point
give an algebra isomorphism
\[
 A^1_{\mathrm{alg}}(C\times C)
    \simeq\mathbb Q\times\operatorname{End}(J(C))_{\mathbb Q}
                        \times\mathbb Q.
\]
The three factors act on $H^0,H^1,H^2$, respectively. The
divisor cycle map is injective, so there are no further
off-diagonal components. The middle identification is the
classical Picard-correspondence identification used in C.4a.
Poincar\'e complete reducibility makes
$\operatorname{End}(J(C))_{\mathbb Q}$ a finite product of matrix
algebras over division algebras. Hence the displayed algebra
is finite and semisimple.

For a specified surface counterexample, let the fourteen
$t_\alpha$, indexed by $|\alpha|=4$ and
$\alpha\ne(4,0,0)$, be algebraically independent over
$\mathbb Q$, and put
\[
 K=\mathbb Q(t_\alpha),\qquad
 F=x_0^4+\sum_{\substack{|\alpha|=4\\\alpha\ne(4,0,0)}}
                      t_\alpha x^\alpha.
\]
Let $C\subset\mathbb P^2_{\overline K}$ be $F=0$.
It is smooth, since the discriminant is not identically zero
on this coefficient chart, as the Fermat quartic shows.
Choose a root $\rho\in\overline K$ of $F(T,1,0)$ and put
$p=(\rho:1:0)$. Embed $\overline K$ into $\mathbb C$.
We claim that the modified diagonal
\[
 \gamma=\Delta_{123}-\Delta_{12}-\Delta_{13}-\Delta_{23}
                     +\Delta_1+\Delta_2+\Delta_3
             \quad\text{in }\mathrm{CH}^{2}(C^3)
\]
has infinite order modulo algebraic equivalence. Here the
coordinates indexed by a subscript vary together and the other
coordinates equal $p$. The opened primary source
\citet[Proposition 3.1, p.~654, and the paragraph immediately
after Corollary 4.9, p.~659]{cyGrossSchoen1995} proves its
homological triviality and records the existence of such
genus-three curves, following the Ceresa and Colombo--van Geemen
results. These examples are nonhyperelliptic: the hyperelliptic
modified diagonal is torsion by their Corollary 4.9.
Thus they occur among smooth plane quartics. We use that
classical nonvanishing input explicitly;
we do not claim an independent proof of Ceresa's theorem or
that its original Annals PDF was obtained. The formula for
$\gamma$ is equation (0.1) of the same opened paper.

We justify the passage from that existence input to this
specified geometric generic curve. The class of $\gamma$
modulo algebraic equivalence is independent of $p$, since
the displayed seven-term formula varies algebraically as
$p$ ranges over the connected curve. Suppose that a nonzero
multiple of the geometric generic $\gamma$ were algebraically
trivial. Its finite cycle, parameter-curve, and rational-function
data descend to a finite extension of the generic coefficient
field. Norm and spreading then make a nonzero integer multiple
algebraically trivial over a nonempty open of the smooth
plane-quartic parameter space. Such a relation specializes to
every smooth fibre: choose an integral curve through the
desired parameter point and meeting that open, normalize,
make the finite base change needed for the witnesses, and
specialize at a point over the desired fibre. Algebraic
equivalence is preserved in this specialization. Concretely,
one closes the witnessing parameter curves and correspondences
over the resulting discrete valuation ring; the proper
connected special parameter curves, component by component,
give the specialized algebraic-equivalence relation. Multiplying
by the finite base-change degree has no effect on the
contradiction with infinite order. Every smooth plane quartic
can be put in the displayed coefficient chart by choosing
coordinates with nonzero $x_0^4$ coefficient. We would therefore
force the cited nonvanishing example to be torsion, a contradiction.
This proves the asserted non-torsion class for the specified
$C$, for any chosen embedding into $\mathbb C$; extension of
algebraically closed fields does not create a new finite
algebraic-equivalence relation, by the same finite-data
specialization argument.

Put $S=C\times C$, and view $\gamma$ as a degree-zero
correspondence $S\to C$. If $\pi:S\to C$ is the first
projection, put $a=\pi^*\circ\gamma$. In coordinates on
$S\times S=C^4$ this is the class $\gamma\times C$.
Pullback to a fixed point in the fourth factor recovers
$\gamma$, so $a\ne0$ modulo algebraic equivalence with
rational coefficients.
For every $b\in A^2_{\mathrm{alg}}(S\times S)$, the composition
\[
       \gamma\circ b\circ\pi^*
          \in A^1_{\mathrm{alg}}(C\times C)
\]
is homologically trivial. It is a divisor class, hence zero
modulo algebraic equivalence. Consequently
\[
 aba=\pi^*(\gamma b\pi^*)\gamma=0
                         \quad\text{for every }b.
\]
The radical argument in C.1l proves $0\ne a\in\operatorname{Jac}(R_S)$.
No nilpotence conjecture or finite-dimensional-motive hypothesis
is used. In particular this surface satisfies both of the
fixed-variety assertions $A(S)$ and $B(S)$, while its
correspondence algebra is not semisimple. The universal
auxiliary-product hypotheses in the letter's motivic argument
cannot be replaced by hypotheses on $S$ alone.

Finally take $X=S\times\mathbb P^{n-2}$ for $n\ge2$, with
projection $f:X\to S$ and a constant section $s:S\to X$.
The class $a_X=f^*a s^*$ is nonzero because
$s^*a_X f^*=a$. For every $b\in R_X$,
\[
 a_X b a_X=f^*a(s^*b f^*)a s^*=0.
\]
Thus it also lies in the nonzero Jacobson radical. This gives
the counterexample in every dimension $n\ge2$ and proves
the stated universal range. Notice that the Griffiths groups
of the surface $S$ itself vanish; it is the Griffiths group
of $S\times S$ that is detected here.

% Exact endpoint-image reduction, owned by cycles. No preamble.
\paragraph{C.1n (Proved: the endpoint image criterion is exactly the
universal standard conjecture $D$).}
We first work over $\mathbb C$ with rational coefficients. Fix integers
$n\ge0$ and $0\le r\le n$. Let $\mathcal I_{n,r}$ mean the
following assertion: for every smooth projective connected
$n$-fold $X$, every smooth projective connected $d$-fold $Y$,
every $0\le j\le d$, and every correspondence
\[
 \Gamma\in\mathrm{CH}^{d+r-j}(Y\times X)_{\mathbb Q},
\]
an element $z\in A^r_{\mathrm{alg}}(X)$ belongs to the image
of $\Gamma_*:A^j_{\mathrm{alg}}(Y)\to A^r_{\mathrm{alg}}(X)$
if and only if its cycle class belongs to the image of
\[
 \Gamma_*:H^{2j}(Y,\mathbb Q(j))
                      \longrightarrow H^{2r}(X,\mathbb Q(r)).
\]
These are the source and degree-shift quantifiers of the image
criterion on manuscript page $14{:}46$: its source degree is
arbitrary, and the correspondence need not have degree zero.
Write $D_{\mathbb C}$ for equality of numerical and Betti
homological equivalence on rational cycles of every smooth
projective complex variety. Then
\begin{equation}\label{eq:endpoint-image-D}
 \mathcal I_{n,r}\quad\Longleftrightarrow\quad D_{\mathbb C}
       \qquad\text{if }r\in\{0,1,n\}.
\end{equation}
For every $2\le r\le n-1$, the assertion $\mathcal I_{n,r}$
is false by C.1k. If we restrict the target to
$X=\mathbb P^n$, its image criterion is equivalent to
$D_{\mathbb C}$ for \emph{every} $0\le r\le n$.
More precisely, for a \emph{fixed} $X$ and $r$, let
$\mathcal I(X,r)$ require the displayed criterion for every
source $Y$, source degree $j$, and correspondence $\Gamma$.
Then
\begin{equation}\label{eq:fixed-target-image-D}
 \mathcal I(X,r)\quad\Longleftrightarrow\quad
             \bigl(D_{\mathbb C}\ \text{and}\ G^r(X)=0\bigr).
\end{equation}

\emph{Necessity, including the restricted projective-space target.}
Suppose that $D_{\mathbb C}$ fails. Then some smooth projective
$d$-fold $Y$ has a numerically trivial rational cycle
$a\in\mathrm{CH}^s(Y)_{\mathbb Q}$ with
$\mathrm{cl}(a)\ne0$. Put $j=d-s$, let $h$ be the hyperplane
class on $\mathbb P^n$, and set
\[
 \Gamma=a\times h^r\in
       \mathrm{CH}^{s+r}(Y\times\mathbb P^n)_{\mathbb Q}.
\]
The codimension is $s+r=d+r-j$, as required. The action on
cycles is
\[
 \Gamma_*b=\deg(a\cdot b)h^r=0
       \qquad\bigl(b\in A^j_{\mathrm{alg}}(Y)\bigr),
\]
by numerical triviality. In contrast, Poincar\'e duality supplies
$v\in H^{2d-2s}(Y,\mathbb Q(d-s))$ with
$\int_Y\mathrm{cl}(a)v=1$. Hence
$\Gamma_*v=\mathrm{cl}(h^r)$, although the nonzero cycle class
$h^r\in A^r_{\mathrm{alg}}(\mathbb P^n)$ is not in the cycle
image. This contradicts the proposed criterion. Notice that
we neither need nor assert that this cohomological preimage
$v$ is algebraic: the criterion itself would force an
algebraic preimage, which numerical triviality excludes.
The same construction works for any fixed target $X$: take
$h$ ample on $X$; its power $h^r$ has nonzero cohomology
class, since $\int_Xh^n>0$. Thus $\mathcal I(X,r)$ also
implies $D_{\mathbb C}$. Taking $\Gamma=0$ shows separately
that $\mathcal I(X,r)$ implies $G^r(X)=0$.

\emph{Sufficiency.}
Assume $D_{\mathbb C}$. By \citet[Theorem~1, pp.~448--449]
{Jannsen1992}, the category of rational numerical motives
is semisimple abelian. Betti realization on this category is
well defined and faithful: $D_{\mathbb C}$ applies on every
product, and K\"unneth together with Poincar\'e duality
identifies a correspondence's total cohomology class with
its action on all cohomological degrees. Here K\"unneth is
only the cohomological decomposition; we do not assume
algebraic K\"unneth projectors.

We use the convention
\[
 \begin{split}
 \operatorname{Hom}(\mathbf1,h(X)(r))
       &=A^r_{\mathrm{num}}(X),\\
 \operatorname{Hom}(h(Y)(j),h(X)(r))
       &=A^{d+r-j}_{\mathrm{num}}(Y\times X).
 \end{split}
\]
Thus $\Gamma$ induces a morphism
$f:M=h(Y)(j)\to N=h(X)(r)$, and the numerical class of
$z$ is a morphism $\beta:\mathbf1\to N$. Semisimplicity
gives $g:N\to M$ for which $p=fg$ is a projector onto
$\operatorname{im}(f)$ and $pf=f$. If $\mathrm{cl}(z)$ is
in the stated cohomological image, then
$H((1-p)\beta)=0$. Indeed realization in degree zero of
$h(Y)(j)$ and $h(X)(r)$ is precisely the pair of groups in
the displayed image criterion. Faithfulness gives
$(1-p)\beta=0$, so
\[
                     \beta=f(g\beta).
\]
Choose a rational cycle $b\in\mathrm{CH}^j(Y)_{\mathbb Q}$
representing $g\beta$. We have proved that $\Gamma_*b$ and
$z$ are numerically equivalent. In target codimensions
$r=0,n$, numerical and algebraic equivalence with rational
coefficients are both determined by degree. In codimension
$r=1$, the same equality follows from the N\'eron--Severi
theorem and the divisor intersection pairing, as recalled
in C.1k. Therefore $\Gamma_*b=z$ in
$A^r_{\mathrm{alg}}(X)$ in every asserted endpoint case.
The converse implication is always given by functoriality
of the cycle class. For $X=\mathbb P^n$, the map
$A^r_{\mathrm{alg}}(X)\to A^r_{\mathrm{num}}(X)$ is an
isomorphism in every codimension, so the same proof gives
the final restricted-target assertion.
For the fixed-target assertion, $D_{\mathbb C}$ identifies
numerical and homological equivalence, and $G^r(X)=0$
identifies homological and algebraic equivalence in the
specified target group. The same numerical preimage therefore
proves \eqref{eq:fixed-target-image-D}.

These two equivalences are not restricted to complex geometry.
Let $k$ be any algebraically closed field and let $H$ be a
fixed Weil cohomology with characteristic-zero coefficient
field $F$. Keep rational cycle coefficients, use $F$ in the
cohomological image, and define $D_{k,H}$ by equality of
numerical and $H$-homological equivalence for rational cycles
on every smooth projective $k$-variety. The identical argument
proves
\[
 \mathcal I_{n,r}^{k,H}\ \Longleftrightarrow\ D_{k,H}
         \qquad(r\in\{0,1,n\}),
\]
\[
 \mathcal I^{k,H}(X,r)\ \Longleftrightarrow\
       \bigl(D_{k,H}\ \text{and}\ G_H^r(X)=0\bigr),
\]
where $G_H^r(X)$ is the kernel of the rational-cycle class map
into $H^{2r}(X)(r)$. The source functional is evaluated using
Poincar\'e duality over $F$; the numerical-motive category
and its image splitting still have rational coefficients.
Numerical and algebraic equivalence agree rationally for
divisors over any algebraically closed field, and degree
handles codimensions zero and top. The intermediate-codimension
universal counterexamples cited above remain assertions over
$\mathbb C$; this extension does not transfer their
nonvanishing to positive characteristic.

Equations~\eqref{eq:endpoint-image-D}--\eqref{eq:fixed-target-image-D}
isolate an exact remaining
standard-conjecture gate. It does not prove $D_{\mathbb C}$
and does not replace homological equivalence by algebraic
equivalence in intermediate target codimension. It also
explains why the absence of a target Griffiths group by
itself was insufficient for the image argument in C.1k:
the unrestricted source can still detect a failure of $D$.

% Exact positivity reductions; owned by cycles. No preamble.
\paragraph{C.1o (Proved: projective-space stabilization identifies the exact positivity kernel).}
Let $X/\mathbb C$ be smooth projective of dimension $n$, let $h$ be an
ample rational divisor class, and put $t=c_1(\mathcal O_{\mathbb P^d}(1))$.
Suppose that $n+d=2m$ is even. For $0\le d\le n$ set
$i=(n-d)/2$. There is an explicit isomorphism
\begin{equation}\label{eq:C1o-radical}
 \begin{split}
 \ker\bigl(h^{d+1}:G^i(X)\longrightarrow G^{i+d+1}(X)\bigr)
 &\xrightarrow{\ \sim\ }
 \operatorname{rad}\bigl(q_{X\times\mathbb P^d,h+t}\bigr),\\
 z&\longmapsto
 \sum_{j=0}^{d}(-1)^j h^jz\,t^{d-j}.
 \end{split}
\end{equation}
Here $q$ is the middle primitive form of C.1j, with cycles modulo
algebraic equivalence and rational coefficients. If $d>n$, that radical
is zero. The statement uses the actual product polarization $h+t$;
it is unchanged after a positive rational rescaling of either factor.

\begin{proof}
The projective-bundle formula descends to both algebraic and
homological equivalence. It gives
\[
 G^m(X\times\mathbb P^d)=
       \bigoplus_{b=0}^{d}G^{m-b}(X)t^b.
\]
Write $\Gamma=\sum_{b=0}^d a_bt^b$ in this direct sum.
The equation $(h+t)\Gamma=0$ is exactly
\[
 ha_0=0,\qquad ha_b+a_{b-1}=0\quad(1\le b\le d).
\]
Hence all coefficients are determined by $z=a_d$:
$a_{d-j}=(-h)^jz$, and the remaining equation is $h^{d+1}z=0$.
Its degree is $m-d=(n-d)/2=i$. Conversely those formulas give
\[
 (h+t)\sum_{j=0}^{d}(-1)^j h^jz\,t^{d-j}
                         =(-1)^d h^{d+1}z,
\]
because $t^{d+1}=0$. The coefficient of $t^d$ recovers $z$, so
the map is injective as well as surjective. If $d>n$, its putative
last coefficient has negative codimension and vanishes; the same
recursion then kills every coefficient. C.1j identifies this primitive
Griffiths kernel with the radical, proving the claim. For a polarization
$h+bt$ with $b>0$, divide by $b$ and apply the same argument to
$h/b+t$. Positive rescaling does not change the kernel or strict
positivity.
\end{proof}

It follows that the following two \emph{universal} assertions over
$\mathbb C$ are equivalent:
\begin{enumerate}
\item the letter's strict middle positivity $B$ holds for every smooth
projective variety of even dimension and every ample class;
\item for every smooth projective $n$-fold $X$, every ample $h$, and
every $0\le 2i\le n$, the map
\begin{equation}\label{eq:C1o-shifted}
 h^{n-2i+1}:G^i(X)\longrightarrow G^{n-i+1}(X)
 \quad\text{is injective}.
\end{equation}
\end{enumerate}
For the forward implication take $d=n-2i$ in
\eqref{eq:C1o-radical}; for the reverse implication take $d=0$.
The forward implication also gives the injectivity part of the
letter's $A$: if $h^{n-2i}x=0$ modulo algebraic equivalence,
cohomological hard Lefschetz first puts $x$ in $G^i(X)$, and
\eqref{eq:C1o-shifted} then kills it. No surjectivity assertion follows
from this argument.

There is an unconditional and useful test family. For every smooth
projective threefold $V$, every odd $d\ge1$, and every ample class on
$V\times\mathbb P^d$, the letter's $B$ holds. Indeed, the only
nonnegative degrees $i=(3-d)/2$ in the formula are $i=1$ and $i=0$,
and $G^1(V)=G^0(V)=0$. Any ample class on this product has the form
$h+bt$ with $h$ ample on $V$ and $b>0$, by the projective-bundle
description of its Picard group and restriction to its two factors.
Taking the specified quintic threefold of C.1g therefore gives, in
every even dimension at least four, a variety satisfying strict
middle positivity while possessing a nonzero Griffiths group and
a nonzero radical in its correspondence algebra by C.1l.
With a very general quintic, Clemens's theorem in C.1k strengthens
the Griffiths group to infinite dimension. This refutes the
\emph{individual-variety} inference from its own middle positivity
to the corresponding conclusions. It does not refute
an implication whose premise is positivity for \emph{all} varieties
and products; that universal premise remains precisely the
unproved assertion \eqref{eq:C1o-shifted}.

\paragraph{C.1p (Proved: the exact Beauville-graded positivity test on abelian stabilizations).}
Fix a complex abelian variety $A$ of dimension $g$. The following
assertions are equivalent:
\begin{enumerate}
\item the letter's $B$ holds for every even-dimensional
$A\times\mathbb P^d$ and every ample class on it;
\item $G^r(A)_{(s)}=0$ for every $r$ and every $s\le0$.
\end{enumerate}
Equivalently, the quotient
$B^*(A)=\mathrm{CH}^*(A)_{\mathbb Q}/
\mathrm{CH}^*(A)_{\mathrm{alg},\mathbb Q}$ has no negative
Beauville pieces, and its grade-zero cycle map is injective.
This is a statement modulo algebraic equivalence. We do not identify
it with the corresponding conjectures for the ordinary Chow group.

\begin{proof}
We use Beauville's algebraic $\mathfrak{sl}_2$-action, as in C.3h
\citep[Theorem 4.2 and Proposition 5.2, pp.~260--261]{rootBeauville2010}.
His footnote 1 on p.~254 permits any adequate equivalence.
The correspondences preserve the homologically trivial subgroup and
algebraic equivalence, so $G^*(A)$ is a direct sum of finite
dimensional strings
\[
 z,\theta z,\ldots,\theta^{g+s-2a}z,
 \qquad z\in G^a(A)_{(s)},\quad Fz=0,
 \quad g+s-2a\ge0,
\]
where $\theta$ is the symmetric ample divisor and $F$ is the lowering
operator. If a nonzero piece with $s\le0$ exists, choose a nonzero
lowest vector in one of its strings. Then $2a\le g$, and with
$d=g-2a$ we have
\[
                   \theta^{d+1}z=0,
       \qquad g+s-2a\le d.
\]
Formula \eqref{eq:C1o-radical} supplies a nonzero primitive radical
on $A\times\mathbb P^d$, disproving (1).

Conversely suppose (2). Every remaining string has $s\ge1$.
For a degree-$r$ element $\theta^{r-a}z$, the power
$\theta^{g-2r+1}$ could kill it only if
\[
 r-a+g-2r+1>g+s-2a,
 \quad\text{that is,}\quad a>r+s-1.
\]
But $a\le r$ and $s\ge1$, so this is impossible. The direct sum of
strings therefore gives all injectivities
\eqref{eq:C1o-shifted} for this $A$, and C.1o proves (1).
An ample divisor and its symmetric component have the same class
modulo algebraic equivalence, so this also handles nonsymmetric
ample classes. The description of ample classes on a product with
projective space handles all the polarizations in (1).

For $s\ne0$, a Beauville eigenclass has zero cohomology class:
multiplication by an integer $q>1$ acts on its cycle class both by
$q^{2r-s}$ and by $q^{2r}$. Thus $G^r(A)_{(s)}=B^r(A)_{(s)}$
for $s\ne0$, while $G^r(A)_{(0)}$ is exactly the kernel of the
grade-zero cycle map. This proves the equivalent formulation.
\end{proof}

The same string argument gives precise first tests without assuming
any Beauville conjecture. For $g\le3$, every string with $s\le0$
would have lowest degree $a\le1$, where the Griffiths group is zero;
hence (1) holds unconditionally. For an abelian fourfold,
\[
 \operatorname{rad}(q_{A,\theta})=G^2(A)_{(0)}.
\]
For an abelian fivefold,
\[
 \operatorname{rad}(q_{A\times\mathbb P^1,\theta+t})
       \simeq G^2(A)_{(-1)}\oplus G^2(A)_{(0)}.
\]
Indeed the only possible lowest degree is $a\ge2$: in dimension
four, $s\le0$ forces $(a,s)=(2,0)$ and string length zero; in
dimension five it permits only $(a,s)=(2,-1)$, of length zero,
or $(2,0)$, of length one. Positive-index strings contribute no
kernel to the displayed maps. The two graded groups in the last
display and the grade-zero group in dimension four are not proved
to vanish here. This exact reduction sharpens the description of
the outstanding question; it is recorded as an equivalence, not as
a closure of the universal positivity problem.

% Exact audit of the blow-up route; owned by cycles. No preamble.
\paragraph{C.1q (Proved: the exact exceptional polynomial and the ample-pencil test).}
Let $M/\mathbb C$ be smooth projective of dimension $2m$, let
$i:Z\hookrightarrow M$ be a smooth connected center of codimension
$c\ge2$, and put $\pi:X=\operatorname{Bl}_Z(M)\to M$.
Write $j:E\hookrightarrow X$ and $p:E\to Z$ for the exceptional
divisor and its projection, and set
\[
 N=N_{Z/M},\qquad \xi=c_1(\mathcal O_{\mathbb P_{\mathrm{lines}}(N)}(1)).
\]
Assume $G^m(M)=G^{m+1}(M)=0$; in particular this holds when
$M$ has a Tate Chow motive. Let
$L=\pi^*D-bE$ be an ample rational class, so $b>0$, and put
$u=i^*D/b$. Define
\[
 q_N(T)=\sum_{k=0}^{c-1}c_k(N)T^{c-1-k},\qquad c_0(N)=1.
\]
Then extraction of the last exceptional coefficient gives an
isomorphism
\begin{equation}\label{eq:C1q-kernel}
 \operatorname{rad}(q_{X,L})\simeq
 \ker\left(q_N(-u):G^{m-c+1}(Z)\longrightarrow G^m(Z)\right).
\end{equation}
Thus Griffiths classes strictly above the middle degree of a
center, by themselves, do not create a primitive middle radical:
the controlling degree is
\[
             m-c+1=\frac{\dim Z-c+2}{2}\le\frac{\dim Z}{2}.
\]
No condition on the Chern classes of the normal bundle can change
that degree constraint.

\begin{proof}
We use the opened blow-up formula
\citep[Theorem 2.8(3), and equation (2.2)]{cyJiang2023}.
Its maps and inverse are algebraic correspondences, so the
formula descends to homological equivalence and to algebraic
equivalence, and hence to the rational Griffiths quotient.
Every $\Gamma\in G^m(X)$ has a unique expression
\[
 \Gamma=\sum_{j=1}^{c-1}f_j(x_j),\qquad
 f_j(x)=j_*(\xi^{j-1}p^*x),\quad x_j\in G^{m-j}(Z).
\]
The excess bundle is $Q=p^*N/\mathcal O_E(-1)$. Its top Chern
class is
\[
 c_{c-1}(Q)=\sum_{k=0}^{c-1}p^*c_k(N)\xi^{c-1-k}.
\]
The excess identity therefore gives
\[
 f_c(z)=\pi^*i_*z-\sum_{k=1}^{c-1}f_{c-k}(c_k(N)z).
\]
For the class $z$ used below, the ambient term lies in
$G^{m+1}(M)$ and is zero. Also $j^*E=-\xi$, so
\[
 b^{-1}L\,f_j(x)=f_j(ux)+f_{j+1}(x).
\]
These formulas fix the signs and include the top exceptional
relation; ignoring that relation would give an incorrect kernel.

Set $z=x_{c-1}$ and $x_0=0$. The coefficient of $f_k$ in
$b^{-1}L\Gamma$ is
\[
       ux_k+x_{k-1}-c_{c-k}(N)z\qquad(1\le k\le c-1).
\]
The equations for $k=c-1,\ldots,2$ determine every coefficient:
\[
 x_{c-1-j}=\sum_{k=0}^{j}c_k(N)(-u)^{j-k}z
                    \qquad(0\le j\le c-2).
\]
The equation for $k=1$ is exactly $q_N(-u)z=0$.
Conversely these coefficients give a primitive class whenever
that last equation holds. Uniqueness of the blow-up decomposition
makes extraction of $z$ inverse to this construction. C.1j now
identifies the kernel with the positivity radical. This proves
\eqref{eq:C1q-kernel}, including the cases of negative or zero
controlling degree, where the relevant Griffiths group vanishes.
\end{proof}

There is a genuine ample test, not just a formal exceptional
class. Suppose $Z$ is the smooth transverse base locus of two
sections of an ample line bundle $H$ on $M$, so $c=2$ and
$N=H|_Z\oplus H|_Z$. The pencil defines $f:X\to\mathbb P^1$,
and its closed graph realization gives
\[
 L=2\pi^*H-E=\pi^*H+f^*\mathcal O_{\mathbb P^1}(1)
                         \quad\text{ample}.
\]
Indeed this is the restriction of the ample external product
on $M\times\mathbb P^1$ to the closed graph. Here $u=2H|_Z$
and $q_N(-u)=c_1(N)-u=0$. Consequently
\begin{equation}\label{eq:C1q-pencil}
       \operatorname{rad}(q_{X,L})\simeq G^{m-1}(Z).
\end{equation}
A nonzero middle Griffiths class on such a pencil base would
therefore refute universal $B$. This is a criterion, not an
assertion that such a class has been constructed. If $M$ has
Tate cohomology, weak Lefschetz gives
$H^{2m-3}(Z,\mathbb Q)=0$; the required class would have zero
ordinary Abel--Jacobi invariant. The above-middle Nori and
quintic classes already established in this paper do not
supply this middle class.

For completeness, we can rule out a substantial stabilized
quintic attempt with its actual ample classes. Let $V\subset
\mathbb P^4$ be any smooth quintic threefold, let $r\ge0$,
$c\ge2$, and assume $r+c+3=2m$. Set
\[
 M=\mathbb P^4\times\mathbb P^{r+c-1},\qquad
 Z=V\times\mathbb P^r,
\]
using a linear inclusion for the second factor. Let $H,t$
be the two hyperplane classes. Every ample rational class on
$X=\operatorname{Bl}_Z(M)$ has the form
$L=\alpha\pi^*H+\beta\pi^*t-bE$ with $b>0$.
It also has $\beta>b$: take a line in a fibre over a point
of $V$ meeting the linear $\mathbb P^r$ once, without being
contained in it. Its strict transform has $L$-degree $\beta-b$.

The only possibly nonzero Griffiths group of $V$ is $G^2(V)$,
and $H G^2(V)=0$ because $G^3(V)=0$. Thus the controlling
group in \eqref{eq:C1q-kernel} is zero unless
\[
 e=\frac{r-c+1}{2}\ge0,
 \quad\text{in which case it is }G^2(V)t^e.
\]
Here
$N=\mathcal O_V(5)\boxplus
\mathcal O_{\mathbb P^r}(1)^{\oplus(c-1)}$.
On the Griffiths summand the terms containing $H$ vanish,
and the binomial theorem gives
\[
 q_N(-u)(g t^e)
      =(1-\beta/b)^{c-1}g t^{e+c-1}.
\]
Since $e\ge0$ implies $e+c-1\le r$, the last projective
coefficient is a genuine direct summand. Its coefficient
is nonzero because $\beta>b$. Therefore this map is
injective, and every such blow-up satisfies its own strict
middle positivity assertion. The apparent kernel at
$\beta=b$ lies outside the ample cone. This explicitly
records why this blow-up attempt does not refute universal
$B$; it neither proves that universal assertion nor rules
out other centers satisfying \eqref{eq:C1q-pencil}.


\paragraph{Remaining work for C.1.}
We have rigorously deduced the quintic refutations from the opened
Clemens theorem, proved the exact divisor criterion, and classified
the universal truth of $A$ in every $(n,i)$ over $\mathbb C$ by
C.1f. C.1g supplies the full line argument and a specified geometric
generic counterexample, also sufficient for C.1f without the stronger
infinite-dimensional input. C.1h--i give the exact geometric-generic
Gysin image and classify all the universal $A',A''$ cases in their
stated range, with explicit complete-intersection counterexamples.
C.1j identifies the exact primitive Griffiths-kernel obstruction to
$B$; C.1k classifies universal cycle-class injectivity, equality
with numerical equivalence, and finite generation. C.1l--m give
explicit nonzero radicals, including a specified geometric generic
plane-quartic surface example, and classify universal finite
semisimplicity of the correspondence algebra: precisely dimensions
zero and one. C.1n identifies the remaining endpoint image
criterion with the universal standard conjecture $D_{\mathbb C}$,
even if the target is restricted to projective space. C.1o identifies
universal $B$ with the precise shifted Lefschetz injectivity on
Griffiths groups, and gives unconditional threefold stabilizations.
C.1p identifies the corresponding abelian test with vanishing of
the nonpositive Beauville pieces of the Griffiths quotient, including
the exact first fourfold and fivefold tests. These equivalences do
not prove the conjectural vanishings. C.1q computes the full
exceptional normal-bundle polynomial, identifies the genuinely
ample pencil-base candidate, and proves that linear stabilized
quintic blow-ups cannot furnish the desired primitive class.
The higher
universal $B$ assertion and its universal positivity-only
consequences remain undecided here, as does the truth of that
standard conjecture.
The full task therefore remains incomplete.

\subsection{C.2: what changes modulo homological equivalence}

\paragraph{Homological definitions.}
The homological versions are
\[
\begin{array}{ll}
A_H:& L^{n-2r}:A^r_{\mathrm{hom}}(X)\xrightarrow{\sim}
                  A^{n-r}_{\mathrm{hom}}(X),\quad 2r\le n,\\
A'_H:&i_*:A^r_{\mathrm{hom}}(Y)\twoheadrightarrow
                  A^{r+1}_{\mathrm{hom}}(X),\quad 2r\ge n-1,\\
A''_H:&A^r_{\mathrm{hom}}(X\setminus Y)=0,\quad 2r>n.
\end{array}
\]
Here the group on the open variety uses its ordinary cohomology.
The letter's $B_H$ is the positive-definiteness of
$(-1)^m\deg(xx')$ on primitive algebraic middle classes of a
$2m$-fold. We denote its extension to all primitive algebraic degrees
by $H_{\mathrm{std}}$, following the modern standard-conjecture
terminology. We reserve $B_{\mathrm{std}}$ for algebraicity of the inverse
Lefschetz operator. The two letters $B$ must not be conflated.

\paragraph{C.2b (Proved: affine vanishing; the missing localization step).}
For a smooth affine complex variety $U$ of dimension $n$,
$H^j(U,\mathbb Q)=0$ for $j>n$. The algebraic de Rham complex has terms
only in degrees $0,\ldots,n$; all its terms are quasi-coherent and thus
acyclic on $U$. Its hypercohomology is consequently computed by a complex
of global sections of length $n$. Comparison with analytic de Rham
cohomology proves the assertion. In arbitrary characteristic the
corresponding $\ell$-adic vanishing is Artin's affine theorem
\citep[Expos\'e XIV, Corollaire 3.2]{cySGA4XIV}; passing from finite
coefficients to $\mathbb Q_\ell$ gives the claimed range. Hence $A''_H$
is a theorem.

The cohomological localization sequence yields, in that range, a
surjection
\[
 H^{2r-2}(Y)(-1)\longrightarrow H^{2r}(X).
\]
Its surjectivity does not say that an \emph{algebraic} class on $X$ has
an \emph{algebraic} preimage on $Y$. Chow localization applies only when
the restricted cycle is zero in the Chow group, whereas $A''_H$ merely
says that its restricted \emph{cohomology class} is zero. Accordingly
the proof of \eqref{eq:alg-localization} cannot be reused for homological
equivalence. This identifies a precise logical gap in the proposed
deduction $A''_H\Rightarrow A'_H$; it is not a geometric counterexample
to $A'_H$ itself.

\paragraph{C.2c (Known comparison and proved positivity implication).}
Hard Lefschetz proves injectivity in $A_H$; its surjectivity is the
standard conjecture $A$. The global equivalence with algebraicity of
the inverse Lefschetz operator is recorded in
\citet[Corollary 2.14]{cyKleiman1968}. More precisely,
$A(X\times X)$ implies $B_{\mathrm{std}}(X)$, whereas
$B_{\mathrm{std}}(X)$ implies $A(X)$ and the algebraicity of the
K\"unneth projectors
\citep[Theorem 4-1]{cyKleiman1994}. The product hypothesis is material.

Assume $A_H(X)$ and $H_{\mathrm{std}}(X)$ in all the degrees used in the
Lefschetz decomposition. Each algebraic class decomposes into algebraic
primitive classes: the inverse powers of $L$ needed in the recursive
decomposition preserve algebraic classes by $A_H$. Define the usual
Lefschetz star by reversing each $L$-string with its Hodge--Riemann sign.
It again preserves algebraic classes. Positivity says that, for a
nonzero algebraic class $x$, the complementary algebraic class $*x$
satisfies $\deg(x\,*x)>0$. Such an $x$ cannot be numerically zero.
Thus homological and numerical equivalence agree under these two
hypotheses. This is the mechanism behind the implication
$A+H_{\mathrm{std}}\Rightarrow D$; removing algebraicity of the
Lefschetz decomposition removes a used premise.

Conversely, if $D(X)$ holds in complementary degrees, the numerical
pairing identifies the two finite-dimensional numerical cycle spaces
as duals. They therefore have the same dimension. Hard Lefschetz
injectivity then makes $A_H(X)$ an isomorphism. This argument is about
the full numerical pairing, not the asserted positivity of a single
primitive middle space.

\paragraph{C.2d (Known semisimplicity theorem; audited quantifiers).}
\citet[Theorem 1]{cyJannsen1992} proves semisimplicity of pure numerical
motives with rational coefficients and characterizes numerical
equivalence by semisimplicity among adequate equivalence relations.
If homological and numerical equivalence agree on the products defining
all morphism spaces, the homological category is identified with that
numerical category and is therefore semisimple. The letter's archive
wording, ``A et B'', is now checked directly in the original typed
scan 46, archive sheet 45
\citep[scan 46]{GrothendieckCote14Scan}. It is consistent with the
positivity argument above. The opened online English translation
instead says ``B also implies'' at printed p.~258
\citep[p.~258]{rootGrothendieckSerreEnglish}; we distinguish that
witness from the original and from the French printed volume.
That shorter wording does not justify deleting the
Lefschetz hypothesis when $B$ means only the positivity assertion that
the letter actually states. The category modulo algebraic equivalence
cannot be declared semisimple by that argument: C.1a supplies a genuine
difference of the equivalence relations.

The conditional use of standard conjectures for the Weil estimates is
set out in \citet[Sections 2--4]{cyGrothendieck1969} and
\citet[Propositions 3.6, 3.8 and Corollary 3.9]{cyKleiman1968};
C.2h--j below reconstruct the conditional estimates and the image
criterion. We do not present the now established Weil conjectures as a proof
of the still conjectural algebraicity of general Lefschetz inverses.

% Independent pencil argument by /root; checked by /root/cycles and /root/frobenius.
\paragraph{C.2e (Proved under the stated field-extension hypothesis): $A'_H$ implies $A_H$.}
Fix an algebraically closed field $k$ of arbitrary characteristic and a
prime $\ell\ne\operatorname{char}(k)$. We use $\ell$-adic cohomology,
with rational cycle coefficients, and write $A^r_H$ for the image of
$\mathrm{CH}^r(-)_{\mathbb Q}$ in $H^{2r}(-,\mathbb Q_\ell(r))$.
Assume $A'_H$ for all smooth projective varieties and smooth ample
hyperplane sections over every algebraically closed extension of $k$.
Then $A_H$ holds over all those extensions. In particular the universal
implication in the letter is valid. The extension hypothesis is used
at the geometric generic fibre of a pencil; we do not replace it by
an unsupported assertion about specializing algebraic preimages.

\begin{proof}
We prove $A_H(X,i)$ first by induction on $i$, and, for fixed $i$,
by induction on $n=\dim X\ge2i$. The induction includes all the
specified base fields simultaneously. For $i=0$, both cycle spaces are
one-dimensional by degree, and $L^n$ has positive nonzero degree.
For any $i$, the case $n=2i$ is the identity. Replacing an ample line
bundle by a sufficiently high very ample power does not change the
assertion: it multiplies each relevant power of $L$ by a nonzero
rational scalar. We therefore work with a projective embedding.

Let $n\ge2i+1$, set $\kappa=n-2i$, and take
$z\in A_H^{n-i}(X)$. Choose a pencil whose base locus $B$ is a smooth
codimension-two complete intersection and whose generic member is
smooth. Bertini supplies this over the algebraically closed base.
No assertion about ordinary double points of the other members is
needed. Write
\[
 \pi:\widetilde X=\operatorname{Bl}_B X\longrightarrow X,
 \qquad f:\widetilde X\longrightarrow\mathbb P^1,
 \qquad K=k(t),
\]
and let $Y_\eta$ be the generic fibre, viewed as a hyperplane section of
$X_K$. Over $\overline K$, the hypothesis $A'_H$ gives a class
$b\in A_H^{n-i-1}(Y_{\overline\eta})$ whose Gysin image is $z$;
its degree lies in the required range since
$2(n-i-1)\ge n-1$. The induction on dimension gives
\[
 b=L_Y^{\kappa-1}\gamma,
 \qquad \gamma\in A_H^i(Y_{\overline\eta}).
\]
Here $n-1\ge2i$, including the critical case $n=2i+1$, where this
induction step is just the identity map. Thus
\begin{equation}\label{eq:C2-generic-push}
             z=L^{\kappa-1}(i_\eta)_*\gamma.
\end{equation}

A rational cycle representing $\gamma$ is defined over a finite
extension $K'/K$. Transfer divided by $[K':K]$ gives a rational cycle
on $Y_\eta$ for which \eqref{eq:C2-generic-push} still holds: the
class $z$ is fixed under base-field automorphisms, and pushforward,
cycle class and transfer commute. Purely inseparable factors multiply
pushforward by their degrees and induce the usual isomorphism on
prime-to-characteristic cohomology, so cause no exception. We keep
$\gamma$ for this averaged cycle and close it in $\widetilde X$.
The Chow blow-up formula, with the explicit maps of
\citet[Theorem 2.8(3)]{cyJiang2023}, writes its closure as
\[
 \Gamma=\pi^*a+j_*q^*\beta,
 \qquad a\in\mathrm{CH}^i(X)_{\mathbb Q},\quad
 \beta\in\mathrm{CH}^{i-1}(B)_{\mathbb Q},
\]
where $j:E\hookrightarrow\widetilde X$ is the exceptional divisor
and $q:E\to B$ its projection. The pushforward of the restriction to
one pencil member is
\begin{equation}\label{eq:C2-exceptional}
                  (i_\eta)_*\gamma=L a+(i_B)_*\beta.
\end{equation}
For clarity, the fibre class is $F=\pi^*L-E$.
We have $\pi_*(F\pi^*a)=La$, since $\pi_*E=0$.
Moreover $F|_E$ has degree one on each exceptional $\mathbb P^1$,
so $q_*(F|_E)=1$ and
$\pi_*(Fj_*q^*\beta)=(i_B)_*\beta$.
This proves both the sign and the coefficient in
\eqref{eq:C2-exceptional}.

Weak Lefschetz for the two hyperplane sections identifies
\[
 H^{2i-2}(X,\mathbb Q_\ell(i-1))
       \xrightarrow{\sim}
 H^{2i-2}(B,\mathbb Q_\ell(i-1)),
\]
since $2i-2<n-2=\dim B$. Let $b_0$ be the preimage of the cycle
class of $\beta$. The projection formula gives
$(i_B)_*\beta=L^2b_0$. The class
\[
 L^{\kappa+2}b_0=L^\kappa(i_B)_*\beta
\]
is algebraic. The already proved induction step with codimension
$i-1$ on $X$ makes $L^{\kappa+2}$ surjective from
$A_H^{i-1}(X)$ onto $A_H^{n-i+1}(X)$. Hard Lefschetz on the full
cohomology makes it injective. Consequently $b_0$ itself is algebraic.
Combining the two displayed pushforward formulas now gives
\[
                 z=L^\kappa(a+L b_0),
\]
which proves surjectivity. Hard Lefschetz supplies injectivity and
completes both inductions.
\end{proof}

This proof explains the role of a nonclosed intermediate field in the
letter. The actual cycle to be closed lives over $k(t)$; existence of
its preimage is applied over $\overline{k(t)}$, then finite transfer
returns it to $k(t)$. If one assumes $A'_H$ only for varieties over a
single fixed algebraically closed field, extension stability is an
additional premise in this argument. An assertion for every closed
member of a pencil is not substituted for the needed generic one.

\paragraph{C.2f (Proved): the converse and the exact affine-vanishing gate.}
For a smooth ample hyperplane section $i:Y\hookrightarrow X$,
$A_H(X)$ already implies $A'_H(X,Y)$. Indeed, for $2r\ge n-1$,
write an algebraic class $z\in A_H^{r+1}(X)$ as
\[
 z=L^{2r+2-n}a,\qquad a\in A_H^{n-r-1}(X),
\]
using $A_H(X,n-r-1)$. Since $2r+1-n\ge0$, the algebraic class
$L_Y^{2r+1-n}i^*a$ lies in $A_H^r(Y)$, and its Gysin image is $z$
by the projection formula. Thus, with the extension-stable universal
quantifiers of C.2e, $A_H$ and $A'_H$ are equivalent.

Since $A''_H$ is already a theorem by C.2b, the universal implication
$A''_H\Rightarrow A'_H$ would prove this standard conjecture $A$;
conversely that standard conjecture proves the implication. Together
with the product-level equivalences in
\citet[Corollary 2.14]{cyKleiman1968}, its exact remaining gate is the
universal Lefschetz standard conjecture. This is a reduction, not a
refutation of the implication and not a proof of that conjecture.

\paragraph{C.2g (Proved): reduction of $A'_H$ to the critical dimension.}
It suffices in the universal $A'_H$ hypothesis to assume the cases
$n=2r+1$. The following proof also works with algebraic equivalence
in place of homological equivalence. We retain rational cycle
coefficients and the smooth very ample hyperplane setting of the
letter.

\begin{proof}
Let $Y\subset X$ have defining section $s_0$ of the very ample line
bundle $L$, let $n=\dim X$, and let $(n-1)/2\le r\le n-1$.
Set $d=2r+1-n$. If $d=0$, this is the assumed case. Otherwise put
\[
 T=X\times\mathbb P^d,\qquad
 H=\left\{\sum_{j=0}^d s_j(x)t_j=0\right\}\subset T,
\]
where $s_1,\ldots,s_d$ are general sections of $L$.
The line bundle defining $H$ is $L\boxtimes\mathcal O(1)$, which
is very ample. We choose the sections so that their common zero
locus with $s_0$ is smooth of codimension $d+1$ (or empty), by
successive Bertini on the already smooth $Y$. Then $H$ is smooth:
away from this base locus a derivative in a projective-coordinate
direction is nonzero; at a base point the independent differentials
$ds_0,\ldots,ds_d$ have no zero linear combination with coefficients
$[t_0:\cdots:t_d]$. Now $\dim T=2r+1$, so the critical case applies.

Given $z\in A_H^{r+1}(X)$, it supplies an actual rational cycle
$\gamma\in\mathrm{CH}^r(H)_{\mathbb Q}$ whose pushforward has
cohomology class $\operatorname{pr}_X^*z$. The Chow projective-bundle
formula gives unique classes
\[
 (i_H)_*\gamma=\sum_{j=0}^d a_j h^j,
 \qquad a_j\in\mathrm{CH}^{r+1-j}(X)_{\mathbb Q},
 \quad h=c_1(\mathcal O_{\mathbb P^d}(1)).
\]
Its cohomological counterpart implies $\operatorname{cl}(a_0)=z$.
We now extract a stronger fact about $a_0$ than cohomological
vanishing. On $U=X\setminus Y$, the section $s_0$ trivializes $L$,
and the triangular coordinate change
\[
 t_0\longmapsto t_0+\sum_{j=1}^d(s_j/s_0)t_j
\]
identifies $H_U$ with the coordinate hyperplane
$\mathbb P^{d-1}_U\subset\mathbb P^d_U$.
This regular projective-bundle automorphism preserves
$\mathcal O(1)$ and every base class. Its Gysin image on Chow groups
is therefore precisely the ideal generated by $h$.
Consequently $a_0|_U=0$ in the \emph{Chow group}, not only in
cohomology. Chow localization provides
$\beta\in\mathrm{CH}^r(Y)_{\mathbb Q}$ with $a_0=(i_Y)_*\beta$.
Taking classes proves $A'_H(X,Y)$.

With algebraic equivalence, the same projective-bundle decomposition
is respected by families of cycles, so its zeroth coefficient is
algebraically equivalent to $z$. The Chow-localization step is
unchanged. This proves the claimed critical-case reduction there
as well, without asserting that the homological pencil proof
C.2e automatically proves its algebraic-equivalence analogue.
\end{proof}

Thus the precise phrase in 14:45, ``il suffit sauf erreur de faire
$n=2i+1$'', has a proof in the stated hyperplane setting. Combined
with C.2e--f, the critical homological $A'_H$ cases already have the
full strength of the extension-stable universal standard conjecture
$A_H$. The product stabilization does not assume that an arbitrary
cohomologically vanishing cycle is rationally equivalent to zero.

\paragraph{C.2h (Proved): middle positivity and the positive trace form.}
For homological cycle classes, the universal middle-dimensional
positivity assertion $B_H$ in 14:45 implies the Hodge standard
positivity assertion in every algebraic degree. We give the reduction,
then specify exactly the additional Lefschetz hypothesis needed for
its consequences on correspondences.

Let $x\in A_H^r(X)$ be primitive, with $\dim X=n\ge2r$, and let
$i:Y\hookrightarrow X$ be a smooth complete intersection of dimension
$2r$ formed from the chosen hyperplanes. Restriction in degree $2r$
is injective by weak Lefschetz. Moreover $i^*x$ is primitive: the
Gysin map out of $H^{2r+2}(Y)$ is injective, by duality with weak
Lefschetz in degree $2r-2$, and
\[
 i_*(L_Y i^*x)=L^{n-2r+1}x=0.
\]
Thus middle positivity on $Y$ gives, for $x\ne0$,
\[
          (-1)^r\deg_X(L^{n-2r}x^2)
          =(-1)^r\deg_Y((i^*x)^2)>0.
\]
This proves the reduction in
\citet[Proposition 3.7]{cyKleiman1968}; its converse is its middle case.
No algebraicity of a Lefschetz inverse was used in this paragraph.

Now assume the Lefschetz standard conjecture $B_{\rm std}$ for the
varieties involved and Hodge standard positivity on their products.
The Lefschetz primitive projectors, inverse operators, and star
operators are algebraic. If $u$ is an algebraic cohomological
correspondence, its adjoint
\[
                  u^\dagger=*_{X}\,{}^t u\,*_{Y}
\]
(with the standard degree-wise Lefschetz signs) is therefore algebraic.
The ordinary trace $\operatorname{Tr}(u^\dagger u)$ is rational and
strictly positive for $u\ne0$. Here is the checked computation behind
this input, as in \citet[Theorem 3.11]{cyKleiman1968}.
Decompose $u$ by the algebraic orthogonal projectors onto Lefschetz
strings. The trace norm is the sum of the trace norms of these blocks.
Transporting each block to its primitive factors reduces to a
primitive class in $P^a(X)\otimes P^b(Y)$, where $a+b$ is even.
The product Lefschetz operator is $L_X\otimes1+1\otimes L_Y$;
expanding its power contributes the positive binomial factor
\[
                   \binom{n-a+m-b}{n-a},
                   \qquad n=\dim X,\quad m=\dim Y.
\]
The trace formula contributes $(-1)^a$, and exchanging the two
primitive factors contributes $(-1)^{ab}$. Their product is $+1$
because $a+b$ is even. Up to that positive binomial factor, the block
norm is therefore its primitive Hodge intersection norm on $X\times Y$,
which is positive by the stated hypothesis. Rationality follows from
the correspondence trace formula, an intersection number of rational
cycles. This verifies the hypotheses of the finite algebra arguments
below; positivity on just $X$, without the relevant products, is not
substituted for them.

The resulting $\mathbb Q$-algebra $E$ of algebraic endomorphisms is
finite dimensional, even before one invokes semisimplicity. To see
this directly, choose a maximal linearly independent list
$u_1,\ldots,u_s$ over the cohomology coefficient field in the image
algebra. The map
\[
 E\longrightarrow\mathbb Q^s,
 \qquad u\longmapsto
       (\operatorname{Tr}(u^\dagger u_1),\ldots,
                         \operatorname{Tr}(u^\dagger u_s))
\]
is injective. If all coordinates vanish, express $u$ in the coefficient-
field span of the $u_j$ to obtain
$\operatorname{Tr}(u^\dagger u)=0$, whence $u=0$.
The $u_j$ are also rationally independent, so the injection proves
$\dim_{\mathbb Q}E=s$. They are consequently a rational basis as well
as a basis after extension to the cohomology coefficient field.
In particular that scalar extension of the correspondence algebra
remains faithful. This proves finite dimensionality and a positive
definite rational trace form. Its extension to $E\otimes\mathbb R$ is positive definite:
a negative real direction would give a negative rational direction,
and a degenerate rational matrix would have a nonzero rational kernel.

\paragraph{C.2i (Proved under the preceding standard-conjecture hypotheses): semisimplicity and algebraic preimages.}
Under C.2h's hypotheses, $E$ is semisimple. Indeed its Jacobson radical
is preserved by the anti-involution $\dagger$. If $u$ belonged to that
radical, then $u^\dagger u$ would be nilpotent and would have zero trace
in cohomology. Positivity implies $u=0$. Thus the radical is zero.
This proves the finite-algebra assertion of
\citet[Corollary 3.12 and Lemma 3.13]{cyKleiman1968} in the present
trace realization. With the hypotheses on all products, the same
argument, or the numerical identification and
\citet[Theorem 1]{Jannsen1992}, gives the semisimplicity of the
idempotent-complete pure homological motive category.

We also prove the full image criterion in 14:46, without assuming
that a cohomological image automatically has an algebraic preimage.
Let $u:H^a(X)\to H^b(Y)$ be induced by a rational algebraic
correspondence, with the Tate twist matching its codimension.
Work in the algebra of blocks on the direct sum of these two
realizations. Put $x=u^\dagger u$. Left multiplication by $x$ on
this finite-dimensional real algebra is self-adjoint for the trace
inner product, so $x$ is semisimple. Its minimal polynomial has
rational coefficients and no repeated root. By the Chinese remainder
theorem there is $p(T)\in\mathbb Q[T]$ which is zero at $T=0$ and
agrees with $1/T$ at every nonzero root of that minimal polynomial.
Let
\[
                    e_0=1-p(x)x,\qquad v=p(x)u^\dagger.
\]
The element $e_0$ is a self-adjoint idempotent with $xe_0=0$.
Therefore
\[
 (ue_0)^\dagger(ue_0)=e_0xe_0=0,
\]
and trace positivity forces $ue_0=0$. It follows that
\[
                         uvu=u.
\]
Every operator used is algebraic. If an algebraic class $z$ lies
in the cohomological image of $u$, then $u(vz)=z$, and $vz$ is an
algebraic class. The converse is immediate. This proves exactly the
surjectivity-on-algebraic-images consequence, including maps that are
neither injective nor surjective. It also constructs the projectors
onto the image and kernel; the faithful realization verifies their
idempotent identities. The argument is the explicit pseudo-inverse
form of \citet[Corollary 3.14]{cyKleiman1968}.

These assertions require both algebraic Lefschetz operators and the
product positivity used above. Since universal middle $B_H$ gives
$H_{\rm std}$, the remaining implication ``$B_H$ alone implies $D$''
has, under weak and hard Lefschetz, the exact further requirement
$A_H$: C.2c gives $A_H+H_{\rm std}\Rightarrow D$ and
$D\Rightarrow A_H$. Over $\mathbb C$, where Hodge--Riemann positivity
already proves $B_H$, that requested implication has the full force
of the universal standard conjecture $A_H$. We have not proved it by
silently strengthening the meaning of the letter's $B$.

\paragraph{C.2j (Proved under explicit hypotheses): the conditional Weil consequences.}
Let $X/\mathbb F_q$ be smooth projective and geometrically connected
of dimension $n$. Use geometric Frobenius on $\ell$-adic cohomology,
$\ell\nmid q$, equivalently the pullback operator of the $q$-power
Frobenius endomorphism. Suppose $B_{\rm std}(X)$ and Hodge standard
positivity on $X\times X$ for the product polarization. Then its
operators $F_i$ on $H^i(X_{\overline{\mathbb F}_q},\mathbb Q_\ell)$
are semisimple, their characteristic polynomials have integer
coefficients, and every complex conjugate of every eigenvalue on
$H^i$ has absolute value $q^{i/2}$. We give the argument in order to
separate these conditional consequences from the already proved
Weil conjectures themselves.

First, $B_{\rm std}$ makes the cohomological projectors $\pi_i$
algebraic. Choose an integer $N>0$ clearing the denominators of a
cycle representing $\pi_i$. For every $m\ge1$, the trace formula gives
\[
             \operatorname{Tr}(F_i^m)\in N^{-1}\mathbb Z.
\]
Indeed the supertrace of the composition of $\pi_i$ with the integral
Frobenius correspondence of power $m$ is an intersection number;
only degree $i$ survives, so the ordinary trace differs by $(-1)^i$.
The same fixed $N$ works for every $m$. Newton's identities imply
$P_i(t)=\det(1-tF_i)\in\mathbb Q[t]$.

We include the bounded-denominator argument that removes every
prime denominator, including the characteristic prime. If a finite
multiset of algebraic numbers $\lambda_j$ has power sums in
$N^{-1}\mathbb Z$ for all positive powers, each $\lambda_j$ is an
algebraic integer. Otherwise, at some nonarchimedean place one of
them has absolute value greater than one. The power series
\[
 \sum_{m\ge1}\left(\sum_j\lambda_j^m\right)t^m
       =-\frac{tP'(t)}{P(t)},
 \qquad P(t)=\prod_j(1-\lambda_jt),
\]
has bounded coefficients at that place and is analytic on $|t|<1$.
The rational function on the right has a pole at the reciprocal
of such a $\lambda_j$, which lies in that disk. There is no cancellation:
a nonzero root of $P$ of multiplicity $e$ contributes a logarithmic
pole with nonzero residue proportional to $e$ in characteristic zero.
This contradicts the power-series identity. Thus all eigenvalues are
integral, and the rational elementary symmetric functions are integers.
This is the mechanism in \citet[Proposition 2.6]{cyKleiman1968}; it
requires no unproved removal of $p$-denominators.

Let $L$ be an ample class defined over $\mathbb F_q$, and extend the
scalar algebra from $\mathbb Q$ to $\mathbb Q(\sqrt q)\subset\mathbb R$.
The normalized degree-wise operator
\[
                   g=\sum_i q^{-i/2}F_i\pi_i
\]
is multiplicative, fixes $L$ because $F^*L=qL$, and acts as the
identity on top cohomology because the Frobenius morphism has degree
$q^n$. Poincar\'e duality therefore gives its transpose as $g^{-1}$.
Since $g$ commutes with $L$, with the primitive projectors, and with
Lefschetz star, its adjoint in C.2h is also $g^{-1}$:
\[
                             g^\dagger g=1.
\]
Left multiplication by $g$ on the finite real correspondence algebra
is an isometry for the positive trace inner product, since
\[
 \operatorname{Tr}((ga)^\dagger gb)
                   =\operatorname{Tr}(a^\dagger b).
\]
It is consequently diagonalizable over $\mathbb C$, with eigenvalues
on the unit circle. The left regular representation and the faithful
cohomological representation have the same minimal polynomial for an
algebra element: a polynomial in that element vanishes in one exactly
when it vanishes in the algebra. Hence the same semisimplicity and
absolute-value statement holds for $g$ in cohomology. On the invariant
left ideal $\pi_iE_{\mathbb R}$, left multiplication by $g$ equals
$q^{-i/2}$ times left multiplication by the rational algebra element
$F_i\pi_i$. The latter has a rational minimal polynomial; every complex
root occurs in its faithful left regular representation on that ideal.
The isometry therefore controls every complex conjugate of $F_i$, not
just one chosen embedding. Multiplying back by the positive real
number $q^{i/2}$ proves the claimed weights.
This supplies the linear algebra of
\citet[Lemma 4.6 and Theorem 4.7]{cyKleiman1968}, with the product
positivity and algebraic-projector premises visible.

Finally the fixed-point formula and the formal identity
\[
-\log\det(1-tF_i)=\sum_{m\ge1}\operatorname{Tr}(F_i^m)t^m/m
\]
give
\[
 Z(X,t)=\exp\left(\sum_{m\ge1}\#X(\mathbb F_{q^m})\frac{t^m}{m}\right)
       =\prod_{i=0}^{2n}P_i(t)^{(-1)^{i+1}}.
\]
Thus the numerator and denominator factors lie in $\mathbb Z[t]$;
$P_0(t)=1-t$, $P_{2n}(t)=1-q^nt$, and the zeros of $P_i$ have
absolute value $q^{-i/2}$. Poincar\'e duality pairs eigenvalues in
complementary degrees by $\alpha\leftrightarrow q^n/\alpha$, yielding
\[
 Z\left(X,\frac1{q^nt}\right)
        =\epsilon\,q^{n\chi/2}t^\chi Z(X,t),
 \qquad \chi=\sum_i(-1)^i\dim H^i(X),\quad \epsilon\in\{1,-1\}.
\]
More precisely, with
$\Delta=\prod_i\det(F_i)^{(-1)^{i+1}}$, the sign is
$\epsilon=(-1)^\chi q^{n\chi/2}\Delta$;
duality gives $\Delta^2=q^{-n\chi}$.
We retain this determinant sign instead of assuming that it is
fixed independently of Frobenius. Also $n\chi$ is even: when $n$
is odd, middle Poincar\'e duality is alternating and its dimension
is even, while the other terms occur in complementary pairs.
Rationality and the functional equation follow formally from the
trace formula and duality; the extra standard-conjecture input above
supplies the positive-involution proof of the absolute values and
semisimplicity. It is not a proof of the general standard conjectures.

% Classical low-degree standard-D cases with a complete derivation.
\paragraph{C.2k (Known with derivation: the exact low-degree standard-$D$ range used here).}
Let $X/\mathbb C$ be smooth projective of dimension $n$ and let
$A^r_H(X)\subset H^{2r}(X,\mathbb Q)$ be the space of algebraic
cycle classes. Numerical and homological equivalence, with rational
coefficients, agree in every codimension
\[
                        r\in\{0,1,2,n-1,n\}.
\]
Repeated or out-of-range entries are omitted. In particular the
full standard conjecture $D(X)$ holds for every complex variety
of dimension at most four. This is not an assertion of the
Lefschetz standard conjecture for all correspondences on those
varieties.

\begin{proof}
We derive these classical cases directly from the opened
Lefschetz $(1,1)$ and Hodge--Riemann inputs
\citep[Remark 3.4, p.~379, and Remark 3.10, p.~380]{cyKleiman1968}.
The degree map treats $r=0,n$. For divisors, the Hodge index
form
\[
       (d,e)\longmapsto\int_X deh^{n-2}
\]
is nondegenerate on the rational Neron--Severi space; for curves
of dimension one we instead use degree directly. Thus a
numerically trivial divisor has zero cohomology class.

For $r=n-1$ and $n\ge2$, cohomological hard Lefschetz is an
isomorphism of rational Hodge structures
\[
 h^{n-2}:H^2(X,\mathbb Q)\xrightarrow{\sim}
                         H^{2n-2}(X,\mathbb Q).
\]
The inverse of an algebraic class is a rational $(1,1)$ class,
and therefore is a divisor class. We have
$A_H^{n-1}(X)=h^{n-2}A_H^1(X)$. Pairing with divisors and
the preceding nondegenerate form proves the assertion.

It remains to treat $r=2$ with $n\ge4$.
Take $a\in A_H^2(X)$ which is numerically trivial.
There is a unique rational class $d\in H^2(X,\mathbb Q)$ such
that
\[
                      h^{n-2}d=h^{n-3}a.
\]
Hard Lefschetz and Hodge type make $d$ a rational $(1,1)$
class, hence an algebraic divisor class. For every divisor
class $e$ we have
\[
             \int_X deh^{n-2}=
             \int_X aeh^{n-3}=0.
\]
The divisor Hodge index form gives $d=0$. Thus
$h^{n-3}a=0$: the class $a$ is primitive in degree four.
Numerical triviality also gives
\[
                         \int_X a^2h^{n-4}=0.
\]
The primitive Hodge--Riemann form is positive definite on
rational $(2,2)$ classes, so $a=0$. These arguments prove
the stated range. For $n\le4$ that range includes every
codimension.
\end{proof}

There is a precise remaining quotient in degree two of the
homological Lefschetz problem. For every $n\ge4$, the pairing
\[
       Q_2(a,b)=\int_X abh^{n-4}
                  \quad(a,b\in A_H^2(X))
\]
is nondegenerate. To verify this without assuming a standard
conjecture, use the same rational algebraic projection
$a\mapsto d$ above. It gives the orthogonal decomposition
\[
 A_H^2(X)=
 \bigl(A_H^2(X)\cap\ker h^{n-3}\bigr)\oplus hA_H^1(X).
\]
The first summand has the positive Hodge--Riemann form and
the second has the nondegenerate divisor form. Define
\[
 R_H^{n-2}(X)=
 \left\{z\in A_H^{n-2}(X):
       \int_X za=0\text{ for every }a\in A_H^2(X)\right\}.
\]
Then there is an exact direct-sum decomposition of rational
cycle-class spaces
\begin{equation}\label{eq:C2k-first-obstruction}
 A_H^{n-2}(X)=h^{n-4}A_H^2(X)\oplus R_H^{n-2}(X).
\end{equation}
Indeed $Q_2$ represents the functional $a\mapsto\int_X za$
by a unique $b\in A_H^2(X)$. Subtracting $h^{n-4}b$ gives the
second summand; the same form shows that the intersection
of the two summands is zero. This is a finite-dimensional
rational linear-algebra argument on algebraic cycle classes,
without assuming that every rational Hodge class is algebraic.

Consequently $A_H(X,2)$ is equivalent to
$R_H^{n-2}(X)=0$, which is exactly standard $D$ in
codimension $n-2$ on this $X$. On a fivefold this reads
\[
 A_H^3(X)/hA_H^2(X)\simeq R_H^3(X).
\]
Thus the first dimension and codimension not covered by
the preceding classical theorem are $(n,r)=(5,3)$.
No vanishing of this last space is proved here.

Finally, C.1n's fixed-target endpoint image criterion still
quantifies over \emph{all source varieties and all source
cycles}. Its necessity argument can insert a numerically
trivial cycle of any source dimension and codimension.
The theorem just proved therefore does not turn that
universal image criterion, even for a point or projective
space as target, into an unconditional assertion.


\paragraph{Remaining work for C.2.}
We have proved affine vanishing, isolated its exact failure to prove
cycle localization, and given the correct conditional positivity
argument. C.2e--f prove the universal equivalence
$A'_H\Longleftrightarrow A_H$ under hypotheses stable under the
algebraically closed field extensions used by the pencil construction.
We have not proved or refuted either conjecture in full.
A proof of $A''_H$ alone does not prove either one.
The universal implication from the letter's
positivity assertion alone has not been decided here; the fixed-variety
counterexample C.1e has a deliberately narrower conclusion.
C.2h--j also supply the conditional Weil estimates, the positive
involution argument, and the correspondence image criterion.
C.2k proves the classical complex standard-$D$ cases in
codimensions $0,1,2,n-1,n$, hence every degree in dimensions
at most four. Its first remaining obstruction is the explicitly
defined numerical radical in codimension three on a fivefold;
these low-degree cases do not settle C.1n's all-source criterion.

\subsection{C.3: the maps on algebraic classes and their algebraically trivial part}

\paragraph{C.3a (Proved exact obstruction sequence).}
For $2r\le n$, put $q=n-2r$. The cycle class maps give a commutative
diagram with exact rows
\[
\begin{array}{ccccccccc}
0&\to&G^r(X)&\to&A^r_{\mathrm{alg}}(X)&\to&A^r_{\mathrm{hom}}(X)&\to&0\\
 &&\downarrow L^q&&\downarrow L^q&&\downarrow L^q\\
0&\to&G^{n-r}(X)&\to&A^{n-r}_{\mathrm{alg}}(X)&\to&A^{n-r}_{\mathrm{hom}}(X)&\to&0.
\end{array}
\]
The last vertical map is injective by hard Lefschetz. The snake lemma
therefore gives
\[
\begin{split}
\ker(L^q|A^r_{\mathrm{alg}})&=\ker(L^q|G^r),\\
0\longrightarrow\operatorname{coker}(L^q|G^r)&\longrightarrow
\operatorname{coker}(L^q|A^r_{\mathrm{alg}})
\longrightarrow\operatorname{coker}(L^q|A^r_{\mathrm{hom}})
\longrightarrow0.
\end{split}
\]
Thus the general question contains two separate gates: the Lefschetz
map on Griffiths groups, and standard-conjecture surjectivity on
algebraic cohomology. Neither is disposed of by hard Lefschetz on the
ambient cohomology. C.1c proves injectivity for $r=1$ over $\mathbb C$;
C.1a identifies the requested quintic cokernel exactly. The higher
injectivity problem has not been settled here.

\paragraph{C.3b (Proved: $J^1$ and its exact cokernel).}
For a smooth projective complex $n$-fold and an ample $L$,
\[
 L^{n-1}:\operatorname{Pic}^0(X)(\mathbb C)\otimes\mathbb Q
 \longrightarrow\mathrm{CH}_0(X)^0\otimes\mathbb Q
\]
is injective. Its cokernel is naturally isomorphic to
$\ker(\operatorname{alb}_X)\otimes\mathbb Q$. Indeed, its composition
with the Albanese map is the isogeny
$\operatorname{Pic}^0(X)\to\operatorname{Alb}(X)$ attached to a complete
intersection curve, proved in C.4a below. On complex points this
isogeny becomes an isomorphism after tensoring with $\mathbb Q$:
it is surjective and has finite kernel. Every zero-cycle has a unique
Lefschetz image with the same rational Albanese class. Subtracting that
image gives the asserted identification of the quotient with the kernel.
In particular, the $J^1$ surjectivity assertion is equivalent to
vanishing of the rational Albanese kernel.

\paragraph{C.3c (Refuted: the unrestricted $J$ and punctured-cone claims).}
Let $S/\mathbb C$ be the Fermat quartic surface
$x_0^4+x_1^4+x_2^4+x_3^4=0$. Its partial derivatives have no common
projective zero, so it is smooth; adjunction gives
$K_S\simeq\mathcal O_S$, and the hypersurface cohomology sequence gives
$H^1(S,\mathcal O_S)=0$. Thus $S$ is a K3 surface with
$\operatorname{Pic}^0(S)=\operatorname{Alb}(S)=0$ and
$H^{2,0}(S)\ne0$. Mumford's theorem gives an infinite-dimensional
degree-zero Chow group \citep[Section 2, main theorem, p.~200]{cyMumford1969}.
This refutes the surjectivity of $J^1(S)\to J^2(S)$.

Let $E$ be the complement of the zero section in
$\mathcal O_S(-1)$. It is the punctured affine cone over this projectively
normal quartic, and has dimension $3$, not $5$. Vector-bundle invariance
and localization for the zero section give
\begin{equation}\label{eq:punctured-line-bundle}
 \mathrm{CH}^r(E)=
 \mathrm{CH}^r(S)/\bigl(h\,\mathrm{CH}^{r-1}(S)\bigr).
\end{equation}
For $r=2$, the denominator is the image of the finitely generated group
$\mathrm{Pic}(S)=\mathrm{NS}(S)$, while the rational degree-zero numerator
is not finite dimensional. Consequently $\mathrm{CH}^2(E)\otimes\mathbb Q$
is not finite dimensional. This is a counterexample to the global
punctured-cone finiteness reading of $\text{\textcircled{3}}'$.

We have checked the codimension shift in
\eqref{eq:punctured-line-bundle}; writing $\mathrm{CH}^r(Y)$ in place of
$\mathrm{CH}^{r-1}(Y)$ in divisor localization is incorrect. A result for
this punctured cone does not automatically pass to the punctured
spectrum of the strict henselization at its vertex. C.3e below
proves the needed separate injection, and C.3f--g give strict local
counterexamples with all the required hypotheses.
The primitive exponents compatible with the two proposed Lefschetz
strings are $n-2i+1$ on $A^i_{\mathrm{alg}}$ and $n-2i+2$ on $J^i$.

\paragraph{C.3d (Proved: the complete universal surjectivity classification for $J$).}
In the nonempty cycle range $1\le i\le(n+1)/2$, the map
\[
 L^{n-2i+1}:J^i(X)\longrightarrow J^{n+1-i}(X)
\]
is surjective for every smooth projective complex $n$-fold exactly
when $2i=n+1$, when it is the identity. For every $n\ge2i\ge2$
we construct $X$ with $J^i(X)=0$ and $J^{n+1-i}(X)\ne0$.
Thus every nonidentity pair fails universal surjectivity, independently
of the choice of ample class.

We apply the incidence construction of C.1f with the Fermat quartic
$W_2=S\subset T_3=\mathbb P^3$ as its initial divisor. Set
\[
\begin{split}
 W_{2i}&=\operatorname{Bl}_{W_{2i-2}\times\{0\}}
                   (T_{2i-1}\times\mathbb P^1),\\
 T_{2i+1}&=\mathbb P_{\mathrm{lines}}
 \bigl(\operatorname{pr}_1^*\mathcal O_{T_{2i-1}}(W_{2i-2})
       \oplus\operatorname{pr}_2^*\mathcal O_{\mathbb P^1}(1)\bigr)
       \qquad(i\ge2).
\end{split}
\]
Exactly the same local equation proves that $W_{2i}$ is a smooth
Cartier divisor in $T_{2i+1}$, and $T_{2i+1}$ has Tate motive.
The initial surface has $J^1(S)=0$, while
$J^2(S)=\mathrm{CH}_0(S)^0\otimes\mathbb Q\ne0$ by C.3c.
All other $J^r(S)$ vanish. The explicit blow-up isomorphism in C.1f
restricts to the algebraically trivial subgroups, so induction gives
\[
 J^r(W_{2i})=0\quad(r\ne i+1),\qquad
 J^{i+1}(W_{2i})\simeq J^2(S)\ne0.
\]
Now take $X=W_{2i}\times\mathbb P^{n-2i}$. The projective-bundle
formula places the copies of $J^2(S)$ in precisely the degrees
$i+1,\ldots,n+1-i$. Its source $J^i$ is zero and its target is
nonzero. This proves every claimed failure. It does not refute
injectivity, whose source is zero in these examples.

% Root proof; independent cone-lemma audit by /root/frobenius.
\paragraph{C.3e (Proved): an integral injection into the strict local puncture.}
Let $V$ be a smooth projective integral variety of dimension $n$ over
an algebraically closed field $k$. Let $L$ be a very ample line bundle
whose section ring is generated in degree one and defines a
projectively normal embedding. Write $C$ for its affine cone, $v$ for
the vertex, and
\[
 R=\mathcal O_{C,v},\qquad
 X^{\mathrm{sh}}=\operatorname{Spec}(R^{\mathrm{sh}})
                    \setminus\{\mathfrak m\}.
\]
We use codimension Chow groups, with codimension inherited from the
local cone; removing the closed point does not lower the heights of
the remaining primes. There are natural injections
\begin{equation}\label{eq:C3-strict-cone-injection}
\begin{aligned}
 \frac{\mathrm{CH}^r(V)}{c_1(L)\mathrm{CH}^{r-1}(V)}
       &\lhook\joinrel\longrightarrow\mathrm{CH}^r(X^{\mathrm{sh}}),\\
 \frac{B^r(V)}{c_1(L)B^{r-1}(V)}
       &\lhook\joinrel\longrightarrow B^r(X^{\mathrm{sh}}).
\end{aligned}
\end{equation}
Here $B^r$ denotes cycles modulo algebraic equivalence with integral
coefficients. These statements concern strict henselization, not
completion, and assert injections rather than surjections.

\begin{proof}
The blow-up of $C$ at $v$ is the total space
$W=\operatorname{Tot}(L^{-1})$, with exceptional zero section
$i:V\hookrightarrow W$ and projection $q:W\to V$.
Indeed, on the blow-up chart belonging to a degree-one coordinate
$x_j$, the degree-zero ratios generate the corresponding affine
chart of $V$, and $x_j$ is the fibre coordinate. The chart ring is
the polynomial ring in that fibre coordinate over the chart ring
of $V$; their transition functions are those of $L^{-1}$.
The degree-one generation hypothesis is used here.
Blow-up commutes with flat base change, and its exceptional divisor
is Cartier \citep[Lemmas 31.33.3--31.33.4]{rootStacksBlowup}.

Take a pointed affine \emph{\'etale} neighbourhood $T\to C$ of $v$.
After removing the other finitely many points over $v$, its fibre over
$v$ is a single reduced $k$-point $t$. Put $W_T=W\times_C T$.
The morphism $W_T\to W$ is \emph{\'etale}, so
$q_T:W_T\to V$ is smooth of relative dimension one. Its exceptional
divisor $i_T:V\hookrightarrow W_T$ is a section of $q_T$, its normal
bundle is $L^{-1}$, and
\[
 W_T\setminus i_T(V)\simeq T\setminus\{t\},\qquad
 \mathcal O_{W_T}(i_T(V))\simeq q_T^*L^{-1}.
\]
Suppose $q_T^*\alpha$ restricts to zero on this complement.
Chow localization gives $\beta\in\mathrm{CH}^{r-1}(V)$ such that
$q_T^*\alpha=(i_T)_*\beta$. Apply the refined pullback $i_T^!$.
The smooth-section identity $i_T^!q_T^*=1$ and the
self-intersection formula give
\begin{equation}\label{eq:C3-local-self-intersection}
          \alpha=c_1(L^{-1})\beta=-c_1(L)\beta.
\end{equation}
For the cited formulas see \citet[Lemma 42.59.2]{rootStacksGysin}
and \citet[Proposition 1.8 and Corollary 6.3]{frC4Fulton1998}.
Conversely, the projection formula and the displayed line-bundle
identity give
\[
 (i_T)_*\beta
  =c_1(\mathcal O_{W_T}(i_T(V)))q_T^*\beta
  =-q_T^*(c_1(L)\beta),
\]
whose restriction to the complement is zero. The kernel at
\emph{every} such neighbourhood is therefore exactly
$c_1(L)\mathrm{CH}^{r-1}(V)$.

The same proof works modulo algebraic equivalence. Proper pushforward,
flat pullback, refined Gysin maps, and Chern-class operations preserve
that equivalence; localization is still exact
\citep[Proposition 10.3 and Example 10.3.4]{frC4Fulton1998}.
One can check the last assertion directly: close a curve-parameter
correspondence from the open complement in $W_T$ times the parameter
curve, take the difference of its refined fibres, and then use
ordinary localization for the remaining boundary-supported class.
Thus the same application of $i_T^!$ gives
\eqref{eq:C3-local-self-intersection} in $B^r(V)$.

Finally, the strict henselization is the filtered limit of these
pointed neighbourhoods. Any cycle, finite list of rational functions
giving a rational equivalence, or curve and correspondence giving an
algebraic equivalence is specified by finitely many equations and
descends to one neighbourhood, after shrinking it. In particular,
a relation killing a pulled-back class on $X^{\mathrm{sh}}$ already
kills it at one stage. The kernel calculation there proves both
injections in \eqref{eq:C3-strict-cone-injection} integrally.
\end{proof}

\paragraph{C.3f (Refuted): the three-dimensional local finiteness,
Gysin, and open-restriction assertions.}
We give equations and a non-torsion witness. Fix the Fermat quartic
\[
 V:\ x_0^4+x_1^4+x_2^4+x_3^4=0\ \subset\mathbb P^3_{\mathbb C},
 \qquad H=V\cap\{x_3=0\},\quad L=\mathcal O_V(1).
\]
Choose $u,v\in\mathbb C$ algebraically independent over
$k_0=\overline{\mathbb Q}$, choose
$w$ with $w^4=-1-u^4-v^4$, and choose
$\zeta\in k_0$ with $\zeta^4=-1$. Set
\[
 P=(u:v:w:1),\qquad a=(1:\zeta:0:0),\qquad
                 \alpha=[P]-[a]\in\mathrm{CH}_0(V)^0.
\]
Then $\alpha$ has infinite order in
$\mathrm{CH}^2(V)/(L\mathrm{CH}^1(V))$.

\begin{proof}
The point $P$ is a geometric generic point of the $k_0$-surface $V$.
We prove the stronger assertion that no nonzero multiple of
$[P]-[a]$ is rationally equivalent to a zero-cycle supported on $H$.
Suppose such a relation exists over $\mathbb C$, and put
$K=k_0(V)$. Its cycles and rational-equivalence functions spread over
a finite-type integral $\overline K$-scheme. Shrink to a smooth
nonempty open on which all these data specialize; a closed point of
that open has residue field $\overline K$. Refined specialization
there gives the relation over $\overline K$, and its finite data
descend to a finite extension $K'/K$. Norm from $K'$ to $K$ then
gives a relation
\[
                   m[\eta]=z_H\quad\text{in }\mathrm{CH}_0(V_K)
\]
for an integer $m\ne0$ and a zero-cycle $z_H$ supported on $H_K$;
the term involving $a$ has been absorbed in $z_H$.
Spread this equality from the generic point of the first factor of
$V\times V$ and apply Chow localization. It yields
\begin{equation}\label{eq:C3-diagonal-support}
 m[\Delta_V]=\Gamma_H+\Gamma_D
                   \quad\text{in }\mathrm{CH}^2(V\times V),
\end{equation}
where $\Gamma_H$ is supported on $V\times H$ and $\Gamma_D$ on
$D\times V$ for a proper divisor $D\subset V$.
Enlarging a proper closed support to a divisor causes no problem.

Adjunction gives $K_V\simeq\mathcal O_V$, so
$H^{2,0}(V)\ne0$. A correspondence supported on $D\times V$
annihilates this space because a two-form restricts to zero on a
resolution of the curve $D$. A correspondence supported on
$V\times H$ also annihilates it: pushforward from a curve raises
both Hodge indices by one and cannot produce type $(2,0)$.
The diagonal acts as the identity, contradicting
\eqref{eq:C3-diagonal-support}. This is the generic-point form of
Mumford's differential-form obstruction; the support argument here
proves the particular assertion we need
\citep[Section 2, Lemmas 2--3 and main theorem]{cyMumford1969}.
Since every product $L\beta$ has a representative supported on $H$,
the claimed quotient class is non-torsion.
\end{proof}

Now take
\[
 R=\left(\mathbb C[x_0,x_1,x_2,x_3]/
                   (x_0^4+x_1^4+x_2^4+x_3^4)\right)_{(x_0,x_1,x_2,x_3)},
 \quad S=\operatorname{Spec}(R^{\mathrm{sh}}),\quad X=S\setminus\{s\}.
\]
This is an excellent strictly Henselian noetherian local scheme of
dimension $N=3$. Its puncture is regular: the four partial derivatives
$4x_j^3$ vanish simultaneously only at the vertex, and regularity
is preserved by the \emph{\'etale} limit. By C.3e, the image
$\gamma\in\mathrm{CH}^2(X)$ of $\alpha$ is non-torsion.
It is algebraically trivial, since degree-zero cycles on the smooth
connected projective surface are algebraically trivial and the maps
used here preserve that equivalence. Thus the rational-equivalence
reading of the boxed assertion on 14:108,
``$A^2(X)$ finite'', is false, and so is its asserted vanishing of
$J^2(X)$. We distinguish this from cycles \emph{modulo} algebraic
equivalence: $\gamma$ is zero in that quotient.

The same equations refute the local Gysin claim without changing
the required regularity hypotheses. Put
\[
 f=x_2-wx_3\in\mathfrak m_R,\qquad
 T=V\cap\{f=0\},\qquad Y=X\cap V(f),\quad U=X\setminus Y.
\]
The plane quartic $T$ has equation
\[
                    x_0^4+x_1^4+(1+w^4)x_3^4=0.
\]
Its three coefficients are nonzero, since
$1+w^4=-u^4-v^4\ne0$; hence $T$ is smooth, and $Y$ is regular.
Both $P$ and $a$ lie on $T$. Pull back
$\delta=[P]-[a]\in\operatorname{Pic}^0(T)$ along its punctured
line bundle and then along the strict local base change. This gives
$\beta\in\mathrm{CH}^1(Y)_{\mathrm{alg}}$, and flat compatibility
of the closed immersion's pushforward gives
\[
                        (Y\hookrightarrow X)_*\beta=\gamma.
\]
Equivalently, the two corresponding cone rays occur with
multiplicity one. This image is non-torsion. Consequently the Gysin
map on 14:107 is neither zero nor zero modulo finite groups, even
at $i=1\le(N-1)/2$. By localization, $\gamma$ maps to zero in
$\mathrm{CH}^2(U)$; hence the restriction
$\mathrm{CH}^2(X)\to\mathrm{CH}^2(U)$ is not an isomorphism modulo
finite groups. All assertions here retain integral Chow groups.
The symbolic certificate \texttt{code/check\_local\_cone.py}
checks the equations, the points, the smoothness coefficients and
the blow-up substitution; the non-torsion assertion is proved by
the preceding diagonal argument, not by that finite computation.

\paragraph{C.3g (Refuted in every nonempty proposed high range):
universal strict local finiteness.}
For every $N\ge3$ and every integer
\[
                  \left\lceil\frac{N+1}{2}\right\rceil
                              \le r\le N-1,
\]
there is an excellent strictly Henselian local complex cone of
dimension $N$, regular away from its closed point, for which
$\mathrm{CH}^r$ of the puncture contains a non-torsion element.
If $r<N-1$, the example can moreover be chosen nonzero after
tensoring cycles modulo algebraic equivalence with $\mathbb Q$.
Thus no nonempty universal high-degree finiteness range displayed
on 14:106 and 14:109 survives under these hypotheses.

\begin{proof}
Write $n=N-1$. If $r<n$, then $r<n<2r$.
Apply C.1i to the smooth $n$-dimensional complete intersection in
the split quadric of dimension $2r$, with all cutting degrees
$4r+2$ and the explicit geometric generic coefficients specified
there. It gives a nonzero class
\[
 z\in B^r(V)_{\mathbb Q}/
                         \bigl(LB^{r-1}(V)_{\mathbb Q}\bigr).
\]
This complete intersection in projective space is projectively
normal and has section ring generated in degree one. Its affine
cone has an isolated singularity, since its puncture is the
nonzero total space of $L^{-1}$ over smooth $V$.
The second injection of C.3e sends $z$ to a nonzero rational
class on the strict local puncture. Clearing denominators gives
the required integral non-torsion class there and in its Chow group.
The range used here satisfies $2r\ge n+2$, including the precise
strict local bound in the manuscript.

If $r=n$, take instead the smooth Fermat hypersurface of degree
$n+2$ in $\mathbb P^{n+1}_{\mathbb C}$, defined over
$\overline{\mathbb Q}$. It has $K_V\simeq\mathcal O_V$ and
$H^{n,0}(V)\ne0$. Choose a fixed hyperplane section $H$, a
$\overline{\mathbb Q}$-point $a\in H$, and a geometric generic
point $P$ of $V$ over $\overline{\mathbb Q}$.
The proof of C.3f, with $n$ in place of $2$, shows that no nonzero
multiple of $[P]-[a]$ is supported on $H$ modulo rational
equivalence: a decomposition of the diagonal supported on
$V\times H$ and $D\times V$ annihilates $H^{n,0}$ on the right
and acts nontrivially on the left. Therefore
\[
 [P]-[a]\ne0\quad\text{in}\quad
 \bigl(\mathrm{CH}^n(V)/L\mathrm{CH}^{n-1}(V)\bigr)
                                      \otimes\mathbb Q.
\]
Apply the first injection of C.3e. The equations for a smooth
homogeneous hypersurface again give all the stated local hypotheses.
\end{proof}

These results settle the strict local passage left open in C.3c.
They do not decide every low-degree restriction or higher injectivity
assertion in the local tables. In particular, a failure of finiteness
or of one Gysin map is not a proof about a different restriction map
whose source and degree lie outside the examples just given.

\paragraph{C.3h (Known mechanism and proved exact graded gates on abelian varieties).}
Let $A$ be a $g$-dimensional abelian variety, $g\ge1$, over an algebraically
closed field, of arbitrary characteristic. All groups in this paragraph
have rational coefficients. Write
\[
 B^r(A)=\mathrm{CH}^r(A)_{\mathbb Q}/
                  \mathrm{CH}^r(A)_{\mathrm{alg},\mathbb Q},
 \qquad J^r(A)=\mathrm{CH}^r(A)_{\mathrm{alg},\mathbb Q}.
\]
Both inherit the Beauville grading, with index $s$ characterized by
$[m]^*x=m^{2r-s}x$. For an ample divisor class $h$, the following
are exact equivalences for this fixed abelian variety:
\begin{equation}\label{eq:C3-abelian-B}
 \begin{gathered}
 h^{g-2r}:B^r(A)\longrightarrow B^{g-r}(A)
 \text{ injective for every }2r\le g\\
 \Longleftrightarrow\quad B^*(A)_{(s)}=0\quad(s<0).
 \end{gathered}
\end{equation}
\begin{equation}\label{eq:C3-abelian-J}
 \begin{gathered}
 h^{g-2r+1}:J^r(A)\longrightarrow J^{g-r+1}(A)
 \text{ injective for every }2r\le g+1\\
 \Longleftrightarrow\quad J^*(A)_{(s)}=0\quad(s\le0).
 \end{gathered}
\end{equation}
The displayed implications do not assert these vanishing conditions
in general. The first is the adequate-equivalence version of
\citet[Proposition 5.11; footnote 1, p.~254]{rootBeauville2010};
the second applies the same representation to its invariant
algebraically trivial subgroup.

\begin{proof}
We give the string argument, including the converse directions.
Choose the symmetric ample class $\theta$ associated to $h$.
\citet[Theorem 4.2 and Proposition 5.2, pp.~260--261]{rootBeauville2010}
construct algebraic correspondences acting by
\[
 \begin{gathered}
 E(x)=\theta x,\qquad
 F(x)=\frac{\theta^{g-1}}{d(g-1)!}*x,\qquad
 d=\frac{\deg\theta^g}{g!},\\
 H(x)=(2r-g-s)x\quad(x\in\mathrm{CH}^r(A)_{(s)}).
 \end{gathered}
\]
where $*$ is the Pontryagin product. They form an
$\mathfrak{sl}_2$-representation that is a direct sum of finite
dimensional representations. All operators preserve algebraic
equivalence, since they are correspondences. Therefore they induce
the same kind of representation on $B^*(A)$ and on $J^*(A)$.
Each is a direct sum of strings
\[
        z,\theta z,\ldots,\theta^{g+s-2a}z,
 \qquad z\in M^a_{(s)},\quad Fz=0,\quad g+s-2a\ge0,
\]
where $M$ denotes either $B$ or $J$. The semisimplicity used here
is ordinary characteristic-zero representation theory over
$\mathbb Q$, independent of the characteristic of the base field.
It concerns these locally finite representations, not a
semisimplicity assertion about the cycle category.

In degree $r$, a nonzero term of this string is
$\theta^{r-a}z$ with $a\le r$ and $r-a\le g+s-2a$.
Multiplication by $\theta^k$ kills it exactly when
\[
                         r-a+k>g+s-2a.
\]
For $k=g-2r$, this inequality becomes $a>r+s$.
It is impossible when $s\ge0$. Conversely, if a negative
graded piece is nonzero, take a nonzero lowest element
$z\in B^a_{(s)}$ in one of its strings. Then
$2a\le g+s<g$ and $g-2a>g+s-2a$, so the asserted
map in degree $a$ kills $z$. This proves
\eqref{eq:C3-abelian-B} for $\theta$.

For $k=g-2r+1$, the corresponding inequality is
$a>r+s-1$, impossible when $s\ge1$. A nonzero string
with $s\le0$ has $2a\le g+s\le g$ and
$g-2a+1>g+s-2a$; its lowest element is therefore a
nonzero kernel element in an allowed degree. This proves
\eqref{eq:C3-abelian-J} for $\theta$.

Finally $h=T_b^*\theta$ for some $b\in A(k)$, as a rational
divisor class: their difference belongs to $\operatorname{Pic}^0(A)$,
and the polarization isogeny associated to $\theta$ is surjective
on points over the algebraically closed field. Clearing the
possible denominator in the symmetric class and using the
surjectivity of multiplication on $\operatorname{Pic}^0(A)(k)$
gives the rational-class assertion as well. The pullback
$T_b^*$ is an automorphism of both $B$ and $J$, and
\[
             h^k T_b^*x=T_b^*(\theta^k x).
\]
It conjugates the relevant maps, proving both claims for $h$.
On $B$ one may also use directly that translations act trivially
modulo algebraic equivalence.
\end{proof}

These formulas specify small remaining tests. For $g=5$, the only
possibly negative piece in degree two is $B^2(A)_{(-1)}$:
the Fourier correspondence identifies the $s$-piece there with
$B^{3+s}(A)_{(s)}$; for $s\le-2$ this is a divisor,
a codimension-zero group, or zero, and its negative part vanishes.
The degree-two $s=-1$ piece is entirely primitive, since its
lowering image is in the zero group $B^1(A)_{(-1)}$.
Hence the first nonidentity higher map in this dimension satisfies
\[
 \ker\bigl(h:B^2(A)\to B^3(A)\bigr)
                 \simeq B^2(A)_{(-1)}.
\]
For $g=4$ and $J^2$, the same Fourier argument eliminates
all negative pieces. Its degree-zero piece is primitive because
$J^1(A)_{(0)}=0$. Consequently
\[
 \ker\bigl(h:J^2(A)\to J^3(A)\bigr)
                 \simeq J^2(A)_{(0)}.
\]
For nonsymmetric $h$ the indicated identifications use a translation
as above and need not be canonical. Neither group on the right
is shown to vanish here. We have reduced these particular
injectivity problems to the stated graded groups, rather than
confusing them with the already refuted surjectivity assertions.

\paragraph{C.3i (Refuted in every proposed degree-one range): strict local Picard restriction without a depth hypothesis.}
For every $N\ge4$ there is an excellent strictly Henselian
normal local complex ring $R$ of dimension $N$, and a
nonzero $f\in\mathfrak m_R$, such that
\[
 X=\operatorname{Spec}R\setminus\{\mathfrak m_R\},\qquad
 Y=X\cap V(f)
\]
are regular, while
\begin{equation}\label{eq:C3-picard-kernel}
 \ker\bigl(\operatorname{Pic}(X)\longrightarrow
                  \operatorname{Pic}(Y)\bigr)
             \simeq (\mathbb C,+)^3.
\end{equation}
The map is surjective in this example. Thus the assertion
on 14:107 that $A^1(X)\to A^1(Y)$ is injective for $N\ge4$,
and is an isomorphism for $N\ge5$, is false on the
rational-equivalence reading $A^1=\mathrm{CH}^1=\operatorname{Pic}$.
The failure persists after tensoring with $\mathbb Q$ and
after discarding finite groups.

\begin{proof}
Put $m=N-2\ge2$ and take the fixed smooth plane quartic
\[
 C:\ z_0^4+z_1^4+z_2^4=0\subset\mathbb P^2_{\mathbb C}.
\]
Its partial derivatives have no simultaneous projective
zero, and adjunction gives $K_C=\mathcal O_C(1)$ and
genus $3$. Set
\[
 W=\mathbb P^m\times C,\qquad
 L=\mathcal O_{\mathbb P^m}(1)\boxtimes K_C,
 \qquad f=x_0z_0+x_1z_1+x_2z_2\in H^0(W,L).
\]
Here unused coordinates $x_3,\ldots,x_m$ are allowed.
The morphism
\[
 \mathcal O_C^{m+1}\longrightarrow K_C,
       \qquad (a_0,\ldots,a_m)\longmapsto
                         a_0z_0+a_1z_1+a_2z_2
\]
is surjective. Its kernel $E$ is a vector bundle of rank
$m$, and $H=V_W(f)$ is its projective bundle of lines.
Consequently $H$ is smooth and connected, of dimension
$m$. Both the source and target varieties used below
therefore have dimension at least two.

The full section ring
\[
 B=\bigoplus_{a\ge0}H^0(W,L^a)
\]
is a normal domain, generated in degree one. Indeed it
is the Segre product of a polynomial ring with the
homogeneous coordinate ring of the plane quartic.
The complete section rings of these projectively normal
embeddings give the asserted normality and generation.
Let $v$ denote the cone vertex, let
$A=B_{B_+}$, and set $R=A^{\mathrm{sh}}$ with the
residue field $\mathbb C$ fixed. This agrees with the
henselization, since the residue field is algebraically
closed. The ring is excellent, normal and of dimension
$m+2=N$. The punctured affine cone over $(W,L)$ is the
complement of the zero section in $\operatorname{Tot}(L^{-1})$.
It is smooth, and the section $f$ cuts out the corresponding
punctured line bundle over $H$. Passing to the strict
henselization preserves regularity, so $X,Y$ have the
required regularity.

We first check that the possible nonnormality of the
closed hyperplane cone at its vertex is harmless here.
Let
\[
 B_H=\bigoplus_{a\ge0}H^0(H,L_H^a),\qquad L_H=L|_H.
\]
The exact sequence
\[
 0\longrightarrow L^{a-1}\xrightarrow{\ f\ }L^a
        \longrightarrow L_H^a\longrightarrow0
\]
gives an injection $B/(f)\hookrightarrow B_H$ whose
cokernel is supported in finitely many degrees: the
next coherent cohomology vanishes for $a$ sufficiently
large. Thus $B_H$ is a finite $B/(f)$-module. It is a
normal domain and the map is an isomorphism away from
the vertex. In particular it is the normalization and
the $S_2$-hull of the hyperplane cone; it has just one
point above the vertex, with residue field $\mathbb C$.
Finite base change and henselization commute in this
one-point situation. Moreover henselization commutes
with quotient by $f$. It follows that the puncture
$Y$ is also the puncture of the strict henselization of
$(B_H)_{(B_H)_+}$. We may compute its local Picard group
with this normal cone. No depth hypothesis on $R/(f)$
has been inserted.

Here is the actual group calculation that will be used.
In characteristic zero, for a smooth projective variety
$V$ of dimension at least two and an ample embedding
line bundle $M$, the local Picard \emph{scheme} of its
cone has the natural exact sequence
\begin{equation}\label{eq:C3-Kollar-cone}
 0\longrightarrow\bigoplus_{a>0}H^1(V,M^a)
  \longrightarrow\mathbf{Pic}^{\mathrm{loc}}(v,C(V,M))
  \longrightarrow\mathbf{Pic}(V)/\mathbb Z[M]
  \longrightarrow0.
\end{equation}
The left term is a vector group, not merely a tangent
space. The sequence is functorial for the hyperplane
restriction under consideration. These are
\citet[Example 4.2, equations (4.2.3)--(4.2.4),
p.~470]{D_Kollar2016LocalPicard}.
Its elementary mechanism is the proper cone contraction
from $\operatorname{Tot}(M^{-1})$:
on the formal line bundle the characteristic-zero
logarithm identifies $1+I$ with $I$, and
$I=\prod_{a>0}M^a$. Restriction to the zero section has
the splitting supplied by the projection to $V$.
Formal existence for this proper contraction gives its line
bundles, and removing the exceptional divisor divides
out by its class $[M]$. In particular this splitting
is compatible with the inclusion $H\hookrightarrow W$ and the
restriction maps. The local Picard-scheme calculation
in the cited result makes these actual vector groups.
For the full section ring $B_H$, this contraction is
the blow-up of the ideal generated by the degree-one
sections; we do not assert that it is the blow-up of
the maximal ideal. Indeed $B_H$ need not be generated
in degree one. In the present example the section-ring
defect is concentrated in degree two and is
$H^1(W,L)\simeq\mathbb C^{m+1}$. On a chart where a
degree-one section $t$ is chosen, the contraction has
ring $\mathcal O(H_t)[t]$, proving this description
directly. Its exceptional divisor and formal normal
bundle are still exactly $(H,L_H^{-1})$.

Also, over an algebraically closed ground field,
\begin{equation}\label{eq:C3-henselian-Picard}
 \mathbf{Pic}^{\mathrm{loc}}(v,C(V,M))(\mathbb C)
   =\operatorname{Pic}
       \bigl(\operatorname{Spec}\mathcal O^{\mathrm{sh}}_{C(V,M),v}
                                                     \setminus\{v\}\bigr).
\end{equation}
This is \citet[Section 3.5, p.~469]{D_Kollar2016LocalPicard}.
It is the point that passes from a local Picard scheme
to the requested strict henselian puncture. We are
not inferring a line-bundle counterexample from a
tangent vector, nor replacing strict henselization
by completion without a comparison theorem.

The cohomology and its restriction can now be calculated
explicitly. Kunneth and Serre duality give
\[
 \begin{aligned}
 H^1(W,\mathcal O_W)&=H^1(C,\mathcal O_C)
                           =H^0(C,K_C)^*,\\
 H^1(W,L)&=H^0(\mathbb P^m,\mathcal O(1))
                            \otimes H^1(C,K_C)
                              \simeq\mathbb C^{m+1},\\
 H^1(W,L^a)&=0\quad(a\ge2),\qquad
 H^2(W,L^a)=0\quad(a\ge0).
 \end{aligned}
\]
Since $z_0,z_1,z_2$ form a basis of $H^0(C,K_C)$,
multiplication by $f$ on the first of these groups is
\[
 H^0(C,K_C)^*\longrightarrow
                H^0(\mathbb P^m,\mathcal O(1)),\qquad
 \lambda\longmapsto\sum_{j=0}^{2}\lambda(z_j)x_j.
\]
It is injective with image
$V_0=\langle x_0,x_1,x_2\rangle\simeq\mathbb C^3$.
The hyperplane exact sequence therefore yields
\[
 H^1(H,L_H)=H^1(W,L)/V_0,
 \qquad H^1(H,L_H^a)=0\quad(a\ge2).
\]
Thus restriction on the vector groups in
\eqref{eq:C3-Kollar-cone} is the quotient
$\mathbb C^{m+1}\to\mathbb C^{m-2}$ with kernel $V_0$.

Finally the projective-bundle formula gives
\[
 \operatorname{Pic}(W)=\mathbb Z[\mathcal O_{\mathbb P^m}(1)]
                              \oplus\operatorname{Pic}(C),
 \qquad
 \operatorname{Pic}(H)=\mathbb Z[\mathcal O_{\mathbb P(E)}(1)]
                              \oplus\operatorname{Pic}(C).
\]
Restriction is an isomorphism between these groups and
carries $L$ to $L_H$; it is the same isomorphism on the
corresponding Picard schemes. Applying
\eqref{eq:C3-Kollar-cone} to $W$ and $H$, with the
compatible projection splittings, shows that the full
kernel on local Picard schemes is precisely the
vector group $V_0$, and that the map is surjective.
Equation \eqref{eq:C3-henselian-Picard} now proves
\eqref{eq:C3-picard-kernel} for actual line bundles.
Every nonzero point of this additive group is
non-torsion. On the regular punctures,
$\operatorname{Pic}=\mathrm{CH}^1$, giving the asserted
refutations.
\end{proof}

\emph{The necessary distinction concerning $J^1$.}
The same calculation refutes injectivity of the
\emph{identity components of the local Picard schemes}:
the displayed vector kernel is already in their
identity components. In this specific example those
components are
\[
 \operatorname{Pic}^0(C)\times\mathbb G_a^{m+1}
       \longrightarrow
 \operatorname{Pic}^0(C)\times\mathbb G_a^{m-2},
\]
with identity on the abelian factor and the displayed
vector-space quotient on the other factor.
However, the algebraic equivalence used in that
local Picard functor permits an etale neighbourhood
of the product with the parameter scheme that varies
with the parameter
\citep[Definition 3.3, clauses (2.a)--(2.c),
pp.~467--468]{D_Kollar2016LocalPicard}.
Ordinary $\mathrm{CH}^1(X)_{\mathrm{alg}}$ uses
families on the fixed scheme $X$. These conventions
must not be identified without proof. Accordingly,
\eqref{eq:C3-picard-kernel} is a refutation of the
full rational-equivalence $A^1$ statements, and of
the local-Picard-identity-component reading of $J^1$;
it is not asserted here to refute
$J^1=\mathrm{CH}^1_{\mathrm{alg}}\otimes\mathbb Q$.
The already visible abelian factor in this example
maps isomorphically.

The missing depth condition is substantive.
Indeed $H^1(W,\mathcal O_W)\ne0$ makes
$\operatorname{depth}A=2$, whereas $f$ is a
nonzerodivisor, so $\operatorname{depth}A/(f)=1$.
This lies outside the characteristic-zero injection
theorem of \citet[Theorem 0.1, p.~833]{D_BhattDeJong2014},
which requires dimension at least four and
$\operatorname{depth}A/(f)\ge2$.
Their Examples 1.31 and 1.33--1.34, pp.~843--844,
also explain why a cohomological or analytic example
requires an additional algebraization argument.
The group-scheme calculation and
\eqref{eq:C3-henselian-Picard} supply the requisite
strict-local argument here, with a smooth projective-bundle
hyperplane section explicitly fixed.
The exact-arithmetic certificate
\texttt{code/verify\_local\_picard.py} verifies the
three affine-chart smoothness and section-generation
certificates, and the contraction/quotient matrices for
$4\le N\le20$. The proof of the assertions for every
$N$ is the uniform calculation above.

\paragraph{C.3j (Proved in the stated algebraizable characteristic-zero scope): ordinary algebraically trivial line bundles restrict injectively.}
Let $k$ be algebraically closed of characteristic zero,
let $(A,\mathfrak m,k)$ be a normal local ring essentially
of finite type over $k$, and put $R=A^{\mathrm{sh}}$.
Assume $N=\dim A\ge4$ and that the punctured spectrum
\[
 X=\operatorname{Spec}R\setminus\{\mathfrak m_R\}
\]
is regular. Let $0\ne f\in\mathfrak m_R$ and assume
that $Y=X\cap V(f)$ is regular. For ordinary algebraic
equivalence, defined by line-bundle families on the
fixed scheme times a connected smooth parameter curve,
the restriction
\begin{equation}\label{eq:C3-ordinary-alg-injection}
 \operatorname{Pic}(X)_{\mathrm{alg}}
           \longrightarrow\operatorname{Pic}(Y)_{\mathrm{alg}}
\end{equation}
is injective integrally. Consequently it is injective
after tensoring with $\mathbb Q$, with the convention
$J^1=\mathrm{CH}^1_{\mathrm{alg}}\otimes\mathbb Q$.
No depth condition on $R/(f)$ is required in this
statement. We do not assert surjectivity, the case
$N=3$, mixed characteristic, or the same theorem for
every nonalgebraizable excellent local ring.

\begin{proof}
First pass to a pointed etale neighbourhood so that
$f$ is defined over the chosen algebraic local ring.
This does not change $R$. The local Picard functor is
also unchanged by this replacement: its definition
already sheafifies the pointed etale neighbourhoods.
Let $G$ denote its representing group scheme, which
exists by the hypotheses on the normal algebraic
singularity and its dimension. It is locally of finite
type, and
\[
                    G(k)=\operatorname{Pic}(X)
\]
by \citet[Definition 3.3 and Section 3.5,
pp.~467--469]{D_Kollar2016LocalPicard}.
The subgroup scheme $K$ representing those line
bundles that become trivial on the punctured
hyperplane section exists even if $A/(f)$ is not
$S_2$; it is defined using its $S_2$-hull, which has
the same puncture. By
\citet[Theorem 1.7(i), p.~463, and Definition 3.3,
p.~468]{D_Kollar2016LocalPicard}, $K$ is a unipotent
algebraic group. All hypotheses of that theorem hold:
the algebraic germ is $S_2$, pure of dimension at least
four, in characteristic zero, and $f$ defines an
effective Cartier divisor.

We prove the separate fact needed to apply this
theorem to \emph{ordinary} algebraic equivalence:
every element of $\operatorname{Pic}(X)_{\mathrm{alg}}$
belongs to the $k$-points of an abelian-subvariety image
in $G$. Suppose first it is the difference of the two
fibres of a line bundle $\mathcal M$ on $X\times T$,
where $T$ is a smooth connected curve with marked
$k$-points $t_0,t_1$.

The entire family, not only its two fibres, descends
to a single pointed etale algebraic neighbourhood of
the original germ. Here are the finite-presentation
details. Choose finitely many generators of the
maximal ideal of $A$; their principal opens cover
$X$. The strict henselization is the filtered limit
of pointed affine etale neighbourhoods, and these
same principal opens describe the corresponding
punctures at each stage. Cover $T$ by finitely many
affine opens. The line bundle $\mathcal M$ is given
on a finite refinement of these product covers by
finitely many generators, transition units and their
inverses, satisfying finitely many cocycle identities.
All coefficients and identities descend to one stage
of the filtered system. The same applies to any
specified fibre isomorphisms. Local freeness of rank
one descends as well, since the finite trivializing
cover and its unit identities descend. After
shrinking an algebraic representative around its
marked point, we therefore have a smooth
quasiprojective variety $U$ and a line bundle
$\mathcal M_U$ on $U\times T$ that pull back to
$X$ and the given family. The shrinking is uniform
in $T$: it is the same finite-descent argument on
the finite affine cover of $T$, with finitely many
denominators from the local algebraic neighbourhood.

Choose a smooth projective compactification
$U\hookrightarrow\overline U$. To make the input
precise, embed a quasiprojective algebraic model in
projective space, take its closure, and apply
characteristic-zero embedded resolution, which is
an isomorphism over its smooth open $U$;
\citet[Theorem 1.0.2(b)--(c), p.~781]{D_Wlodarczyk2005}
supplies this statement. Also compactify $T$ to its
smooth projective completion $\overline T$.

The line bundle extends to
$\overline U\times\overline T$. Indeed choose a
nonzero rational section of $\mathcal M_U$ on the
integral variety $U\times T$. Its divisor has a
closure on $\overline U\times\overline T$. This
product is smooth, hence regular and locally
factorial, so the closure is a Cartier divisor.
The associated invertible sheaf $\overline{\mathcal M}$
restricts to $\mathcal M_U$. Its two marked fibres
satisfy
\[
 [\overline{\mathcal M}_{t_1}\otimes
             \overline{\mathcal M}_{t_0}^{-1}]
                           \in\operatorname{Pic}^0(\overline U)(k),
\]
since $\overline T$ is connected. The usual
Poincare line bundle on
$\overline U\times\operatorname{Pic}^0(\overline U)$,
restricted to the chosen neighbourhood and then its
strict henselization, defines a homomorphism of
group schemes
\[
              \operatorname{Pic}^0(\overline U)\longrightarrow G.
\]
Its image $B$ is a proper connected algebraic subgroup:
the source is an abelian variety, and the image of
this proper group in the separated algebraic group
$G^0$ is closed and proper. In characteristic zero
it is an abelian subvariety. Our original difference
of fibres lies in $B(k)$.
For a finite sum of algebraic-equivalence relations,
apply the same argument to the product of their
Picard varieties and the sum of the homomorphisms.
Its image is again an abelian subvariety of $G$
containing the proposed line bundle.

Now suppose that this line bundle restricts to the
trivial bundle on $Y$. Its point lies in $B\cap K$.
This intersection is both proper, as a closed
subgroup of $B$, and affine, as a closed subgroup
of the unipotent group $K$. It is therefore finite.
A unipotent group over a characteristic-zero field
has no nonidentity torsion point: in a faithful
unipotent matrix representation, the minimal
polynomial of a torsion element divides both a power
of $T-1$ and the separable polynomial $T^r-1$, so
the element is the identity. Thus $(B\cap K)(k)=\{1\}$,
and the proposed line bundle on $X$ is trivial.
This proves \eqref{eq:C3-ordinary-alg-injection}.
\end{proof}

\emph{Exact ordinary-equivalence calculation in C.3i.}
In the Fermat/projective-bundle example, the ordinary
algebraically trivial subgroup of each puncture's
Picard group is exactly the visible
$\operatorname{Pic}^0(C)(\mathbb C)$ factor. For the
inclusion in this factor, apply the preceding
finite-descent and compactification argument: every
such family maps from an abelian variety, whose
map to the affine vector factor is zero. For the
reverse inclusion, pull the Poincare family on
$C\times\operatorname{Pic}^0(C)$ back along the
punctured line bundle's projection to $C$.
These projections exist for the source and the
normalized target cone and therefore on their
strict local punctures. It follows that
\[
 J^1(X)\longrightarrow J^1(Y)
\]
is an isomorphism in this example with ordinary
Chow algebraic equivalence, even though the
restriction of the local Picard identity components
has the additive kernel from C.3i.

Moreover put
$B^1(Z)=\mathrm{CH}^1(Z)/\mathrm{CH}^1(Z)_{\mathrm{alg}}$
with ordinary algebraic equivalence. Since the
restriction of the algebraically trivial subgroups
is an isomorphism and the full Picard restriction
is surjective, the exact kernel calculation also
gives
\[
 \ker\bigl(B^1(X)\longrightarrow B^1(Y)\bigr)
                       \simeq(\mathbb C,+)^3,
\]
and this quotient restriction is surjective.
Thus the same explicit classes refute the proposed
unrestricted degree-one cycle restriction modulo
ordinary algebraic equivalence as well. The
different local Neron--Severi group, formed by
dividing by $\mathbf{Pic}^{\mathrm{loc},0}$,
does not retain this vector kernel. These are
three distinct quotients; the notation must specify
which one is meant.

\paragraph{C.3k (Proved): torsion injectivity in dimension three and the ordinary algebraic-equivalence consequence.}
The cyclic-cover mechanism used in
\citet[Section 2.5, p.~466]{D_Kollar2016LocalPicard}
also gives the following precise statement in dimension
three. We give the normalization, unit and connectedness
arguments, since they are indispensable at this boundary.

\emph{(i) Prime-to-residue-characteristic torsion.}
Let $(R,\mathfrak m,k)$ be an excellent, normal, strictly
henselian local ring of dimension $N\ge3$. Let
$0\ne f\in\mathfrak m$, and put
\[
 X=\operatorname{Spec}R\setminus\{\mathfrak m\},
 \qquad Y=X\cap V(f).
\]
Assume that $Y$ is regular. For every positive integer
$n$ invertible in $R$, restriction induces an injection
\begin{equation}\label{eq:C3-prime-to-p-Pic-injection}
              \operatorname{Pic}(X)[n]
                    \longrightarrow\operatorname{Pic}(Y)[n].
\end{equation}
Here the line bundles and the assertion are integral.
The result applies in equal and mixed characteristic,
and does not require algebraizability. Normality of
$R$ suffices on the source; regularity of $X$ is an
allowed stronger hypothesis from the manuscript.

\begin{proof}
We first determine the global units that enter Kummer
theory. Since $R$ is normal and $N\ge2$, the local
Hartogs property gives
\[
                  \Gamma(X,\mathcal O_X)=R.
\]
For completeness, the extension statement follows by
expressing a normal noetherian domain as the intersection
of its height-one localizations inside its fraction
field: every height-one point lies in $X$.
Every unit of $R$ has an $n$th root. Its reduction has
a root in the separably closed residue field, and
the derivative of $T^n-u$ is invertible at that root;
the henselian property supplies the lift.

We need the corresponding statement on $Y$, where
$R/(f)$ need not be normal or $S_2$. Put $C=R/(f)$.
This ring is reduced. Indeed normality gives
$\operatorname{depth}R\ge2$, and $f$ is a non-zero-divisor,
so $\operatorname{depth}C\ge1$. At every point except
its closed point, $C$ is regular by hypothesis.
It therefore satisfies $S_1$ everywhere and is
generically reduced, which together imply reducedness.
Each component has dimension $N-1\ge2$: the minimal
primes over $(f)$ have height one, and excellence
gives the catenary dimension formula in the local
domain $R$.

Let $\widetilde C$ be the normalization of $C$ in
its total ring of fractions. It is finite over $C$
because $C$ is excellent, hence Nagata; see
\citet[Lemma 29.54.14, Tag 0AVK]{D_StacksLocalPicard}.
Normalization is an isomorphism above the regular
open $Y$. A finite algebra over a henselian local
ring is a product of henselian local rings
\citep[Lemma 10.153.4, Tag 04GH]{D_StacksLocalPicard},
so
\[
                    \widetilde C=\prod_{a=1}^{s} C_a,
\]
where each $C_a$ is a normal henselian local domain
of dimension $N-1$. Its residue field is a finite
extension of the separably closed field $k$, hence
is purely inseparable and still separably closed.
Thus each $C_a$ is strictly henselian. Removing the
closed points of $\operatorname{Spec}\widetilde C$
gives exactly $Y$. Applying the same height-one
intersection argument to the normal rings $C_a$
gives the equality, not merely an inclusion,
\begin{equation}\label{eq:C3-section-finite-normalization}
             \Gamma(Y,\mathcal O_Y)=\widetilde C.
\end{equation}
In particular this ring is finite over $C$, and
its group of units is $n$-divisible, by Hensel's
lemma in each factor. This supplies a global
$n$th root for every unit that can occur in a
trivialization on $Y$.

The Kummer sequences on $X$ and $Y$, together with
these two divisibility statements, now give
compatible isomorphisms
\begin{equation}\label{eq:C3-Kummer-Pic-identifications}
\begin{aligned}
 H^1_{\mathrm{et}}(X,\mu_n)&\simeq\operatorname{Pic}(X)[n],\\
 H^1_{\mathrm{et}}(Y,\mu_n)&\simeq\operatorname{Pic}(Y)[n].
\end{aligned}
\end{equation}
Suppose there were a nontrivial kernel element of
order $d$ dividing $n$. Use $d$ in place of $n$.
All $d$th roots of unity lie in $R$, so $\mu_d$ is
the constant cyclic group on the connected normal
scheme $X$. Its Kummer class has exact order $d$.
Under
\[
 H^1_{\mathrm{et}}(X,\mu_d)
       =\operatorname{Hom}_{\mathrm{cont}}
                       (\pi_1^{\mathrm{et}}(X),\mu_d),
\]
the order of a character equals the order of its
image. Thus this class defines a connected finite
etale $\mu_d$-torsor $X'\to X$ of degree $d$.
Its restriction to $Y$ is the trivial torsor by
\eqref{eq:C3-Kummer-Pic-identifications}. In particular
\begin{equation}\label{eq:C3-split-restricted-cover}
                       X'\times_XY\simeq\coprod_{j=1}^{d}Y.
\end{equation}

The connected normal scheme $X'$ is integral.
Let $S$ be the normalization of $R$ in its function
field. This is a finite normal $R$-algebra, again
by excellence and the finite field extension.
As a finite algebra over the henselian ring $R$,
it is a product of local rings; since it is a
domain, there is only one factor. Consequently
$S$ is a normal henselian local domain, of
dimension $N$. Its punctured spectrum is $X'$:
both are the normalization of $X$ in the same
field. There is one closed point above
$\mathfrak m$, and therefore
\[
 X'\times_XY
       =\operatorname{Spec}(S/fS)\setminus\{\mathfrak m_S\}.
\]

We show that this punctured divisor is connected.
The ring $S$ is excellent. Its completion
$\widehat S$ is a normal local domain
\citep[Lemma 15.53.6, Tag 0C23]{D_StacksLocalPicard}.
The spectrum of this complete domain is irreducible
of dimension $N$, and
$\operatorname{cd}(\widehat S,(f))\le1$, as is seen
from the two-term Cech complex for a principal
ideal. The connectedness theorem
\citep[Theorem 1.6, p.~2497]{D_Varbaro2009},
in the equivalent punctured form of
\citet[Lemma 52.14.1, Tag 0ECR]{D_StacksLocalPicard},
applies because
\[
              1<N-1=\min\{N,N-1\}.
\]
It gives that
$\operatorname{Spec}(\widehat S/f\widehat S)
\setminus\{\mathfrak m_{\widehat S}\}$ is connected.
Completion is faithfully flat and its fibre at
the closed point is the residue field. Hence
the induced morphism from this puncture onto
$\operatorname{Spec}(S/fS)\setminus\{\mathfrak m_S\}$
is surjective. A continuous image of a connected
space is connected, proving the claim before
completion as well. Since $Y$ is nonempty and
$d>1$, this contradicts
\eqref{eq:C3-split-restricted-cover}.
Thus the proposed torsion kernel is zero.
\end{proof}

\emph{(ii) Ordinary algebraic equivalence in
algebraizable characteristic zero.}
Let $k$ be algebraically closed of characteristic
zero, let $A$ be a normal local ring essentially
of finite type over $k$ with residue field $k$,
and put $R=A^{\mathrm{sh}}$. In the setting above,
assume also that $X$ is regular. Then for every
$N\ge3$ the restriction
\begin{equation}\label{eq:C3-ordinary-alg-N3}
          \operatorname{Pic}(X)_{\mathrm{alg}}
                  \longrightarrow\operatorname{Pic}(Y)_{\mathrm{alg}}
\end{equation}
is injective integrally, where algebraic equivalence
is defined by families on the fixed scheme times
a smooth connected parameter curve. It follows
that the rational ordinary $J^1$ restriction is
injective in this exact scope, including $N=3$.

\begin{proof}
After descent of $f$ to one pointed etale
neighbourhood, use that neighbourhood as the
algebraic germ; this leaves $R$ unchanged.
The source local Picard scheme $G$ exists for
the normal algebraic germ of dimension at least
three, with $G(k)=\operatorname{Pic}(X)$.
The finite family-descent, smooth projective
compactification and Poincare-family argument
in C.3j uses this source representability and
characteristic-zero resolution, but does not
use $N\ge4$ or unipotence. Applied here, it shows
that every ordinarily algebraically trivial
line bundle belongs to an abelian subvariety
$B\subset G$.

The kernel of restriction on $B$ is a closed
subgroup scheme of $B$, including when $N=3$.
To justify this without assuming a finite-type
Picard scheme for a two-dimensional target,
take the finite normalization of the algebraic
hyperplane germ. It is an isomorphism away
from the closed point because that puncture
is regular. Its finitely many local factors
are normal of dimension $N-1\ge2$, and hence
have depth at least two. The unit section of
each local Picard functor is represented by
a finitely presented closed immersion by
\citet[Theorem 1(i), Introduction, p.~VII]{D_Boutot1978};
the exact statement is also recorded in
\citet[Definition 3.3, p.~468]{D_Kollar2016LocalPicard}.
Pulling back the product of these unit
sections along the morphism of group functors
\[
 B\lhook\joinrel\longrightarrow G
       \longrightarrow
       \prod_a\mathbf{Pic}^{\mathrm{loc}}(y_a,\widetilde D_a)
\]
gives a closed subgroup scheme $K_B\subset B$.
Here $\widetilde D_a$ are the normal local
factors of the finite normalization and
$y_a$ their closed points; the second arrow
is restriction of local Picard functors.
This construction does not require a
universal line bundle on $B$.
Its $k$-points are precisely
the line bundles in $B(k)$ trivial on $Y$.
Etale base change of the finite normalization
and then the strict henselization preserve
its puncture, so this describes the required
restriction, not a different algebraic model.

If $K_B$ had a nonidentity $k$-point, it would
have a nonidentity torsion $k$-point. Indeed
its reduced identity component is an abelian
subvariety: if it has positive dimension,
its $\ell$-torsion is nontrivial for any prime
$\ell$; if its dimension is zero, the group
is finite and every $k$-point is torsion.
Here group schemes of finite type in
characteristic zero are smooth, so no
infinitesimal qualification is needed.
The inclusion $B\subset G$ would make this
a nontrivial torsion line bundle on $X$
restricting trivially to $Y$, contrary to
part (i), since every positive integer is
invertible in characteristic zero. Thus
$K_B(k)=\{1\}$, proving
\eqref{eq:C3-ordinary-alg-N3}.
\end{proof}

Part (ii) strengthens the injectivity statement
of C.3j, but does not assert surjectivity in
dimension at least four. Part (i) is valid
for arbitrary excellent normal strictly
henselian rings in the stated characteristics;
the ordinary-algebraic-equivalence conclusion
uses the separate algebraizable
characteristic-zero hypotheses explicitly.

\paragraph{C.3l (Proved): the two universal abelian injectivity assertions have one Beauville gate.}
Fix an algebraically closed field $k$ of arbitrary characteristic.
All cycle groups in this paragraph have rational coefficients.
The following two assertions are equivalent:
\begin{enumerate}
\item For every abelian variety $A/k$, both injectivity assertions
      of C.3h hold in every indicated degree.
\item For every abelian variety $V/k$, the ordinary Chow group
      has no negative Beauville pieces:
      $\mathrm{CH}^r(V)_{(t)}=0$ for every $r$ and $t<0$.
\end{enumerate}
The second assertion is the usual Beauville vanishing conjecture,
recalled in \citet[Equation (5.9), p.~262]{rootBeauville2010}.
We prove the equivalence, not the conjecture. Its quantifier over
all abelian varieties is essential: a family on $A$ introduces
the Jacobian of its parameter curve.

\begin{proof}
The additional point beyond C.3h is the following parameterization.
For each $z\in J^r(A)_{(s)}$, there are finitely many smooth
projective pointed curves $C_j$, points $a_j\in J(C_j)(k)$,
and cycles
\[
 Z_j\in\mathrm{CH}^r(J(C_j)\times A)_{(s-1)}
 \quad\text{such that}\quad
 z=\sum_j\bigl((Z_j)_{a_j}-(Z_j)_0\bigr).
\]
Here the subscript on $Z_j$ denotes the Beauville grading for
simultaneous multiplication on both abelian factors; its
evaluation is the refined pullback to the indicated point.

To prove this, express ordinary algebraic equivalence by curve
families: $z$ is a finite sum of $W_{j*}\delta_j$, where
$W_j\in\mathrm{CH}^r(C_j\times A)$ and
$\delta_j\in\mathrm{CH}_0(C_j)^0_{\mathbb Q}$.
Clear the finitely many rational coefficients, incorporating
them in $W_j$, so that each $\delta_j$ is an integral
degree-zero divisor class. It determines a point
$a_j\in J(C_j)(k)=\operatorname{Pic}^0(C_j)(k)$.
Let $D_j=c_1(\mathcal P_j)$ for the Poincar\'e line bundle on
$J(C_j)\times C_j$, normalized along zero and the base point
of $C_j$. On $J(C_j)\times C_j\times A$ put
\[
 \widetilde Z_j=(p_{13})_*
               \bigl(p_{12}^*D_j\cdot p_{23}^*W_j\bigr).
\]
The codimension is $1+r-1=r$. The Poincar\'e bundle's
defining universal property and its rigidification give
$([m]\times1)^*D_j=mD_j$: tensor power $m$ classifies
multiplication by $m$ on $\operatorname{Pic}^0(C_j)$,
and the rigidification removes a pullback from the base.
Fix $m>1$ prime to the characteristic. Flat base change
therefore gives $([m]\times1)^*\widetilde Z_j=m\widetilde Z_j$.
Decompose $\widetilde Z_j$ by the ordinary Beauville grading
of the abelian variety $J(C_j)\times A$. Its projectors are
polynomials in simultaneous multiplication and commute with
the first-factor operator. Thus each piece $Z_{j,(t)}$
retains first-factor eigenvalue $m$, while its total
eigenvalue is $m^{2r-t}$. The first-factor pullback is
injective on rational Chow groups: its finite-flat
pushforward composed with pullback is its nonzero degree.
Cancelling this pullback consequently gives
\[
 (1\times[m])^*Z_{j,(t)}=m^{2r-t-1}Z_{j,(t)}.
\]
Its fibre at any point of $J(C_j)$ therefore has $A$-index
$t+1$. The unprojected fibre at $a_j$ is $W_{j*}\delta_j$,
and the zero fibre is zero, by normalization of the
Poincar\'e bundle. These statements follow from commutation
of proper pushforward with refined pullback in the product
square. Every graded piece also has zero fibre at zero:
evaluation there commutes with first-factor multiplication,
so this fibre equals $m$ times itself. Taking the $s$-piece
of the sum of the evaluated classes now proves the required
formula with $Z_j=Z_{j,(s-1)}$. This uses only the ordinary
grading on an abelian product, without a relative Chow
projector on $C_j\times A$. No characteristic-zero argument
or isogeny denominator is involved; rational coefficients
are used for the Beauville decomposition itself.

Assume (2). If $s\le0$, all of these parameterizing cycles
vanish, since $s-1<0$ on the abelian varieties
$J(C_j)\times A$. Hence $J^r(A)_{(s)}=0$ for $s\le0$.
The quotient $B^r(A)$ also has no negative piece, directly
from (2). C.3h now proves both assertions in (1).

Conversely (1) and C.3h give $B^r(A)_{(s)}=0$ for $s<0$
and $J^r(A)_{(s)}=0$ for $s\le0$. The degreewise exact
sequence for the invariant subgroup and its quotient is
\[
 0\longrightarrow J^r(A)_{(s)}
 \longrightarrow\mathrm{CH}^r(A)_{(s)}
 \longrightarrow B^r(A)_{(s)}\longrightarrow0.
\]
For $s<0$ both ends vanish, proving (2).
\end{proof}

This uniform equivalence does not replace the smaller,
fixed-dimension tests in C.3h by a proved vanishing theorem.
In particular, it supplies no dimension bound on the
Jacobians of parameter curves for a proposed class in
$J^2(A)_{(0)}$.

\paragraph{C.3m (Refuted, with an explicit counterexample): ordinary $J^1$ surjectivity fails in every dimension at least four.}
For every $N\ge4$ we construct an excellent normal
strictly henselian local complex ring $R$ of dimension
$N$, with regular puncture $X$ and a regular punctured
principal divisor $Y$, such that
\begin{equation}\label{eq:C3-ordinary-nonsurjection}
 \operatorname{Pic}(X)_{\mathrm{alg}}=0,
 \qquad
 \operatorname{Pic}^0(C)(\mathbb C)
       \lhook\joinrel\longrightarrow
       \operatorname{Pic}(Y)_{\mathrm{alg}},
\end{equation}
where $C$ is the explicitly specified Fermat plane
quartic. The right-hand injection has nontorsion
image. Thus ordinary $J^1(X)\to J^1(Y)$ is not
surjective, even with rational coefficients.
All algebraic equivalences here use families on the
fixed scheme times a smooth connected parameter
curve. The argument below works with actual line
bundles throughout.

\emph{1. The algebraic contraction and its divisor.}
Put $m=N-2\ge2$, and define
\[
\begin{aligned}
 F(z)&=z_0^4+z_1^4+z_2^4, & C&=V(F)\subset\mathbb P^2_{\mathbb C},\\
 H(z)&=z_0^2z_1^2+z_0^2z_2^2+z_1^2z_2^2,&
 \ell(x,z)&=x_0z_0+x_1z_1+x_2z_2.
\end{aligned}
\]
Here $x_0,\ldots,x_m$ are homogeneous coordinates on
$\mathbb P^m$. Set
\[
 B=\mathbb P^m_x\times\mathbb P^2_z,
 \qquad \mathcal N=\mathcal O_B(-1,-1),
 \qquad T=\operatorname{Tot}(\mathcal N).
\]
If $t$ denotes the tautological fibre vector, the
expression
\[
                         F+t\ell H
\]
is a well-defined section of the pullback of
$\mathcal O_B(0,4)$ to $T$. Define $Z$ as its zero
scheme. Its zero-section divisor is
\[
 V=\mathbb P^m\times C,
 \qquad
 \mathcal N|_V=L^{-1},
 \qquad L=\mathcal O_{\mathbb P^m}(1)\boxtimes K_C.
\]
The curve $C$ is smooth, and $Z$ is smooth in a
neighbourhood of $V$, since its differential there
has the nonzero $dF$ component. The projection of
$Z$ onto $B$ is the graph
$t=-F/(\ell H)$ on the dense open $\ell H\ne0$.
The locus $F=\ell H=0$ has codimension two in $B$,
so its inverse image cannot be another component
of $Z$. Thus $Z$ is integral.

The standard Segre cone contraction
\[
      T\longrightarrow\mathbb A^{3(m+1)},
           \qquad w_{ij}=t x_i z_j
\]
is proper, contracts the zero section, and is an
isomorphism outside it. Let $W$ be the reduced
closed image of $Z$, and write $v$ for the vertex.
This gives a completely algebraic finite-type
construction. Explicit equations for $W$ are the
radical of the following ideal: all $2$ by $2$
minors of $(w_{ij})$, together with, for every
ordered quadruple $I=(i_1,i_2,i_3,i_4)$,
\begin{equation}\label{eq:C3-deformed-cone-equations}
 \sum_{j=0}^{2}\prod_{r=1}^{4}w_{i_rj}
 +f\sum_{0\le a<b\le2}
       w_{i_1a}w_{i_2a}w_{i_3b}w_{i_4b}=0,
 \qquad f=w_{00}+w_{11}+w_{22}.
\end{equation}
Indeed on the rank-one locus the left side is
$t^4x_{i_1}\cdots x_{i_4}(F+t\ell H)$;
choosing a nonzero row proves the converse away
from $v$. The vertex lies in the closure of this
integral open, so the reduced image is integral.

Let $W^{\mathrm n}$ be the normalization of the
original reduced affine scheme $W$; it is finite.
For the proof of its local properties choose an
open neighbourhood $W^\circ$ of $v$ such that
$Z^\circ=Z\times_W W^\circ$ is smooth.
Such a neighbourhood exists: the image of the
complement of the smooth locus is closed by
properness and does not contain $v$.
The normal smooth scheme $Z^\circ$ factors
through $(W^{\mathrm n})^\circ$, the inverse image
of $W^\circ$; the factor is proper, birational
and an isomorphism away from the vertex.
There is only one point $v^{\mathrm n}$ above $v$.
In fact the factorization is surjective, and the
connected fibre $V$ cannot map onto more than one
point in the finite fibre of the normalization.
Its residue field is $\mathbb C$. Define
\[
 A=\mathcal O_{W^{\mathrm n},v^{\mathrm n}},
 \qquad R=A^{\mathrm{sh}},
 \qquad
 \mathcal Z=Z^\circ\times_{(W^{\mathrm n})^\circ}
                         \operatorname{Spec}R.
\]
In particular $R$ is specified by the original
polynomial equations and normalization, followed
by localization at its unique vertex and strict
henselization. The auxiliary shrinking only
proves regularity and introduces no choice in $R$.
Then $R$ has all the asserted excellence, normality
and dimension properties. The regular scheme
$\mathcal Z$ is proper over $R$, its exceptional
fibre is the irreducible smooth divisor $V$, and
its complement is the regular puncture $X$.
These statements are preserved under the pointed
etale limit defining the strict henselization.

The displayed $f$ is a nonzero nonunit in $R$, and
on $Z$ it equals $t\ell$. The global equation of
the strict transform of its divisor is
\[
 D=V(\ell,F)\subset T,
\]
the total space of $L_H^{-1}$ over
\[
 H_0=V(\ell)\subset V,
 \qquad L_H=L|_{H_0}.
\]
Indeed setting $\ell=0$ in the equation of $Z$
leaves $F=0$. The surjection
\[
 \mathcal O_C^{m+1}\longrightarrow K_C,
 \qquad (a_i)\longmapsto a_0z_0+a_1z_1+a_2z_2
\]
identifies $H_0$ with the projective bundle of
lines in its rank-$m$ kernel. In particular
$H_0$ and $D$ are smooth. Set $D^\circ=D\cap Z^\circ$;
its morphism to $(W^{\mathrm n})^\circ$ is the
restriction of the factorization already constructed.
The base change
\[
 \mathcal D=D^\circ\times_{(W^{\mathrm n})^\circ}
                              \operatorname{Spec}R
\]
is regular and proper over $R$. Its exceptional
divisor is $H_0$, with normal bundle $L_H^{-1}$,
and
\[
            \mathcal D\setminus H_0=Y=X\cap V(f).
\]
This proves regularity of the required punctured
divisor, without asserting any projective-normality
property of its section ring.

\emph{2. A full-rank first-neighbourhood deformation.}
Let $V_2\subset\mathcal Z$ be defined by the square
of the ideal of $V$. It is a square-zero thickening
of $V$ by $L$. Its class
\[
                   \kappa\in H^1(V,T_V\otimes L)
\]
is obtained by taking the boundary of $\ell H$
in the normal sequence of $V\subset B$, twisted
by $L$. Its curve component is
\begin{equation}\label{eq:C3-KS-components}
       \kappa=\sum_{i=0}^{2}x_i\otimes\delta(z_iH),
 \qquad
 \delta:H^0(C,\mathcal O_C(5))\longrightarrow H^1(C,\mathcal O_C),
\end{equation}
where $\delta$ is the boundary of
\[
 0\longrightarrow T_C\otimes K_C=\mathcal O_C
 \longrightarrow T_{\mathbb P^2}|_C\otimes K_C
 \longrightarrow\mathcal O_C(5)\longrightarrow0.
\]
To verify the class directly, choose local vector
fields lifting the normal displacement $\ell H$.
Their differences on overlaps are the derivations
which identify the local retractions of $V_2$
onto $V$; this is precisely the boundary cocycle.
The sign depends on the choice of retraction and
does not affect the injectivity below.

The Euler sequence twisted by $K_C=\mathcal O_C(1)$
and Serre duality give
$H^1(C,T_{\mathbb P^2}|_C\otimes K_C)=0$.
Thus $\delta$ is surjective. More explicitly its
cokernel description is the degree-five part of
the Fermat Jacobian ring:
\[
 H^1(C,\mathcal O_C)
   \simeq
 \left(\frac{\mathbb C[z_0,z_1,z_2]}
                  {(z_0^3,z_1^3,z_2^3)}\right)_5.
\]
Here the image of quadratic vector fields consists
of the degree-five Jacobian ideal, modulo the
three multiples $z_iF$ already zero on $C$.
It has dimension $15$, while
$h^0(\mathcal O_C(5))=18$; the remaining global
tangent section is the kernel
$H^0(T_C\otimes K_C)$, so introduces no additional
relation. The three images of $z_iH$ are
\[
             z_0z_1^2z_2^2,\qquad
             z_0^2z_1z_2^2,\qquad
             z_0^2z_1^2z_2.
\]
They form a basis of that degree-five part.

Contraction with $\kappa$ consequently gives an
injective linear map
\begin{equation}\label{eq:C3-KS-full-rank}
 c_\kappa:H^0(C,K_C)
       =H^0(V,\Omega_V^1)
       \longrightarrow H^1(V,L).
\end{equation}
Indeed $H^1(V,L)=H^0(\mathbb P^m,\mathcal O(1))
\otimes H^1(C,K_C)$ and the $i$th coefficient
of this contraction is the perfect Serre-duality
pairing with $\delta(z_iH)$. Its rank is exactly
three for every $m\ge2$.

\emph{3. The actual Poincare-family obstruction.}
We record the precise first-order calculation.
For a smooth variety $V$, a square-zero thickening
by $L$ with class $\kappa$, and a line bundle
$\mathcal Q$ on $V\times P$, its obstruction to
extension to $V_2\times P$ is
\begin{equation}\label{eq:C3-line-obstruction}
        \kappa\mathbin{\lrcorner}c_1^{V}(\mathcal Q)
                   \in H^2(V\times P,\operatorname{pr}_V^*L),
\end{equation}
where $c_1^{V}$ is the component of its Atiyah
class in the $V$-differentials. This is
\citet[Theorem 2.6 and formula (2.7.2),
pp.~166--167]{D_Illusie1972Torsors}, applied to
$\mathbb G_m$-torsors. In this line-bundle case
it follows directly from transition functions:
local splittings of the thickening differ by
derivations $d_{ij}$ with values in $L$; on triple
intersections the product of lifted transition
units differs from one by the evaluation of
$d_{ij}$ on $d\log g_{jk}$. The resulting Cech
two-cocycle is the cup product in
\eqref{eq:C3-line-obstruction}. Thus the formula
uses the actual extension obstruction, not just
a tangent-space dimension.

Let $J=\operatorname{Pic}^0(C)$ and let $P$ be an
abelian variety. If a line bundle on $V\times P$
defines a morphism $\psi:P\to J$ carrying $0$ to
$0$, its $H^0(C,K_C)\otimes H^1(P,\mathcal O_P)$
mixed Atiyah component is that of the normalized
Poincare bundle pulled back along $\psi$.
The corresponding linear map is
\[
                  \psi^*:H^1(J,\mathcal O_J)
                                  \longrightarrow H^1(P,\mathcal O_P).
\]
One can verify the normalization by pulling the
Poincare bundle back along the Abel map
$C\to J$: on $C\times C$ it is
$\mathcal O(\Delta-C\times c_0-c_0\times C)$,
whose mixed diagonal class gives the identity
under Serre duality. The Abel map induces an
isomorphism on $H^1(\mathcal O)$, so this
determines the asserted component. The Poincare
universal property and this cohomological
isomorphism are recorded in
\citet[Theorem 1.2 and Lemma 9.5,
pp.~3 and 31--32]{D_MilneJacobians2021}.

For our $V_2$, the component of
\eqref{eq:C3-line-obstruction} in
\[
                H^1(V,L)\otimes H^1(P,\mathcal O_P)
\]
is obtained by applying $c_\kappa$ to the first
factor of this Poincare tensor. Since
$c_\kappa$ is injective, vanishing of the
obstruction forces $\psi^*=0$. The morphism
$\psi$ is a homomorphism by rigidity of abelian
varieties. Over $\mathbb C$, a nonzero
homomorphism of abelian varieties induces a
nonzero map on $H^1(\mathcal O)$: complex
conjugation gives the corresponding map on
holomorphic one-forms, and zero differential
everywhere forces a homomorphism to be constant.
Consequently an actual extension to
$V_2\times P$ forces $\psi=0$.

\emph{4. Ordinary algebraic families on $X$ vanish.}
Take an ordinarily algebraically trivial line
bundle $M$ on $X$. The finite-descent and smooth
compactification proof in C.3j supplies an
abelian variety $P$, a point $a\in P(\mathbb C)$,
and an actual line bundle $\mathcal M$ on
$X\times P$ such that
\[
                    \mathcal M_0\simeq\mathcal O_X,
                    \qquad \mathcal M_a\simeq M.
\]
For a sum of relations one takes a product of
the Picard varieties and the tensor product
of their pulled-back Poincare bundles. This
uses actual Poincare bundles on those source
Picard varieties, not on their possibly
non-fine abelian images.

The scheme $\mathcal Z\times P$ is regular:
$\mathcal Z$ is regular and its projection
from this product is smooth. Extend
$\mathcal M$ to a line bundle
$\widetilde{\mathcal M}$ on this product by
closing the divisor of a rational section.
The product is noetherian and locally factorial,
so the closure is Cartier. Its fibre at $0$
is trivial outside the irreducible divisor
$V$, hence is $\mathcal O_{\mathcal Z}(bV)$
for an integer $b$. Tensoring by the pullback
of its inverse normalizes that fibre to
$\mathcal O_{\mathcal Z}$, without changing
the original family on $X$.

Restrict to $V\times P$. Since the fibre at
$0$ is trivial and $P$ is connected, its
classifying morphism takes values in
$\operatorname{Pic}^0(V)=J$. It extends to
$V_2\times P$, so step 3 shows this morphism
is zero. By the Poincare universal property
or the see-saw principle the restriction is
therefore pulled back from a line bundle on
$P$. Twist by that pullback. This changes
none of the individual fibre isomorphism
classes and makes the restriction to
$V\times P$ trivial.

The unit sequence for the square-zero
thickening now shows that the class on
$V_2\times P$ comes from
\[
\begin{aligned}
 H^1(V\times P,\operatorname{pr}_V^*L)
  ={}&H^1(V,L)\otimes H^0(P,\mathcal O_P)\\
    &\oplus H^0(V,L)\otimes H^1(P,\mathcal O_P).
\end{aligned}
\]
Restriction to a point of $P$ kills the second
summand and is the same map on the first
summand for every point, since
$H^0(P,\mathcal O_P)=\mathbb C$. Thus all
its individual fibres have the same class
in $\operatorname{Pic}(V_2)$. The fibre at
$0$ is trivial, so the fibre at $a$ is
trivial as well.

For every $r\ge2$,
\[
                  H^1(V,L^r)=H^2(V,L^r)=0.
\]
This follows from the projective-space
cohomology and $H^1(C,K_C^r)=0$ for $r\ge2$.
If $V_r$ is defined by $\mathcal I_V^r$,
the exact unit sequences therefore give
isomorphisms
$\operatorname{Pic}(V_{r+1})\to\operatorname{Pic}(V_r)$
for $r\ge2$. They also make the transition
maps on global units surjective. It follows
inductively that $\widetilde{\mathcal M}_a$
admits compatible trivializations on every
$V_r$, and is formally trivial along $V$.

We spell out why this proves triviality on
the strict local puncture. On $\mathcal Z$,
\[
                        \mathfrak m_R\mathcal O_{\mathcal Z}
                                          =\mathcal I_V.
\]
The old cone coordinates generate the ideal
$(t)$ on every Segre chart; every additional
function from normalization vanishes on the
reduced fibre and hence also belongs to
$(t)$. These two inclusions give the equality.
After completing $R$, full faithfulness of
formal completion for a proper scheme
\citep[Lemma 30.24.1, Tag 0883]{D_StacksFormalPicard}
turns the compatible formal trivializations
into an isomorphism with the trivial line
bundle on $\mathcal Z_{\widehat R}$. In
particular $M$ becomes trivial on the
puncture of $\operatorname{Spec}\widehat R$.

Finally let $Q=j_*M$, the finite rank-one
reflexive $R$-module extending $M$ across
the closed point. Its local cohomology
$H^0_{\mathfrak m}(Q)$ and
$H^1_{\mathfrak m}(Q)$ vanish. Flat base
change of these groups, and the fact that
the closed fibre of completion is a point,
show that $Q\otimes_R\widehat R$ is the
extension of the pulled-back line bundle
from the completed puncture. It is therefore
$\widehat R$. Faithfully flat descent of
finite-presentation flatness gives that $Q$
is a projective rank-one $R$-module, hence
is free since $R$ is local. Thus $M$ is
trivial. This proves the first equality in
\eqref{eq:C3-ordinary-nonsurjection} without
inferring algebraization from tangent data.

\emph{5. Actual nontrivial target families.}
The projection $\mathcal D\to H_0\to C$
pulls back the Poincare bundle on $C\times J$
to an actual family on $Y\times J$.
It induces an injection
\[
                     \operatorname{Pic}^0(C)(\mathbb C)
                              \longrightarrow\operatorname{Pic}(Y).
\]
Indeed if the pullback of a line bundle $Q_C$
is trivial on $Y$, its extension to the
regular $\mathcal D$ is
$\mathcal O_{\mathcal D}(bH_0)$ for some
integer $b$. Restricting to $H_0$ gives
\[
                    \pi_C^*Q_C\simeq L_H^{-b}.
\]
Its degree on a line of the fibre
$\mathbb P^{m-1}$ is zero on the left and
$-b$ on the right. Thus $b=0$, and the
projective-bundle formula for the Picard
group gives $Q_C=\mathcal O_C$. The image
is ordinarily algebraically trivial:
points of the connected projective variety
$J$ can be joined by a chain of smooth
curves, to which the displayed family can
be restricted.

For an explicit nontorsion witness choose
$u\in\mathbb C$ transcendental over
$\overline{\mathbb Q}$, choose $v$ with
$v^4=-1-u^4$, and choose
$\zeta\in\overline{\mathbb Q}$ with
$\zeta^4=-1$. The points
\[
                         p=(u:v:1),\qquad c_0=(1:\zeta:0)
\]
lie on $C$. The line bundle
$\mathcal O_C(p-c_0)$ has infinite order.
Otherwise its Abel--Jacobi point would be
a torsion point of the Jacobian defined
over $\overline{\mathbb Q}$, hence itself
defined over that field. The Abel map is
a closed immersion, also defined over
$\overline{\mathbb Q}$
\citep[Proposition 2.3, p.~7]{D_MilneJacobians2021},
forcing $p$ to be
algebraic, contrary to the choice of $u$.
Its target pullback is therefore a nonzero
element of $J^1(Y)$, while $J^1(X)=0$.

The finite polynomial and Jacobian-ring
certificates are supplied separately. The
obstruction, proper-family and completion
arguments above, not those finite checks,
establish the actual counterexample.

\paragraph{C.3n (Proved): ordinary degree-one injectivity with a coefficient field, in every characteristic.}
Let $k$ be an algebraically closed field, and let
$(R,\mathfrak m,k)$ be an excellent normal strictly
henselian local $k$-algebra whose residue map restricts
to the identity on the specified coefficient field $k$.
Assume $N=\dim R\ge3$. Put
\[
 X=\operatorname{Spec}R\setminus\{\mathfrak m\},
 \qquad Y=X\cap V(f),\qquad 0\ne f\in\mathfrak m,
\]
and assume $X$ and $Y$ regular. Ordinary algebraic
equivalence is defined by actual line bundles on
the fixed $k$-scheme times a smooth connected
$k$-curve. Then
\begin{equation}\label{eq:C3-general-ordinary-injectivity}
 \ker\bigl(\operatorname{Pic}(X)_{\mathrm{alg}}
             \longrightarrow
            \operatorname{Pic}(Y)_{\mathrm{alg}}\bigr)
 \begin{cases}
   =0,&\operatorname{char}k=0,\\
   \text{is $p$-primary torsion},&\operatorname{char}k=p>0.
 \end{cases}
\end{equation}
In particular $J^1(X)\to J^1(Y)$ is injective with
rational coefficients in every characteristic;
in characteristic $p$ the integral map also becomes
injective after tensoring with $\mathbb Z[1/p]$.
Neither algebraizability nor a resolution of
singularities is assumed. The coefficient-field
hypothesis is retained explicitly; this statement
does not reinterpret ordinary algebraic equivalence
over an unspecified mixed-characteristic base.

\begin{proof}
\emph{1. An actual abelian-parameter family without
compactifying the puncture.}
We first prove the required family lemma for any
connected regular noetherian $k$-scheme $U$.
Let $T$ be a smooth connected $k$-curve with two
marked points $t_0,t_1$, and let $\mathcal M$ be
a line bundle on $U\times T$. Its smooth projective
completion is a curve $C$, with the same marked
points, and $c_0=t_0$. The product $U\times C$ is
regular and noetherian: its projection to $U$ is
smooth and of finite type. It is integral, since
$U$ is connected and regular and $C$ is
geometrically integral. Closing the divisor of a
rational section of $\mathcal M$ therefore extends
it to an invertible sheaf $\overline{\mathcal M}$
on $U\times C$. Local factoriality makes this
Weil-divisor closure Cartier.

The degree of its restriction to each curve fibre
over $U$ is locally constant. For example, this
follows from constancy of the Hilbert polynomial
for the invertible sheaf in the proper flat family
$C\times U\to U$. Since $U$ is connected, it is
one integer $d$. After the two twists
\[
 \mathcal M'=
 \overline{\mathcal M}\otimes
 \operatorname{pr}_U^*
       (\overline{\mathcal M}_{c_0}^{-1})\otimes
 \operatorname{pr}_C^*\mathcal O_C(-d c_0),
\]
the degree is zero and the restriction to
$U\times\{c_0\}$ is trivial. The last assertion
includes a choice of basis of the fixed
one-dimensional $k$-vector space
$\mathcal O_C(-dc_0)|_{c_0}$; it imposes no
condition on $U$ and requires no $k$-point of $U$.

Write $J=\operatorname{Pic}^0(C)$, and let
$\mathcal P_C$ be the Poincare bundle on $C\times J$
normalized along $c_0\times J$ and $C\times0$.
The Picard universal property gives a morphism
\[
                         g:U\longrightarrow J
\]
whose pullback of $\mathcal P_C$, with the factors
exchanged, is $\mathcal M'$.
The relevant form is
\citet[Theorem 1.1 and Lemma 1.3, pp.~2--3]{D_MilneJacobians2021}:
the point $c_0$ makes the Picard functor represented,
and normalization at $c_0$ removes the possible
line-bundle factor pulled back from the base.
Although those notes conventionally discuss
finite-type schemes, the same universal property
applies here. On each affine open of $U$, the
finite presentation of a line bundle and its
trivialization on the projective curve descend
to a finitely generated $k$-subalgebra. The
finite-type universal property gives the map
there; its uniqueness makes these maps glue.
Equivalently, this is the usual extension of
the represented Picard functor to arbitrary
$k$-schemes by affine limits.

Let $P=J^\vee$ and let $\mathcal P_J$ be the
normalized Poincare bundle on $J\times P$.
These actual universal line bundles and their
tensor-product group law are described in
\citet[Sections 9--10, especially the universal
property on p.~16]{D_MilneAbelian2022}.
Transposing $\mathcal P_C$ gives a family of
algebraically trivial line bundles on $J$
parametrized by $C$, trivial at $c_0$. It therefore
defines a morphism $h:C\to P$ with $h(c_0)=0$.
Now set
\[
             \mathcal N=(g\times1_P)^*\mathcal P_J
                    \quad\text{on }U\times P.
\]
This is an actual line bundle, trivial at $0$,
and its fibres satisfy
\begin{equation}\label{eq:C3-dual-Jacobian-family}
       [\mathcal N_{h(c)}]
        =[\overline{\mathcal M}_{c}\otimes
                     \overline{\mathcal M}_{c_0}^{-1}]
                    \quad(c\in C(k)).
\end{equation}
The degree-normalizing factor contributes only
a one-dimensional constant vector space to each
such fibre and hence does not change its line-bundle
class on $U$. This argument uses no autoduality
identification or choice of a principal polarization.
The Poincare tensor law shows that the family
$\mathcal N$ induces a homomorphism to any Picard
functor to which its line bundles map. Products
of the abelian varieties and tensor products
of these families handle finite sums of relations.

\emph{2. The exact general local Picard input.}
Apply this lemma to $U=X$. We use Boutot's local
Picard functor $G=\mathbf{Pic}^{\mathrm{loc}}_{R/k}$,
defined by sheafifying, on $k$-algebras $A$, the
Picard groups of the puncture of the henselization
of the pair
\[
                  (R\otimes_k A,\mathfrak m(R\otimes_k A)).
\]
The hypotheses and definition in
\citet[Introduction, Theorem 1(i)--(ii) and the
following paragraph, p.~VII]{D_Boutot1978}
start with an arbitrary noetherian local
$k$-algebra with residue field $k$.
They have no finite-type-over-$k$ requirement.
Depth at least two gives a finitely presented
closed unit section. Normal completion of
dimension at least three gives the additional
finite-dimensional first punctured coherent
cohomology and hence representability by a
group scheme locally of finite type over $k$.
Our $\widehat R$ is normal by excellence, so
all these hypotheses hold. The closed unit
section makes this group scheme separated.
Since $R$ is already henselian and $k$ is
algebraically closed, its $k$-points are exactly
\[
                         G(k)=\operatorname{Pic}(X).
\]
Indeed the pair henselization at $A=k$ is $R$,
and etale sheafification over the algebraically
closed field adds no new point classes.

An actual family on $X\times P$ restricts
functorially to these henselized products.
The family from step 1 therefore gives a
homomorphism $P\to G$. Its scheme-theoretic
image $B$ is a proper connected reduced
algebraic subgroup of $G$. To check the
finiteness and properness here, the morphism
from the proper finite-type scheme $P$ to
the separated scheme $G$ is proper, its
quasicompact image lies in a finite-type
open of the locally finite-type scheme $G$,
and its closed image is proper. The image
is reduced because $P$ is reduced. Over
the perfect field $k$ a reduced algebraic
group is smooth, so this proper connected
group is an abelian variety. Its inclusion
in $G$ is a closed immersion.
Thus every ordinary algebraically trivial
class on $X$ belongs to $B(k)$ for some
abelian subvariety $B\subset G$.

\emph{3. Closed restriction kernel and the
prime-to-$p$ obstruction.}
As proved in C.3k(i), $R/(f)$ is reduced,
its finite normalization is a product
\[
                            \prod_a C_a
\]
of normal strictly henselian local domains
of dimension $N-1\ge2$, and the disjoint
union of their punctures is $Y$.
Each residue field is $k$, since $k$ is
algebraically closed, and each $C_a$ has
the inherited coefficient field $k$.
Boutot's theorem therefore gives a finitely
presented closed unit section for each
functor $H_a=\mathbf{Pic}^{\mathrm{loc}}_{C_a/k}$,
even when that whole functor is not
represented by a finite-type scheme.

Restriction induces a morphism of group
functors
\[
                  G\longrightarrow\prod_a H_a.
\]
The maps on henselized tensor products
exist by the universal property of
henselization. In each finite local
factor $C_a$, the ideals
$\mathfrak m C_a$ and $\mathfrak m_{C_a}$
differ by a nilpotent ideal in the
quotient. This remains true after
tensoring with any $k$-algebra, so they
give the same henselization of the pair.
Pullback of the closed
unit section along
$B\to G\to\prod_a H_a$ defines a closed
subgroup scheme $K\subset B$. Its
$k$-points are precisely those classes
in $B(k)$ which restrict trivially to $Y$.
No universal line bundle on $B$ is used
in defining this subgroup.

For every prime $\ell$ different from
$\operatorname{char}k$, C.3k(i) gives
\[
                             K(k)[\ell]=0.
\]
If $K$ had positive dimension, its
reduced identity component would be
a positive-dimensional abelian variety.
Such a variety has nontrivial
$\ell$-torsion for every such $\ell$,
a contradiction. Hence $K$ is finite.
In characteristic zero its finite group
of $k$-points has no torsion of any
prime order and is therefore trivial.
In characteristic $p$ the finite
abelian group $K(k)$ is a $p$-group.
Every ordinary kernel element lies in
one of these groups $K(k)$, proving
\eqref{eq:C3-general-ordinary-injectivity}.
No uniform exponent for the entire
ordinary kernel is asserted.
\end{proof}

This strengthens the positive statement C.3k(ii)
in both its algebraizability and characteristic
restrictions. It is compatible with C.3m, which
refutes surjectivity, and with C.3i, whose full
Picard kernel lies outside ordinary algebraic
equivalence.

\paragraph{C.3o (Proved/refuted): the remaining ordinary degree-one boundary dimensions over $\mathbb C$.}
In the coefficient-field setting of C.3n, with ordinary
algebraic equivalence and rational coefficients, the
universal injectivity threshold over $\mathbb C$ is
$N\ge3$. Surjectivity is automatic for $N\le2$ and
fails in every dimension $N\ge3$. The assertions
for $N\ge4$ follow from C.3m--n. Here are explicit
examples and proofs for the remaining dimensions;
they use actual line bundles on strictly henselian
punctures.

\begin{center}
\begin{tabular}{lll}
\toprule
Dimension $N$ & Always injective & Always surjective \\
\midrule
$1$ & Yes, both groups vanish & Yes, both groups vanish \\
$2$ & No, the cone example below & Yes, the target vanishes \\
$\ge3$ & Yes, by C.3n & No, below and C.3m \\
\bottomrule
\end{tabular}
\end{center}

\begin{proof}
Let
\[
 F=z_0^4+z_1^4+z_2^4,\qquad
 C=V_+(F)\subset\mathbb P^2_{\mathbb C},\qquad
 L=\mathcal O_C(1).
\]
This is the smooth plane quartic already used in
C.3i and C.3m, with genus three and $\deg L=4$.
Define the two-dimensional excellent normal
strictly henselian local ring
\[
 R_2=
 \left(\mathbb C[z_0,z_1,z_2]/(F)\right)_{(z_0,z_1,z_2)}^{\mathrm{sh}},
 \qquad
 X_2=\operatorname{Spec}R_2\setminus\{\mathfrak m_2\}.
\]
The affine cone is smooth away from its vertex,
so $X_2$ is regular. Normality follows, for example,
from the hypersurface Cohen--Macaulay property and
regularity in codimension one; excellence preserves
normality under strict henselization. As throughout,
the algebraically closed residue field is identified
with the coefficient field $\mathbb C$.

The cone projection gives an actual morphism
$\rho:X_2\to C$. We claim that
\begin{equation}\label{eq:C3-low-dimension-Jacobian}
 \rho^*:\operatorname{Pic}^0(C)(\mathbb C)
       \lhook\joinrel\longrightarrow
       \operatorname{Pic}(X_2)_{\mathrm{alg}}
\end{equation}
is injective. Pulling back the usual Poincare
families on $C$ shows that its image consists of
ordinary algebraically trivial classes.
For injectivity, base change the proper cone
resolution $\operatorname{Tot}(L^{-1})\to
\operatorname{Spec}\mathbb C[z_0,z_1,z_2]/(F)$
to $\operatorname{Spec}R_2$. The resulting regular
scheme $Z_2$ has exceptional divisor $C$, is
isomorphic to $X_2$ off $C$, and satisfies
\[
                 \mathcal O_{Z_2}(C)|_C=L^{-1}.
\]
If $Q\in\operatorname{Pic}^0(C)$ pulls back to
zero on $X_2$, its pullback to $Z_2$ has a rational
trivialization with divisor supported on the single
prime divisor $C$. Thus it is $\mathcal O_{Z_2}(bC)$
for an integer $b$. Restriction to $C$ gives
$Q\simeq L^{-b}$. Degrees give $0=-4b$, hence
$b=0$ and $Q=0$. This proves
\eqref{eq:C3-low-dimension-Jacobian} on the strict
henselization itself.

An explicit infinite-order class makes the
rational-coefficient failure just as definite.
Choose $u\in\mathbb C$ transcendental over
$\overline{\mathbb Q}$, choose $v$ with
$v^4=-1-u^4$, put $p=(u:v:1)$, and fix
$c_0=(1:\zeta:0)$ with $\zeta^4=-1$. Then
\[
                         Q=\mathcal O_C(p-c_0)
\]
has infinite order. Indeed $C$, $c_0$, its
Jacobian, and the Abel closed immersion are
defined over $\overline{\mathbb Q}$. Every
torsion point of that Jacobian over $\mathbb C$
is already defined over $\overline{\mathbb Q}$,
since each torsion subgroup scheme is finite
over this algebraically closed field. A torsion
$Q$ would therefore force $p$ to be defined
over $\overline{\mathbb Q}$, contrary to the
choice of $u$. The Abel closed immersion and
the finite torsion property are the standard
Jacobian facts recorded in
\citet[Propositions 2.1--2.3]{D_MilneJacobians2021}
and \citet[Section 12]{D_MilneAbelian2022}.

\emph{Failure of injectivity for $N=2$.}
Take $f=z_2\in R_2$. The divisor puncture
\[
                         Y_2=X_2\cap V(z_2)
\]
is regular and zero dimensional: before
henselization its four branches are the
distinct lines $z_0^4+z_1^4=0$, meeting only
at the removed vertex. Equivalently, the
finite normalization of $R_2/(z_2)$ is a
product of four strictly henselian discrete
valuation rings, and their punctures are
four spectra of fields. Hence
$\operatorname{Pic}(Y_2)=0$. The infinite-order
class $\rho^*Q$ gives a nonzero element of
$J^1(X_2)$ restricting to zero in $J^1(Y_2)$.

\emph{Failure of surjectivity for $N=3$.}
Now take
\[
 R_3=\mathbb C[z_0,z_1,z_2]_{(z_0,z_1,z_2)}^{\mathrm{sh}},
 \qquad X_3=\operatorname{Spec}R_3\setminus\{\mathfrak m_3\},
 \qquad f=F.
\]
The source is a regular local ring, so its
puncture has trivial Picard group
\citep[Lemma 52.27.5, Tag 0F2H]{D_StacksRegularPuncture}.
Strict henselization commutes with this quotient,
and therefore $Y_3=X_3\cap V(F)$ is exactly $X_2$.
The nonzero class $\rho^*Q\in J^1(Y_3)$ has no
preimage from $J^1(X_3)=0$.

Finally, for $N=1$ a normal local domain is a
discrete valuation ring: its puncture is the
spectrum of its fraction field, and the punctured
divisor is empty. Both Picard groups vanish.
For $N=2$, every regular punctured divisor is
a zero-dimensional noetherian scheme and hence
a finite disjoint union of spectra of fields.
Its Picard group is zero, so degree-one
surjectivity is automatic. Together with C.3m--n,
these facts give exactly the dimension claims
stated above.
\end{proof}

These are statements over $\mathbb C$, not a
claim that each failure persists over every
fixed positive-characteristic field. For example,
over $k=\overline{\mathbb F}_p$ every point of an
abelian variety is torsion. The actual
abelian-parameter construction in step 1 of C.3n
then shows that ordinary $\operatorname{Pic}_{\mathrm{alg}}$
of any connected regular noetherian $k$-scheme
is torsion: every class lies in the image of
the point group of an abelian variety. Thus
the rational ordinary $J^1$ groups in this
setting vanish. The characteristic-zero
counterexamples above must not be used to
infer the corresponding nonvanishing over
$\overline{\mathbb F}_p$.

The same torsion conclusion holds for ordinary
$\operatorname{CH}^r_{\mathrm{alg}}$ on these
excellent regular punctures in every codimension.
Here we use Chow groups for schemes of finite
type over the universally catenary base $X$
with its codimension dimension function, so
$X$ need not be of finite type over $k$;
this is precisely the setup of
\citet[Situation 42.7.1 and Lemma 42.20.3]{D_StacksChowProper}.
For a generator represented by a cycle family
over a smooth connected curve $T$ and points
$t_0,t_1\in T(k)$, complete $T$ to a smooth
projective curve $C$ and close each horizontal
integral component in $X\times C$. The closure
is still dominant, hence flat, over the regular
curve. Since $[t_1]-[t_0]$ is torsion in
$\operatorname{Pic}^0(C)(\overline{\mathbb F}_p)$,
there are $a>0$ and $u\in k(C)^\times$ with
\[
                 \operatorname{div}_C(u)
                         =a([t_1]-[t_0]).
\]
On each horizontal component, the pulled-back
rational function has divisor equal to $a$
times the difference of its two fibre cycles,
with their Cartier-divisor multiplicities.
The projection of that closure to $X$ is
proper and of finite type, so its pushforward
is rationally equivalent to zero by the cited
lemma. Thus $a$ kills the given fibre difference.
Finite sums are torsion as well. This proves
$\operatorname{CH}^r_{\mathrm{alg}}(X)\otimes\mathbb Q=0$
in this ordinary curve-family convention,
without asserting any relation to the
identity component of a local cycle functor.

\paragraph{C.3p (Known with derivation): local etale Lefschetz, Picard torsion, and completion.}
Let $(R,\mathfrak m)$ be an excellent normal strictly
henselian local ring, let $0\ne f\in\mathfrak m$,
and write
\[
 X=\operatorname{Spec}R\setminus\{\mathfrak m\},
 \qquad Y=X\cap V(f),\qquad N=\dim R.
\]
Assume $X$ regular. Restriction of finite etale covers
is fully faithful for $N\ge3$ and is an equivalence
for $N\ge4$:
\begin{equation}\label{eq:C3-etale-Lefschetz}
 \operatorname{FEt}(X)\xrightarrow{\ \sim\ }
 \operatorname{FEt}(Y)\qquad(N\ge4).
\end{equation}
If also $Y$ is regular, then, for every positive
integer $n$ invertible in $R$,
\begin{equation}\label{eq:C3-Picard-torsion-Lefschetz}
 \operatorname{Pic}(X)[n]\xrightarrow{\ \sim\ }
 \operatorname{Pic}(Y)[n]\qquad(N\ge4).
\end{equation}
This assertion includes mixed characteristic and
has no coefficient-field or algebraizability
hypothesis.

Separately, in the coefficient-field setting of
C.3n, completion induces a canonical isomorphism
of the represented local Picard functors
\begin{equation}\label{eq:C3-Picard-completion}
 \mathbf{Pic}^{\mathrm{loc}}_{R/k}
       \xrightarrow{\ \sim\ }
 \mathbf{Pic}^{\mathrm{loc}}_{\widehat R/k}
                     \qquad(N\ge3).
\end{equation}
The functors in this last assertion use
henselized tensor products as defined in C.3n.
It does not identify their arbitrary parameter
families with line bundles on the fixed product
$X\times T$.

\begin{proof}[Proof of the finite-etale and torsion assertions]
Put $A=\widehat R$ and use hats for the base changes
of $X,Y$. These are their punctured completions:
the closed fibre of $\operatorname{Spec}A\to
\operatorname{Spec}R$ has just the one closed
point. Excellence makes $A$ a normal complete
local domain of dimension $N$. It also makes
$R\to A$ regular. Thus $\widehat X$ is regular:
each point above the puncture is obtained by
localizing a regular morphism over a regular
local ring. The completed ring has a dualizing
complex.

Elkik's exact input is
\citet[Section III, Theorem 5, pp.~577--579, and
its corollary on p.~579]{D_Elkik1973}.
For a noetherian henselian pair $(B,I)$, completion
induces an equivalence of the categories of
finite etale covers of
$\operatorname{Spec}B\setminus V(I)$ and its
completed counterpart. Apply this both to
$(R,\mathfrak m)$ and to $(R/(f),\mathfrak m/(f))$.
The resulting vertical functors in
\[
\begin{array}{ccc}
 \operatorname{FEt}(X)&\longrightarrow&\operatorname{FEt}(Y)\\
 \big\downarrow\rlap{$\scriptstyle\sim$}&&
 \big\downarrow\rlap{$\scriptstyle\sim$}\\
 \operatorname{FEt}(\widehat X)&\longrightarrow&
 \operatorname{FEt}(\widehat Y)
\end{array}
\]
are equivalences. The hypotheses here concern
noetherian henselian pairs; no principal-ideal
restriction from Elkik's later Theorem 7 is
being added or silently removed.

It remains to check the complete-ring restriction.
The ring $A$ is $f$-adically complete. Indeed
an $f$-adically Cauchy sequence is
$\mathfrak m_A$-adically Cauchy, and each
ideal $(f^r)$ is closed for the latter topology
by the Krull intersection theorem applied to
$A/(f^r)$. Its limit therefore has the required
residues modulo every $f^r$.
The ring $A_f$ is regular and hence satisfies
$(S_2)$. For $N\ge3$, all the hypotheses of
\citet[Lemma 58.19.5, Tag 0EK5]{D_StacksEtaleLefschetz}
hold, giving full faithfulness.

For essential surjectivity with $N\ge4$, apply
\citet[Lemma 58.24.1, Tag 0EK9]{D_StacksEtaleLefschetz}.
Its additional condition is purity at every
maximal ideal of $A_f$. Here is the dimension
check for that condition. The punctured spectrum
of a noetherian local ring is Jacobson
\citep[Lemma 28.6.4, Tag 02IM]{D_StacksEtaleLefschetz}.
Hence a point closed in the open $\operatorname{Spec}A_f$
is closed in the puncture. If its prime in $A$
is $\mathfrak p$, this says
$\dim A/\mathfrak p=1$. The complete local domain
$A$ is catenary and equidimensional, so
\[
                \dim A_{\mathfrak p}=N-1\ge3.
\]
This is a regular local ring of dimension at
least two. Purity therefore holds by
\citet[Lemma 58.21.3(1), Tag 0BMA]{D_StacksEtaleLefschetz}.
This is the categorical purity clause (1);
no assertion about finite algebras with extra
components supported at the closed point is
needed. Lemma 58.24.1 now proves essential
surjectivity. Together with full faithfulness
and the Elkik comparisons it proves
\eqref{eq:C3-etale-Lefschetz}.

For \eqref{eq:C3-Picard-torsion-Lefschetz}, use
the Kummer exact sequence with $n$ invertible.
The global units on $X$ are $R^\times$ by
normal Hartogs extension, and every such unit
has an $n$th root in $R$ by strict henselianity.
For $Y$, use the finite normalization from
C.3k(i): the normal strictly henselian factors
$C_a$ have dimension $N-1\ge3$, their punctures
are the disjoint components of $Y$, and
\[
                 \Gamma(Y,\mathcal O_Y)=\prod_a C_a.
\]
Their units likewise have $n$th roots.
Consequently Kummer gives natural isomorphisms
\[
 H^1_{\mathrm{et}}(X,\mu_n)=\operatorname{Pic}(X)[n],
 \qquad
 H^1_{\mathrm{et}}(Y,\mu_n)=\operatorname{Pic}(Y)[n].
\]
The finite-etale equivalence preserves torsors
under the finite etale group $\mu_n$ and their
group actions. It thus gives the required
isomorphism. In dimension three the same
argument with full faithfulness gives the
injectivity already proved independently in
C.3k(i).
\end{proof}

\begin{proof}[Separate proof of the local Picard functor comparison]
Let $B$ be a finite-type $k$-algebra and form the
two henselian pairs
\[
 S_B=(R\otimes_k B)^h_{\mathfrak m},\qquad
 T_B=(\widehat R\otimes_k B)^h_{\mathfrak m\widehat R}.
\]
Both rings are noetherian: their unhenselized
rings are finite type over noetherian rings,
and henselization of a noetherian pair is
noetherian. The natural map $S_B\to T_B$
induces the canonical identification of their
completions along the displayed ideals,
\begin{equation}\label{eq:C3-common-tensor-completion}
 \widehat S_B=\widehat T_B
       =\varprojlim_r (R/\mathfrak m^r)\otimes_k B.
\end{equation}
In particular the completed complements of
these ideals are the same scheme.

The exact source now is
\citet[Section III, Theorem 3, p.~569, and its
unnumbered Picard corollary on p.~573]{D_Elkik1973}:
for every noetherian henselian pair, the Picard
group of the ideal complement maps bijectively
to that of its completion. Apply it to $S_B$
and $T_B$. Both Picard groups map bijectively
to the Picard group of the common completed
complement in \eqref{eq:C3-common-tensor-completion}.
The natural map between them is therefore
a bijection. These maps are functorial in $B$,
so the two local Picard presheaves, and then
their etale sheafifications, agree on finite-type
$k$-algebras.

Boutot's representability theorem applies to
both $R$ and $\widehat R$ in the hypotheses
of C.3n. The resulting group schemes are
locally of finite type over $k$. The canonical
morphism between them is an isomorphism
because it is a natural bijection on every
finite-type $k$-algebra: the affine finite-type
open subschemes of these group schemes test
the inverse by Yoneda, and the inverses glue
by their uniqueness. Equivalently one may
extend from finite-type parameter algebras
by their filtered-colimit property.
This proves \eqref{eq:C3-Picard-completion}.

Every comparison in this proof occurs after
henselizing the tensor pair. None removes
that henselization from an arbitrary family
on a proper parameter scheme. Therefore
\eqref{eq:C3-Picard-completion}, even combined
with \eqref{eq:C3-Picard-torsion-Lefschetz},
does not by itself prove surjectivity for
ordinary algebraic equivalence on fixed
products. This precise remaining distinction
is compatible with the characteristic-zero
counterexample C.3m.
\end{proof}

\paragraph{C.3q (Proved under the stated geometric hypotheses): lifting abelian subgroups of the local Picard functor in characteristic $p$.}
Let $k$ be algebraically closed of characteristic
$p>0$. Let $R$ be the strict henselization at a
$k$-rational point of a normal local scheme
essentially of finite type over $k$. Suppose
$N=\dim R\ge4$ and that
\[
 X=\operatorname{Spec}R\setminus\{\mathfrak m\},
 \qquad Y=X\cap V(f),\qquad0\ne f\in\mathfrak m,
\]
are regular. Put
\[
 G=\mathbf{Pic}^{\mathrm{loc}}_{R/k},\qquad
 H=\prod_a\mathbf{Pic}^{\mathrm{loc}}_{C_a/k},
 \qquad r:G\longrightarrow H,
\]
where $\prod_a C_a$ is the finite normalization
of $R/(f)$ and $r$ is the restriction morphism
of local Picard group schemes. These group
schemes are locally of finite type, since
$\dim C_a=N-1\ge3$; as in C.3n,
$G(k)=\operatorname{Pic}(X)$ and
$H(k)=\operatorname{Pic}(Y)$.

For every closed abelian subvariety $B\subset H$
there exist a closed abelian subvariety $A\subset G$,
an integer $e\ge0$, and isogenies of abelian
varieties over $k$
\begin{equation}\label{eq:C3-abelian-local-lift}
 s:B\longrightarrow A,\qquad
 r_A:A\longrightarrow B,\qquad
 r_A\circ s=[p^e]_B,
\end{equation}
where $r_A$ is induced by $r$. Both isogenies
have $p$-primary kernels. Thus every target
abelian subgroup lifts up to a power-of-$p$
isogeny at the level of the specified local
Picard functors.

\begin{proof}
We state separately the two geometric inputs.
First, in the present normal essentially-finite-type
scope, the local Neron--Severi group
\[
                  \operatorname{NS}^{\mathrm{loc}}
                         =G/G^0
\]
has finitely generated point group. This is
the positive-characteristic normal case stated
in \citet[Definition 3.4 and the following
paragraph, p.~468]{D_Kollar2016LocalPicard};
the cited argument combines Boutot's method
with alterations. We only need the weaker
consequence that its prime-to-$p$ torsion
has bounded exponent. Second,
\citet[Theorem 1.7(ii), p.~463, with Definition
3.3]{D_Kollar2016LocalPicard} makes
$\ker(r)$ a unipotent algebraic group.
The target in that theorem is interpreted by
its closed unit section; replacing the divisor
by its finite normalization does not change
the kernel, since its puncture is the same
$Y$. Henselization does not change the local
Picard functor, by its definition using henselian
tensor pairs. If $f$ is only given in $R$, it
descends to a pointed etale neighbourhood of
the original algebraic model. The cited
theorem applies there, giving the same local
functor and restriction map. These inputs are
not asserted here for an arbitrary excellent
ring outside that essentially-finite-type
geometric scope.

Choose a prime $\ell\ne p$ not dividing the
order of the finite prime-to-$p$ torsion
subgroup of $\operatorname{NS}^{\mathrm{loc}}(k)$.
By C.3p every point of $B[\ell^n](k)$ has a
unique preimage in $G(k)[\ell^n]$. Its image
in $\operatorname{NS}^{\mathrm{loc}}(k)$ is zero,
so the preimage belongs to $G^0(k)$. Since
$k$ is perfect, the reduced identity component
\[
                          G_0=(G^0)_{\mathrm{red}}
\]
is a smooth connected commutative algebraic
group with $G_0(k)=G^0(k)$.
The image $I=r(G_0)$ is a closed connected
algebraic subgroup of $H$. It contains all
$\ell$-power torsion points of $B$.
Their union is Zariski dense in $B$: otherwise
its closed subgroup closure would have smaller
identity-component dimension, contradicting
the $\ell^{2n\dim B}$ torsion points in
$B[\ell^n]$ and the corresponding growth in
that smaller-dimensional subgroup. Thus
\[
                              B\subset I.
\]

Let $E$ be the reduced identity component
of $r^{-1}(B)\cap G_0$. Then $E\to B$ is
surjective. Indeed the whole preimage maps
onto $B$, and the images of its finitely
many components are translates of the
closed connected image of $E$; a finite
union of proper closed subvarieties cannot
cover the irreducible variety $B$. A
surjective homomorphism of algebraic groups
is faithfully flat; here this also follows
directly from generic flatness followed
by translations. Write its kernel as $U$.
It is a closed subgroup scheme of the
unipotent group $\ker(r)$, hence is
unipotent and of finite type.

There is an integer $e$ such that
$[p^e]_U=0$ as a morphism of group schemes.
For example, a finite-type unipotent group
over the algebraically closed field embeds
in a group of upper unitriangular matrices;
if the matrices have size $d$ and $p^e\ge d$,
then $(1+M)^{p^e}=1+M^{p^e}=1$ in every
coefficient algebra. This argument includes
nonreduced kernels. The commutativity of $E$
now makes $[p^e]_E$ constant on the $U$-cosets.
Faithfully flat descent through $E\to B=E/U$
therefore gives a unique homomorphism
\[
                         \widetilde s:B\longrightarrow E
\]
with $\widetilde s\circ r_E=[p^e]_E$.
Composing with $r_E$ and using its faithful
flatness yields
\[
                            r_E\circ\widetilde s=[p^e]_B.
\]

Let $A$ be the scheme-theoretic image of
$\widetilde s$ in $E\subset G$. It is
proper, connected and reduced, since $B$
has those properties. A reduced algebraic
group over the perfect field $k$ is smooth,
so $A$ is an abelian variety. The kernel
of $\widetilde s$ is contained in
$B[p^e]$, hence is finite and $p$-primary.
Thus $s:B\to A$ is an isogeny. Its
composition with $r_A$ is $[p^e]_B$, so
$r_A$ is also surjective with finite
$p$-primary kernel, proving
\eqref{eq:C3-abelian-local-lift}.
\end{proof}

The exact ordinary-equivalence consequence
is conditional on an additional family
assertion. If the points of the resulting
$A$ are represented by ordinary algebraic
equivalences on the fixed product $X$ times
curves, then $r_A$ makes the corresponding
target classes lift in ordinary
$\operatorname{Pic}_{\mathrm{alg}}$, and hence
in $J^1$ with rational coefficients.
The theorem itself gives a morphism into
the henselized-product local Picard functor;
that alone is not this additional family
assertion. In particular the proof has not
used a universal actual line bundle on
$X\times A$, a blanket henselian GAGA
essential-surjectivity statement, or an
uncontrolled norm from an etale neighbourhood.
The further ordinary-equivalence step is
the realization of the lifted points by
actual product families. C.3u supplies
a direct torsion argument for smooth-cone
sources in every characteristic. C.3w
settles point realization for the present
algebraizable scope using alterations,
proper curve families, and transfer.
It consequently proves ordinary
surjectivity in the characteristic-$p$
scope of this paragraph, while retaining
the distinction between point realization
and a universal actual family on the
whole local Picard subgroup.

% Root construction; independent correspondence audit by /root/cycles.
\paragraph{C.3r (Refuted): ordinary higher-cycle restriction at every boundary dimension $N=2r+1$.}
For every integer $r\ge1$, we give an excellent normal strictly
henselian local $\mathbb C$-algebra of dimension $N=2r+1$, a nonzero
element $f$ of its maximal ideal, and regular punctures $X$ and
$Y=X\cap V(f)$ for which
\[
 J^r(X)=0,\qquad J^r(Y)\ne0.
\]
Here $J^r$ means the subgroup of ordinary algebraically trivial
codimension-$r$ Chow classes, tensored with $\mathbb Q$; the family
in the construction is an actual family on the fixed puncture.
In particular, restriction $J^r(X)\to J^r(Y)$ is not surjective.
The $r=1$ case has the same dimension as C.3o; the higher cases
give explicit failures beginning with $(N,r)=(5,2)$.

\begin{proof}
Put
\[
 V=\left\{\sum_{i=0}^{2r}z_i^3=0\right\}\subset\mathbb P^{2r}_{\mathbb C},
 \qquad H=\mathcal O_V(1),
 \qquad E=\{a^3+b^3+c^3=0\}\subset\mathbb P^2_{\mathbb C}.
\]
The smooth cubic $V$ has dimension $2r-1$, and $E$ is embedded in
its last three coordinates. Fix a primitive cube root of unity
$\zeta$. For $\lambda\in\{1,\zeta\}$, let $R_\lambda\subset V$
be cut out by
\[
 z_{2j}+\lambda z_{2j+1}=0\qquad(0\le j\le r-2).
\]
For $r\ge2$, this is the join of $E$ with the linear vertex
$\mathbb P^{r-2}$ in the first $2r-2$ coordinates. Its standard
resolution is the projective bundle with fibres the linear
$\mathbb P^{r-1}$ spanned by that vertex and a point of $E$.
Let $Z_\lambda\subset E\times V$ be this incidence projective
bundle, and put
\[
 \Gamma_\lambda=[Z_\lambda]\in\mathrm{CH}^r(E\times V).
\]
Equivalently, over $[a:b:c]\in E$, its points are
\[
 [-\lambda x_0:x_0:\cdots:-\lambda x_{r-2}:x_{r-2}:ta:tb:tc].
\]
This notation is homogeneous in $[a:b:c]$: it describes the
projective bundle, and does not choose global affine coordinates
on $E$. For $r=1$, take $V=E$ and $\Gamma_1=\Gamma_\zeta=\Delta_E$.

We have the integral correspondence identity
\begin{equation}\label{eq:C3-fermat-left-inverse}
             {}^t\Gamma_\zeta\circ\Gamma_1=\Delta_E
                         \quad\text{in }\mathrm{CH}^1(E\times E).
\end{equation}
Indeed the two linear equations in each paired set of coordinates
force those two coordinates to vanish, since $\zeta-1\ne0$.
Consequently $R_1\cap R_\zeta=E$ scheme-theoretically, and this
intersection avoids both vertices. At a point of $E$, the tangent
space of $V$ is the direct sum of $T_E$ and the $2r-2$ coordinate
directions belonging to the pairs. The two tangent spaces of the
joins use complementary $(r-1)$-dimensional subspaces of those
directions. Thus they meet transversely. The fibre product of
the two incidence resolutions over $V$ is exactly the diagonal
copy of $E$, with intersection multiplicity one. The composition
formula for Chow correspondences proves
\eqref{eq:C3-fermat-left-inverse}.

It follows that $(\Gamma_1)_*$ injects
$\operatorname{Pic}^0(E)(\mathbb C)$ into
$\mathrm{CH}^r(V)_{\mathrm{alg}}$. We next check that no nonzero
rational degree-zero class in this image is killed by passage to
$\mathrm{CH}^r(V)_{\mathbb Q}/H\mathrm{CH}^{r-1}(V)_{\mathbb Q}$.
For $r\ge2$, weak Lefschetz gives
\[
 H^{2r-3}(V,\mathbb Q)=0,\qquad
 H^{2r-2}(V,\mathbb Q)=\mathbb Q\,\operatorname{cl}(H^{r-1}).
\]
Write $b=cH^{r-1}+b_0$ for any
$b\in\mathrm{CH}^{r-1}(V)_{\mathbb Q}$, with $b_0$ homologically
trivial and $c\in\mathbb Q$. The intermediate Jacobian governing
the Abel--Jacobi class of $b_0$ is zero, by the first displayed
vanishing. Functoriality of Abel--Jacobi under multiplication by
$H$ and under $ {}^t\Gamma_\zeta$ therefore shows
\[
 ({}^t\Gamma_\zeta)_*(H b_0)=0
                            \quad\text{in }\operatorname{Pic}(E)_{\mathbb Q}.
\]
Here the resulting divisor is homologically trivial, and its
Abel--Jacobi map is the identity on $\operatorname{Pic}^0(E)$;
thus the conclusion is an equality of divisor classes, not merely
an equality of their cycle classes. Moreover
\[
 \deg\bigl(({}^t\Gamma_\zeta)_*(H^r)\bigr)
       =\int_{R_\zeta}H^r=3,
\]
because $R_\zeta$ is a linear section of the cubic $V$.
If $(\Gamma_1)_*\alpha=Hb$ for a degree-zero
$\alpha\in\operatorname{Pic}^0(E)_{\mathbb Q}$, applying
$ {}^t\Gamma_\zeta$ yields
$\alpha=c({}^t\Gamma_\zeta)_*H^r$. Degrees give $3c=0$,
and then $\alpha=0$. For $r=1$, the same conclusion follows
directly from $\deg H=3$ and $\mathrm{CH}^0(E)=\mathbb Z[E]$.

Choose $u\in\mathbb C$ transcendental over $\overline{\mathbb Q}$,
choose $v$ with $v^3=-1-u^3$, and put
\[
 p=(u:v:1),\quad o=(1:-1:0),\quad
                    \alpha=[p]-[o]\in\operatorname{Pic}^0(E).
\]
This class is not torsion. The curve, origin and Abel isomorphism
$E\simeq\operatorname{Pic}^0(E)$ are defined over
$\overline{\mathbb Q}$, and all torsion points of that abelian
variety are defined over $\overline{\mathbb Q}$. Torsion of
$\alpha$ would contradict the choice of $u$.

Finally set
\[
 R=\mathbb C[z_0,\ldots,z_{2r}]_{(z_0,\ldots,z_{2r})}^{\mathrm{sh}},
 \quad f=\sum_{i=0}^{2r}z_i^3,
 \quad X=\operatorname{Spec}R\setminus\{\mathfrak m\},
 \quad Y=X\cap V(f).
\]
The source is regular of dimension $2r+1$. The divisor is the
strictly henselian local cone on the smooth projectively normal
cubic $V$, and is regular away from its vertex. Strict
henselization commutes with the indicated quotient. Thus all
local hypotheses are verified. C.3e injects the nonzero quotient
class of $(\Gamma_1)_*\alpha$ into $\mathrm{CH}^r(Y)_{\mathbb Q}$.
It is ordinarily algebraically trivial, because the displayed
incidence family on $E\times V$ pulls back to an actual family
on $E\times Y$.

For the source, \citet[Corollary~6.4(i), author p.~27]
{rootSkalit2017} applies directly to the geometrically regular
map $\mathbb C\to R$ with henselian target. Its base has
dimension zero, so it gives $\mathrm{CH}^r(R)=0$ for every
$r>0$, integrally. The corollary derives this from exactness of
the Gersten complex: its degree-$r$ cohomology is the group of
codimension-$r$ cycles modulo principal divisors. Chow localization then gives
$\mathrm{CH}^r(X)=0$, proving the assertion.
\end{proof}

The exact symbolic checks in
\texttt{code/check\_fermat\_boundary.py} verify the cubic
substitution, the complementary paired-coordinate matrices and
the non-torsion witness equations. The transversality,
correspondence identity and non-torsion conclusions are proved
above, rather than inferred from finite examples. The construction
uses the familiar joins in Fermat hypersurfaces; we make no
priority claim. It does not assert the corresponding failure
for every larger $N$ at fixed $r\ge2$.

% Root construction; independent quotient audit by /root/frobenius.
\paragraph{C.3s (Refuted in every available positive codimension): universal local Chow and ordinary-$J$ finiteness.}
Let $N\ge2$ and $1\le r\le N-1$ be arbitrary integers.
There exists an explicitly specified excellent normal strictly
henselian local $\mathbb C$-algebra of dimension $N$, with regular
puncture $U$, such that
\[
       \mathrm{CH}^r(U)_{\mathrm{alg}}\otimes\mathbb Q\ne0.
\]
Consequently neither $\mathrm{CH}^r(U)$ nor its ordinary
algebraically trivial subgroup is universally finite in any
available positive codimension. For $r\ge N$ these groups
are zero for every such puncture, since it has no point of
codimension $r$. Thus this is the exact universal
finiteness/vanishing range for these two groups over $\mathbb C$.
The groups modulo algebraic equivalence remain separately
governed by C.3g and its stated ranges.

\begin{proof}
Set $d=r+2$, $m=N-1-r\ge0$, and take the smooth Fermat
Calabi--Yau $r$-fold
\[
 T=\{x_0^d+\cdots+x_{r+1}^d=0\}\subset\mathbb P^{r+1}_{\mathbb C},
 \qquad H=\mathcal O_T(1),\qquad D=T\cap\{x_{r+1}=0\}.
\]
Choose algebraically independent
$u_1,\ldots,u_r\in\mathbb C$ over $k_0=\overline{\mathbb Q}$,
choose $w$ with $w^d=-1-\sum_{i=1}^r u_i^d$, and choose
$\eta\in k_0$ with $\eta^d=-1$. Define
\[
 P=(u_1:\cdots:u_r:w:1),\qquad
 a=(1:\eta:0:\cdots:0)\in D(k_0),\qquad
 \alpha=[P]-[a]\in\mathrm{CH}^r(T)_{\mathrm{alg}}.
\]
All degree-zero zero-cycles on a smooth connected projective
variety are algebraically trivial. More explicitly, a smooth
connected complete-intersection curve through $P$ and $a$
supplies an ordinary curve family when $r\ge2$; when $r=1$
we use $T$ itself.

We recall the precise non-support argument, including its
direction of transfer. The point $P$ is a geometric generic
point of the $k_0$-variety $T$. If a nonzero multiple of
$\alpha$ were rationally equivalent to a cycle supported on
$D$, its finite cycle-and-function data over $\mathbb C$
could be spread and specialized to $\overline{k_0(T)}$,
then descended to a finite extension of $k_0(T)$.
The norm to $k_0(T)$ would give a nonzero multiple of its
generic point supported on $D_{k_0(T)}$. Closing this relation
in $T\times T$ and applying localization would yield
\[
       b\Delta_T=\Theta_D+\Theta_Z\quad(b\ne0),
\]
with supports $T\times D$ and $Z\times T$ for a proper
divisor $Z\subset T$. These are the same finite-extension,
norm and closure operations justified in C.3f. By adjunction
$K_T=\mathcal O_T$, so $H^{r,0}(T)\ne0$. Both right-hand
correspondences annihilate this space: restriction to a
resolution of $Z$ kills an $r$-form, and pushforward from
$D$ cannot produce Hodge bidegree $(r,0)$. The left-hand
side acts as multiplication by $b$, a contradiction.
Every class in $H\mathrm{CH}^{r-1}(T)$ has a representative
supported on $D$. We have therefore proved
\begin{equation}\label{eq:C3-CY-quotient-nonzero}
  \alpha\ne0\quad\text{in}\quad
  \mathrm{CH}^r(T)_{\mathbb Q}/
                      H\mathrm{CH}^{r-1}(T)_{\mathbb Q}.
\end{equation}

Now let
\[
           V=\mathbb P^m\times T,\qquad
           L=\mathcal O_{\mathbb P^m}(1)\boxtimes H.
\]
For $m=0$, this means $V=T$ and $L=H$.
The Segre embedding is projectively normal and has section
ring generated in degree one: the two factor embeddings have
these properties, their section multiplication maps are
surjective in every degree, and the normal full section ring
of the smooth product is their diagonal Segre product.
Its affine cone is normal and is smooth away from its vertex.
We take the local ring at that vertex and then its strict
henselization; this specifies the asserted ring and puncture
$U$. Equivalently its affine equations are the usual rank-one
Segre minors together with the homogeneous equations obtained
from the displayed Fermat equation. Its dimension is
$\dim V+1=m+r+1=N$.

Write $h=c_1(\mathcal O_{\mathbb P^m}(1))$ and also write
$H$ for the pullback of $c_1(H)$ from $T$. The projective-bundle
formula gives
\[
 \mathrm{CH}^*(V)
       =\mathrm{CH}^*(T)[h]/(h^{m+1}),\qquad c_1(L)=h+H.
\]
Eliminating $h$ by $h=-H$ gives an isomorphism of graded groups
\begin{equation}\label{eq:C3-Segre-quotient}
 \frac{\mathrm{CH}^r(V)}{c_1(L)\mathrm{CH}^{r-1}(V)}
 \simeq
 \frac{\mathrm{CH}^r(T)}
       {H^{m+1}\mathrm{CH}^{r-m-1}(T)}.
\end{equation}
Negative Chow degrees in the denominator mean zero.
Because $m+1\ge1$, that denominator is contained in
$H\mathrm{CH}^{r-1}(T)$. Equation
\eqref{eq:C3-CY-quotient-nonzero} therefore shows that the
pullback of $\alpha$ to $V$ survives rationally in
\eqref{eq:C3-Segre-quotient}. The integral injection C.3e
sends it to a non-torsion class in $\mathrm{CH}^r(U)$.
Its ordinary algebraic triviality is preserved by the
displayed smooth pullbacks and strict-local base change;
the curve family above pulls back on the fixed $U$.
This proves the assertion for every pair $(N,r)$ in scope.
\end{proof}

The program \texttt{code/check\_segre\_nonvanishing.py} checks
the graded quotient elimination and the Fermat witness equations
in a finite grid of dimensions. The proof works for all
dimensions by the exact projective-bundle identity; its
non-torsion step is the differential-form argument, not a
computational extrapolation.

% Owned by /root/frobenius. Explicit extension of the local cone Gysin example.
\paragraph{C.3t (Refuted in every odd critical dimension): ordinary local Gysin vanishing.}
Let $r\ge2$ and
\[
                         r+1\le N\le2r-1.
\]
We construct an excellent normal strictly henselian local complex
ring of dimension $N$, a nonzero element $f$ of its maximal ideal,
and regular punctures $X$ and $Y=X\cap V(f)$, such that the Gysin map
\[
             J^{r-1}(Y)\longrightarrow J^r(X)
\]
has a non-torsion integral image. Here the $J$ groups mean ordinary
algebraically trivial Chow classes, with rational coefficients when
testing nonvanishing. In particular, taking $N=2r-1$ refutes the
Gysin-vanishing assertion on manuscript $14{:}107$ in every odd
critical dimension, with $i=r-1=(N-1)/2$. It also gives a non-torsion
kernel for the open restriction
$\mathrm{CH}^r(X)\to\mathrm{CH}^r(X\setminus Y)$.

\begin{proof}
Put $m=N-1-r$, so $0\le m\le r-2$, and $d=r+2$. We use
\[
 T=\left\{\sum_{j=0}^{r+1}z_j^d=0\right\}
          \subset\mathbb P^{r+1}_{\mathbb C},\qquad
 V=\mathbb P^m\times T,\qquad
 L=\mathcal O_{\mathbb P^m}(1)\boxtimes\mathcal O_T(1).
\]
Choose $u_0,\ldots,u_{r-1}$ algebraically independent over
$\overline{\mathbb Q}$, choose $u_r$ satisfying
$u_r^d=-1-\sum_{j=0}^{r-1}u_j^d$, and choose
$\eta\in\overline{\mathbb Q}$ with $\eta^d=-1$. Set
\[
       P=(u_0:\cdots:u_r:1),\qquad
       a=(1:\eta:0:\cdots:0),\qquad
       \alpha=[P]-[a]\in\mathrm{CH}^r(T).
\]
The generic-point and differential-form argument in C.3s shows that
$\alpha$ is nonzero in
$\mathrm{CH}^r(T)_{\mathbb Q}/H\mathrm{CH}^{r-1}(T)_{\mathbb Q}$,
where $H=c_1(\mathcal O_T(1))$. Its pullback to $V$ consequently
survives modulo $c_1(L)\mathrm{CH}^{r-1}(V)_{\mathbb Q}$ by the
projective-bundle quotient in C.3s.

We give a smooth divisor carrying a specific actual family whose
pushforward is this pullback. For $2\le j\le r$ put
\[
                       s_j=z_j-u_jz_{r+1},\qquad j_0=r-m\ge2.
\]
The common zero locus on $T$ of the $m+1$ sections
$s_{j_0},\ldots,s_r$ is the hypersurface
\[
 Z:\quad \sum_{j=0}^{j_0-1}z_j^d+Cz_{r+1}^d=0
                 \quad\text{in }\mathbb P^{j_0},\qquad
 C=1+\sum_{j=j_0}^{r}u_j^d
        =-\sum_{j=0}^{j_0-1}u_j^d\ne0.
\]
It is smooth, since all the displayed coefficients are nonzero
and the characteristic is zero. Its dimension is
$j_0-1=r-m-1\ge1$. In particular, the sections cut $Z$ transversely
on $T$: their differentials are linearly independent at every point
of $Z$.

Write $[x_0:\cdots:x_m]$ for the coordinates of $\mathbb P^m$ and
define the divisor
\begin{equation}\label{eq:C3-uniform-incidence-divisor}
       W=\left\{\sum_{j=j_0}^{r}x_{j-j_0}s_j=0\right\}\subset V.
\end{equation}
This is a member of $|L|$. Away from $Z$, its equation cuts a
hyperplane in the projective-space fibre; when $m=0$ that part is
empty. At a point above $Z$, the derivative along $T$ is the
nonzero linear combination $\sum x_{j-j_0}ds_j$ of independent
covectors. Thus $W$ is smooth everywhere, including above $Z$.
We also have $\mathbb P^m\times Z\subset W$.

For the parameter curve we use the smaller common zero locus
\[
       B=T\cap\{s_2=\cdots=s_r=0\}\subset Z.
\]
It has the explicit plane equation
\[
 z_0^d+z_1^d+C_Bz_{r+1}^d=0,\qquad
 C_B=1+\sum_{j=2}^{r}u_j^d=-u_0^d-u_1^d\ne0.
\]
Hence $B$ is a smooth connected projective plane curve containing
both $P$ and $a$. The incidence
\[
 \mathbb P^m\times B\longrightarrow B\times W,
               \qquad(x,b)\longmapsto(b,(x,b)),
\]
is an actual flat family of codimension-$(r-1)$ cycles on $W$.
Its two marked fibres give the ordinary algebraically trivial class
\[
 \beta=[\mathbb P^m\times\{P\}]
                   -[\mathbb P^m\times\{a\}]
                         \in\mathrm{CH}^{r-1}(W)_{\mathrm{alg}}.
\]
If $i:W\hookrightarrow V$ is the inclusion, then
\begin{equation}\label{eq:C3-uniform-gysin-projective}
                         i_*\beta=\operatorname{pr}_T^*\alpha.
\end{equation}
All coefficients here are integral and equal to one.

We now pass to the local construction. Let $C(V,L)$ be the affine
cone and $v$ its vertex, and put
\[
 R=\mathcal O_{C(V,L),v}^{\mathrm{sh}},\qquad
 X=\operatorname{Spec}R\setminus\{\mathfrak m\}.
\]
This is the normal Segre cone of C.3s, of dimension
$m+r+1=N$, and its puncture is regular. To specify its coordinate
ring explicitly, it is the diagonal Segre product
\[
 \bigoplus_{n\ge0}\mathbb C[x_0,\ldots,x_m]_n
   \otimes_{\mathbb C}
 \left(\mathbb C[z_0,\ldots,z_{r+1}]/
                         \left(\sum z_j^d\right)\right)_n.
\]
It is generated by $w_{ij}=x_i z_j$. The required local element is
the nonzero degree-one element
\begin{equation}\label{eq:C3-uniform-local-f}
             f=\sum_{j=j_0}^{r}
                 \bigl(w_{j-j_0,j}-u_jw_{j-j_0,r+1}\bigr).
\end{equation}
Before localizing, the divisor $f=0$ on the punctured cone is the
nonzero line bundle over $W$ obtained by restricting $L^{-1}$.
Since $W$ is smooth, this divisor is regular. Flat etale base change
and the strict-local limit preserve regularity, so
$Y=X\cap V(f)$ is regular as required. No normality assertion at the
vertex of the divisor cone is needed.

Pull back $\beta$ to $Y$ and call the resulting class $\beta_Y$.
The displayed family over $B$ pulls back to an actual flat family
on the fixed product $B\times Y$, so
$\beta_Y\in\mathrm{CH}^{r-1}(Y)_{\mathrm{alg}}$ in the ordinary
sense. Flat compatibility of proper pushforward in the Cartesian
square of the two punctured line bundles, and then at each etale
stage, transforms \eqref{eq:C3-uniform-gysin-projective} into
\[
                   (Y\hookrightarrow X)_*\beta_Y=\gamma,
\]
where $\gamma$ is the pullback of $\operatorname{pr}_T^*\alpha$.
The integral cone injection C.3e and the quotient nonvanishing
already proved show that $\gamma$ is non-torsion. This proves
the Gysin assertion, including its failure modulo finite groups.
Localization sends $\gamma$ to zero on $X\setminus Y$, proving
the final assertion about open restriction.
\end{proof}

For $N=2r-1$ we have $m=r-2$ and $Z=B$, giving the odd critical
case directly. The larger formal choice $m=r-1$ would make $Z$
zero-dimensional and would remove the curve family used in this
proof; we make no claim for that choice. The symbolic certificate
\texttt{agents/frobenius/code/check\_uniform\_gysin.py} checks the
Fermat substitutions, nonzero coefficient formulas, and derivative
rank in a finite grid. Smoothness and non-torsion in all dimensions
are established by the preceding proof. This is an explicit
extension of C.3f; we make no priority claim.

\paragraph{C.3u (Proved under the stated cone hypotheses): ordinary degree-one surjectivity in every characteristic.}
Let $k$ be an algebraically closed field of arbitrary
characteristic. Let $V$ be a smooth projective integral
$k$-variety of dimension $d\ge3$, and let $L$ be a
very ample invertible sheaf such that its full section
ring
\[
 S=\bigoplus_{a\ge0}H^0(V,L^a)
\]
is generated in degree one. Write $\mathfrak v=S_{>0}$,
and form the normal strictly henselian cone local ring
and its puncture
\[
 R=(S_{\mathfrak v})^{\mathrm{sh}},\qquad
 X=\operatorname{Spec}R\setminus\{\mathfrak m\},
 \qquad N=d+1\ge4.
\]
The geometric residue field in this strict
henselization is the given field $k$. Let
$0\ne f\in\mathfrak m$, and suppose
$Y=X\cap V(f)$ is regular. The element $f$ need
not be homogeneous. Let $\rho:X\to V$ be the
cone projection, and put
$A=(\operatorname{Pic}^0_{V/k})_{\mathrm{red}}$.
For ordinary algebraic equivalence defined by
actual line bundles on products with smooth
connected $k$-curves, the following statements hold:
\begin{align}
 A(k)&\xrightarrow[\rho^*]{\ \sim\ }
                    \operatorname{Pic}(X)_{\mathrm{alg}},
                    \label{eq:C3-cone-ordinary-source}\\
 \operatorname{Pic}(X)_{\mathrm{alg}}
       &\longrightarrow\operatorname{Pic}(Y)_{\mathrm{alg}}
                    \quad\text{is surjective}.
                    \label{eq:C3-cone-ordinary-surjectivity}
\end{align}
In characteristic zero the map in
\eqref{eq:C3-cone-ordinary-surjectivity} is an
isomorphism. In characteristic $p>0$ its kernel
is a finite $p$-primary group. Consequently
\[
                       J^1(X)\xrightarrow{\ \sim\ }J^1(Y)
\]
with rational coefficients in every characteristic.
These are statements about the specified smooth
cone source and arbitrary regular divisor punctures
inside it; they do not impose a cone structure
on the divisor or assert the same conclusion
for a deformation of the source cone.

\begin{proof}
\emph{1. The proper resolution and its actual
Picard family.}
The full section ring of the normal projective
variety is normal. Its degree-one generation
identifies the standard contraction
\[
 Z_0=\operatorname{Tot}(L^{-1})
          \longrightarrow\operatorname{Spec}S
\]
with the blow-up of the vertex. One can check
this on a degree-one homogeneous chart: the
blow-up chart is the affine-line bundle obtained
by retaining the radial coordinate over the
corresponding affine open of $\operatorname{Proj}S=V$.
Thus this contraction is proper, its exceptional
divisor is the zero section $V$, and it is an
isomorphism away from $V$. Base change gives
\[
 Z=Z_0\times_{\operatorname{Spec}S}\operatorname{Spec}R
                  \longrightarrow\operatorname{Spec}R.
\]
The scheme $Z$ is regular and integral. Indeed
it is the blow-up of the domain $R$, and its
regularity follows from that of $Z_0$ by the
localization and ind-etale base change defining
$R$. It is proper over the henselian ring $R$,
has the single prime exceptional divisor $V$,
and has reduced closed fibre exactly $V$.
Its complement is $X$, which is consequently
regular. The projection of the line bundle
induces $\pi:Z\to V$, extending $\rho$, with
\begin{equation}\label{eq:C3-cone-exceptional-restriction}
 \pi|_V=1_V,\qquad
             \mathcal O_Z(V)|_V=L^{-1}.
\end{equation}

The variety $A$ is an abelian variety. Moreover
there is a Poincare line bundle on $V\times A$,
normalized at a chosen $k$-point of $V$ and at
$0\in A$. Here the precise standard inputs are
\citet[Theorem 2.5(2), Exercise 4.3 with its answer,
and Theorem 5.4, author-version pp.~18, 27,
37--38, 65]{D_KleimanPicard2005v1}.
The section provided by the chosen point makes
the universal line bundle exist; the normality
of $V$ makes the reduced identity component
proper. Pulling this bundle back along $\rho$
gives an actual family on $X\times A$ and a
homomorphism
\[
           j:A\longrightarrow G,
           \qquad G=\mathbf{Pic}^{\mathrm{loc}}_{R/k}.
\]
The local Picard scheme is the separated group
scheme locally of finite type described in
C.3n, with $G(k)=\operatorname{Pic}(X)$.
Its closed reduced image $P=j(A)$ is a proper
connected algebraic subgroup, hence an abelian
variety, and $P(k)=j(A(k))$.

The map $j$ is injective on $k$-points.
In fact, if $Q\in A(k)$ has $\rho^*Q$ trivial,
the line bundle $\pi^*Q$ is trivial off $V$.
A rational trivialization therefore gives
\[
                         \pi^*Q\simeq\mathcal O_Z(bV)
                              \quad(b\in\mathbb Z).
\]
Restricting to $V$ gives $Q\simeq L^{-b}$.
An algebraically trivial line bundle has degree
zero on every curve, whereas a nonzero power
of the ample bundle $L$ does not. Thus $b=0$
and $Q$ is trivial. This argument asserts
injectivity on points; it does not assume
that $j$ has no infinitesimal kernel in
positive characteristic.

Every point of $A(k)$ gives an ordinary
algebraically trivial class on $X$. For example,
restrict the actual Poincare family to curves
joining the relevant points in the smooth
projective connected variety $A$, and pull
back to $X$. Such points can be joined by a
chain of integral curves, whose normalizations
are smooth projective over $k$. Conversely,
the equality between the points of
$\operatorname{Pic}^0_{V/k}$ and algebraically
trivial line bundles is
\citet[Proposition 5.10, pp.~41--42]{D_KleimanPicard2005v1}.

\emph{2. All source torsion at one suitable prime
comes from this actual family.}
The Neron--Severi group $\operatorname{NS}(V)$
is finitely generated in every characteristic;
see \citet[Theorem 7.3, p.~18]{D_ChengJiLarsonOlander2021}.
Choose a prime $\ell$ invertible in $k$ and not
dividing the order of the finite torsion subgroup
of
\[
                         \operatorname{NS}(V)/\mathbb Z[L].
\]
We claim that for every $n\ge1$,
\begin{equation}\label{eq:C3-cone-prime-torsion}
 A[\ell^n](k)\xrightarrow[\rho^*]{\ \sim\ }
                         \operatorname{Pic}(X)[\ell^n].
\end{equation}
Only surjectivity remains to be proved.

Let $M\in\operatorname{Pic}(X)[\ell^n]$.
The closure of the divisor of a rational
section extends $M$ to an invertible sheaf
$\widetilde M$ on the regular noetherian
integral scheme $Z$. The chosen triviality
of $M^{\ell^n}$ on $X$ gives
\[
                 \widetilde M^{\ell^n}\simeq\mathcal O_Z(bV)
                             \quad(b\in\mathbb Z).
\]
Put $Q=\widetilde M|_V$. By
\eqref{eq:C3-cone-exceptional-restriction},
$Q^{\ell^n}\simeq L^{-b}$. Hence the class
of $Q$ modulo $\mathbb Z[L]$ in
$\operatorname{NS}(V)$ is killed by $\ell^n$.
Our choice of $\ell$ makes that class zero.
Choose $a\in\mathbb Z$ with $[Q]=a[L]$,
and replace $\widetilde M$ by
$\widetilde M\otimes\mathcal O_Z(aV)$.
Its restriction, still denoted $Q$, now belongs
to $A(k)$. Its $\ell^n$th power is still of
the form $\mathcal O_Z(b'V)$.
Restricting this identity to $V$ shows
$\ell^n[Q]=-b'[L]$ in $\operatorname{NS}(V)$.
The left side is zero. Ampleness makes
$[L]$ of infinite order, so $b'=0$.
Consequently
\[
       \widetilde M^{\ell^n}\simeq\mathcal O_Z,
       \qquad Q^{\ell^n}\simeq\mathcal O_V.
\]

The bundle
$T=\widetilde M\otimes\pi^*Q^{-1}$ is
$\ell^n$-torsion and restricts trivially to
$V$. A trivialization of $T^{\ell^n}$ defines
a finite etale $\mu_{\ell^n}$-torsor over $Z$.
Its restriction to $V$ is trivial: after a
trivialization of $T|_V$, its defining unit
lies in $\Gamma(V,\mathcal O_V)^\times=k^\times$,
and has an $\ell^n$th root in $k$.
Now
\citet[Lemma 58.9.1, Tag 0A48]{D_StacksProperEtale}
gives an equivalence
\[
                  \operatorname{FEt}(Z)
                       \xrightarrow{\ \sim\ }
                  \operatorname{FEt}(V),
\]
because $Z$ is proper over the henselian
local ring $R$ and $V$ is its closed fibre.
Full faithfulness lifts a section of this
torsor over $V$ to a section over $Z$;
the torsor is therefore trivial. Hence
$T\simeq\mathcal O_Z$ and
$M\simeq\rho^*Q$, proving
\eqref{eq:C3-cone-prime-torsion}.
This uses the equivalence for finite etale
covers of the proper scheme $Z$, not a
claim that arbitrary line bundles or
parameter families algebraize from the
henselized product.

\emph{3. Ordinary source classes and target
surjectivity.}
We use the actual-family lemma proved in
C.3n: every ordinary algebraically trivial
class on the regular puncture lies in the
point group of a closed abelian subvariety
of its local Picard scheme. This lemma
starts with the given actual curve family
and constructs a dual-Jacobian family;
it does not infer an actual family from
an arbitrary subgroup of a local Picard
functor.

Apply it first to $X$, and let $B\subset G$
be such an abelian subvariety. Every point
of $B[\ell^n](k)$ is in $P(k)$ by
\eqref{eq:C3-cone-prime-torsion}.
The union of these torsion groups is Zariski
dense in $B$. Indeed its closed subgroup
closure, if of smaller dimension, could
have at most a fixed finite multiple of
$\ell^{2n\dim C}$ points of order dividing
$\ell^n$, where $C$ is its identity
component; this contradicts the
$\ell^{2n\dim B}$ such points of $B$.
Since $P$ is closed, $B\subset P$.
Together with the actual family in step 1
and pointwise injectivity of $j$, this
proves \eqref{eq:C3-cone-ordinary-source}.

For the target, write $\prod_a C_a$ for
the finite normalization of $R/(f)$ and
put
\[
 H=\prod_a\mathbf{Pic}^{\mathrm{loc}}_{C_a/k},
                  \qquad r:G\longrightarrow H.
\]
As in C.3n--p, the punctures of the
normal strictly henselian factors are the
disjoint components of $Y$. They have
dimension $N-1\ge3$, so $H$ is a
separated group scheme locally of finite
type and $H(k)=\operatorname{Pic}(Y)$.
The same actual-family lemma, applied to
the finitely many components, puts each
ordinary algebraically trivial class of
$Y$ in a closed abelian subvariety
$B\subset H$.

By C.3p, restriction gives isomorphisms
\[
 \operatorname{Pic}(X)[\ell^n]
       \xrightarrow{\ \sim\ }
 \operatorname{Pic}(Y)[\ell^n]
                      \qquad(n\ge1).
\]
Every point of $B[\ell^n](k)$ therefore
comes from $P(k)$ by step 2.
The image $r(P)$ is closed and proper,
since $P$ is proper and $H$ separated.
It contains the dense union of these
torsion groups, and hence contains $B$.
Every $k$-point of this image lifts to a
$k$-point of $P$: the nonempty fibre is
of finite type over the algebraically
closed field $k$. The points of $P$
are ordinary classes on $X$ by step 1.
This proves
\eqref{eq:C3-cone-ordinary-surjectivity}.

\emph{4. The integral kernel.}
The closed subgroup scheme
$K=\ker(r\circ j)\subset A$ is of
finite type. For every prime $\lambda$
invertible in $k$, C.3p and the pointwise
injectivity of $j$ imply
$K(k)[\lambda^n]=0$ for all $n\ge1$.
If $K$ had positive dimension, its reduced
identity component would be a
positive-dimensional abelian variety and
would have nonzero $\lambda$-torsion,
a contradiction. Thus $K$ is finite.
Its finite point group has no torsion
of order prime to the characteristic.
It is therefore zero in characteristic
zero and $p$-primary in characteristic
$p>0$. By
\eqref{eq:C3-cone-ordinary-source}, this
is exactly the ordinary restriction
kernel asserted above.
\end{proof}

The proof requires neither a splitting
of the whole local Picard identity
component nor the positive-characteristic
unipotent-kernel theorem used in C.3q.
The actual source family comes from
$\operatorname{Pic}^0(V)$, and proper
henselian finite-etale comparison detects
all the torsion needed for the density
argument. This closes the ordinary-family
issue for these cone sources in every
characteristic. C.3w addresses the more
general algebraizable characteristic-$p$
scope of C.3q by a separate alteration
and proper-curve construction.


% Root proof, independently checked by /root/cycles and /root/real_weil.
\paragraph{C.3v (Refuted already in dimension two): finiteness of the ordinary divisor quotient.}
For every $N\ge2$, there is an explicitly specified excellent
normal strictly henselian local $\mathbb C$-algebra $R_N$ of
dimension $N$, with regular puncture $X_N$, for which
\begin{equation}\label{eq:C3v-ordinary-quotient}
 B^1(X_N):=\frac{\operatorname{Pic}(X_N)}
                   {\operatorname{Pic}(X_N)_{\mathrm{alg}}}
 \simeq
 \begin{cases}
  (\mathbb C,+)\oplus\mathbb Z/4,&N=2,\\
  (\mathbb C,+)^{N-1}\oplus\mathbb Z,&N\ge3.
 \end{cases}
\end{equation}
Here algebraic equivalence is ordinary algebraic equivalence
through line bundles on the fixed product with a smooth
connected curve. In particular the rational quotient is
nonzero and infinite-dimensional over $\mathbb Q$, already
in dimension two. Thus a finiteness assertion for the
degree-one quotient on 14:107 is false on this reading.
The quotient using the connected local Picard functor is a
different relation and is not refuted by this calculation.

\begin{proof}
Let
\[
 C=\{z_0^4+z_1^4+z_2^4=0\}\subset\mathbb P^2_{\mathbb C},
 \qquad m=N-2,\qquad V=\mathbb P^m\times C,
 \qquad L=\mathcal O_{\mathbb P^m}(1)\boxtimes K_C.
\]
For $m=0$, the first factor and its line bundle mean a
point and its trivial line bundle. The partial derivatives
$4z_i^3$ have no common projective zero. Adjunction gives
$K_C=\mathcal O_C(1)$, $\deg K_C=4$ and $g(C)=3$.
The canonical embedding of this plane quartic and the
standard projective embedding are projectively normal.
Their Segre section ring
\[
 S=\bigoplus_{a\ge0}H^0(V,L^a),\qquad
 R_N=(S_{S_+})^{\mathrm{sh}}
\]
is therefore a normal domain generated in degree one.
We use the strict henselization with its specified residue
field $\mathbb C$; it equals the henselization. Excellence,
normality and dimension $\dim R_N=m+2=N$ follow from
this finite-type construction. The cone puncture is the
complement of the zero section in $\operatorname{Tot}(L^{-1})$,
and is smooth. Its strict-local base change $X_N$ is regular.

Let $Z\to\operatorname{Spec}R_N$ be the base change of the
blow-up of the vertex. It is regular and proper, with sole
exceptional divisor $V$, projection $\pi:Z\to V$, and
\[
 I_V=\mathfrak m_{R_N}\mathcal O_Z,
 \qquad \mathcal O_Z(V)=\pi^*L^{-1}.
\]
Degree-one generation proves the first identity. Completion
along $V$ is the formal completion of the line bundle along
its zero section. On its $n$th infinitesimal neighborhood,
the projection splits the units onto $\mathcal O_V^*$,
while the characteristic-zero logarithm identifies the
remaining units with the additive nilpotent ideal
\[
 I_V/I_V^n=\bigoplus_{a=1}^{n-1}L^a.
\]
The logarithm is a finite polynomial at each level, so
these are actual compatible isomorphisms of sheaves of
abelian groups, not merely tangent-space identifications.
Consequently
\[
 \operatorname{Pic}(V_n)
 \simeq\operatorname{Pic}(V)\oplus
                   \bigoplus_{a=1}^{n-1}H^1(V,L^a).
\]
For $a>0$, the K\"unneth formula and Serre duality on $C$
give
\begin{align*}
 H^1(V,L)&=H^0(\mathbb P^m,\mathcal O(1))
                         \otimes H^1(C,K_C)
                  \simeq\mathbb C^{m+1},\\
 H^1(V,L^a)&=0\qquad(a\ge2).
\end{align*}
Indeed $H^1(C,K_C^a)$ is dual to
$H^0(C,K_C^{1-a})$, whose degree is negative for $a\ge2$;
the higher cohomology of positive projective-space line
bundles contributes nothing in degree one.

There is no inverse-limit ambiguity here. Global units
on $V_n$ split as $\mathbb C^*$ times the finite direct
sum of the $H^0(V,L^a)$ under the logarithm. Their
transition maps are surjective, so the corresponding
$\varprojlim^1$ is zero. Equivalently, isomorphisms of
successive line bundles can be chosen compatibly by
these surjections. Thus the formal Picard group is
\[
 \operatorname{Pic}(\widehat Z)
       \simeq\operatorname{Pic}(V)\oplus\mathbb C^{m+1}.
\]
Grothendieck existence, in the precise form of
\citet[Theorem 30.27.1, Tag 088E]{rootStacksExistence},
algebraizes every such compatible line bundle to the
proper scheme $Z_{\widehat R_N}$ and is fully faithful.
The algebraized coherent sheaf is invertible: it is
invertible near the closed fibre by Nakayama's lemma,
and the image of its closed nonfree locus would be a
nonempty closed subset of the local base missing its
closed point, which properness excludes.

The excellent regular base changes used here keep
$Z_{\widehat R_N}$ and its puncture regular. Cartier
divisor localization therefore identifies
\[
 \operatorname{Pic}(X_{\widehat R_N})
   \simeq
   \bigl(\operatorname{Pic}(V)/\mathbb Z[L]\bigr)
                                      \oplus\mathbb C^{m+1}.
\]
The exceptional divisor contributes exactly $[L^{-1}]$
and has zero additive coordinate. Finally,
\citet[Section III, Theorem 3, p.~569, and its unnumbered
Picard corollary, p.~573]{D_Elkik1973}
gives a bijection
$\operatorname{Pic}(X_N)\to\operatorname{Pic}(X_{\widehat R_N})$.
Its hypotheses are the noetherian henselian pair
$(R_N,\mathfrak m_{R_N})$ and its completion; there is no
dimension-at-least-three assumption. We have proved the
same direct-sum formula for the actual punctured
henselian scheme $X_N$.

We now identify ordinary algebraic equivalence without
replacing it by formal-family equivalence. Suppose
$\mathcal M$ is an actual line bundle on $X_N\times T$,
where $T$ is a smooth connected curve. Compactify $T$
to a smooth proper connected curve $\overline T$.
The schemes $X_N\times\overline T$ and
$Z\times\overline T$ are regular and noetherian.
Closing a Cartier divisor first across
$X_N\times(\overline T\setminus T)$ and then across
$V\times\overline T$ extends $\mathcal M$ to a line
bundle on $Z\times\overline T$. Restrict it to
$V_2\times\overline T$. The split square-zero
description above is functorial in a parameter scheme;
it gives the Picard group scheme decomposition
\[
 \mathbf{Pic}_{V_2/\mathbb C}
   \simeq \mathbf{Pic}_{V/\mathbb C}
                       \times\mathbb G_a^{m+1}.
\]
The resulting map from the proper connected curve
$\overline T$ to the affine vector group is constant.
Its two fibre differences therefore have zero additive
coordinate. Their $\operatorname{Pic}(V)$ coordinates
differ by an element of $\operatorname{Pic}^0(V)(\mathbb C)$.
The ambiguity of extension by exceptional divisors has
zero additive coordinate and vanishes modulo $\mathbb Z[L]$.
It follows, for every actual family and hence every sum
of family differences, that
\[
 \operatorname{Pic}(X_N)_{\mathrm{alg}}
       \subseteq\rho^*\operatorname{Pic}^0(V)(\mathbb C),
\]
where $\rho:X_N\to V$ is the radial projection.
Conversely a Poincar\'e line bundle on
$V\times\operatorname{Pic}^0(V)$, pulled back by $\rho$,
gives these ordinary algebraic classes. Any two points
of the abelian parameter can be joined by a smooth
connected curve (or a chain of such curves). Also
$\mathbb Z[L]\cap\operatorname{Pic}^0(V)=0$, since a
nonzero power of an ample line bundle is not numerically
trivial. Thus this subgroup is an actual embedded copy
of $\operatorname{Pic}^0(C)(\mathbb C)$.

Taking the quotient gives
\[
 B^1(X_N)\simeq
       \bigl(\operatorname{NS}(V)/\mathbb Z[L]\bigr)
                                      \oplus\mathbb C^{m+1}.
\]
If $m=0$, $\operatorname{NS}(C)=\mathbb Z$ by degree and
$[K_C]$ has degree four. If $m\ge1$, the projective
bundle formula gives $\operatorname{NS}(V)=\mathbb Z^2$,
and $[L]=(1,4)$ is primitive, leaving one free copy
of $\mathbb Z$. This proves \eqref{eq:C3v-ordinary-quotient}.

For completeness, we specify one non-torsion divisor
class in the smallest counterexample, rather than only
its ambient group. Choose $\eta\in\mathbb C$ with
$\eta^4=-1$ and let $c_0=(1:\eta:0)\in C$.
The function $t=z_2/z_0$ is a uniformizer at $c_0$,
since the derivative of $1+(z_1/z_0)^4+t^4$ with respect
to $z_1/z_0$ is $4\eta^3\ne0$ there. Take
$U=C\setminus\{c_0\}$ and an affine neighborhood
$W\subset\{z_0\ne0\}$ of $c_0$ excluding the other
three zeros of $t$. The regular differential
$dt/t$ on $U\cap W$ represents a nonzero class
$\alpha\in H^1(C,K_C)$: its Serre-duality residue
pairing with $1\in H^0(C,\mathcal O_C)$ is one.
On the formal total space of $K_C^{-1}$, contract
$dt/t$ with the tautological radial vector to obtain
$a\in I_V(U\cap W)$. Glue the trivial formal line
bundles by the transition $\exp(a)$; this is a
compatible formal unit, and its additive Picard
coordinate is exactly $\alpha$. Proper formal existence
algebraizes it, and Elkik's bijection specifies a unique
line-bundle class $M$ on $X_2$. This construction fixes
its class independently of all auxiliary algebraization
choices. Its image in $B^1(X_2)$ has additive coordinate
$\alpha\ne0$, so every nonzero integer multiple remains
nonzero. A Cartier divisor representing $M$ is the
required ordinary cycle counterexample.
\end{proof}

The finite certificate
\texttt{code/check\_ordinary\_divisor\_quotient.py}
checks the smooth quartic equations, the chosen local
parameter, the residue coefficient, the indicated
cohomology dimensions in its stated finite input range,
and the two N\'eron--Severi lattice quotients. The
formal descent, the all-dimension formula and the
non-torsion conclusion follow from the proof, not
from extrapolation of that computation.


\paragraph{C.3w (Proved in the stated algebraizable scope): proper local Picard images and actual families via alterations.}
Let $k$ be algebraically closed. Let $(R,\mathfrak m,k)$
be the strict henselization at a $k$-rational point of
a normal local scheme essentially of finite type over
$k$. Assume $N=\dim R\ge3$ and that its puncture
$X=\operatorname{Spec}R\setminus\{\mathfrak m\}$ is
regular. Write
\[
                   G=\mathbf{Pic}^{\mathrm{loc}}_{R/k}
\]
for the henselized-tensor-product local Picard scheme
of C.3n. Then the following point-realization assertion holds:
\begin{equation}\label{eq:C3-proper-local-image-ordinary}
 B(k)\subset\operatorname{Pic}(X)_{\mathrm{alg}}
              \quad\text{for every closed abelian subvariety }B\subset G.
\end{equation}
This asserts ordinary algebraic equivalence of each
point, not the existence of a universal actual line
bundle on $X\times B$ representing the inclusion.

Consequently, if $\operatorname{char}k=p>0$,
C.3q closes the ordinary-surjectivity step:
for $N\ge4$,
$0\ne f\in\mathfrak m$, and regular
$Y=X\cap V(f)$, restriction
\begin{equation}\label{eq:C3-charp-ordinary-isomorphism}
 \operatorname{Pic}(X)_{\mathrm{alg}}
       \longrightarrow\operatorname{Pic}(Y)_{\mathrm{alg}}
\end{equation}
is surjective, with $p$-primary torsion kernel.
In particular $J^1(X)\to J^1(Y)$ is a rational
isomorphism. The algebraizability and specified
coefficient/residue field hypotheses in this
assertion are retained.

\begin{proof}[Proof of the point-realization assertion]
\emph{1. Fixing one regular alteration.}
Choose a normal affine algebraic model near the
specified point. The exact geometric input is
\citet[Theorem 4.1 and Remark 4.2, p.~66]{D_DeJong1996Alterations}:
it has a proper generically finite alteration
whose source is regular and quasiprojective.
Base change to $R$ preserves regularity, since
strict henselization is a localization and
ind-etale base change. Choose an integral
connected component dominating
$\operatorname{Spec}R$. We obtain
\[
             h:Z\longrightarrow\operatorname{Spec}R,
             \qquad d=[k(Z):\operatorname{Frac}R]>0,
\]
with $Z$ regular, integral, and proper over $R$.
No assertion that this alteration is finite or
is an isomorphism over the puncture is made.
Put $Z^\circ=h^{-1}(X)$ and $E=h^{-1}(\mathfrak m)$.

The Stein algebra
\[
                         R'=\Gamma(Z,\mathcal O_Z)
\]
is finite over $R$, is its normalization in
$k(Z)$, and is a normal domain. Indeed an
element integral over $R$ is regular in every
normal local ring of $Z$, and proper coherence
makes the global algebra finite. Since $R$ is
henselian, the finite domain $R'$ is local;
its residue field is $k$. Stein factorization
gives a proper birational map
$Z\to\operatorname{Spec}R'$ with connected
fibres. We use
\citet[Theorem 37.53.4]{D_StacksAlterationFamilies}
for this factorization. In particular $E$ is
connected. Choose a $k$-point $e\in E$.
Let $E_1,\ldots,E_s$ denote the prime divisors
of $Z$ contained in $E$; components of $E$
of larger codimension do not enter this list.
Each $E_i$ is Cartier because $Z$ is regular.

\emph{2. Extending a proper parameter curve as a
Picard-functor class.}
We use exactly the etale sheafification of
henselized-tensor-product Picard groups in
\citet[Introduction, Theorem 1 and preceding
definition, p.~VII]{D_Boutot1978}, as specified
in C.3n. In particular its sections have
representatives etale locally on the parameter.
Let $C$ be a smooth projective connected
$k$-curve, and let $g:C\to G$ be a morphism.
For an affine etale $U=\operatorname{Spec}B\to C$,
set
\[
 T_B=(R\otimes_kB)^h_{\mathfrak m},\quad
 Z_B^h=Z\times_R\operatorname{Spec}T_B,
 \quad (Z_B^h)^\circ=Z_B^h\setminus h^{-1}V(\mathfrak mT_B).
\]
All these rings and schemes are noetherian.
The scheme $Z_B^h$ is regular: $B/k$ is smooth,
so $R\to R\otimes_kB$ is smooth, and the
subsequent henselization is ind-etale.
Let $\mathcal F$ and $\mathcal H$ be the etale
sheaves on $C$ obtained from
\[
 U\longmapsto\operatorname{Pic}(Z_B^h),\qquad
 U\longmapsto\operatorname{Pic}((Z_B^h)^\circ).
\]
Pullback along $h$ maps the local Picard family
$g$ to a section $g'\in\mathcal H(C)$.

There is an exact sequence of sheaves on
$C_{\mathrm{et}}$
\begin{equation}\label{eq:C3-exceptional-lattice-sheaves}
 0\longrightarrow\underline{\mathbb Z}^{s}
       \longrightarrow\mathcal F\longrightarrow\mathcal H
       \longrightarrow0,
\end{equation}
where the first map sends a tuple to
$\mathcal O(\sum_i a_iE_i)$.
Surjectivity follows by closing the divisor
of a rational section on the regular scheme
$Z_B^h$, component by component.
The kernel consists of Cartier divisors
supported on the boundary. Over each connected
component of $U$, those prime boundary divisors
are exactly $E_i\times_k U$, so their coefficients
are locally constant integers.
There are no relations in this lattice.
To see this, flat base change identifies the
Stein algebra of $Z_B^h$ with
$R'\otimes_RT_B$, which is normal.
The boundary in this normal affine base has
codimension $N\ge2$. A principal divisor on
$Z_B^h$ supported on that boundary comes from
a rational function on this Stein base.
The proper birational map is an isomorphism
at its codimension-one points; thus that
function and its inverse are regular away
from a subset of codimension at least two.
Normal Hartogs extension makes it a global
unit, so its divisor is zero. This proves
injectivity in \eqref{eq:C3-exceptional-lattice-sheaves}.

One has $H^1_{\mathrm{et}}(C,\underline{\mathbb Z}^{s})=0$.
For completeness, if $j:\operatorname{Spec}k(C)\to C$
is the generic point, normality of every etale
curve over $C$ gives
$j_*\underline{\mathbb Z}^{s}=\underline{\mathbb Z}^{s}$:
its connected components and generic components
coincide. The first Leray edge map therefore
injects this $H^1$ into
\[
 H^1(k(C),\mathbb Z^s)
   =\operatorname{Hom}_{\mathrm{cont}}
              (\operatorname{Gal}(k(C)^{\mathrm{sep}}/k(C)),\mathbb Z^s)=0.
\]
A continuous homomorphism from a profinite
group to the discrete torsion-free group
$\mathbb Z^s$ has finite and hence zero image.
The exact cohomology sequence of
\eqref{eq:C3-exceptional-lattice-sheaves} now
lifts $g'$ to a section $a\in\mathcal F(C)$.
This is still a sheaf-valued Picard class;
the next step constructs its actual family.

\emph{3. Formal units on a curve have the required
cohomology after rigidification.}
Write $\widehat R'$ for the completion of $R'$;
the $\mathfrak mR'$-adic and maximal-ideal-adic
topologies agree. Let $\mathfrak n$ be its
maximal ideal. On $C_{\mathrm{et}}$ consider
the sheaves
\[
 \mathcal R=\varprojlim_r
       (\widehat R'/\mathfrak n^r)\otimes_k\mathcal O_C,
 \qquad \mathcal K=\ker(\mathcal R^\times\to\mathcal O_C^\times).
\]
Thus $\mathcal K=1+\mathfrak n\mathcal R$,
with complete separated filtration
$\mathcal K_r=1+\mathfrak n^r\mathcal R$ and
\begin{equation}\label{eq:C3-formal-unit-quotients}
 \mathcal K_r/\mathcal K_{r+1}
       = (\mathfrak n^r/\mathfrak n^{r+1})\otimes_k\mathcal O_C
\end{equation}
as additive sheaves. This description is
valid in every characteristic and uses no
logarithm.

For every affine etale $U\to C$,
\begin{equation}\label{eq:C3-formal-unit-affine-acyclicity}
                      H^i_{\mathrm{et}}(U,\mathcal K)=0
                                \quad(i>0).
\end{equation}
Here is an explicit proof avoiding a passage
from continuous cohomology to ordinary
cohomology. Take any finite affine etale
cover of $U$. Its disjoint union is affine
and faithfully flat over $U$. The augmented
Amitsur complex for each additive quotient
in \eqref{eq:C3-formal-unit-quotients} is
exact, by faithfully flat descent of modules.
Given a positive-degree Cech cocycle with
values in $\mathcal K$, solve the coboundary
equation modulo $\mathcal K_2$, lift the
cochain to $\mathcal K_1$, and divide by
its coboundary. The resulting cocycle lies
in $\mathcal K_2$. Repeat at every level.
All cochains lift on these affine intersections,
since the quotient maps of the coefficient
vector spaces split over $k$. The infinite
product of the correcting cochains converges
in the complete filtration and solves the
original coboundary equation on this same
cover. Thus every such cover has zero
positive Cech cohomology. Finite affine
etale covers form a cofinal system, stable
under fibre products. The exact criterion
\citet[Lemma 21.10.9]{D_StacksAlterationFamilies}
proves \eqref{eq:C3-formal-unit-affine-acyclicity}.

Cover $C$ by the complements of two distinct
$k$-points. Both opens and their intersection
are affine. Mayer--Vietoris and
\eqref{eq:C3-formal-unit-affine-acyclicity} give
\begin{equation}\label{eq:C3-formal-unit-curve-H2}
                           H^2_{\mathrm{et}}(C,\mathcal K)=0.
\end{equation}
In particular a gerbe banded by $\mathcal K$
on $C$ has an object: the affine etale covers
and their intersections satisfy the $H^1$
conditions, and
\citet[Lemma 21.11.1, Tag 0CK0]{D_StacksAlterationFamilies}
applies. This is the precise gerbe-neutralization
statement used below.

\emph{4. Constructing the actual formal family.}
Complete the local representatives of $a$
along $\mathfrak m$. They give a locally
represented Picard class on the formal
completion of $Z_{\widehat R}\times_k C$
along its closed fibre. Consider the stack
of its invertible-sheaf representatives,
rigidified along the closed subscheme
$\{e\}\times C$. It is locally nonempty:
after obtaining a local representative,
tensor it by the inverse of its restriction
to $e\times U$, pulled back from $U$.
That tensor factor is zero in the sheafified
Picard functor. Any two such representatives
are locally isomorphic. An isomorphism can
be made compatible with their rigidifications
by multiplying it by the inverse of its
restriction at $e$; the coefficient field
$k$ supplies the lift of this base unit.
Thus this stack is a gerbe.

Its band is exactly $\mathcal K$ of step 3.
Indeed flat base change gives
\[
 \Gamma(Z_B^h,\mathcal O)=R'\otimes_RT_B,
\]
by \citet[Lemma 30.5.2, Tag 02KH]{D_StacksAlterationFamilies}.
The theorem on formal functions then gives
the full formal global-functions algebra
\[
 \Gamma(\widehat{Z_B^h},\mathcal O)
    =\varprojlim_r (R'/\mathfrak m^rR')\otimes_kB,
\]
using \citet[Theorem 30.20.5, Tag 02OC]{D_StacksAlterationFamilies}.
This is $\mathcal R(U)$, with the equivalent
$\mathfrak n$-adic filtration. Evaluation at
$e$ is its residue map to $B$, so the
rigidification-preserving automorphisms are
$\mathcal K(U)$.
This argument identifies the inverse limit
of global functions. It does \emph{not}
claim that global functions on each individual
thickening are $R'/\mathfrak m^rR'$; such
a claim could fail through higher-cohomology
base-change terms.

Step 3 therefore neutralizes this gerbe,
producing a compatible actual formal line
bundle $\widehat{\mathcal M}$ on
$\widehat{Z_{\widehat R}\times C}$.
The scheme $Z_{\widehat R}\times C$ is
proper over the complete noetherian ring
$\widehat R$. Grothendieck existence,
\citet[Proposition 30.25.4, Tag 088C]{D_StacksAlterationFamilies},
algebraizes $\widehat{\mathcal M}$ to an
invertible sheaf $\mathcal M$ on
$Z_{\widehat R}\times C$.
Invertibility follows either by algebraizing
its formal inverse and the two evaluation
isomorphisms, or by properness and the
closed non-locally-free locus.
For every $t\in C(k)$, its restriction to
$Z^\circ_{\widehat R}\times\{t\}$ represents
the pullback of $g(t)\in G(k)=\operatorname{Pic}(X)$
after completion. This follows from the
definition of the lifted class $a$; at a
$k$-point etale sheafification adds no new
isomorphism classes.

\emph{5. Rigidity of proper-curve differences
under approximation.}
We need a precise fact about a proper scheme
$W$ over a complete noetherian local
$k$-algebra with residue $k$ and connected
closed fibre. Suppose two line bundles
$\mathcal L_1,\mathcal L_2$ on $W\times C$
agree on the closed fibre times $C$.
Then for every $t_0,t_1\in C(k)$ their
fibre differences agree in $\operatorname{Pic}(W)$:
\begin{equation}\label{eq:C3-proper-curve-rigidity}
 [\mathcal L_{1,t_1}\otimes\mathcal L_{1,t_0}^{-1}]
   =[\mathcal L_{2,t_1}\otimes\mathcal L_{2,t_0}^{-1}].
\end{equation}

To prove this, let $W_r$ be the closed
infinitesimal thickenings and put
$\mathcal D=\mathcal L_1\otimes\mathcal L_2^{-1}$.
Each $W_r$ is a proper finite-type $k$-scheme,
so its Picard functor is a separated group
scheme locally of finite type
\citep[Proposition 1.3, p.~4]{D_ChengJiLarsonOlander2021}.
The kernel of
$\operatorname{Pic}_{W_r/k}\to\operatorname{Pic}_{W_1/k}$
is quasi-affine of finite presentation.
For this precise property apply
\citet[Proposition 1.4, p.~4, from Theorem 1.2,
p.~3]{D_ChengJiLarsonOlander2021}
to the proper surjective closed immersion
$W_1\hookrightarrow W_r$: its base is
$\operatorname{Spec}k$, so the theorem's
nonempty base open is the whole base.
No smoothness or reducedness of the
thickenings is required.

The family $\mathcal D$ defines a morphism
$C\to\operatorname{Pic}_{W_r/k}$ whose
image is in this quasi-affine kernel.
Every morphism from the proper connected
curve $C$ to a quasi-affine $k$-scheme is
constant: compose with an open immersion
into an affine scheme and use
$\Gamma(C,\mathcal O_C)=k$.
Thus $\mathcal D_{t_1}\otimes\mathcal D_{t_0}^{-1}$
is trivial on every $W_r$.
These trivializations can be chosen compatibly.
For fixed $r$, the images of
$\Gamma(W_m,\mathcal O)\to\Gamma(W_r,\mathcal O)$
stabilize, since they are descending
subspaces of a finite-dimensional $k$-vector
space. These global-function algebras are
local Artinian with residue $k$, because
the common underlying closed fibre is
connected. Hence images of their unit
groups stabilize as well: an element of
an image algebra is a unit precisely when
its residue is nonzero, and then any
preimage is a unit. The nonempty torsors
of trivializations therefore satisfy the
Mittag--Leffler condition and admit a
compatible inverse-limit point.
Formal full faithfulness
\citep[Lemma 30.24.1, Tag 0883]{D_StacksFormalPicard}
now proves \eqref{eq:C3-proper-curve-rigidity}.

The line bundle $\mathcal M$ of step 4
descends, with its invertibility, to
$Z_B\times C$ for a finitely generated
$R$-subalgebra $B\subset\widehat R$:
this is finite-presentation descent for
line bundles on the fixed quasi-compact
quasi-separated scheme $Z\times C$.
Since $R$ is an excellent henselian local
ring, Artin approximation in the precise
form
\citet[Theorem 16.13.1, Tag 07QY]{D_StacksAlterationFamilies}
supplies an $R$-map $B\to R$ agreeing
modulo $\mathfrak m$ with $B\to\widehat R$.
Pulling back gives an actual invertible
sheaf $\mathcal M_R$ on $Z\times C$
whose completion agrees with $\mathcal M$
on the closed fibre times $C$.
Apply \eqref{eq:C3-proper-curve-rigidity}
to $W=Z_{\widehat R}$. The completed fibre
differences of $\mathcal M_R$ are exactly
those of $\mathcal M$, hence on
$Z^\circ_{\widehat R}$ are the pullbacks
of $g(t_1)-g(t_0)$ after completion.

\emph{6. Transfer of the actual family.}
There is an open $U\subset X$ containing
all points of codimension at most one
such that $h^{-1}(U)\to U$ is finite flat
of degree $d$. Indeed, under the finite
normalization over the normal domain $R$,
primes above a codimension-one prime have
codimension one. Their local rings in
$R'$ are discrete valuation rings, and
the proper birational map
$Z\to\operatorname{Spec}R'$ is an
isomorphism there. Over that codimension-one
point of the regular $X$, the finite
normalization algebra is torsion-free
over a discrete valuation ring and is
therefore flat. Properness makes
the finite locus open, and flatness then
gives the required $U$.

On the actual product $U\times C$, form
the finite-flat norm of
$\mathcal M_R|_{h^{-1}(U)\times C}$.
Explicitly, for a finite flat map $v$,
\[
                 \operatorname{Nm}_v(\mathcal L)
                  =\det(v_*\mathcal L)\otimes\det(v_*\mathcal O)^{-1}.
\]
This norm commutes with arbitrary base
change and satisfies
$\operatorname{Nm}_v(v^*M)=M^d$.
The complement of $U\times C$ has
codimension at least two in the regular
scheme $X\times C$. Closure of divisors
therefore extends the norm uniquely to
an invertible sheaf $\mathcal N$ on
$X\times C$. The same uniqueness shows
that its fibres are the corresponding
extensions of the fibre norms.

Completion is flat and regular here.
The complement of $U_{\widehat R}$ still
has codimension at least two, by going
down, and $X_{\widehat R}$ is regular.
Thus the extended norm also commutes
with completion: this can be checked on
the finite-flat open and extended
uniquely. Step 5 and the norm identity
therefore give
\[
 [\mathcal N_{t_1}\otimes\mathcal N_{t_0}^{-1}]_{\widehat R}
                  =d\bigl(g(t_1)-g(t_0)\bigr)_{\widehat R}.
\]
Elkik's Picard-group comparison for a
noetherian henselian pair,
\citet[Section III, unnumbered corollary,
p.~573]{D_Elkik1973}, makes
$\operatorname{Pic}(X)\to\operatorname{Pic}(X_{\widehat R})$
injective. Consequently
\begin{equation}\label{eq:C3-alteration-curve-multiple}
 d\bigl(g(t_1)-g(t_0)\bigr)
      =[\mathcal N_{t_1}\otimes\mathcal N_{t_0}^{-1}]
             \in\operatorname{Pic}(X)_{\mathrm{alg}}.
\end{equation}
Every line bundle used in this last
ordinary relation is on the fixed actual
product $X\times C$.

\emph{7. Removing the degree and applying C.3q.}
Let $b\in B(k)$ for a closed abelian
subvariety $B\subset G$. Multiplication
by $d$ on $B(k)$ is surjective, including
when $d$ is divisible by the characteristic.
Choose $b'\in B(k)$ with $db'=b$.
Join $0$ and $b'$ by a chain of integral
curves in the projective variety $B$,
and normalize them. These are smooth
projective curves over the perfect field
$k$. Apply \eqref{eq:C3-alteration-curve-multiple}
to their maps to $G$ and add the
resulting actual ordinary relations.
Their sum is $db'=b$. This proves
\eqref{eq:C3-proper-local-image-ordinary}.

Finally suppose $\operatorname{char}k=p>0$,
$N\ge4$, and $Y$ is regular. C.3n places
every ordinary target class in a closed
abelian subvariety of the target local
Picard scheme. C.3q gives a surjective
isogeny from a closed source abelian
subvariety onto that target subvariety.
Every $k$-point lifts under this isogeny,
and \eqref{eq:C3-proper-local-image-ordinary}
makes the lifted point ordinary on $X$.
Thus \eqref{eq:C3-charp-ordinary-isomorphism}
is surjective integrally. C.3n identifies
its kernel as $p$-primary torsion, giving
the rational isomorphism asserted above.
\end{proof}

\emph{Scope of the construction.}
The proof uses an actual regular
alteration of an algebraic model, a
constant exceptional-divisor lattice,
neutralization of a specifically
rigidified formal Picard gerbe, proper
formal existence, and Artin approximation
of a whole proper-curve family. None of
these steps is a blanket essential
surjectivity assertion for henselian
GAGA. The transfer is taken only after
the actual family on $Z\times C$ has
been constructed. No all-excellent-ring
extension beyond the stated model is
claimed here.


\paragraph{C.3x (Proved/refuted): the complete ordinary rational degree-one dimension table in positive characteristic.}
Fix an algebraically closed field $k$ of characteristic $p>0$.
Let $R$ be the strict henselization at a $k$-rational point of a
normal local scheme essentially of finite type over $k$, with the
given coefficient field inducing the residue-field identification.
Write
\[
 N=\dim R,\qquad
 X=\operatorname{Spec}R\setminus\{\mathfrak m\},\qquad
 Y=X\cap V(f),\quad 0\ne f\in\mathfrak m,
\]
and assume $X$ and $Y$ regular. Use ordinary algebraic equivalence,
defined by actual line-bundle families on products with smooth
connected curves, and put
\[
 J^1(X)=\operatorname{Pic}(X)_{\mathrm{alg}}\otimes_{\mathbb Z}\mathbb Q,
 \qquad
 J^1(Y)=\operatorname{Pic}(Y)_{\mathrm{alg}}\otimes_{\mathbb Z}\mathbb Q.
\]
If $k=\overline{\mathbb F}_p$, both groups vanish in every dimension.
If $k\ne\overline{\mathbb F}_p$, the universal assertions in this
algebraizable scope have exactly the following answers.
Here ``No'' means that the explicit example below over the fixed
field $k$ disproves the universal assertion.

\begin{center}
\begin{tabular}{lll}
\toprule
Dimension $N$ & Always injective & Always surjective \\
\midrule
$1$ & Yes, both groups vanish & Yes, both groups vanish \\
$2$ & No, the elliptic cone below & Yes, the target vanishes \\
$3$ & Yes, by C.3n & No, the same cone as divisor below \\
$\ge4$ & Yes, by C.3n & Yes, by C.3w \\
\bottomrule
\end{tabular}
\end{center}

In particular, the two boundary failures hold in characteristic
$2$, in characteristic $3$, and in every other positive
characteristic. The last row uses the algebraizability hypothesis
of C.3w. It is not an assertion for arbitrary nonalgebraizable
excellent germs, and it differs from the characteristic-zero
surjectivity table in C.3o.

\begin{proof}
\emph{1. A smooth constant elliptic curve in every characteristic.}
Use homogeneous coordinates $(x:y:z)$ and define
\[
 F_p=
 \begin{cases}
   y^2z+yz^2+x^3,&p=2,\\
   y^2z-x^3+xz^2,&p\ne2,
 \end{cases}
 \qquad
 E=V_+(F_p)\subset\mathbb P^2_k,
 \qquad O=(0:1:0).
\]
For $p=2$, the three partial derivatives are $x^2,z^2,y^2$.
They have no common projective zero. For $p\ne2$, the derivatives
are
\[
 (F_x,F_y,F_z)=(z^2-3x^2,\ 2yz,\ y^2+2xz).
\]
At $z=0$ the equation forces $x=0$, and then $F_z=y^2\ne0$.
On $z=1$, a singular point would have $y=0$ and
$x(x^2-1)=0$. The value of $F_x$ is $1$ when $x=0$ and $-2$
when $x=\pm1$, so it cannot vanish. This also checks $p=3$.
An exact identity over $\mathbb Z[1/2][x,y]$, with
$G=y^2-x^3+x$, certifies the same affine assertion:
\begin{equation}\label{eq:C3x-odd-smoothness}
 1=
 \left(1-\frac{3x^2}{2}\right)G_x
 +\frac{9x}{2}G
 -\frac{9xy}{4}G_y.
\end{equation}
Thus $E$ is a smooth plane cubic with the specified point $O$,
and hence a genus-one curve with its elliptic-curve structure.
All these data are defined over $\mathbb F_p$.

\emph{2. A specified infinite-order actual class over the fixed field.}
Let $k_0\subset k$ be the algebraic closure of $\mathbb F_p$ in $k$.
The assumption $k\ne k_0$ gives an element $t\in k$
transcendental over $k_0$. Choose $v\in k$ satisfying
\[
 v^2+v=t^3\quad(p=2),\qquad
 v^2=t^3-t\quad(p\ne2),
\]
and set
\[
 P=(t:v:1)\in E(k),\qquad Q=\mathcal O_E(P-O).
\]
The Abel map $E\to\operatorname{Pic}^0(E)$, $P\mapsto
\mathcal O_E(P-O)$, is an isomorphism over $k_0$.
For example, Riemann--Roch shows that each degree-one line
bundle on a genus-one curve has a unique effective divisor;
this gives the inverse. For every integer $a>0$, the group
scheme $E[a]$ is finite over $k_0$, including when $p\mid a$:
multiplication by $a$ is a finite flat isogeny of degree $a^2$
\citep[Proposition 8.1 and Theorem 8.2, pp.~12--14]{D_MilneAbelian2022}.
Every $k$-point of a finite $k_0$-scheme is already a $k_0$-point,
since its residue field is finite over the algebraically closed
field $k_0$. Consequently all torsion points of this constant
elliptic curve over $k$ lie in $E(k_0)$. The point $P$ does not,
because its $x$-coordinate is $t$. Hence $Q$ has infinite order.
This argument includes the non-etale $p$-primary torsion; an
etale-only torsion calculation would not suffice.

\emph{3. The actual Picard injection on the strict local cone.}
Put $L=\mathcal O_E(1)$, so $\deg L=3$, and define
\[
 R_2=\left(k[x,y,z]/(F_p)\right)_{(x,y,z)}^{\mathrm{sh}},
 \qquad
 X_2=\operatorname{Spec}R_2\setminus\{\mathfrak m_2\}.
\]
The cone is integral and smooth away from its vertex.
It is normal: the hypersurface is Cohen--Macaulay and is
regular in codimension one. Its local ring is excellent,
and its strict henselization remains an excellent normal
local domain. The usual cone projection induces
$\rho:X_2\to E$. Base change the proper cone resolution
$\operatorname{Tot}(L^{-1})\to\operatorname{Spec}k[x,y,z]/(F_p)$
to $\operatorname{Spec}R_2$. The resulting regular integral
scheme $Z_2$ is isomorphic to $X_2$ away from its single
exceptional prime divisor $E$, and
\[
                     \mathcal O_{Z_2}(E)|_E=L^{-1}.
\]
If $Q'\in\operatorname{Pic}^0(E)(k)$ pulls back to a trivial
line bundle on $X_2$, its pullback to $Z_2$ has a rational
trivialization whose Cartier divisor is supported on $E$.
It is therefore $\mathcal O_{Z_2}(bE)$ for some integer $b$.
Restriction gives $Q'\simeq L^{-b}$, and degrees give
$0=-3b$, so $b=0$ and $Q'=0$. This proves an actual injection
\begin{equation}\label{eq:C3x-elliptic-injection}
 \rho^*:\operatorname{Pic}^0(E)(k)
       \lhook\joinrel\longrightarrow
       \operatorname{Pic}(X_2)_{\mathrm{alg}}.
\end{equation}
The target subscript is justified directly, not by an
identity-component convention for a local Picard functor:
the line bundle $\mathcal O_{E\times E}
(\Delta-\{O\}\times E)$ has fibres
$\mathcal O_E(P'-O)$ as $P'\in E(k)$ varies. Pull it back
along $\rho\times\mathrm{id}_E$ to obtain an actual family
on $X_2\times E$. In particular, $\rho^*Q$ has infinite
order and gives a nonzero element of $J^1(X_2)$.

\emph{4. Failure of injection in dimension two.}
Take the specified nonzero divisor
\[
 s=
 \begin{cases}x,&p=2,\\y,&p\ne2\end{cases}
 \quad\text{in }R_2,
 \qquad Y_2=X_2\cap V(s).
\]
Its equation on the plane section factors as
\[
 F_2|_{x=0}=yz(y+z),\qquad
 F_p|_{y=0}=-x(x-z)(x+z)\quad(p\ne2).
\]
In each case these are three distinct lines meeting only
at the vertex. Removing that vertex gives a regular
zero-dimensional puncture, and strict henselization
preserves this regularity. More explicitly, normalization
has three strictly henselian discrete valuation local
factors, so $Y_2$ is a disjoint union of three spectra
of fields. Thus $\operatorname{Pic}(Y_2)=0$.
The nonzero rational class $\rho^*Q\otimes1$ lies in the
kernel of $J^1(X_2)\to J^1(Y_2)$.

\emph{5. Failure of surjection in dimension three.}
Define instead
\[
 R_3=k[x,y,z]_{(x,y,z)}^{\mathrm{sh}},\qquad
 X_3=\operatorname{Spec}R_3\setminus\{\mathfrak m_3\},
 \qquad f=F_p.
\]
The ring $R_3$ is regular local, and its puncture has
trivial Picard group
\citep[Lemma 52.27.5, Tag 0F2H]{D_StacksRegularPuncture}.
Strict henselization commutes with the quotient by
$F_p$, so $Y_3=X_3\cap V(F_p)$ is the preceding $X_2$.
It is regular, and the nonzero element
$\rho^*Q\otimes1\in J^1(Y_3)$ has no preimage from
$J^1(X_3)=0$.

\emph{6. Completion of the universal table and the finite-field case.}
For $N=1$, a normal noetherian local domain is a discrete
valuation ring; $X$ is a field spectrum and $Y$ is empty,
so both groups vanish. For $N=2$, regularity makes $Y$
a zero-dimensional noetherian scheme, hence a finite
disjoint union of spectra of fields. Its Picard group
vanishes, proving universal surjectivity in that row.
For $N\ge3$, C.3n proves that ordinary Picard restriction
has only $p$-primary torsion in its kernel, and therefore
is injective after tensoring with $\mathbb Q$. For
$N\ge4$, C.3w proves integral ordinary surjectivity in
exactly the algebraizable scope stated here. Tensoring
with $\mathbb Q$ gives the asserted isomorphism.

When $k=\overline{\mathbb F}_p$, every point of every
abelian variety over $k$ is torsion: the variety, its
group law, and the given point descend to some finite
field, whose point group is finite. The actual
dual-Jacobian-family construction of C.3n places
each ordinarily algebraically trivial line-bundle
class in an abelian point-group image. Thus every
such class is torsion and both rational groups vanish.
This last conclusion does not require algebraizability;
it has the regular noetherian coefficient-field scope
of that actual-family argument.
\end{proof}

The exact certificate in
\texttt{code/verify\_positive\_characteristic\_boundary.py}
checks the polynomial identity
\eqref{eq:C3x-odd-smoothness} over $\mathbb Q$, its sole
denominator prime $2$, the characteristic-two derivative
identities, and the two squarefree hyperplane factorizations.
It also independently enumerates the projective singular
and section points over the prime fields of characteristic
$2,3,5,7,11,13$. The universal smoothness and infinite-order
assertions are proved above; they are not inferred from
this finite sample.

\paragraph{C.3y (Proved/refuted): an explicit multiplicative ordinary divisor quotient in every characteristic.}
For every algebraically closed field $k$, there is an explicitly
specified excellent normal strictly henselian two-dimensional local
$k$-algebra $R$, with regular puncture $X$, such that
\begin{equation}\label{eq:C3y-answer}
 \operatorname{Pic}(X)_{\mathrm{alg}}=0,
 \qquad
 B^1(X)=\operatorname{Pic}(X)
       \simeq k^\times\oplus(\mathbb Z/4\mathbb Z)^2.
\end{equation}
The last isomorphism is one of abstract abelian groups and is
noncanonical; the injection of $k^\times$ below is specified by
gluing on the exceptional triangle. In particular $B^1(X)$ is
infinite in every characteristic. Its rationalization is zero
when $k=\overline{\mathbb F}_p$, and is infinite dimensional over
$\mathbb Q$ in characteristic zero and whenever $k$ has positive
transcendence degree in positive characteristic. Algebraic
equivalence in \eqref{eq:C3y-answer} is ordinary equivalence by
actual product families, not membership in a local Picard
functor's identity component.

Here are equations defining the ring, including the exceptional
characteristic choice:
\[
 a=\begin{cases}2,&\operatorname{char}k=229,\\1,&\text{otherwise},\end{cases}
 \qquad
 D=xyz,\qquad
 G=x^4+y^4+z^4+a(x^3y+y^3z+z^3x),
\]
\begin{equation}\label{eq:C3y-ring}
 R=\left(k[x,y,z]/(xyz-G)\right)_{(x,y,z)}^{\mathrm{sh}},
 \qquad
 X=\operatorname{Spec}R\setminus\{\mathfrak m\}.
\end{equation}
No classification of cusp singularities and no contraction
theorem for algebraic spaces is needed for this construction.

\begin{proof}
\emph{1. Twelve transverse points and an explicit projective resolution.}
Let $D_0\subset\mathbb P^2_k$ be the coordinate triangle $xyz=0$.
The quartic $G$ is nonzero at all three vertices. On each side,
after choosing the cyclic affine coordinate with the two
vertices at zero and infinity, its restriction is
\[
                       g_a(t)=t^4+a t^3+1.
\]
The resultant with its derivative is
\[
                \operatorname{Res}(g_a,g_a')=256-27a^4.
\]
For $a=1$ this is $229$, which is nonzero in every allowed
characteristic except the explicitly excluded characteristic
$229$. There $a=2$ gives $-176\equiv53\pmod{229}$.
Thus $D_0$ and $G=0$ meet transversely in twelve distinct
points, four on each side and none at a vertex. This assertion
also holds in characteristics $2$ and $3$. Exact integer
Bezout identities, valid before reduction, are
\begin{align*}
 229&=(48t^2+229)g_1
       +(-12t^3-3t^2-55t-16)g_1',\\
 -176&=(192t^2-176)g_2
       +(-48t^3-24t^2+80t-32)g_2'.
\end{align*}

Blow up these twelve reduced points in $\mathbb P^2_k$, obtaining
a smooth projective surface $S$. Let $E_1,E_2,E_3$ be the strict
transforms of the three sides, and put $E=E_1+E_2+E_3$.
Each $E_i\simeq\mathbb P^1_k$ has self-intersection $-3$,
and each pair of different components meets transversely
in its unchanged triangle vertex. Consequently
\begin{equation}\label{eq:C3y-matrix}
 M=(E_i\cdot E_j)_{i,j}
 =\begin{pmatrix}-3&1&1\\1&-3&1\\1&1&-3\end{pmatrix},
 \qquad \det M=-16,
 \qquad \deg\mathcal O_E(-E)|_{E_i}=1.
\end{equation}

The four quartics $G,xD,yD,zD$ have base ideal $(G,D)$
as an ideal sheaf on $\mathbb P^2$. Its base points are
exactly the twelve simple points just blown up, so they
define a morphism
\[
 S\longrightarrow W\subset\mathbb P^3_k,
 \qquad [G:xD:yD:zD],
\]
where, in coordinates $(u_0:u_1:u_2:u_3)$,
\begin{equation}\label{eq:C3y-projective-image}
             W=V_+\bigl(u_0u_1u_2u_3-G(u_1,u_2,u_3)\bigr).
\end{equation}
The equation follows from homogeneity:
$G(xD,yD,zD)=D^4G(x,y,z)$ and
$(xD)(yD)(zD)=D^4$. It is irreducible, since it is
primitive of degree one in $u_0$ and
$\gcd(u_1u_2u_3,G(u_1,u_2,u_3))=1$.

This morphism contracts exactly $E$ to
$v=(1:0:0:0)$ and is an isomorphism
\begin{equation}\label{eq:C3y-resolution-open}
                           S\setminus E\simeq W\setminus\{v\}.
\end{equation}
Here is a chart verification, including the exceptional curves
of the twelve blowups. On the chart $u_i\ne0$ of the right-hand
side, the ratios $(u_1:u_2:u_3)$ give a point of $\mathbb P^2$
in its corresponding coordinate chart. After putting its
$i$th coordinate equal to one, the remaining equation is
$G=tD$, where $t=u_0/u_i$. This is precisely the $D$-chart
of the blowup of the ideal $(D,G)$. These charts exhaust
$S\setminus E$: at a base point the omitted $G$-chart
has coordinate $D/G$, and its zero locus is the strict
transform of the triangle. Away from the base points the
same assertion is immediate. Thus there are no additional
contracted curves or unidentified puncture branches.
The image is all of the integral surface $W$, since the
morphism is proper and has the dense open image just
identified. The hypersurface $W$ is Cohen--Macaulay and
regular away from $v$ by \eqref{eq:C3y-resolution-open};
it is therefore normal. Its chart $u_0=1$ is exactly
the affine equation in \eqref{eq:C3y-ring}.

Base change $S\to W$ to $\operatorname{Spec}R$, and call
the resulting regular integral proper surface $Z$.
Strict henselization is ind-etale over the local algebraic
model; it preserves regularity of this resolution and the
specified reduced exceptional triangle. The open complement
of $E$ is $X$. Moreover
\begin{equation}\label{eq:C3y-maximal-ideal}
                        \mathfrak m\mathcal O_Z=\mathcal O_Z(-E).
\end{equation}
Indeed, near $E$, the section $G$ representing $u_0$ is
nonvanishing in the resolved linear system. The three
functions $u_i/u_0$ are $(D/G)x,(D/G)y,(D/G)z$ in local
trivializations, and $x,y,z$ generate the trivialized
tautological line bundle. The ratio $D/G$ has a simple
zero on each strict triangle component; the common
exceptional factors at the twelve base points cancel.
Equivalently, its blowup $G$-chart is exactly the local
coordinate $D/G$ cutting out the strict triangle. At
a triangle node the equation is the product of the two
branch parameters, which cuts out the reduced Cartier
divisor $E$. This proves \eqref{eq:C3y-maximal-ideal}
scheme-theoretically.

\emph{2. The formal Picard group is exactly that of the triangle.}
Let $E_n\subset Z$ be defined by $\mathcal O_Z(-nE)$,
so $E_1=E$. For every $n\ge1$,
\begin{equation}\label{eq:C3y-coherent-vanishing}
             H^1(E,\mathcal O_E(-nE))=0.
\end{equation}
To see this directly, normalize the triangle. The line bundle
restricts to $\mathcal O_{\mathbb P^1}(n)$ on each of its
three components. Each component has two attaching points,
and, for $n\ge1$, evaluation of its sections at those two
points is surjective. In the normalization exact sequence
the map from the three spaces of sections to the three
node fibres is therefore surjective. The component
$H^1$ groups vanish, proving \eqref{eq:C3y-coherent-vanishing}.
Also $H^2$ of every coherent sheaf on this curve vanishes.

Apply these facts to the square-zero unit sequence
\[
 1\longrightarrow\mathcal O_E(-nE)
 \longrightarrow\mathcal O_{E_{n+1}}^\times
 \longrightarrow\mathcal O_{E_n}^\times\longrightarrow1.
\]
Restriction gives isomorphisms
$\operatorname{Pic}(E_{n+1})\simeq\operatorname{Pic}(E_n)$,
and the maps on global units are surjective. Hence compatible
formal line bundles, including their isomorphisms, give
exactly $\operatorname{Pic}(E)$; there is no residual inverse-limit
term from units. Write $\widehat R$ for the maximal-ideal
completion and $\widehat Z=Z\times_R\widehat R$.
Proper formal existence and full faithfulness, applied using
\eqref{eq:C3y-maximal-ideal}, now give
\begin{equation}\label{eq:C3y-formal-pic}
                         \operatorname{Pic}(\widehat Z)\simeq\operatorname{Pic}(E).
\end{equation}
The exact primary theorem is Grothendieck existence for a
proper scheme over a complete noetherian ring
\citep[Theorem 30.27.1, Tag 088E]{rootStacksExistence}.
A formal invertible sheaf algebraizes to an invertible
sheaf: local freeness holds near the closed fibre by
Nakayama, and the non-locally-free closed locus has closed
image under the proper morphism, so cannot be nonempty
while avoiding the closed point of the local base.

\emph{3. Passage to the actual puncture and computation of the group.}
Excellence makes $R\to\widehat R$ regular, so $\widehat Z$
is regular. Every line bundle on
$\widehat X=\operatorname{Spec}\widehat R\setminus\{\widehat{\mathfrak m}\}$
extends to $\widehat Z$ by Cartier closure. The kernel
of restriction consists exactly of the divisors supported
on $E_1,E_2,E_3$. Thus \eqref{eq:C3y-formal-pic} gives
\[
 \operatorname{Pic}(\widehat X)
  \simeq\operatorname{Pic}(E)\big/
        \langle\mathcal O_Z(E_1)|_E,
                  \mathcal O_Z(E_2)|_E,
                  \mathcal O_Z(E_3)|_E\rangle.
\]
The Picard comparison for a noetherian henselian pair and
its completion identifies this with $\operatorname{Pic}(X)$
\citep[Section III, unnumbered Picard corollary, p.~573]{D_Elkik1973}.
Here the pair is $(R,\mathfrak m)$, and its complementary
open is exactly the puncture. Thus this calculation gives
actual line bundles on the strict henselization, not only
formal line bundles.

Normalization and gluing at the three nodes give an exact
sequence of abelian groups
\begin{equation}\label{eq:C3y-triangle-pic}
 0\longrightarrow k^\times\longrightarrow
 \operatorname{Pic}(E)\xrightarrow{\deg}
 \mathbb Z^3\longrightarrow0.
\end{equation}
Explicitly, trivial line bundles on the three components
are glued by scalars at the three nodes. Changing the
three component trivializations removes two independent
scalars, and the oriented product of the three gluing
scalars is the remaining element of $k^\times$.
For arbitrary specified component degrees, gluing the
three component line bundles similarly shows surjectivity
onto $\mathbb Z^3$. The degrees of the three exceptional
divisor classes are the columns of $M$ in
\eqref{eq:C3y-matrix}. Since $\det M\ne0$, their lattice
has zero intersection with the displayed $k^\times$.
The Smith form of $M$ is $(1,4,4)$. Consequently
\begin{equation}\label{eq:C3y-puncture-extension}
 0\longrightarrow k^\times\longrightarrow
 \operatorname{Pic}(X)\longrightarrow
 (\mathbb Z/4\mathbb Z)^2\longrightarrow0.
\end{equation}
The group $k^\times$ is divisible because $k$ is
algebraically closed. Divisible abelian groups are
injective $\mathbb Z$-modules, so this sequence splits
as abstract abelian groups. Its multiplicative subgroup
has a concrete specification: for $\lambda\in k^\times$,
glue the three trivial component bundles with scalars
$(\lambda,1,1)$ in a fixed cyclic orientation, extend
uniquely through \eqref{eq:C3y-formal-pic}, restrict to
$\widehat X$, and use Elkik's bijection to obtain the
actual class in $\operatorname{Pic}(X)$.

\emph{4. Ordinary algebraic equivalence is zero.}
Let $T$ be a smooth connected parameter curve and let
$\mathcal M$ be an actual line bundle on $X\times T$.
Complete $T$ to a smooth projective curve $C$.
The scheme $Z\times C$ is regular and noetherian, so
Cartier closure extends $\mathcal M$ to a line bundle
on $Z\times C$. Its restriction to $E\times C$ has
constant isomorphism class on the fibres over $C$.
One can verify this without any local Picard representability:
on each $E_i\times C\simeq\mathbb P^1\times C$ the
line bundle has the form
\[
 \operatorname{pr}_{E_i}^*\mathcal O_{\mathbb P^1}(d_i)
       \otimes\operatorname{pr}_C^*N_i
\]
with constant integer $d_i$ and $N_i\in\operatorname{Pic}(C)$.
Gluing at a node identifies the corresponding $N_i$ and
$N_j$, after fixed one-dimensional fibre factors are
chosen. Make these identifications along two edges of
the triangle. The last gluing is an automorphism of
one line bundle on $C$, hence multiplication by an
element of $\Gamma(C,\mathcal O_C^\times)=k^\times$.
All component degrees and the remaining gluing scalar
are therefore constant on $C$. This is the explicit
proper-curve-to-$\mathbb G_m$ constancy argument.

For $t_0,t_1\in T(k)$, the quotient of the two extended
fibre line bundles is trivial on $E$. It is trivial on
$\widehat Z$ by \eqref{eq:C3y-formal-pic}, hence on
$\widehat X$, and then on $X$ by Elkik's injectivity.
Every ordinary algebraic-equivalence generator is thus
zero. This proves $\operatorname{Pic}(X)_{\mathrm{alg}}=0$
integrally and completes \eqref{eq:C3y-answer}.

\emph{5. Precise finiteness consequences.}
Every algebraically closed field has infinite multiplicative
group, so $B^1(X)$ is infinite in every case. If
$k=\overline{\mathbb F}_p$, this multiplicative group is
torsion, and \eqref{eq:C3y-puncture-extension} shows
$B^1(X)_{\mathbb Q}=0$. In positive characteristic with
$k\ne\overline{\mathbb F}_p$, choose $t\in k$
transcendental over $k_0=\overline{\mathbb F}_p$.
The classes of $t-c$, $c\in k_0$, are linearly independent
in $k^\times\otimes\mathbb Q$: after clearing a putative
rational relation, a finite product
$\prod_c(t-c)^{n_c}$ would be a root of unity, hence a
constant in $k_0$. Its divisor on $\mathbb P^1_{k_0}$
forces every $n_c=0$. In characteristic zero, the
distinct positive rational primes give the analogous
infinite independent set by unique factorization in
$\mathbb Q^\times$. The injection in
\eqref{eq:C3y-puncture-extension} preserves this
independence after tensoring with $\mathbb Q$.
\end{proof}

The exact script
\texttt{code/verify\_multiplicative\_divisor\_quotient.py}
checks the two integer Bezout certificates, the universal
resultant and projective-image identities, the intersection
matrix and its Smith invariants, and the component evaluation
ranks for the finite sample $1\le n\le12$. Its elapsed
runtime and exact output are saved. The all-$n$ vanishing,
formal existence, strict-local descent, and infinite
multiplicative independence are proved above rather than
inferred from those computations.

\paragraph{C.3z (Proved): removing the algebraizability restriction for the ordinary degree-one theorem.}
Let $k$ be an algebraically closed field and let $(R,\mathfrak m,k)$
be an excellent henselian normal noetherian local $k$-algebra whose
specified coefficient map induces the identity on the residue field.
Assume that $X=\operatorname{Spec}R\setminus\{\mathfrak m\}$ is regular.
There is a normal local $k$-algebra $A_0$ essentially of finite type,
with residue field $k$, and a local regular $k$-homomorphism
\[
                    A_0\longrightarrow R
\]
which induces an isomorphism of maximal-ideal completions.
Consequently $\operatorname{Spec}R$ admits a proper generically finite
surjective morphism from a regular integral scheme.

If $N=\dim R\ge3$, the point-realization conclusion of C.3w holds
for this $R$: every $k$-point of every closed abelian subvariety
of $\mathbf{Pic}^{\mathrm{loc}}_{R/k}$ belongs to
$\operatorname{Pic}(X)_{\mathrm{alg}}$ for ordinary algebraic
equivalence on the fixed product.
In positive characteristic $p$, if $N\ge4$, $0\ne f\in\mathfrak m$,
and $Y=X\cap V(f)$ is regular, then
\[
 \operatorname{Pic}(X)_{\mathrm{alg}}
       \longrightarrow\operatorname{Pic}(Y)_{\mathrm{alg}}
\]
is surjective with $p$-primary torsion kernel, and hence is an
isomorphism after tensoring with $\mathbb Q$.
Thus the entire ordinary degree-one table in C.3x extends to
these possibly nonalgebraizable excellent henselian germs.
No statement about mixed characteristic or about higher Chow
groups is included in this extension.

\begin{proof}
\emph{1. Algebraizing the completion with the given coefficient field.}
If $R=k$, we take $A_0=k$ and the identity alteration.
We may therefore assume that the embedding dimension below is positive.
The map $R\to\widehat R$ is regular by excellence.
Normality ascends along this map, so the complete local ring
$\widehat R$ is a normal domain of dimension $N$.
Its puncture is regular: a nonmaximal prime of $\widehat R$
cannot lie over $\mathfrak m$, because
$\widehat R/\mathfrak m\widehat R=k$; at its contraction
the source ring is regular, and regularity ascends along
a regular ring map.

Choose finitely many generators of $\mathfrak m/\mathfrak m^2$.
Their lifts and the specified copy of $k$ give a presentation
\[
             \widehat R\simeq k[[x_1,\ldots,x_e]]/I.
\]
Surjectivity follows by successively lifting homogeneous
polynomials in the generators and using completeness.
Since $k$ is perfect, the formal-power-series Jacobian
criterion identifies regularity of each localization of
this quotient with the required full-rank Jacobian minor.
The exact version we use is
\citet[Proposition 2.5, with the definition of $\mathcal B$
in Section 2]{rootHauserWoblistin2021}: take
$\mathcal B=k[[\mathrm t,\mathrm y]]$, simply relabeling
the $e$ variables. Thus the presentation is formally
smooth away from its closed point in the precise sense
of \citet[Theorem 3.8, p.~32]{rootArtin1969}.
That theorem supplies a finite-type $k$-scheme $V$, a
$k$-rational point $v$, and a $k$-isomorphism
\begin{equation}\label{eq:C3z-formal-algebraization}
                  \widehat{\mathcal O}_{V,v}\simeq\widehat R.
\end{equation}
The local ring $A_0=\mathcal O_{V,v}$ is normal by
faithfully flat descent from its normal completion.
After shrinking $V$ around $v$ we may and do take $V$
normal and integral. Here we use openness of the normal
locus for a finite-type algebra over a field; no
regularity assertion away from $v$ is needed for the
alteration below.

\emph{2. Obtaining an actual map to the original henselian ring.}
Write an affine neighborhood of $v$ in the form
$\operatorname{Spec}(k[t_1,\ldots,t_a]/(F_1,\ldots,F_b))$,
translating the coordinates so that $v$ is the origin.
The formal isomorphism \eqref{eq:C3z-formal-algebraization}
gives a solution $(\widehat t_1,\ldots,\widehat t_a)$
of these polynomial equations in $\widehat R$.
The general henselian G-ring approximation theorem
\citep[Theorem 16.13.1, Tag 07QY]{rootStacksRegularApproximation}
provides a solution $(t'_1,\ldots,t'_a)$ in $R$ with
\[
                   t'_j-\widehat t_j\in\mathfrak m^2\widehat R.
\]
This gives an actual local $k$-map $A_0\to R$.
It extends to a continuous map
$\widehat A_0\to\widehat R$.
The coordinate images generate the cotangent space
of $\widehat R$, since they agree modulo
$\mathfrak m^2$ with those of the given isomorphism.
The same successive approximation by polynomials now
shows that $\widehat A_0\to\widehat R$ is surjective.
Its kernel is prime, both rings are domains of the
same finite dimension $N$, and a nonzero prime in a
domain strictly lowers the maximal possible chain
length in its quotient. The kernel is therefore zero.
We have constructed a map inducing a genuine
isomorphism of completions, not merely matching
some chosen finite jet.

The composite $A_0\to R\to\widehat R$ is regular:
under this new isomorphism it is the completion map
of the excellent ring $A_0$.
The second arrow $R\to\widehat R$ is faithfully flat.
It follows that $A_0\to R$ is regular, by
\citet[Lemma 15.42.7, Tag 07NT]{rootStacksRegularApproximation}.
For clarity, that descent first gives flatness;
on each fibre, faithful flatness descends geometric
regularity. This proves the first assertion.

\emph{3. A regular proper alteration over $R$.}
Apply \citet[Theorem 4.1 and Remark 4.2, p.~66]{D_DeJong1996Alterations}
to the normal algebraic neighborhood $V$.
It gives a regular integral scheme $Z_0$ and a proper
generically finite surjective map $Z_0\to V$.
Base change first to $A_0$ and then by the actual
regular map $A_0\to R$.
Regular maps remain regular under finite-type base
change \citep[Lemma 15.42.3, Tag 07BY]{rootStacksRegularApproximation},
so the resulting scheme $Z_R$ is regular.
It is noetherian and proper over $R$; its generic
fibre is finite and nonempty because $A_0\to R$
is faithfully flat and the original alteration is
generically finite. A regular noetherian scheme
has disjoint integral connected components.
Choose one component meeting this generic fibre.
It is proper and generically finite over the
integral local base; its image is closed and
contains the generic point, hence it is surjective.
This is the claimed regular integral alteration.
We neither assert regular alterations for all
excellent schemes nor use altered local uniformization
as if it supplied a single global regular alteration.

\emph{4. Transfer to ordinary algebraic equivalence.}
The only use of an algebraic model in C.3w is its
step 1, which supplies one regular integral proper
alteration. Step 3 above supplies exactly that input
over the present $R$. Every subsequent hypothesis
is unchanged: the Stein algebra is a finite normal
local algebra with residue $k$; tensor products with
smooth parameter curves are noetherian and regular;
proper formal existence applies over $\widehat R$;
and the whole-family approximation in C.3w uses the
same general theorem 07QY, already valid for this
excellent henselian $R$.
The exceptional-divisor lattice, formal-unit gerbe,
proper-curve rigidity and finite-flat norm arguments
of steps 2--7 therefore prove the stated
point-realization assertion with this replacement of step 1.

\emph{5. Lifting target abelian subgroups without algebraizing the divisor.}
We supply the two replacements needed to use the argument of C.3q;
we do not assume that $f$ descends through $A_0\to R$.
Put $G=\mathbf{Pic}^{\mathrm{loc}}_{R/k}$ and let
$H=\prod_a\mathbf{Pic}^{\mathrm{loc}}_{C_a/k}$, where
$\prod_a C_a$ is the finite normalization of $R/(f)$.
These are locally finite-type group schemes in dimension $N\ge4$
by the precise Boutot input in C.3n. Restriction gives $r:G\to H$.

First, the functor comparison proved separately in C.3p gives
\[
 G\simeq\mathbf{Pic}^{\mathrm{loc}}_{\widehat R/k}
   \simeq\mathbf{Pic}^{\mathrm{loc}}_{A_0^h/k}.
\]
Here this is a comparison of group schemes, including identity
components, rather than only a bijection on $k$-points.
Indeed, for every finite-type $k$-algebra $D$, the henselized
tensor pairs for $R$ and $A_0^h$ have the same ideal-adic completion
\[
             \varprojlim_n(R/\mathfrak m^n)\otimes_k D.
\]
They are noetherian henselian pairs, so the Picard corollary of
\citet[Section III, p.~573]{D_Elkik1973} identifies the Picard
groups of their ideal complements with that of this common
completed complement. The maps are natural in $D$ and survive
etale sheafification. The affine finite-type open subschemes
of the represented group schemes then test the inverse by
Yoneda. The given coefficient maps and the $k$-isomorphism of
completions ensure that this comparison is over the specified $k$.
Thus the finite generation of the local Neron--Severi point
group for the normal algebraic model
\citep[Definition 3.4 and the following paragraph, p.~468]
{D_Kollar2016LocalPicard} gives the same finite generation
for $G/G^0$. Only bounded prime-to-$p$ torsion is needed.

Choose a prime $\ell\ne p$ avoiding that finite torsion group.
For any closed abelian subvariety $B\subset H$, C.3p lifts
every $B[\ell^n](k)$ uniquely into $G(k)[\ell^n]$.
Its component class is zero, so the lift lies in
$G_0=(G^0)_{\mathrm{red}}$. The closed image $r(G_0)$
contains the Zariski-dense union of these torsion groups;
therefore it contains $B$. As in C.3q, taking the reduced
identity component of $r^{-1}(B)\cap G_0$ gives a smooth
connected commutative algebraic group $E$ and a faithfully
flat homomorphism $r_E:E\to B$.

Second, its scheme-theoretic kernel $K$ is of finite type
and C.3k(i) gives $K(k)[\ell']=0$ for every prime
$\ell'\ne p$. We now prove that some $[p^e]$ kills $K$
as a group scheme. This assertion allows an infinitesimal
multiplicative part such as $\mu_p$.
The reduced identity component $U=(K_{\mathrm{red}})^0$
is unipotent. To see this, apply the smooth connected
Chevalley decomposition to $U$. Its connected affine
part is an extension of a torus by a unipotent group.
A nonzero torus already supplies prime-to-$p$ torsion.
If the abelian quotient were positive dimensional,
its prime-to-$p$ torsion would lift to torsion in $U(k)$,
because multiplication by that prime is surjective on
the connected affine part. Both possibilities contradict
the displayed torsion exclusion. This is the argument
of \citet[Lemma 2.6 and proof, p.~466]{D_Kollar2016LocalPicard},
applied to the reduced group.
The same lifting observation, using the invertibility
of prime-to-$p$ multiplication on $U$, shows that the
finite component group of $K_{\mathrm{red}}$ is a $p$-group.
The group $U$ embeds in an upper unitriangular matrix
group and is killed by a power of $p$. Hence a power
$[p^a]$ kills all of $K_{\mathrm{red}}$.

The reduced group is affine: it is a finite disjoint
union of translates of $U$. Its nilpotent thickening
$K$ is affine as well
\citep[Lemma 37.2.3, Tag 06AD]{rootStacksRegularApproximation}.
Choose a faithful closed representation $K\hookrightarrow
\operatorname{GL}_d$
\citep[Theorem 4.9 and Corollary 4.10, p.~87]
{rootMilneAlgebraicGroups2022}.
In the universal matrix of $[p^a]$ on $K$, all entries
of $M=[p^a]-I_d$ lie in the nilradical $J$ of $\mathcal O(K)$,
since this map is the identity matrix on $K_{\mathrm{red}}$.
Choose $s$ with $J^s=0$ and $b$ with $p^b\ge s$.
In characteristic $p$ we have
\[
                    (I_d+M)^{p^b}=I_d+M^{p^b}=I_d.
\]
Faithfulness proves $[p^{a+b}]_K=0$ scheme-theoretically.

Write $e=a+b$. The homomorphism $[p^e]_E$ descends
through the faithfully flat quotient $E\to B$ to a
homomorphism $\widetilde s:B\to E$ with
$r_E\widetilde s=[p^e]_B$. Its image $A\subset E\subset G$
is a closed abelian subvariety; the maps $B\to A\to B$
are isogenies with $p$-primary kernels, exactly as in
the final descent argument of C.3q.
Every ordinary algebraically trivial class of $Y$
belongs to some such $B(k)$ by the actual abelian-family
lemma of C.3n, applied to its finitely many regular
components. The surjection $A(k)\to B(k)$ and step 4
realize a preimage by an ordinary algebraic equivalence
on $X$. This proves integral surjectivity without any
simultaneous algebraization of $(R,f)$.

C.3n already has the general coefficient-field
scope used here and gives the $p$-primary kernel.
The counterexamples in dimensions two and three in
C.3x lie in this larger class; the dimension-one
vanishing and the $\overline{\mathbb F}_p$ torsion
argument also retain their original general scope.
This proves the extension of the entire table.
\end{proof}

\paragraph{C.3aa (Refuted): every non-top positive-codimension local finiteness assertion over $\mathbb C$.}
For every pair of integers
\[
                        N\ge4,\qquad 2\le r\le N-2,
\]
there is a specified excellent normal strictly henselian local
$\mathbb C$-algebra $R$ of dimension $N$, with regular puncture
$X$, such that
\begin{equation}\label{eq:C3aa-nonvanishing}
                         B^r(X)_{\mathbb Q}\ne0.
\end{equation}
Here $B^r$ is cycles modulo ordinary algebraic equivalence,
and codimension is inherited from the local spectrum as in
C.3e. In particular the integral group $B^r(X)$ is not
finite. Together with the divisor examples C.3v and C.3y,
this settles every non-top positive-codimension pair
$1\le r\le N-2$ over $\mathbb C$. This is a nonvanishing
assertion; no infinite-dimensionality claim for these
higher rational groups is made. The remaining top pair
$r=N-1$, $N\ge3$, is not decided by this argument.

\begin{proof}
\emph{1. Fix the explicit Nori witness in the smallest needed dimension.}
Apply the fully proved construction C.1i with its indices
$c=r$ and $n=r+1$. Its strict inequalities hold because
$r\ge2$. More explicitly, take the split quadric
\[
 Q=V_+\left(\sum_{i=0}^{r}x_i y_i\right)
           \subset\mathbb P^{2r+1}_{\mathbb Q},
 \qquad H_Q=c_1(\mathcal O_Q(1)),
\]
with the two specified maximal planes
\[
 \Pi_+=V_+(x_0,\ldots,x_r),\qquad
 \Pi_-=V_+(y_0,x_1,\ldots,x_r).
\]
Set $d=4r+2$. Let $T\subset Q$ be the geometric generic
complete intersection of $r-1$ degree-$d$ sections,
with exactly the independent coefficients and the one
normalized coefficient specified in C.1i. Thus the
defining field is the algebraic closure of that explicit
purely transcendental extension of $\mathbb Q$.
Fix any embedding of this field into $\mathbb C$;
C.1i proves persistence under every such embedding.
Writing $j:T\hookrightarrow Q_{\mathbb C}$ and
$H=j^*H_Q$, the actual codimension-$r$ cycle
\[
              z=j^*([\Pi_+]-[\Pi_-])
\]
satisfies
\begin{equation}\label{eq:C3aa-nori-witness}
    \dim T=r+1,\qquad
    z\notin H B^{r-1}(T)_{\mathbb Q}.
\end{equation}
In particular $z\ne0$ in $B^r(T)_{\mathbb Q}$.
This invokes the exact Nori--Nagel connectivity theorem
and the complete trace/incidence proof already checked
in C.1i; it introduces no new unspecified degree bound.

\emph{2. Raise the local dimension while retaining the obstruction.}
Put
\[
 m=N-r-2\ge0,\qquad V=\mathbb P^m_{\mathbb C}\times T,
 \qquad
 L=\mathcal O_{\mathbb P^m}(1)\boxtimes\mathcal O_T(1).
\]
Then $\dim V=N-1$. If $m=0$, the first factor is a point
and $L=\mathcal O_T(1)$, as intended. Denote the projective
hyperplane class by $h$ and continue to write $H$ for
its pullback from $T$; thus $c_1(L)=h+H$.

The projective-bundle formula, and preservation of
ordinary algebraic equivalence by its pullbacks,
pushforwards and Chern-class operations, give the
graded-ring identity
\begin{equation}\label{eq:C3aa-projective-ring}
 B^*(V)_{\mathbb Q}
        =B^*(T)_{\mathbb Q}[h]/(h^{m+1}).
\end{equation}
For completeness, the inverse to the expression
$\sum_{a=0}^m h^a\alpha_a$ is given by pushing
$h^{m-a}$ times a class to $T$. These operations
preserve ordinary algebraic equivalence, so the
usual Chow decomposition descends to this quotient
without imposing a stronger equivalence relation.
Imposing $h+H=0$ in \eqref{eq:C3aa-projective-ring}
therefore gives exactly
\begin{equation}\label{eq:C3aa-segre-quotient}
 \frac{B^r(V)_{\mathbb Q}}
      {(h+H)B^{r-1}(V)_{\mathbb Q}}
 \simeq
 \frac{B^r(T)_{\mathbb Q}}
      {H^{m+1}B^{r-m-1}(T)_{\mathbb Q}}.
\end{equation}
The map is evaluation at $h=-H$; explicitly it sends
$\sum_a h^a\alpha_a$ to $\sum_a(-H)^a\alpha_a$.
If $r-m-1<0$, the denominator on the right is zero.
Otherwise
\[
 H^{m+1}B^{r-m-1}(T)_{\mathbb Q}
           \subset H B^{r-1}(T)_{\mathbb Q}.
\]
Consequently the pullback of $z$ has nonzero image
in \eqref{eq:C3aa-segre-quotient}, by
\eqref{eq:C3aa-nori-witness}. This argument covers
all $m\ge0$, including $m\ge r$, where the additional
denominator has disappeared entirely.

\emph{3. Verify the cone hypotheses and pass to its strict puncture.}
The hyperplane embedding of $T$ is a smooth complete
intersection in projective space. Its homogeneous
coordinate ring is Cohen--Macaulay, its affine cone
is regular away from the vertex, and the vertex has
codimension $r+2\ge4$. Serre's criterion therefore
makes that ring normal. The complete-intersection
Koszul resolution, equivalently the usual vanishing
of intermediate projective-space cohomology, identifies
it with the full section ring
\[
                    \bigoplus_{n\ge0}H^0(T,\mathcal O_T(n)).
\]
It is generated in degree one. The full section ring
of $(V,L)$ is the Segre product
\[
 A=\bigoplus_{n\ge0}
       H^0(\mathbb P^m,\mathcal O(n))
          \otimes_{\mathbb C} H^0(T,\mathcal O_T(n)).
\]
It is generated in degree one because both factors
are: products of paired degree-one generators span
each displayed tensor product. It is normal, either
by the full-section-ring valuation criterion on the
smooth normal variety $V$, or directly as follows.
The tensor product of the two normal coordinate
rings is the coordinate ring of their product,
which is normal over the perfect algebraically
closed field. The Segre ring is its weight-zero
subring for the action with weights $+1$ and $-1$.
A weight-zero fraction integral over this subring
is integral over the normal tensor ring and hence
belongs to it; being fixed by the action, it lies
in the weight-zero subring. This proves normality.
Thus $L$ is a very ample projectively normal
degree-one-generated embedding of precisely the
kind required in C.3e.

Let $v$ be the vertex of $\operatorname{Spec}A$,
and define the actual local ring and puncture by
\[
 R=\mathcal O_{\operatorname{Spec}A,v}^{\mathrm{sh}},
 \qquad X=\operatorname{Spec}R\setminus\{\mathfrak m\}.
\]
Then $R$ is excellent, normal and strictly henselian,
$\dim R=N$, and $X$ is regular: the cone off its
vertex is the $\mathbb G_m$-bundle of $L^{-1}$
over the smooth variety $V$, and regularity persists
under strict henselization. The integral strict-cone
injection of C.3e, tensored with the flat module
$\mathbb Q$, embeds the left-hand side of
\eqref{eq:C3aa-segre-quotient} into $B^r(X)_{\mathbb Q}$.
The specified pullback of $z$ remains nonzero,
which proves \eqref{eq:C3aa-nonvanishing} and gives
the asserted actual counterexample for each $(N,r)$.
\end{proof}

The algebraic elimination in
\eqref{eq:C3aa-segre-quotient} is also recorded by
\texttt{code/verify\_segre\_nori\_quotient.py}. It checks
the exact identity
\[
 h^a-(-H)^a=(h+H)\sum_{b=0}^{a-1}h^{a-1-b}(-H)^b
\]
for its stated finite sample and verifies the index
and degree inequalities over a finite grid of $(N,r)$.
The universal quotient formula is proved above;
the finite program is an independent sign/index
check and does not compute any Griffiths group.

% Final: full independent proof/source audits by /root and /root/frobenius.
\paragraph{C.3ab (Proved): integral top-cycle finiteness for surface cones with an explicit vanishing hypothesis.}
Let $k$ be an algebraically closed field of characteristic zero,
$V$ a smooth projective integral surface over $k$, and $L$ a very
ample line bundle whose full section ring
\[
 S=\bigoplus_{j\ge0}H^0(V,L^j)
\]
is normal and generated in degree one. Put
\[
 R=(S_{S_+})^{\mathrm{sh}},\qquad
 X=\operatorname{Spec}R\setminus\{\mathfrak m_R\},\qquad
 g=\gcd\{\deg(c_1(L)\cdot C):C\subset V
                         \text{ an integral curve}\}.
\]
Assume the following cohomological condition, with \emph{absolute}
K\"ahler differentials over $\mathbb Q$:
\begin{equation}\label{eq:C3ab-absolute-vanishing}
 H^2(V,\Omega^1_{V/\mathbb Q}\otimes L^j)=0
                         \qquad(j\ge1).
\end{equation}
Then the natural maps from C.3e are isomorphisms, integrally:
\begin{equation}\label{eq:C3ab-answer}
 \mathrm{CH}^2(X)\simeq
 \frac{\mathrm{CH}_0(V)}{c_1(L)\mathrm{CH}^1(V)},
 \qquad
 B^2(X)\simeq\mathbb Z/g\mathbb Z.
\end{equation}
Here $B^2$ uses ordinary algebraic equivalence on actual fixed
products with parameter curves. The integer $g$ is positive.
Condition \eqref{eq:C3ab-absolute-vanishing} holds for every such
surface over $\overline{\mathbb Q}$. Over an arbitrary $k$ as above,
it holds if $H^2(V,L^j)=0$ for all $j\ge1$; in particular this
includes $K3$ surfaces and surfaces with $p_g=0$.
This statement does not settle arbitrary three-dimensional
strict local germs, arbitrary complex surface cones, or the top
groups in local dimension at least four.

\begin{proof}
\emph{1. The resolution and its actual infinitesimal neighbourhoods.}
The degree-one hypothesis identifies the blow-up of
$\operatorname{Spec}S$ at its vertex with
$W=\operatorname{Tot}(L^{-1})$, as proved by charts in C.3e.
Set
\[
 Z=W\times_{\operatorname{Spec}S}\operatorname{Spec}R,
 \qquad i:V\hookrightarrow Z,
 \qquad p:Z\longrightarrow V.
\]
Then $Z$ is a regular integral three-dimensional scheme, proper
over $R$, and $Z\setminus V=X$. The divisor $V$ is Cartier,
its normal bundle is $L^{-1}$, $p\circ i=1$, and
$\mathfrak m_R\mathcal O_Z=\mathcal O_Z(-V)$.
The morphism $p$ is the limit of the smooth relative-dimension-one
maps in C.3e; in particular it is flat and has the indicated
flat pullback on cycles. For $a\ge1$ write
\[
 V_a=Z\times_R R/\mathfrak m_R^a.
\]
Strict henselization does not change the quotients of the local
cone by powers of its maximal ideal, because the residue field
is already algebraically closed. Thus these are the split
thickenings
\begin{equation}\label{eq:C3ab-split-thickenings}
 V_a=\operatorname{Spec}_V
       (\mathcal O_V\oplus L\oplus\cdots\oplus L^{a-1}),
\end{equation}
with the truncated graded multiplication. Their projections to
$V$ split restriction on $K$-theory and on Picard groups.

\emph{2. Pro cdh descent gives detection at a finite thickening.}
We use nonconnective algebraic $K$-theory, with $K_0$ the
Grothendieck group of perfect complexes. Apply
\citet[Theorem 3.7(i), author pp.~21--22]{D_Morrow2016}
to the abstract blow-up square formed by $Z$, $\operatorname{Spec}R$
and their closed fibres. All schemes are noetherian,
quasi-excellent $\mathbb Q$-schemes of finite Krull dimension.
The resulting exact sequence of pro abelian groups includes
\[
 K_0(R)\longrightarrow
 \{K_0(R/\mathfrak m_R^a)\}_a\oplus K_0(Z)
 \longrightarrow\{K_0(V_a)\}_a.
\]
The first groups are all canonically $\mathbb Z$, and the first
arrow is $n\mapsto(n,n[\mathcal O_Z])$. Quotienting by this graph
shows that
\begin{equation}\label{eq:C3ab-pro-monomorphism}
 K_0(Z)\longrightarrow\{K_0(V_a)\}_a
\end{equation}
is a monomorphism in the abelian category of pro abelian groups.
This assertion is stronger than merely separating a fixed
individual element in an inverse limit. Indeed, the kernel pro
group is represented by the decreasing subgroups
$\ker(K_0(Z)\to K_0(V_a))$, whose transition maps are inclusions.
A pro-zero system of this form has a zero term. Consequently
there is $a_0$ such that
\begin{equation}\label{eq:C3ab-finite-detection}
 K_0(Z)\longrightarrow K_0(V_a)\quad\text{is injective for }a\ge a_0.
\end{equation}
No assertion that arbitrary $K$-theory commutes with a formal
inverse limit is used.

\emph{3. The relative $K_2$ calculation retains integral information.}
Let $\mathcal K_q(V_a,V)$ denote the Zariski sheaf of relative
$K_q$ on their common underlying space $V$. Nilpotence gives
$\mathcal K_0(V_a,V)=0$. Determinant and the finite logarithm give
\[
 \mathcal K_1(V_a,V)=1+\mathcal I_a
       \simeq\bigoplus_{j=1}^{a-1}L^j,
 \qquad \mathcal I_a=\ker(\mathcal O_{V_a}\to\mathcal O_V).
\]
For a nilpotent ideal in a $\mathbb Q$-algebra, all relative
$K$-groups are uniquely divisible
\citep[Consequence 1.4, p.~392]{D_Weibel1982}.
Therefore the rational comparison of
\citet[Main Theorem, p.~348]{D_Goodwillie1986}
is an integral isomorphism of the underlying abelian groups here.
In degree two it identifies the relative sheaf with relative
$HC_1$ over $\mathbb Q$.

The latter calculation can be made explicitly. On a trivializing
affine open of $V$, write $A[t]/(t^a)$ for
\eqref{eq:C3ab-split-thickenings}. For a commutative
$\mathbb Q$-algebra $B$, the degree-one cyclic complex gives
$HC_1(B)=\Omega^1_{B/\mathbb Q}/dB$: this follows directly from
the cyclic relation in degree one, or from the low-degree
Connes exact sequence $B\xrightarrow{d}\Omega^1_{B/\mathbb Q}
\to HC_1(B)\to0$.
The positive homogeneous part of degree $j$, $1\le j<a$, is
\[
 \frac{t^j\Omega^1_{A/\mathbb Q}\oplus
                 A t^{j-1}dt}
      {\{t^j\,db+j b t^{j-1}dt:b\in A\}}.
\]
Every element has a unique horizontal representative:
\[
 [t^j\omega+b t^{j-1}dt]
          \longmapsto t^j(\omega-j^{-1}db).
\]
Equivalently, inclusion of the horizontal summand is an
isomorphism onto this quotient. The relation
$a t^{a-1}dt=0$ removes degree $a$, and there are no further
positive degrees. Under a change of fibre coordinate the
horizontal inclusion transforms by the $j$th power of the
line-bundle transition function. Thus these local computations
give the global isomorphism
\begin{equation}\label{eq:C3ab-relative-K2}
 \mathcal K_2(V_a,V)\simeq
       \bigoplus_{j=1}^{a-1}
                      \Omega^1_{V/\mathbb Q}\otimes L^j.
\end{equation}
In particular no extra conormal summand has been retained after
quotienting by exact differentials.

\emph{4. The only relative $K_0$ classes are determinant classes.}
Apply Zariski descent to the fibre of $K(V_a)\to K(V)$.
Theorem 10.3 of \citet[pp.~383--386]{D_ThomasonTrobaugh1990}
applies to both spectra and hence to their fibre; its spectral
sequence has
\[
 E_2^{s,t}=H^s(V,\mathcal K_t(V_a,V))
                       \Longrightarrow K_{t-s}(V_a,V).
\]
The underlying space has dimension two. On the diagonal
$t=s$ the only possibly nonzero terms are
$H^1(\mathcal K_1)$ and $H^2(\mathcal K_2)$.
The second vanishes by
\eqref{eq:C3ab-absolute-vanishing} and
\eqref{eq:C3ab-relative-K2}. The first has no possible incoming
or outgoing differential. Naturality for determinant therefore
gives an isomorphism
\[
 K_0(V_a,V)\xrightarrow{\det}\operatorname{Pic}(V_a,V)
         =H^1(V,1+\mathcal I_a).
\]
All these relative groups inject into their absolute groups
because the thickening has a projection to $V$.
It follows that restriction induces an isomorphism
\[
 SK_0(V_a)\xrightarrow{\sim}SK_0(V),\qquad
 SK_0(T)=\ker\bigl(K_0(T)\xrightarrow{(\operatorname{rk},\det)}
                         \mathbb Z\oplus\operatorname{Pic}(T)\bigr).
\]
Combining injectivity in \eqref{eq:C3ab-finite-detection}
with the section $p^*$ gives
\begin{equation}\label{eq:C3ab-SK0}
                 i^*:SK_0(Z)\xrightarrow{\sim}SK_0(V).
\end{equation}

\emph{5. Zero-cycles on $Z$ vanish by the henselian property.}
Every closed point of $Z$ lies on its closed fibre $V$, since
$Z$ is proper over the local ring $R$. Let $z$ be such a point.
In the regular local ring $\mathcal O_{Z,z}$ choose two parameters
whose common prime avoids a local equation for $V$, and close
the resulting integral one-dimensional germ in $Z$. This gives
an integral closed curve $D\subset Z$ through $z$, not contained
in $V$. To justify the parameter choice, complete a local
equation of the smooth Cartier divisor $V$ to a regular system
of three parameters and use the other two.

The closed fibre of $D\to\operatorname{Spec}R$ is
zero-dimensional. Properness and upper semicontinuity of fibre
dimension imply that every fibre is zero-dimensional: otherwise
the closed image of the non-quasi-finite locus would contain
the closed point of the local base. Thus $D$ is finite over
$R$. Its normalization $\widetilde D$ is also finite over $R$
by excellence. A finite algebra over a henselian local ring is
a product of local algebras; because $\widetilde D$ is integral,
it is local. It is consequently the spectrum of a DVR, with
residue field $k$ and unique closed point over $z$.
A uniformizer has divisor equal to that point. Proper
pushforward of this principal divisor proves $[z]=0$ in
$\mathrm{CH}_0(Z)$. Hence
\begin{equation}\label{eq:C3ab-CH0-resolution}
                          \mathrm{CH}_0(Z)=0.
\end{equation}

\emph{6. Codimension two needs no denominator in Riemann--Roch.}
For a smooth variety $T$, let $F^rK_0(T)$ be the filtration by
codimension of coherent support. The rank and determinant
identifications give $F^2K_0(T)=SK_0(T)$, and
\citet[Example 15.3.6, p.~298]{frC4Fulton1998}
gives
\[
 \mathrm{CH}^2(T)\xrightarrow{[D]\mapsto[\mathcal O_D]}
                 F^2K_0(T)/F^3K_0(T)
\]
as an integral isomorphism, inverse to $-c_2$. The two composites
have multiplier $(-1)^{2-1}(2-1)!=-1$, so rationalization is
unnecessary.

These identities also hold on the particular $Z$ used here.
Indeed, $Z$ is the filtered inverse limit of the smooth
threefolds $W_T$ over pointed affine \emph{\'etale}
neighbourhoods of the cone vertex constructed in C.3e.
The transition maps are flat. Vector bundles, their exact
sequences, cycles, and the finite rational-equivalence data
all descend to one such neighbourhood. The support filtrations,
rank, determinant, Chern classes and the displayed isomorphism
therefore pass to this limit. This supplies the precise scope
needed here without invoking a Riemann--Roch theorem for
arbitrary nonalgebraizable regular schemes.

Since $Z$ has dimension three, $F^3K_0(Z)$ is generated by
classes of its closed points. It vanishes by
\eqref{eq:C3ab-CH0-resolution}. Since $V$ is a surface,
$F^3K_0(V)=0$ as well. Compatibility of Chern classes with
pullback and \eqref{eq:C3ab-SK0} now give
\begin{equation}\label{eq:C3ab-CH2-resolution}
 i^!: \mathrm{CH}^2(Z)\xrightarrow{\sim}\mathrm{CH}_0(V),
                      \qquad (i^!)^{-1}=p^*.
\end{equation}

\emph{7. Localization gives the actual ordinary quotient.}
Chow localization gives
\[
 \mathrm{CH}^1(V)\xrightarrow{i_*}\mathrm{CH}^2(Z)
       \longrightarrow\mathrm{CH}^2(X)\longrightarrow0.
\]
Under \eqref{eq:C3ab-CH2-resolution}, self-intersection
identifies the first arrow with multiplication by
$c_1(N_{V/Z})=-c_1(L)$. This proves the first isomorphism
in \eqref{eq:C3ab-answer}.

Flat pullback by $p$, followed by restriction to $X$, sends
ordinary algebraically trivial cycles on $V$ to ordinary
algebraically trivial cycles on $X$: it pulls back each actual
curve-parameter correspondence on the fixed product. Thus the
first isomorphism induces a surjection
\[
 B^2(V)/c_1(L)B^1(V)\longrightarrow B^2(X).
\]
The integral injection already proved in C.3e is this same map,
so it is an isomorphism. On a smooth projective integral
variety over an algebraically closed field, zero-cycles modulo
ordinary algebraic equivalence are classified by degree.
Consequently the last quotient is $\mathbb Z/g\mathbb Z$.
Ampleness makes every degree in the defining gcd positive,
and any hyperplane curve supplies one such degree; hence
$g>0$ and the group is finite.

\emph{8. The field restriction and examples are explicit.}
Akizuki--Kodaira--Nakano vanishing gives
\[
 H^2(V,\Omega^1_{V/k}\otimes L^j)=0\qquad(j\ge1)
\]
because $2+1>\dim V=2$
\citep[Theorem 1.3(a), p.~4; Corollary 6.4 and its proof,
 pp.~55--56]{D_EsnaultViehweg1992}.
If $k=\overline{\mathbb Q}$, absolute and relative differentials
agree. In general smoothness supplies the exact sequence
\[
 0\longrightarrow\Omega^1_{k/\mathbb Q}\otimes_k\mathcal O_V
 \longrightarrow\Omega^1_{V/\mathbb Q}
 \longrightarrow\Omega^1_{V/k}\longrightarrow0.
\]
Thus $H^2(V,L^j)=0$ is a sufficient additional condition.
Cohomology here commutes with the direct sum defining the
possibly infinite-dimensional constant $k$-vector space
$\Omega^1_{k/\mathbb Q}$, as follows from a finite affine
\v Cech cover. For a $K3$ surface this vanishing follows
from $K_V=\mathcal O_V$ and Serre duality, since $L^{-j}$
has no sections. If $p_g=0$, multiplication by a nonzero
section of $L^j$ injects
$H^0(V,K_V\otimes L^{-j})$ into $H^0(V,K_V)=0$,
and Serre duality again applies. The absolute-differential
hypothesis has therefore been verified in the claimed cases,
without replacing it silently by a relative one.

For example, $V=\mathbb P^2_k$ and $L=\mathcal O(d)$ give
$B^2(X)=\mathbb Z/d\mathbb Z$ for every $d\ge1$.
For the Fermat quartic surface in C.3f, choose
$\zeta^4=\eta^4=-1$. The equations
$x_0=\zeta x_1$ and $x_2=\eta x_3$ define a line on the
surface. Its hyperplane degree is one, so $g=1$ and
$B^2(X)=0$ for that particular complex strict local cone.
This is compatible with the non-torsion class of C.3f in
$\mathrm{CH}^2(X)$: its image is ordinarily algebraically
trivial. The two equivalence relations are kept distinct.
\end{proof}

\paragraph{C.3ac (Known finite-field theorem and proved deductions; the strengthened cycle-class conclusion uses the cited 2026 input).}
Let $k=\overline{\mathbb F}_p$, let $A/k$ be an abelian variety of
dimension $g$, and let $h$ be any ample divisor class. Independently
of the 2026 input in D.4c--e, we have
\begin{equation}\label{eq:C3ac-classical}
 \begin{gathered}
 \mathrm{CH}^r(A)_{\mathbb Q,(s)}=0\quad(s\ne0),\qquad
 J^r(A)=0,\\
 h^{g-2r}:\mathrm{CH}^r(A)_{\mathbb Q}
                 \xrightarrow{\ \sim\ }
                 \mathrm{CH}^{g-r}(A)_{\mathbb Q}
                 \quad(2r\le g).
 \end{gathered}
\end{equation}
Consequently both abelian injectivity assertions of C.3h hold
over this field; the first map is an isomorphism. The letter's
$A$, $A'$ and $A''$ hold for these abelian varieties, with the
same rational coefficients and degree ranges as in C.1.

If we additionally invoke Theorem I of
\citet{D_VanHoftenSempliner2026}, with exactly the source qualification
in D.4c--e, then for every $r$ we also have
\begin{equation}\label{eq:C3ac-stronger}
 \mathrm{CH}^r(A)_{\mathbb Q}
       \xrightarrow{\ \sim\ }
       \mathrm{CH}^r(A)_{\mathrm{num},\mathbb Q},
 \qquad G^r(A)=0.
\end{equation}
Under this same input the letter's strict $B$ holds for $A$.
The analogous rational--numerical comparison holds on every
product of abelian varieties and, more generally, on every
smooth projective variety over a finite field whose \emph{Chow}
motive is of abelian type. In the category generated by these
Chow motives, numerical passage is an equivalence onto the
corresponding semisimple numerical category. This last assertion
does not include correspondences from arbitrary varieties
outside that category.

\begin{proof}
\emph{The unconditional statement.}
Choose a finite-field model $A_0/\mathbb F_q$ of $A$.
\citet[Theorem 7.1, pp.~99--100]{cyKunnemann1993} proves
\[
       \mathrm{CH}^r(h^i(A_0))_{\mathbb Q}=0
                    \qquad(i\ne2r).
\]
The same theorem applies after every finite extension of
$\mathbb F_q$. A cycle and a rational-equivalence witness
over $k$ involve finitely many coefficients, so ordinary Chow
groups commute with this filtered union of finite fields.
The canonical multiplication projectors commute with base
change, and the index convention is $i=2r-s$. This gives the
first assertion in \eqref{eq:C3ac-classical}. The middle
assertion is C.5a: every algebraically trivial class is torsion.

The ample line bundle also descends after a finite extension.
\citet[Theorem 7.3, p.~100]{cyKunnemann1993} proves the
displayed hard Lefschetz isomorphism over each such finite
field. Passing to the same filtered union proves it over $k$.
Theorem 7.2 is \emph{not} this result: the source's item 7.2
is a definition of purity. The original result appears in
\citet[\S3.4, Theorem 7, pp.~339--341]{cySoule1984}.
Since $J^r(A)=0$, the quotient $B^r(A)$ agrees with the
ordinary rational Chow group. This proves both C.3h assertions.

For completeness, let $i:Y\hookrightarrow A$ be any smooth
divisor representing $h$, and take $z\in\mathrm{CH}^{r+1}(A)_{
\mathbb Q}$ with $2r\ge g-1$. Hard Lefschetz writes
\[
 z=h^{2r+2-g}a,\qquad
 a\in\mathrm{CH}^{g-r-1}(A)_{\mathbb Q}.
\]
The exponent $2r+1-g$ is nonnegative, and the projection
formula gives
\[
 z=i_*i^*(h^{2r+1-g}a).
\]
This proves $A'$ in this case. Ordinary Chow localization
then gives $A''$ for $A\setminus Y$, and hence for its
quotient modulo algebraic equivalence as well. There is no
appeal here to a localization sequence for homological
equivalence.

\emph{The stronger conclusion and its additional input.}
Theorem I cited above, together with
\citet[Theorem 7.1, author p.~25]{D_Milne1999}, supplies the
Tate conjecture for every finite-field abelian variety and
every finite-field model of every product. We retain Milne's
cofinal-CM-field correction exactly as in D.4e. Tate
surjectivity passes through direct sums, Tate twists, and
algebraic projectors, so it holds for the stated Chow motives
of abelian type, including Artin factors after finite descent.

Now \citet[Definition 1, p.~978; Theorem 1.10, p.~982]{cyKahn2003}
proves rational equals numerical equivalence on each such
finite-field variety. The theorem requires both abelian type
and Tate; K\"unnemann's grading statement alone supplies
neither the missing grade-zero injectivity nor this conclusion.
The finite-dimensionality and nilpotence inputs used by Kahn
are explicit in his Theorem 1.4 and Proposition 1.5, p.~981.
His Lemma 1.6 and Proposition 1.7 then eliminate all nonunit
simple summands that could contribute a homologically trivial
map from the unit. Thus the application uses his full
rational-equivalence theorem, rather than merely a statement
about numerical and homological equivalence.

Let $z\in\mathrm{CH}^r(A)_{\mathbb Q}$ be numerically trivial.
Represent it over a finite extension $\mathbb F_{q^e}$.
Its intersections with all cycles defined over that finite
field vanish, because the corresponding geometric
intersection numbers vanish. Kahn's theorem makes the
representative zero over that finite field. This proves
the injectivity in \eqref{eq:C3ac-stronger}; surjectivity is
the definition of the numerical quotient. A homologically
trivial class is numerically trivial, so $G^r(A)=0$ follows.
The same argument handles every product. Finite data descend
for Chow projectors and morphisms too, giving full
faithfulness after passing to $k$. Splitting numerical
idempotents now produces the claimed equivalence of the
generated pseudo-abelian categories; numerical semisimplicity
is Jannsen's theorem
\citep[Theorem 1, pp.~448--449]{Jannsen1992}.

Finally, for $g=2m$, the primitive middle space modulo
algebraic equivalence identifies, by
\eqref{eq:C3ac-stronger}, with its space of primitive algebraic
cohomology classes. Under the same Hodge-CM input,
\citet[Theorem 3.3 and Corollary 3.4, pp.~607--608]{D_Milne2002}
give the Hodge standard conjecture for abelian varieties.
Thus $(-1)^m\deg(x^2)$ is strictly positive for every nonzero
primitive class, which is the letter's $B$. This positivity
step is stated separately from the classical hard Lefschetz
statement.
\end{proof}

\emph{Scope of the function-field citation.}
\citet[\S8, pp.~101--102]{cyKunnemann1993} relates the
hypercube assertion to negative Beauville vanishing. The
opened original
\citet[\S3.5.5, Corollary 3, p.~343]{cySoule1984}
has a precise constant-variety hypothesis: its varieties over
the function field are $Y_j=X_j\otimes K$ for varieties
$X_j$ over a finite field. We do not remove that hypothesis
when reading the shorter function-field description in
K\"unnemann's discussion. In particular, the proof here
does not settle C.3h for arbitrary nonconstant abelian
varieties over $\overline{\mathbb F_p(t)}$, nor does its
restricted semisimple category settle C.1n with unrestricted
source varieties.

% Final: full proof and exact-primary-source audits passed by root and frobenius.
\paragraph{C.3ad (Proved in the stated rational smooth-cone scope): the top quotient over algebraic numbers.}
Let $k=\overline{\mathbb Q}$, let $V/k$ be a smooth projective
integral variety of dimension $d\ge2$, and let $L$ be very ample.
Assume its full section ring is normal and generated in degree one.
Let $R$ be the strict henselization of the local ring of its affine
cone at the vertex, and put
\[
             X=\operatorname{Spec}R\setminus\{\mathfrak m_R\}.
\]
Then the natural maps give
\begin{equation}\label{eq:C3ad-top-cone}
 \mathrm{CH}^d(X)_{\mathbb Q}\simeq
 \frac{\mathrm{CH}_0(V)_{\mathbb Q}}
      {c_1(L)\mathrm{CH}^{d-1}(V)_{\mathbb Q}},
 \qquad B^d(X)_{\mathbb Q}=0.
\end{equation}
The equivalence relation defining $B^d$ remains ordinary
algebraic equivalence on actual fixed products. In dimensions
$d>2$ this statement does not assert integral finiteness;
it also does not concern arbitrary strict local singularities.

\begin{proof}
Use the resolution $Z\to\operatorname{Spec}R$, its exceptional
divisor $i:V\hookrightarrow Z$, and the projection $p:Z\to V$
from C.3e and C.3ab. Thus $Z$ is regular of dimension $d+1$,
$p i=1$, $Z\setminus V=X$, and
\[
 V_a=Z\times_R R/\mathfrak m_R^a
     =\operatorname{Spec}_V
           (\mathcal O_V\oplus L\oplus\cdots\oplus L^{a-1}).
\]
The multiplication on the right is truncated graded
multiplication. The pro cdh argument in C.3ab applies without
change in this dimension: $K_0(R)=K_0(R/\mathfrak m_R^a)=\mathbb Z$
and the first map is the graph of $[\mathcal O_Z]$.
Consequently
\begin{equation}\label{eq:C3ad-finite-detection}
               K_0(Z)\longrightarrow K_0(V_a)
               \quad\text{is injective for }a\gg0.
\end{equation}

\emph{The exact cyclic-homology input.}
For an arbitrary commutative $\mathbb Q$-algebra $A$, including
a nonregular algebra, the Eulerian decomposition satisfies
\begin{equation}\label{eq:C3ad-Hodge}
 HC_n^{(r)}(A)=0\quad(r>n\ge0),\qquad
 HC_n^{(n)}(A)=\Omega^n_{A/\mathbb Q}/d\Omega^{n-1}_{A/\mathbb Q}.
\end{equation}
These are \citet[Theorems 4.6.7--4.6.8,
pp.~151--152]{D_Loday1992}. The hypotheses are commutativity
and a ground ring containing $\mathbb Q$; no smoothness
assumption occurs. The displayed top component is obtained
by taking the cokernel of the Hochschild boundary and then
the Connes operator, as in the source's proof.
Thus \eqref{eq:C3ad-Hodge} does not apply the smooth HKR
isomorphism to $V_a$.

The relative Chern character preserves the Adams decomposition,
with the shift
\[
       K_q^{(r)}(A,I)_{\mathbb Q}
                    \simeq HC_{q-1}^{(r-1)}(A,I)
\]
for a nilpotent ideal $I$. We use the scheme-compatible
form of Cathelineau's theorem proved by
\citet[Theorem 6.1 and Corollary 6.2,
author p.~15]{D_CortinasHaesemeyerWeibel2009}.
The shift is explicitly also the one in their Theorem B.10.
The relative groups here are already uniquely divisible,
but we retain rational coefficients throughout this proof.
Since the projection $V_a\to V$ splits restriction,
the relative groups are kernels of the absolute restriction
maps. Equations \eqref{eq:C3ad-Hodge} therefore imply, for
the relative Zariski sheaves on $V$,
\begin{equation}\label{eq:C3ad-relative-sheaves}
 \mathcal K_q^{(d)}(V_a,V)_{\mathbb Q}=0\quad(0\le q<d),
 \qquad
 \mathcal K_d^{(d)}(V_a,V)_{\mathbb Q}
       \simeq\bigoplus_{j=1}^{a-1}
                         \Omega^{d-1}_{V/\mathbb Q}\otimes L^j.
\end{equation}
For $q=0$ use nilpotent invariance of $K_0$ for affine rings.

Here is the local calculation and its gluing in the last
isomorphism. Put $s=d-1$ and trivialize $L$, so the algebra
is $A[t]/t^a$. Its positive internal-degree-$j$ forms, for
$1\le j<a$, have the form
\[
              t^j\omega+t^{j-1}dt\wedge\eta,
 \qquad \omega\in\Omega_A^s,
        \quad\eta\in\Omega_A^{s-1}.
\]
Modulo exact forms this has the unique horizontal
representative
\begin{equation}\label{eq:C3ad-horizontal}
              t^j(\omega-j^{-1}d\eta).
\end{equation}
Indeed subtract $j^{-1}d(t^j\eta)$; the remaining exact
relations from a form containing $dt$ map to zero because
$d^2=0$. The relation $a t^{a-1}dt=0$ removes internal
degree $a$, and there are no higher degrees. Inclusion
of horizontal forms is consequently an isomorphism onto
this quotient. Under $t=u t'$ a horizontal $s$-form
transforms as $t^j\omega=u^j(t')^j\omega$.
The isomorphism therefore glues with precisely $L^j$,
without any additional conormal term.

\emph{The weight-$d$ part on the resolution.}
The relative descent spectral sequence, with the Adams
component just specified, is
\[
 H^b(V,\mathcal K_q^{(d)}(V_a,V)_{\mathbb Q})
                         \Longrightarrow K_{q-b}^{(d)}(V_a,V)_{\mathbb Q}.
\]
It follows either by taking the relative fibre of
\citet[Theorem 10.3]{D_ThomasonTrobaugh1990} and the compatible
Adams summand, or directly from the scheme comparison above.
Since $\dim V=d$, only $0\le b\le d$ occur. On the
diagonal $q=b$ all terms vanish by
\eqref{eq:C3ad-relative-sheaves}, except possibly
\[
 H^d\!\left(V,
      \bigoplus_{j=1}^{a-1}\Omega^{d-1}_{V/\mathbb Q}\otimes L^j\right).
\]
The latter vanishes as well. Indeed
$\Omega^1_{k/\mathbb Q}=0$, so these are relative
differentials over $k$, and Kodaira--Akizuki--Nakano
vanishing applies because
\[
                     (d-1)+d>d\qquad(d\ge2).
\]
The exact characteristic-zero version is
\citet[Theorem 1.3(a) and Section 8]{D_EsnaultViehweg1992}.
It follows that
\[
                     K_0^{(d)}(V_a,V)_{\mathbb Q}=0.
\]
The split projection to $V$, together with
\eqref{eq:C3ad-finite-detection}, now proves
\begin{equation}\label{eq:C3ad-weight-isomorphism}
              i^*:K_0^{(d)}(Z)_{\mathbb Q}
                   \xrightarrow{\sim}K_0^{(d)}(V)_{\mathbb Q},
                        \qquad(i^*)^{-1}=p^*.
\end{equation}

\emph{Chow groups and ordinary algebraic equivalence.}
Rational Chern character identifies the weight-$d$
summands in \eqref{eq:C3ad-weight-isomorphism} with
$\mathrm{CH}^d(Z)_{\mathbb Q}$ and
$\mathrm{CH}_0(V)_{\mathbb Q}$, respectively. Its application
to this $Z$ is justified by the same filtered pointed-etale
models used in C.3ab: each model is a smooth variety, and
perfect complexes, cycles, rational-equivalence data and
Chern classes descend to a finite stage. The usual
rational Chern character isomorphism on those models
\citep[Example 15.2.16, pp.~294--295]{frC4Fulton1998} therefore passes to
the limit and remains compatible with $i^*$ and $p^*$.

Chow localization for $V\subset Z$, and
$i^!i_*=c_1(L^{-1})$, give the first isomorphism
in \eqref{eq:C3ad-top-cone}. The maps are the same actual
flat pullbacks as in C.3e. Thus they preserve ordinary
algebraic equivalence, and C.3e also gives injectivity
on the corresponding ordinary quotients. In the top
degree on the connected smooth projective variety $V$,
ordinary equivalence is degree, so $B^d(V)=\mathbb Z$.
The degree of $c_1(L)\cdot C$ is positive for every
integral curve $C\subset V$. Such a curve exists by
hyperplane sections. Hence the image of
$c_1(L)B^{d-1}(V)$ is nonzero in $\mathbb Z$ and is
all of $\mathbb Q$ after tensoring with $\mathbb Q$.
This proves $B^d(X)_{\mathbb Q}=0$.

The last use of rational Chern character is precisely
why the argument, as stated, does not prove integral
finiteness for $d>2$. For $d=2$, the independent
codimension-two integral calculation in C.3ab gives
the stronger finite cyclic answer.
\end{proof}

% C.3ae: full proof independently audited by /root and /root/frobenius.
\paragraph{C.3ae (Refuted, with an explicit counterexample): ordinary $J^2$ restriction is not surjective in every dimension at least seven.}
For every $N\ge7$ we specify an excellent normal strictly
henselian local complex ring $R$ of dimension $N$, with
regular puncture $X$ and regular punctured principal divisor
$Y$, for which there is a surjection
\begin{equation}\label{eq:C3ae-cokernel}
 \operatorname{coker}\!\left(
 \mathrm{CH}^2(X)_{\mathrm{alg}}
       \longrightarrow\mathrm{CH}^2(Y)_{\mathrm{alg}}\right)
       \longrightarrow \operatorname{Pic}^0(C)(\mathbb C),
 \qquad C:\ z_0^4+z_1^4+z_2^4=0.
\end{equation}
Here the algebraically trivial groups use actual families on
the fixed scheme times a smooth connected parameter curve.
Consequently the ordinary rational map $J^2(X)\to J^2(Y)$
is not surjective. The argument concerns these ordinary
groups; it does not assert that every cycle on $X$ is
annihilated by the invariant constructed below.

\emph{1. The coefficients and the algebraic local scheme.}
Put $m=N-4\ge3$, and let
\[
 W=V(w_0^4+w_1^4+w_2^4+w_3^4)\subset\mathbb P^3,
 \qquad S=W\times C,
 \qquad M=\mathcal O_W(1)\boxtimes K_C.
\]
Write $A_i=(a_{ijk})_{0\le j\le3,\,0\le k\le2}$ for
the following four integer matrices:
\[
\begin{aligned}
 A_0&=\begin{pmatrix}-3&-5&5\\-3&-5&-2\\-5&3&-3\\5&1&-3\end{pmatrix},&
 A_1&=\begin{pmatrix}-4&4&-1\\3&5&2\\-5&0&-1\\0&2&-5\end{pmatrix},\\[2mm]
 A_2&=\begin{pmatrix}0&-1&-1\\0&-2&4\\-1&1&-4\\-5&-1&-3\end{pmatrix},&
 A_3&=\begin{pmatrix}0&-4&1\\-5&-4&-4\\-1&5&1\\4&2&5\end{pmatrix}.
\end{aligned}
\]
Define
\[
 s_i(w,z)=\sum_{j=0}^3\sum_{k=0}^2a_{ijk}w_jz_k,
 \qquad
 \ell(x,w,z)=\sum_{i=0}^3x_i s_i(w,z),
 \qquad
 G(z)=z_0^2z_1^2+z_0^2z_2^2+z_1^2z_2^2.
\]
For $i>3$ we set $s_i=0$. The four $s_i$ have no
common zero on $S$. This is an exact finite verification:
over $\mathbb F_{101}$, in each of the twelve affine charts
$w_j=z_k=1$, the ideal generated by the two Fermat equations
and the four $s_i$ has Groebner basis $[1]$. The common-zero
scheme is projective over $\mathbb Z$, so its image in
$\operatorname{Spec}\mathbb Z$ is closed. A nonempty generic
fibre would force that image to contain every prime, contrary
to the computation at $101$. Thus this particular integer
tuple, without a generic-coefficient qualification, is
base point free over $\mathbb C$.

Set
\[
 B=\mathbb P^m_x\times W\times\mathbb P^2_z,
 \qquad
 \mathcal N=\mathcal O_{\mathbb P^m}(-1)
       \boxtimes\mathcal O_W(-1)\boxtimes\mathcal O_{\mathbb P^2}(-1),
 \qquad T_B=\operatorname{Tot}(\mathcal N).
\]
If $t$ is the tautological fibre vector, define the hypersurface
\[
                   Z=V(F+t\ell G)\subset T_B,
                   \qquad F=z_0^4+z_1^4+z_2^4.
\]
Its zero-section divisor is
\[
 V=\mathbb P^m\times W\times C,
 \qquad N_{V/Z}=L^{-1},
 \qquad L=\mathcal O_{\mathbb P^m}(1)\boxtimes M.
\]
The expression defining $Z$ is a section of the pullback of
$\mathcal O_B(0,0,4)$, so the definition is independent of
trivializations. The nonzero differential $dF$ makes $Z$
smooth near $V$. On $\ell G\ne0$ it is the graph
$t=-F/(\ell G)$. The locus $F=\ell G=0$ has codimension two
in $B$ and produces no additional component of the
hypersurface. It is a Cartier hypersurface in a regular
scheme, hence Cohen--Macaulay, and it is generically
reduced on the graph. It follows that $Z$ is integral.

Contract the zero section of $T_B$ by the Segre coordinates
\[
                         y_{ijk}=t x_iw_jz_k.
\]
This is a proper morphism to the affine cone of the indicated
Segre embedding, and it is an isomorphism off the zero
section. Let $Q$ be the reduced closed image of $Z$, let
$Q^{\mathrm n}$ be its finite normalization, and let
$v^{\mathrm n}$ be the point over the cone vertex. This
point is unique: after shrinking around the vertex, the
smooth proper birational model $Z$ factors through the
normalization and is surjective; its connected exceptional
fibre $V$ cannot map onto more than one point of the finite
normalization fibre. Define
\[
 R=\mathcal O_{Q^{\mathrm n},v^{\mathrm n}}^{\mathrm{sh}}.
\]
This specifies $R$ by the displayed integer equations,
closed image, normalization, and strict henselization.
As in C.3m, properness allows us to shrink only to prove
regularity of the model; this shrinking makes no change
to $R$. Write $\mathcal Z\to\operatorname{Spec}R$ for
the resulting projective regular model. Its exceptional fibre
is $V$, and $X=\mathcal Z\setminus V$ is the regular
puncture of $\operatorname{Spec}R$. The ring has dimension
$m+4=N$ and is excellent and normal. Pointed etale base
change and its limit preserve these assertions and the
first infinitesimal neighbourhood of $V$. More explicitly,
the $y_{ijk}$ generate $(t)$ on each Segre chart, while
every extra normalization function in the maximal ideal
vanishes on the reduced fibre. Thus
$\mathfrak m_R\mathcal O_{\mathcal Z}=\mathcal I_V$,
as in C.3m; the quotients by the squares of these ideals
are unchanged by a pointed etale extension with the same
residue field.

The nonzero nonunit
\[
                      f=\sum_{i,j,k}a_{ijk}y_{ijk}
\]
pulls back to $t\ell$. Its strict transform is
$D=V(F,\ell)\subset T_B$, the total space of $L_H^{-1}$
over
\[
 H=V(\ell)\subset V,
 \qquad L_H=L|_H.
\]
The surjection $\mathcal O_S^{m+1}\to M$ with entries
$s_i$ has a rank-$m$ vector-bundle kernel $E$, and
\[
             H=\mathbb P_{\mathrm{lines}}(E),
             \qquad c_1(L_H)=h+a,
             \qquad a=c_1(M),\quad h=c_1(\mathcal O_H(1)).
\]
Thus $H$ and $D$ are smooth. Their corresponding base
change $\mathcal D\subset\mathcal Z$ is regular, and
$Y=\mathcal D\setminus H=X\cap V(f)$ is the required
regular punctured divisor. The divisors $V$ and
$\mathcal D$ intersect transversely in $H$.

\emph{2. A cycle invariant, with its family compatibility.}
Fix $\eta\in\overline{\mathbb Q}$ with $\eta^4=-1$ and
the line
\[
 D_W=V(w_0-\eta w_1,\,w_2-\eta w_3)\subset W.
\]
Writing $H_W=c_1(\mathcal O_W(1))$, put
$\tau=4[D_W]-H_W$. Adjunction and the degree of the line
give
\begin{equation}\label{eq:C3ae-intersections}
 H_W^2=4,\qquad H_W[D_W]=1,\qquad [D_W]^2=-2,
 \qquad \tau H_W=0,\qquad \tau[D_W]=-9.
\end{equation}

For a smooth projective parameter $P$, Chow localization
and refined pullback along $H\times P\subset\mathcal D\times P$
define a homomorphism
\[
 \sigma_P:\mathrm{CH}^2(Y\times P)
 \longrightarrow
 \frac{\mathrm{CH}^2(H\times P)}
 {(h+a)\mathrm{CH}^1(H\times P)}.
\]
Indeed we lift a class to $\mathrm{CH}^2(\mathcal D\times P)$
and pull back. Two lifts differ by the pushforward of a
class in $\mathrm{CH}^1(H\times P)$, and the
self-intersection formula makes their pullbacks differ
by $c_1(N_{H/\mathcal D})=-(h+a)$ times that class.
This proves independence, including for rational-equivalence
witnesses. These are the same localization and
self-intersection operations used in C.3e.

Since $\operatorname{rank}E=m\ge3$, the first
projective-bundle relation has degree at least three.
Consequently evaluation at $h=-a$ gives an integral
isomorphism
\begin{equation}\label{eq:C3ae-evaluation}
 \frac{\mathrm{CH}^2(H\times P)}
 {(h+a)\mathrm{CH}^1(H\times P)}
                   \simeq \mathrm{CH}^2(S\times P).
\end{equation}
Intersect the right side with $\operatorname{pr}_W^*\tau$
and push forward along $W$. We obtain
\begin{equation}\label{eq:C3ae-rho}
 \rho_P:\mathrm{CH}^2(Y\times P)
       \longrightarrow \mathrm{CH}^1(C\times P)
       =\operatorname{Pic}(C\times P).
\end{equation}
All the operations commute with refined pullback to a
point of $P$. For example, a lift chosen before taking
the fibre is still a valid lift after taking it, and
refined pullbacks commute. Hence $\rho=\rho_{\operatorname{Spec}\mathbb C}$
sends ordinary algebraic-equivalence differences to
$\operatorname{Pic}^0(C)(\mathbb C)$.

If $\Gamma\in\mathrm{CH}^2(X\times P)$ and
$\widetilde\Gamma$ is a closure on $\mathcal Z\times P$,
write its refined restriction to $V\times P$ as
\[
 b=b_0+h b_1+h^2 b_2,
 \qquad b_i\in\mathrm{CH}^{2-i}(S\times P).
\]
Transverse refined pullback to $H$ and
\eqref{eq:C3ae-evaluation} give the exact formula
\begin{equation}\label{eq:C3ae-source-invariant}
 \rho_P(\Gamma|_{Y\times P})
   =\operatorname{pr}_{C\times P,*}
            \bigl(\tau(b_0-a b_1+a^2b_2)\bigr).
\end{equation}
There is no presumed extension of a cycle over a varying
etale neighbourhood in this formula: the parameter $P$
is fixed, and the closure is on the fixed proper model
$\mathcal Z\times P$.

\emph{3. The first-neighbourhood obstruction for a cycle class.}
We give the required Chern-class argument explicitly.
All differential sheaves in this argument are relative
to $\mathbb C$.
Let $V_2\subset\mathcal Z$ be defined by $\mathcal I_V^2$.
Its square-zero ideal is $L$, and its class is
$\kappa\in H^1(V,T_V\otimes L)$. For any smooth
projective $P$, the conormal sequence of $V\times P$
in $\mathcal Z\times P$ and its second exterior power are
\[
\begin{gathered}
 0\longrightarrow L\longrightarrow
 \Omega^1_{\mathcal Z\times P}|_{V\times P}
 \longrightarrow\Omega^1_{V\times P}\longrightarrow0,\\
 0\longrightarrow L\otimes\Omega^1_{V\times P}
 \longrightarrow\Omega^2_{\mathcal Z\times P}|_{V\times P}
 \longrightarrow\Omega^2_{V\times P}\longrightarrow0.
\end{gathered}
\]
The second sequence is exact because the kernel of
the first has rank one. Its cohomology boundary is
contraction with $\kappa$.
This can be checked by taking local splittings of the
first sequence: their differences are derivations
representing $\kappa$, and the induced difference on
a two-form is contraction in either factor with the
usual alternating sign.

For a vector bundle on $\mathcal Z\times P$, its Hodge
Chern character restricts through the middle term of
the second sequence. Concretely, on a trivializing cover
the matrices $g_{uv}^{-1}dg_{uv}$ represent the Atiyah
class, and $\tfrac12\operatorname{tr}(\operatorname{At}^2)$
represents its degree-two Chern character. These matrices
are defined before restriction, so their trace cup
products give the asserted lift. The formula is
compatible with direct sums and exact sequences; on a
flag bundle it reduces to $\tfrac12\sum c_1(L_i)^2$,
which identifies it with the usual Hodge realization
of the Chow Chern character. Thus it also applies to
virtual vector bundles and rational combinations.

For our closure class we may choose such a rational
virtual bundle whose degree-two Chow Chern character
is $\widetilde\Gamma$. Precisely, the regular scheme
$\mathcal Z\times P$ is finite dimensional, has affine
diagonal and an ample invertible sheaf, and the rational
Chern character is an isomorphism onto its Chow group
\citep[Lemma 42.58.1, Tag 0FEY]{cyStacksChernObstruction}.
One may equally descend this finite collection of cycles,
bundles and maps to a pointed etale neighbourhood and
work on a smooth finite-type complex model; its restriction
to $V_2\times P$ is the same. Naturality of Chern classes
under refined pullback
\citep[Remark 42.59.9, Tag 0FF8]{cyStacksChernObstruction}
therefore gives
\begin{equation}\label{eq:C3ae-obstruction-zero}
 \kappa\mathbin{\lrcorner}\operatorname{cl}_{\mathrm{Hdg}}(b)
       =0
 \quad\hbox{in }H^3(V\times P,L\otimes\Omega^1_{V\times P}).
\end{equation}
Only this necessary vanishing is used. We make no
converse assertion about deforming cycles from their
Hodge classes.

The class $\kappa$ has only a curve-tangent component,
since the construction has a morphism to
$\mathbb P^m\times W$. The normal sequence calculation
from C.3m gives
\begin{equation}\label{eq:C3ae-kappa}
 \kappa=\sum_{i,j,k}a_{ijk}x_iw_j\otimes\delta(z_kG),
 \qquad
 \delta:H^0(C,\mathcal O_C(5))\longrightarrow H^1(C,\mathcal O_C).
\end{equation}
Up to the same harmless overall sign, the three values
$\delta(z_kG)$ are a basis. In the Fermat Jacobian ring
they are respectively
$z_0z_1^2z_2^2$, $z_0^2z_1z_2^2$, and
$z_0^2z_1^2z_2$. The source basis of $H^0(C,K_C)$
used below is the Serre-dual basis to these three
values; it need not be identified with a chosen
coordinate basis of that space.

\emph{4. The exact rank calculation and the Hodge projection.}
Put
\[
 \mathcal R=\mathbb C[w_0,w_1,w_2,w_3]/(w_0^3,w_1^3,w_2^3),
 \qquad \dim\mathcal R_3=16,\quad\dim\mathcal R_4=19.
\]
The Euler and normal sequences on the quartic give
\[
 H^1(W,T_W(-1))\simeq\mathcal R_3,
 \qquad
 H^1(W,T_W)_{\mathrm{pol}}\simeq\mathcal R_4.
\]
For completeness, $H^1(T_{\mathbb P^3}(-1)|_W)=0$,
and quotienting $H^0(\mathcal O_W(3))$ by the four
partial derivatives gives the first equality. In degree
zero, quotienting $H^0(\mathcal O_W(4))$ by linear
vector fields gives the $19$-dimensional kernel of
\[
                   H^1(T_W)\longrightarrow H^1(T_{\mathbb P^3}|_W).
\]
The latter
target has dimension one by the Euler sequence.
The K3 equalities $K_W\simeq\mathcal O_W$ and
$h^{1,1}(W)=20$ identify this kernel with the
polarized deformation space. Its Serre-dual annihilator
in $H^1(W,\Omega^1_W)$ is exactly $\mathbb C H_W$:
the polarization extends to every embedded deformation,
and its class is nonzero since $H_W^2=4$.

Serre duality makes
\[
 \frac{H^1(W,\Omega^1_W)}{\mathbb C H_W}
      \simeq\mathcal R_4^*,
 \qquad H^1(W,\Omega^1_W(1))\simeq\mathcal R_3^*.
\]
Multiplication by $w_j$ on these Hodge groups is dual
to multiplication $w_j:\mathcal R_3\to\mathcal R_4$.
In particular every such multiplication kills $H_W$.
The displayed identifications use dual vector spaces,
not an unmentioned self-dual monomial basis.

Order the degree-three and degree-four monomials
lexicographically by their exponent tuples in
$\{0,1,2\}^4$. In the dual bases, the map
\begin{equation}\label{eq:C3ae-rank-map}
 \mathcal R_4^*\otimes\mathbb C^3
       \longrightarrow\mathbb C^4\otimes\mathcal R_3^*
\end{equation}
induced by \eqref{eq:C3ae-kappa} has rows $(i,u)$,
columns $(k,v)$, and entry
\[
                  M_{(i,u),(k,v)}
                     =\sum_{\{j:\,u+e_j=v\}} a_{ijk}.
\]
It is a $64$ by $57$ integer matrix. The determinant of
the submatrix consisting of its first $57$ rows is
$47$ modulo $101$. Hence \eqref{eq:C3ae-rank-map} is
injective over $\mathbb C$. The retained code records
the monomial order, matrices and exact finite-field
elimination, and independent implementations reproduce
both this minor and full column rank.

Now project \eqref{eq:C3ae-obstruction-zero} to
\begin{equation}\label{eq:C3ae-target-component}
 \begin{split}
 H^0(\mathbb P^m,\mathcal O(1))
  \otimes H^1(W,\Omega^1_W(1))
  \otimes H^1(C,K_C)
  \otimes H^1(P,\mathcal O_P).
 \end{split}
\end{equation}
Because $\kappa$ contracts only a curve differential,
the relevant source component of $\operatorname{cl}(b_0)$ is
\begin{equation}\label{eq:C3ae-source-component}
 H^1(W,\Omega^1_W)\otimes H^0(C,K_C)
                          \otimes H^1(P,\mathcal O_P).
\end{equation}
Terms with a projective-space differential cannot
contribute to \eqref{eq:C3ae-target-component}. They
also vanish by the Euler-sequence equality
\[
                       H^1(\mathbb P^m,\Omega^1(1))=0.
\]
The action on
\eqref{eq:C3ae-source-component}, modulo $\mathbb C H_W$
in the first factor, is exactly the injective map
\eqref{eq:C3ae-rank-map} tensored with
$H^1(P,\mathcal O_P)$. Thus the primitive part of
\eqref{eq:C3ae-source-component} is zero.

Pairing its first factor with $\tau$ annihilates the
remaining $H_W$ part by \eqref{eq:C3ae-intersections}.
It follows that the line bundle
$\operatorname{pr}_{C\times P,*}(\tau b_0)$ has zero
mixed component in $H^0(C,K_C)\otimes H^1(P,\mathcal O_P)$.
The correction $-a b_1+a^2b_2$ in
\eqref{eq:C3ae-source-invariant} has no effect on this
conclusion. Indeed $H^1(W,\mathcal O_W)=0$ and
$H^0(W,\Omega_W^1)=0$, so the only varying mixed
component of a divisor on $S\times P$ comes from
$C\times P$. Multiplication by $a$ supplies the
$H_W$ factor, which pairs to zero with $\tau$.
All other divisor components are constant in $P$;
$b_2$ has codimension zero and is constant as well.
This verifies the boundary correction rather than
assuming that a chosen lift is unique.

\emph{5. Actual ordinary families and the nonzero target.}
Every generator of $\mathrm{CH}^2(X)_{\mathrm{alg}}$ is
a difference of fibres of a cycle on $X$ times a
smooth connected curve. Compactify the curve smoothly
and close that cycle in the product. The fibres at
the original points are unchanged. Thus it suffices
to take the parameter $P$ above to be a smooth
projective connected curve and a fixed actual cycle
on $X\times P$.

The line bundle in \eqref{eq:C3ae-source-invariant}
defines, after subtracting its value at a base point,
a morphism $P\to\operatorname{Pic}^0(C)$. Its mixed
Hodge component is zero by step 4. By the universal
property of $J(P)$, this morphism factors through a
homomorphism $J(P)\to\operatorname{Pic}^0(C)$, whose
map on $H^1(\mathcal O)$ is exactly that mixed
component, as in the Poincare calculation in C.3m
\citep[Theorem 1.2 and Lemma 9.5]{D_MilneJacobians2021}.
A homomorphism of complex abelian varieties
with zero such map is zero: complex conjugation gives
zero on holomorphic differentials, hence zero
differential, and a homomorphism with zero differential
is constant. The classifying morphism is therefore
constant. Fibre compatibility of $\rho_P$ proves
\begin{equation}\label{eq:C3ae-source-zero}
                \rho\bigl(\mathrm{CH}^2(X)_{\mathrm{alg}}|_Y\bigr)=0.
\end{equation}
This proof uses fixed proper parameter curves. It
does not require extending a class from a varying
etale neighbourhood across a compactified parameter.

Pulling the divisors $D_W$ and a divisor on $C$ to
$Y$ gives a homomorphism
\[
 \theta:\operatorname{Pic}^0(C)(\mathbb C)
               \longrightarrow \mathrm{CH}^2(Y)_{\mathrm{alg}}.
\]
Its image is ordinarily algebraically trivial,
as is seen from the Poincare line bundle on
$C\times\operatorname{Pic}^0(C)$: two points of that
connected projective abelian variety can be joined
by a chain of smooth connected curves, and we restrict
the actual family to those curves. Taking the corresponding
extension on $\mathcal D$ and applying the definition
of $\rho$ yields the integral identity
\[
                         \rho\theta=-9\,\mathrm{id}.
\]
Multiplication by $9$ is surjective on the complex
points of an abelian variety. Equations
\eqref{eq:C3ae-source-zero} and the last identity prove
\eqref{eq:C3ae-cokernel}.

For one specified nontorsion witness, choose
$u\in\mathbb C$ transcendental over
$\overline{\mathbb Q}$, choose $v$ with
$v^4=-1-u^4$, and set
\[
                       p=(u:v:1),\qquad o=(1:\eta:0).
\]
Then $[p-o]\in\operatorname{Pic}^0(C)(\mathbb C)$ is
nontorsion: otherwise its Abel point would be a
torsion point defined over $\overline{\mathbb Q}$,
and the Abel closed immersion, defined over that
field, would force $p$ to be algebraic
\citep[Proposition 2.3]{D_MilneJacobians2021}. The target
cycle $\theta([p-o])$ therefore has nonzero image
$-9[p-o]$ under $\rho$ and cannot lie in the ordinary
source image, even after tensoring with $\mathbb Q$.

\emph{Finite certificates and proof scope.}
The rank calculation took $0.048689819$ seconds and
the twelve-chart base point calculation took
$4.204962002$ seconds in the retained implementations.
Independent checks used different matrix and Groebner
basis algorithms; in particular the $57$ by $57$
minor was computed independently. Those computations
certify the stated finite linear algebra and absence
of base points. The localization, Chern obstruction,
boundary correction and proper-family argument above
are the separate reasons that these data give the
claimed ordinary Chow-group counterexample.

% C.3af: full proof and compressed section construction independently audited.
\paragraph{C.3af (Refuted, with explicit counterexamples): ordinary restriction fails in every codimension in an explicit unbounded dimension range.}
Let $r\ge2$ and
\[
                         N\ge 4^{r-1}+2r-1.
\]
There is a completely specified excellent normal strictly
henselian local complex ring $R$ of dimension $N$, with
regular puncture $X$ and regular punctured principal divisor
$Y$, such that
\begin{equation}\label{eq:C3af-cokernel}
 \operatorname{coker}\!\left(
 \mathrm{CH}^r(X)_{\mathrm{alg}}
       \longrightarrow\mathrm{CH}^r(Y)_{\mathrm{alg}}\right)
          \longrightarrow\operatorname{Pic}^0(C)(\mathbb C)
\end{equation}
is surjective. Here $C$ is the Fermat plane quartic and
the equivalence relation is ordinary algebraic equivalence
through actual families on the fixed product. In particular,
the ordinary rational $J^r$ restriction is not surjective.
For $r=2$ this recovers the threshold $N\ge7$ of C.3ae.
No conclusion in the intermediate dimensions follows from
the bound in this statement.

\emph{1. An explicit full set of sections.}
We first use the larger bound
$N\ge3\cdot4^{r-1}+2r-1$ and then reduce it in step 5.
Set $s=r-1$, $q=3\cdot4^s$ and $m=N-2r\ge q-1$.
Let $W_1,\ldots,W_s$ be copies of the Fermat quartic K3
surface from C.3ae. On each $W_i$, write
\[
 H_i=c_1(\mathcal O_{W_i}(1)),\qquad
 D_i=V(w^{(i)}_0-\eta w^{(i)}_1,
        w^{(i)}_2-\eta w^{(i)}_3),\qquad
 \tau_i=4[D_i]-H_i,
 \quad\eta^4=-1.
\]
Put
\[
 \mathcal W=\prod_{i=1}^s W_i,\qquad
 S=\mathcal W\times C,\qquad
 M=\left(\boxtimes_{i=1}^s\mathcal O_{W_i}(1)\right)\boxtimes K_C,
 \qquad a=c_1(M).
\]
The $q$ sections
\begin{equation}\label{eq:C3af-sections}
       t_{\mathbf j,k}
          =z_k\prod_{i=1}^s w^{(i)}_{j_i},
 \qquad \mathbf j\in\{0,1,2,3\}^s,\quad k\in\{0,1,2\},
\end{equation}
form a basis of $H^0(S,M)$ and have no common zero.
Indeed at any point we choose one nonzero coordinate in
each projective factor and take their product. Give
$\mathbb P^m$ coordinates $x_{\mathbf j,k}$ for these
$q$ indices and, if needed, additional coordinates which
will not occur in
\[
                      \ell=\sum_{\mathbf j,k}
                            x_{\mathbf j,k}t_{\mathbf j,k}.
\]
Every coefficient is now zero or one; no generic choice
of a coefficient or a degree remains.

Take
\[
 B=\mathbb P^m\times\mathcal W\times\mathbb P^2,
 \qquad
 \mathcal N=\mathcal O_{\mathbb P^m}(-1)
       \boxtimes\left(\boxtimes_i\mathcal O_{W_i}(-1)\right)
       \boxtimes\mathcal O_{\mathbb P^2}(-1).
\]
In $\operatorname{Tot}(\mathcal N)$ define
\[
 Z=V(F+t\ell G),\qquad
 F=z_0^4+z_1^4+z_2^4,\qquad
 G=z_0^2z_1^2+z_0^2z_2^2+z_1^2z_2^2.
\]
Contract by all Segre coordinates
$t x_\alpha z_k\prod_i w^{(i)}_{j_i}$, take the reduced
closed image, normalize it, localize at the unique point
over the vertex, and take its strict henselization.
This defines $R$. The graph, Cohen--Macaulay,
normalization and scheme-theoretic exceptional-fibre
arguments in C.3ae apply verbatim to this explicitly
given smooth product. They give a projective regular
model $\mathcal Z$ with exceptional divisor
\[
 V=\mathbb P^m\times S,
 \qquad N_{V/\mathcal Z}=L^{-1},
 \qquad L=\mathcal O_{\mathbb P^m}(1)\boxtimes M.
\]
Its dimension is $m+2s+2=N$. In particular the
finite normalization has one vertex, and
$\mathfrak m_R\mathcal O_{\mathcal Z}=\mathcal I_V$;
the first neighbourhood is preserved by strict
henselization. The linear cone function $f=t\ell$ has
regular strict transform $\mathcal D$, with
\[
 H=V(\ell)\subset V,
 \qquad
 H=\mathbb P_{\mathrm{lines}}(E),
 \quad E=\ker(\mathcal O_S^{m+1}\to M),
 \quad\operatorname{rank}E=m.
\]
The map defining $E$ is surjective by
\eqref{eq:C3af-sections}. Its strict transform before
the local base change is the total space of $L_H^{-1}$.
Thus $X=\mathcal Z\setminus V$ and
$Y=\mathcal D\setminus H=X\cap V(f)$ have all the
asserted regularity properties.

\emph{2. The invariant in codimension $r$.}
The inequality $m\ge q-1>r$ implies that no
projective-bundle relation has degree at most $r$.
For every smooth projective parameter $P$, localization
and the refined self-intersection formula therefore give
\[
\begin{split}
 \sigma_P:\mathrm{CH}^r(Y\times P)&\longrightarrow
 \frac{\mathrm{CH}^r(H\times P)}
 {(h+a)\mathrm{CH}^{r-1}(H\times P)}\\
 &\simeq\mathrm{CH}^r(S\times P),\qquad h\longmapsto-a.
\end{split}
\]
The construction and independence are exactly those
proved in C.3ae, now with codimension $r$.
Intersect with $\tau=\prod_i\tau_i$ and push along
the $2s$-dimensional variety $\mathcal W$. Since
$r+s-2s=1$, this defines an integral homomorphism
\[
 \rho_P:\mathrm{CH}^r(Y\times P)
                \longrightarrow\operatorname{Pic}(C\times P)
\]
commuting with fibres. In particular $\rho$ maps the
ordinary algebraically trivial subgroup into
$\operatorname{Pic}^0(C)(\mathbb C)$.

For a closure of a cycle $\Gamma$ from $X\times P$,
write its refined restriction to $V\times P$ as
\[
 b=\sum_{j=0}^r h^j b_j,
 \qquad b_j\in\mathrm{CH}^{r-j}(S\times P).
\]
Transversality gives
\begin{equation}\label{eq:C3af-invariant}
 \rho_P(\Gamma|_{Y\times P})
       =\operatorname{pr}_{C\times P,*}
          \left(\tau\sum_{j=0}^r(-a)^j b_j\right).
\end{equation}
These refined pullbacks and pushforwards act on
actual Chow classes on the fixed product.

\emph{3. The only relevant Hodge component.}
All differential sheaves here are relative to
$\mathbb C$. The rank-one conormal sequence gives
\[
 0\longrightarrow L\otimes\Omega^{r-1}_{V\times P}
 \longrightarrow\Omega^r_{\mathcal Z\times P}|_{V\times P}
 \longrightarrow\Omega^r_{V\times P}\longrightarrow0.
\]
Its cohomology boundary is contraction with the
first-neighbourhood class $\kappa$. The rational
Chern-character lift from C.3ae applies in degree $r$:
the trace formula is
$\operatorname{ch}_r=\operatorname{tr}(\operatorname{At}^r)/r!$.
Consequently every closure restriction $b$ satisfies
\[
       \kappa\mathbin{\lrcorner}\operatorname{cl}_{\mathrm{Hdg}}(b)
                  =0\quad\hbox{in }
       H^{r+1}(V\times P,L\otimes\Omega^{r-1}_{V\times P}).
\]
Here the existence and pullback compatibility of the
virtual-bundle lift use the already opened
\citet[Lemma 42.58.1, Tag 0FEY;
Remark 42.59.9, Tag 0FF8]{cyStacksChernObstruction}.

The class $\kappa$ contracts only a $C$ differential;
it is
\[
          \kappa=\sum_{\mathbf j,k}x_{\mathbf j,k}
                   \left(\prod_i w^{(i)}_{j_i}\right)
                     \otimes\delta(z_kG).
\]
The three curve classes form the basis checked in
C.3ae. Project the vanishing to the Kunneth summand
\[
 H^0(\mathbb P^m,\mathcal O(1))\otimes
 \left(\bigotimes_{i=1}^sH^1(W_i,\Omega^1_{W_i}(1))\right)
 \otimes H^1(C,K_C)\otimes H^1(P,\mathcal O_P).
\]
Its cohomology degree is $s+1+1=r+1$ and its
differential degree is $s=r-1$. Since multiplication
by each $w^{(i)}_{j_i}$ preserves the cohomology and
differential degrees on $W_i$, the only source
component that can reach it is the component of $b_0$
in
\begin{equation}\label{eq:C3af-source}
 \left(\bigotimes_{i=1}^sH^1(W_i,\Omega^1_{W_i})\right)
         \otimes H^0(C,K_C)\otimes H^1(P,\mathcal O_P).
\end{equation}
All components involving a projective-space
differential are excluded. On a K3 factor the
$(2,0)$ and $(0,2)$ pieces have the wrong separate
degrees, and cannot enter this projection by
multiplication by a function. This excludes the
additional Kunneth pieces that appear on a product
of several K3 surfaces.

For each factor put
$U_i=H^1(W_i,\Omega^1_{W_i})$ and
$V_i=H^1(W_i,\Omega^1_{W_i}(1))$. C.3ae's Euler,
normal and Serre-duality calculation proves that
the joint multiplication map
\[
 U_i\longrightarrow\mathbb C^4\otimes V_i,
       \qquad u\longmapsto(w^{(i)}_0u,\ldots,w^{(i)}_3u)
\]
has kernel exactly $\mathbb C H_i$. Its quotient
is dual to the surjective multiplication
$\mathcal R_1\otimes\mathcal R_3\to\mathcal R_4$.
Surjectivity is elementary for the Fermat ring:
every degree-four monomial with all exponents at
most two is a variable times a surviving
degree-three monomial.

Because every section index in
\eqref{eq:C3af-sections} has its own independent
coordinate $x_{\mathbf j,k}$, the projected obstruction
is the tensor product of these joint multiplication
maps and the perfect curve pairing. Tensor products
over $\mathbb C$ preserve injections. Its kernel in
\eqref{eq:C3af-source} is therefore the sum of the
subspaces having $\mathbb C H_i$ in at least one
K3 factor. Pairing with $\tau_i$ in every factor
kills this entire kernel, since $\tau_iH_i=0$.
Thus the mixed $H^0(C,K_C)\otimes H^1(P,\mathcal O_P)$
component of $\operatorname{pr}_{C\times P,*}(\tau b_0)$
is zero.

We also check every correction term in
\eqref{eq:C3af-invariant}. The class $a$ is
$\sum_iH_i+c_1(K_C)$. To contribute to the specified
mixed Picard component after intersection with
$\prod_i\tau_i$, the class before that intersection
must have cohomology and differential degrees $(1,1)$
on every $W_i$, and degrees $(0,1)$ on $C$.
A factor $c_1(K_C)$ already has curve cohomology
degree one, so it cannot occur. A factor $H_i$
requires the rest of the tensor to have degrees
$(0,0)$ on that factor; its pairing is then
$\tau_iH_i=0$. Two such factors on one surface
have too large a degree. Thus every term with
$j\ge1$ contributes zero to the mixed component,
including terms whose $b_j$ contains nontrivial
cycles on several K3 factors. The line bundle
in \eqref{eq:C3af-invariant} has zero mixed
component.

\emph{4. Ordinary families and the target quotient.}
Take a generator of the ordinary algebraically
trivial source group, represented by a curve family.
Compactify that parameter smoothly, close the
cycle on the fixed product, and then on
$\mathcal Z$ times the proper curve. The invariant
defines a morphism from the parameter curve to
$\operatorname{Pic}^0(C)$ after normalization at
one point. Step 3 makes its induced homomorphism
from the parameter Jacobian zero on
$H^1(\mathcal O)$, hence zero. The family is
constant in $\operatorname{Pic}(C)$, proving
\[
             \rho\bigl(\mathrm{CH}^r(X)_{\mathrm{alg}}|_Y\bigr)=0.
\]

The product of the pulled-back divisors
$D_1,\ldots,D_s$ with a divisor class from
$\operatorname{Pic}^0(C)$ gives an actual ordinary
target family and a homomorphism
\[
 \theta:\operatorname{Pic}^0(C)(\mathbb C)
                    \longrightarrow\mathrm{CH}^r(Y)_{\mathrm{alg}}.
\]
As before, the Poincare parameter can be restricted
to chains of smooth connected curves. The exact
intersection identity is now
\[
                         \rho\theta=(-9)^s\,\mathrm{id}.
\]
Divisibility of the complex point group proves
the surjection \eqref{eq:C3af-cokernel} under the
temporary larger dimension bound. The same
specified point $p=(u:v:1)$, $v^4=-1-u^4$, with
$u$ transcendental over $\overline{\mathbb Q}$,
and $o=(1:\eta:0)$, gives a nontorsion target
cycle $\theta([p-o])$. This proves failure with
rational coefficients as well.

\emph{5. A smaller explicit set of sections.}
We now replace $3\cdot4^s$ by $4^s$ and retain the
same proof. Use the four matrices $A_i=(a_{ijk})$
from C.3ae, and define the $4$ by $3$ matrix
\[
                       B(w)_{ik}=\sum_{j=0}^3a_{ijk}w_j.
\]
At every point of the Fermat surface $W$, every
pair of columns of $B(w)$ is linearly independent.
This assertion has a separate exact certificate:
for each of the three column pairs and each of
the four charts $w_j=1$, the Fermat equation
together with the six $2$ by $2$ minors has
Groebner basis $[1]$ over $\mathbb F_{101}$.
There are twelve charts in total. Properness
of each projective determinantal locus over
$\mathbb Z$ proves the assertion in characteristic
zero, exactly as in C.3ae. The three coefficient
matrices for the individual columns have
determinants $-36$, $4$, and $200$, which also
verify that none of these columns vanishes at
a projective point. The retained computation
took $0.325551998$ seconds.

For $\mathbf i\in\{0,1,2,3\}^s$ set
\begin{equation}\label{eq:C3af-compressed-sections}
             t'_{\mathbf i}
                  =\sum_{k=0}^2 z_k
                        \prod_{l=1}^s B(w^{(l)})_{i_lk}.
\end{equation}
These $4^s$ sections of $M$ are base point free.
For $s=1$ this is the original base point
certificate in C.3ae. Suppose $s\ge2$. Write
$b_{lk}\in\mathbb C^4$ for the $k$th column
of $B(w^{(l)})$ at the chosen point. The three
vectors
\[
                         v_k=\bigotimes_{l=1}^s b_{lk}
                         \quad(k=0,1,2)
\]
are linearly independent. Here is a direct proof.
The three first factors $b_{1k}$ are pairwise
independent, and so are the three nonzero
remaining tensor factors
$\bigotimes_{l=2}^s b_{lk}$. Two pure tensors
cannot be proportional unless their nonzero
factors are proportional in every factor.
In a supposed dependence of the three $v_k$,
all coefficients must be nonzero, since each
pair is independent. Projecting the first
factor modulo the span of its first two
vectors shows that the third belongs to
that span. Both of its coefficients in that
span are nonzero. Comparing the first two
coefficients in the tensor product would
then make two remaining tensor factors
proportional, a contradiction. Thus the
$v_k$ are independent. Since the coordinate
vector $(z_0,z_1,z_2)$ is nonzero, its linear
combination of the $v_k$ is nonzero. These
are exactly the values of
\eqref{eq:C3af-compressed-sections}.

Put $q'=4^s$, take $m\ge q'-1$, and use
\[
                       \ell'=\sum_{\mathbf i}
                                      x_{\mathbf i}t'_{\mathbf i}
\]
instead of $\ell$ in the geometric construction.
Its coefficients are the specified products
$\prod_l a_{i_lj_lk}$, so the construction is
still completely determined by the four
integer matrices. Base point freeness gives
the same regular divisor model, and
$m\ge4^{r-1}-1>r$ for $r\ge2$ gives the
same projective-bundle quotient.

It remains to check the obstruction, since
the new sections have correlated coefficients.
Let $\overline U=U/\mathbb C H_W$ and let
$E_C$ be the three-dimensional space with
basis $e_k$ dual to $\delta(z_kG)$. The
rank-$57$ certificate of C.3ae says precisely
that
\[
 \iota:\overline U\otimes E_C
       \longrightarrow\mathbb C^4\otimes H^1(W,\Omega_W^1(1)),
 \qquad
 u\otimes e_k\longmapsto
       \left(\sum_j a_{ijk}w_ju\right)_{i=0}^3
\]
is injective. For the new sections, the map
on the primitive quotient of
\eqref{eq:C3af-source} is the composition,
up to reordering tensor factors, of
\[
\begin{split}
 \overline U^{\otimes s}\otimes E_C
   &\longrightarrow(\overline U\otimes E_C)^{\otimes s}\\
   &\xrightarrow{\ \iota^{\otimes s}\ }
       \left(\mathbb C^4\otimes
              H^1(W,\Omega_W^1(1))\right)^{\otimes s}.
\end{split}
\]
The first map sends
$u_1\otimes\cdots\otimes u_s\otimes e_k$
to $\bigotimes_l(u_l\otimes e_k)$. It is
injective because the three tensors
$e_k^{\otimes s}$ are linearly independent.
The second is injective because tensor
products over $\mathbb C$ preserve
injections. Expanding the composition gives
exactly the coefficients
$\prod_l a_{i_lj_lk}$ in
\eqref{eq:C3af-compressed-sections}; there is
no replacement of the actual curve by a
product of curves. Therefore the only
kernel before the quotient still consists
of tensors with at least one $H_i$ factor.
The Hodge projection, all boundary
corrections and the Picard-family argument
of steps 3--4 apply unchanged. This proves
\eqref{eq:C3af-cokernel} for the asserted
smaller bound $N\ge4^{r-1}+2r-1$.

The proof uses the finite Jacobian-ring spanning
certificate in C.3ae and the tensor-product
argument for arbitrary $r$. It does not rely on
testing a finite list of codimensions or on
unbounded matrix computations. The coefficient
choice and the dimension bound remain explicit
for every $r$.

% Root construction C.3ag; independent full audits by frobenius and real_weil.
\paragraph{C.3ag (Refuted): ordinary restriction in dimension six.}
We specify an excellent normal strictly henselian local
complex ring $R$ of dimension six, with regular puncture
$X$ and regular punctured principal divisor $Y$, and a
surjection
\[
 \operatorname{coker}\bigl(
 \mathrm{CH}^2(X)_{\mathrm{alg}}\longrightarrow
 \mathrm{CH}^2(Y)_{\mathrm{alg}}\bigr)
       \longrightarrow \operatorname{Pic}^0(C)(\mathbb C),
 \qquad C=V(z_0^4+z_1^4+z_2^4)\subset\mathbb P^2.
\]
The algebraic-equivalence groups use actual product
families. All differential sheaves below are relative
to $\mathbb C$. The surface in this construction is a
specified very general member with a Hodge-generic
complex choice; unlike C.3ae, we do not assert that it
is the Fermat quartic or that its coefficients are integers.

\emph{1. The surface and the precise genericity in the data.}
Fix the line $D=V(w_0,w_1)\subset\mathbb P^3$. Let
$\mathcal U$ be the open set of smooth quartics in
\[
 \left\{\sum_{\substack{|e|=4\\e_0+e_1>0}}
                 u_e w^e=0\right\}.
\]
There are thirty coefficients. We choose a Hodge-generic
complex point in this family for which
\[
 \operatorname{NS}(W)=\mathbb Z H_W\oplus\mathbb Z[D],
 \qquad \operatorname{End}_{\mathrm{Hdg}}(T_W)=\mathbb Q,
 \qquad T_W=\operatorname{NS}(W)_{\mathbb Q}^{\perp}
                  \subset H^2(W,\mathbb Q).
\]
Here ``Hodge-generic'' is part of the choice, not a
claim about every embedding of its coefficient field.
The existence, exact source inputs, and the consequence
for mixed cycle classes are supplied in the following
lemma. In particular they are not additional conjectures.

% Auxiliary input for the parent's prospective dimension-six construction.
% No complete local-geometric assertion is made in this fragment.
\paragraph{Auxiliary Hodge calculation for a quartic containing a line.}
Fix a line $D_0\subset\mathbb P^3_{\mathbb C}$ and let
$B\subset\mathbb P H^0(\mathcal I_{D_0}(4))$ be the nonempty open subset
parametrizing smooth quartics. We choose a \emph{Hodge-generic complex
point} $b\in B$ as part of the data, in the precise sense described below.
Write $W=W_b$, $H=c_1(\mathcal O_W(1))$, and $D=D_0\subset W$.
We can and do choose $b$ so that
\[
 \begin{gathered}
 \operatorname{NS}(W)=\mathbb ZH\oplus\mathbb ZD,
 \qquad
 \begin{pmatrix}H^2&H\cdot D\\H\cdot D&D^2\end{pmatrix}
 =\begin{pmatrix}4&1\\1&-2\end{pmatrix},\\
 T=\operatorname{NS}(W)_{\mathbb Q}^{\perp}
          \subset H^2(W,\mathbb Q),
 \qquad \operatorname{End}_{\mathrm{Hdg}}(T)=\mathbb Q.
 \end{gathered}
\]
Let $C$ be the Fermat plane quartic and put $S=W\times C$.
Then the largest rational sub-Hodge structure of $H^3(S,\mathbb Q)$
whose Hodge types are contained in $\{(2,1),(1,2)\}$ is
\begin{equation}\label{eq:N6-level-one}
       \operatorname{NS}(W)_{\mathbb Q}\otimes H^1(C,\mathbb Q).
\end{equation}
In particular, for every smooth projective complex curve $P$ and every
$z\in\mathrm{CH}^2(W\times C\times P)$, the component of its rational
cycle class in $H^3(S,\mathbb Q)\otimes H^1(P,\mathbb Q)$ belongs to
\[
 \operatorname{NS}(W)_{\mathbb Q}\otimes
 H^1(C,\mathbb Q)\otimes H^1(P,\mathbb Q).
\]

\emph{Choice of the quartic and its Hodge-generic complex point.}\par
Theorem~1.1 and Corollary~1.3(c) of
\citet[p.~2 of the opened author version]{frBrevikNollet2011}
apply to $D_0$ with $d=4$: $\mathcal I_{D_0}(3)$ is globally generated
and $H^0(\mathcal I_{D_0}(2))\ne0$. They give the asserted Picard group
for a very general member of $B$. A smooth quartic is a K3 surface,
and adjunction gives the displayed intersection matrix.

For completeness, we also justify the endomorphism condition rather
than inferring it solely from the Picard rank. Locally mark the second
cohomology. The Local Torelli Theorem and its line-bundle formulation
in \citet[Proposition~6.2.8 and Section~6.2.4, pp.~107--108 of the
opened author version]{frHuybrechtsK3Book} identify the local deformations
preserving $H,D$ with an open subset of the period quadric in
$(\mathbb QH+\mathbb QD)^\perp_{\mathbb C}$. It has dimension $18$.
These are locally realized by quartics containing a line: very ampleness
of $H$ persists; Riemann--Roch, $D^2=-2$, and $H\cdot D=1$ make $D$
effective, and an effective divisor of hyperplane degree one is a line.
A projective coordinate change puts that line back at $D_0$.

For every nonscalar $a\in\operatorname{End}_{\mathbb Q}(T)$, the
condition that $a$ is a Hodge endomorphism implies that the period line
is an eigenline of $a$. The eigenlines form a finite union of proper
linear subspaces. They cannot contain the nondegenerate irreducible
period quadric, which spans the ambient projective space. Thus each
such Hodge-endomorphism locus is proper. There are countably many
rational endomorphisms. Applied to the polarized weight-zero variation
$\operatorname{End}(T)$, the Hodge-locus theorem of
\citet[Theorem~1.1 and Corollary~1.3, p.~2]{frCDK1994}
shows that these loci are algebraic (with their monodromy translates).
Avoiding their countable union, together with the extra Picard loci,
produces the stipulated complex point. This is the meaning of the
Hodge-generic choice. We do not assert that every embedding of a field
with coefficientwise algebraically independent quartic parameters has
this property.

\emph{The product Hodge group and its representations.}
By \citet[Theorem~2.2.1 and Example~2.2.3, p.~211]{frZarhin1983},
\[
                    \operatorname{Hg}(T)=\operatorname{SO}(T).
\]
The exact original theorem and its proof in Section~2.5, pp.~212--213,
have been inspected. By \citet[Lemma~2, p.~106, and Example~3,
pp.~106--107]{frBeauville2014}, the Jacobian of $C$ is isogenous to
$E^3$, where $E=\mathbb C/\mathbb Z[i]$. Hence, for $V=H^1(E,\mathbb Q)$,
we have $H^1(C,\mathbb Q)\simeq V^{\oplus3}$ and
$\operatorname{Hg}(V)=\operatorname{Res}^{1}_{\mathbb Q(i)/\mathbb Q}
\mathbb G_m=:U$.

Let $G=\operatorname{Hg}(T\oplus V)$. Its two projections are surjective,
and $G$ is connected. The image of $G^{\mathrm{der}}$ on
$\operatorname{SO}(T)$ is the full group, while its image on the torus
$U$ is trivial. Therefore $\operatorname{SO}(T)\times1\subset G$;
surjectivity on $U$ then gives
\[
                   G=\operatorname{SO}(T)\times U.
\]
This argument uses only that the special orthogonal group in rank $20$
is connected and semisimple; no independence assumption on periods
of $W$ and $E$ is being inserted.

Over $\mathbb Q(i)$, the representation $T\otimes V$ is the direct sum
of the two twists of the standard $\operatorname{SO}(T)$ representation
by the two opposite nontrivial characters of $U$. Both summands are
absolutely irreducible, they are nonisomorphic, and complex conjugation
exchanges them. A rational invariant subspace therefore contains both
or neither. Thus $T\otimes V$ is a simple rational Hodge structure.
Moreover
\[
                 (T\otimes V)^{3,0}=T^{2,0}\otimes V^{1,0}\ne0.
\]
It consequently has level three. Semisimplicity of polarizable rational
Hodge structures implies that every nonzero sub-Hodge structure of
$(T\otimes V)^{\oplus3}$ has a simple constituent isomorphic to
$T\otimes V$; in particular none has level at most one.

Since $H^1(W)=H^3(W)=0$, Kunneth gives
\[
 H^3(S,\mathbb Q)=
 \bigl(\operatorname{NS}(W)_{\mathbb Q}\otimes H^1(C,\mathbb Q)\bigr)
 \oplus\bigl(T\otimes H^1(C,\mathbb Q)\bigr).
\]
The first summand has exactly the two indicated Hodge types and the
second has no nonzero sub-Hodge structure with those types. This proves
\eqref{eq:N6-level-one}.

Finally, the indicated Kunneth component of the class of $z$ is a
rational Hodge class of type $(2,2)$. Poincare duality on $P$ turns it
into a morphism of Hodge structures
$H^1(P,\mathbb Q)(-1)\to H^3(S,\mathbb Q)$. The source has types
$(2,1),(1,2)$, so its image lies in \eqref{eq:N6-level-one}, proving
the claimed consequence for every such cycle and parameter curve.


Let $K_0\subset\mathbb C$ be the field generated over
$\overline{\mathbb Q}$ by these thirty coefficients.
Choose thirty-six elements $a_{ijk}$, indexed by
$0\le i\le2$, $0\le j\le3$, $0\le k\le2$, algebraically
independent over $K_0$. Put
\[
 S=W\times C,\quad M=\mathcal O_W(1)\boxtimes K_C,
 \quad a=c_1(M),\quad
 s_i=\sum_{j,k}a_{ijk}w_jz_k.
\]
This specifies the geometric generic ordered triple
of sections over the chosen surface, with any complex
embedding extending the stated data. We prove the
two required open conditions for this triple below.

\emph{2. The section geometry and the dimension-six germ.}
The very ample bundle $M$ embeds the smooth threefold
$S$. Three general hyperplanes in this embedding have
transverse zero-dimensional intersection. This follows
successively from Bertini: after each smooth section,
the remaining very ample system is still without base
points. The last intersection is reduced, and its length is
\[
                        \int_S a^3
              =3\left(\int_W H_W^2\right)\deg K_C=48.
\]
The resulting nonempty Zariski-open condition is
defined over $K_0$, so it holds for the chosen triple.
Write $B=V(s_0,s_1,s_2)\subset S$ for these forty-eight
points, and put
\[
 V=\mathbb P^2_x\times S,\qquad
 L=\mathcal O_{\mathbb P^2}(1)\boxtimes M,
 \qquad \ell=\sum_{i=0}^2 x_i s_i,
 \qquad H=V(\ell)\subset V.
\]
The projection $p:H\to S$ is a $\mathbb P^1$-bundle
over $U=S\setminus B$, and has fibre $\mathbb P^2$
over each point of $B$. The latter locus has codimension
two in $H$. At a point of $B$, the three differentials
$ds_i$ form a basis of the cotangent space of $S$.
For every nonzero $x$, the differential
$\sum x_i ds_i$ is therefore nonzero. Together with
the bundle description over $U$, this proves that
$H$ is smooth. It is integral: its integral open
$\mathbb P^1$-bundle is dense and the exceptional
fibres have too small a dimension to be components.

Set
\[
 B'=\mathbb P^2_x\times W\times\mathbb P^2_z,
 \quad\mathcal N=\mathcal O_{B'}(-1,-1,-1),\quad
 G=z_0^2z_1^2+z_0^2z_2^2+z_1^2z_2^2.
\]
Inside $\operatorname{Tot}(\mathcal N)$ define
$Z=V(F+t\ell G)$, where $F=z_0^4+z_1^4+z_2^4$.
Contract by the Segre coordinates $y_{ijk}=tx_iw_jz_k$,
take the reduced closed image, its finite normalization,
and the strict henselization at the unique point over
the vertex. This is $R$. All these operations and all
their equations have now been specified. The proof of
C.3ae applies without a base-point-free hypothesis on
the $s_i$: $Z$ is integral, smooth near its connected
exceptional divisor $V$, and its regular projective
model $\mathcal Z$ over $R$ has
$\mathfrak m_R\mathcal O_{\mathcal Z}=\mathcal I_V$.
Its complement is the regular puncture $X$.
The dimension is $\dim V+1=6$.

The function $f=\sum a_{ijk}y_{ijk}$ pulls back to
$t\ell$. Its strict transform is the total space of
$L_H^{-1}$ on $H$, before the indicated base changes.
Write $\mathcal D$ for the corresponding regular
divisor on $\mathcal Z$. Then
$Y=\mathcal D\setminus H$ is the regular punctured
principal divisor required in the statement, and
$\mathcal D$ meets $V$ transversely in $H$.

\emph{3. An integral incidence invariant and its boundary.}
Put $h=c_1(\mathcal O_{\mathbb P^2}(1))|_H$ and
$\tau=4[D]-H_W$. Adjunction gives
\[
 H_W^2=4,\quad H_W[D]=1,\quad [D]^2=-2,
 \quad \deg(\tau H_W)=0,\quad\deg(\tau[D])=-9.
\]
For a smooth projective connected parameter $P$,
including a point, localization and self-intersection
give
\[
 \sigma_P:\mathrm{CH}^2(Y\times P)\longrightarrow
 \frac{\mathrm{CH}^2(H\times P)}
             {(h+a)\mathrm{CH}^1(H\times P)}.
\]
Indeed we extend a cycle to $\mathcal D\times P$ and
restrict to $H\times P$; a different extension changes
the restriction by $-(h+a)$ times a divisor class.

We need only a divisor decomposition, not a
projective-bundle description of all of $H$. Removing
$p^{-1}(B)\times P$ changes no Picard group, because
it has codimension two in a regular scheme. Removing
$B\times P$ similarly changes no Picard group of
$S\times P$. On the complement $H_U\times P$ is a
$\mathbb P^1$-bundle, with the line bundle of class
$h$ of relative degree one. Therefore, integrally,
\[
 \mathrm{CH}^1(H\times P)=
 p^*\mathrm{CH}^1(S\times P)
       \oplus h\,p^*\mathrm{CH}^0(S\times P).
\]
These equalities can also be read directly from Weil
divisors: deletion in codimension at least two preserves
divisor classes on a regular scheme, and the generic
fibre degree splits the divisor group of a projective
bundle. The only pullbacks by the possibly nonflat
map $p$ used here are of divisor classes or of their
products.

The hypersurface class $[H]=h+a$ in
$\mathbb P^2\times S$ gives the exact pushforwards
\[
            p_*1=0,\quad p_*h=1,\quad p_*h^2=a,
                      \quad p_*h^3=0.
\]
For $\alpha\in\mathrm{CH}^2(H\times P)$ set
\[
 \Phi_P(\alpha)=\operatorname{pr}_{C\times P,*}
          \left(\tau\,p_*\bigl((h-2a)\alpha\bigr)\right)
          \in\operatorname{Pic}(C\times P).
\]
Write a divisor on $H\times P$ as $\gamma_0+h\gamma_1$
by the preceding decomposition. Projection formula gives
\[
 p_*\bigl((h-2a)(h+a)(\gamma_0+h\gamma_1)\bigr)
                         =-3a^2\gamma_1.
\]
Its product with $\tau$ has zero pushforward to
$C\times P$: the $H_W^2$ term has codimension three
on $W$, and the $2H_WK_C$ term has coefficient
$2\deg(\tau H_W)=0$. Thus $\Phi_P$ descends to the
quotient defining $\sigma_P$, and
$\rho_P=\Phi_P\sigma_P$ is well defined integrally.
Localization, refined restriction, products and proper
pushforward commute with taking a parameter fibre.
Hence these $\rho_P$ are compatible with actual families.

If a cycle on $X\times P$ is closed on
$\mathcal Z\times P$ and restricted to $V\times P$,
write its class as $b=b_0+h b_1+h^2b_2$ using the
projective bundle $V=\mathbb P^2\times S$. Transverse
restriction to $H$, followed by the preceding identities,
gives
\[
 \rho_P(\Gamma|_{Y\times P})=
 \operatorname{pr}_{C\times P,*}
                  \bigl(\tau(b_0-a b_1-2a^2b_2)\bigr).
\]
This formula is independent of the closure. It has
been derived on the fixed scheme times $P$, without
any assertion about spreading a varying etale model.

\emph{4. Why three sections suffice on this surface.}
Let $F_W$ be the equation of $W$, and set
\[
 \mathcal R=\mathbb C[w_0,w_1,w_2,w_3]/
                 (\partial F_W/\partial w_0,\ldots,
                  \partial F_W/\partial w_3).
\]
The Euler and normal sequences used in C.3ae work
for every smooth quartic. They give
\[
 \frac{H^1(W,\Omega_W^1)}{\mathbb C H_W}
        \simeq\mathcal R_4^*,\qquad
 H^1(W,\Omega_W^1(1))\simeq\mathcal R_3^*,
 \qquad \dim\mathcal R_4=19,\quad\dim\mathcal R_3=16.
\]
Under these identifications, multiplication by $w_j$
is dual to $w_j:\mathcal R_3\to\mathcal R_4$ and kills
$H_W$. The sum
$\mathcal R_3^{\oplus4}\to\mathcal R_4$ is surjective:
every polynomial of degree four is a sum of $w_j$
times polynomials of degree three, and this property
survives the homogeneous quotient. Consequently the
four dual multiplication maps have zero common kernel
modulo $\mathbb C H_W$.

The class $\tau=4[D]-H_W$ is nonzero modulo $H_W$;
indeed $\tau H_W=0$ in degree and $\tau^2=-36$.
Thus some linear form $w$ satisfies
$\mu_w(\tau)\ne0$ in $H^1(W,\Omega_W^1(1))$.
The necessary rank condition on our triple is the
injectivity of
\[
 \mathbb C\tau\otimes H^0(C,K_C)
   \longrightarrow
 H^0(\mathbb P^2,\mathcal O(1))
                 \otimes H^1(W,\Omega_W^1(1)),
 \quad
 \tau\otimes v\longmapsto
 \sum_{i,j,k}a_{ijk}x_i\otimes\mu_{w_j}(\tau)
                           \langle\delta(z_kG),v\rangle.
\]
Here $\delta(z_kG)$ are the three basis classes in
$H^1(C,\mathcal O_C)$ calculated in C.3ae, and the
pairing is Serre duality. This is a three-column map.
For the special triple $s_i=wz_i$, it sends the three
dual basis vectors to $x_i\otimes\mu_w(\tau)$ and
is injective. Therefore its rank-three locus is a
nonempty Zariski open. The map, its class $\tau$,
and its rank condition are defined over $K_0$ by
algebraic de Rham cohomology. It follows that our
algebraically independent triple satisfies this condition
as well as the transverse-intersection condition of
step 2. We use the special triple only to prove that
the rank-open set is nonempty; it is not asserted
to have the required finite base locus itself.

\emph{5. The obstruction on an actual source family.}
Let $P$ now be a smooth projective connected curve
and let $\Gamma\in\mathrm{CH}^2(X\times P)$.
The fixed closure on $\mathcal Z\times P$ used in
step 3 has refined restriction $b$ to $V\times P$.
The conormal sequence and the lifted Hodge Chern
character in C.3ae, with exactly the same rational
Chern-character and Gysin inputs, imply
\[
 \kappa\mathbin{\lrcorner}\operatorname{cl}_{\mathrm{Hdg}}(b)
       =0\quad\hbox{in }
 H^3(V\times P,L\otimes\Omega^1_{V\times P}),
 \qquad
 \kappa=\sum_{i,j,k}a_{ijk}x_iw_j\otimes\delta(z_kG).
\]
The first neighbourhood has this class because its
equation is $F+t\ell G$ and its only moving tangent
direction is that of $C$. No converse assertion
about deforming cycles is used.

Project to the Kunneth summand
\[
 H^0(\mathbb P^2,\mathcal O(1))\otimes
 H^1(W,\Omega_W^1(1))\otimes
 H^1(C,K_C)\otimes H^1(P,\mathcal O_P).
\]
The only contributing component of $b_0$ is
\[
 H^1(W,\Omega_W^1)\otimes H^0(C,K_C)
                                  \otimes H^1(P,\mathcal O_P).
\]
By the auxiliary Hodge calculation in step 1, its
$W$ factor lies in
$\operatorname{NS}(W)\otimes\mathbb C
 =\mathbb C H_W+\mathbb C\tau$.
This assertion applies uniformly to every $P$ and
every algebraic cycle $b_0$; it is not an assumption
that the other Hodge components are algebraic.
The $h b_1$ and $h^2b_2$ terms cannot contribute to
the displayed summand, since contraction acts only
on the $C$ differential and cannot remove their
projective-space differential. The $H_W$ component
is killed by each $\mu_{w_j}$. On the $\tau$
component the map is the injective three-column
map from step 4 tensored with $H^1(P,\mathcal O_P)$.
Hence that component is zero.

Pairing with $\tau$ now shows that
$\operatorname{pr}_{C\times P,*}(\tau b_0)$ has
zero mixed component in
$H^0(C,K_C)\otimes H^1(P,\mathcal O_P)$.
The two correction terms in step 3 do not affect
this conclusion. More explicitly, $q(W)=0$ gives
\[
       \operatorname{Pic}(W\times C\times P)
       =\operatorname{Pic}(W)\oplus
                         \operatorname{Pic}(C\times P).
\]
This follows from the Picard product description:
the possible cross homomorphisms involve
$\operatorname{Alb}(W)=\operatorname{Pic}^0(W)=0$.
If $b_1=D'+\lambda$ in this decomposition, projection
formula gives
\[
 \operatorname{pr}_{C\times P,*}(\tau a b_1)
                 =\deg(\tau D')\,\operatorname{pr}_C^*K_C.
\]
The other terms have either insufficient $W$
codimension or coefficient $\deg(\tau H_W)=0$.
Also $\operatorname{pr}_{C\times P,*}(\tau a^2b_2)=0$.
Thus the correction is constant in $P$, and the
line bundle $\rho_P(\Gamma|_{Y\times P})$ has zero
mixed Hodge component.

Subtracting its fibre at a base point gives a
morphism $P\to\operatorname{Pic}^0(C)$. By the
Jacobian universal property it factors through a
homomorphism $J(P)\to\operatorname{Pic}^0(C)$.
The vanishing mixed component is its zero differential.
Over $\mathbb C$ a homomorphism of abelian varieties
with zero differential is zero. The morphism is
therefore constant. Fibre compatibility of $\rho_P$
proves that $\rho$ kills every source-family difference.
Compactifying a smooth parameter curve and closing
its cycle in the fixed product reduces every ordinary
algebraic-equivalence generator to this case. Hence
\[
      \rho\bigl(\mathrm{CH}^2(X)_{\mathrm{alg}}|_Y\bigr)=0.
\]

\emph{6. A nonzero target and the integral cokernel.}
Pull back the divisors $D\subset W$ and a divisor
class on $C$ through $Y\to H\to W\times C$ and
take their product. This defines
\[
 \theta:\operatorname{Pic}^0(C)(\mathbb C)
                      \longrightarrow\mathrm{CH}^2(Y)_{\mathrm{alg}}.
\]
The Poincare family, restricted to smooth curves
joining the required points in its connected parameter,
proves ordinary algebraic triviality. The map is a
product of divisor pullbacks, so it uses no false
flatness assumption on $p$. Taking the same product
on $\mathcal D$ and using $p_*(h-2a)=1$ yields
\[
                              \rho\theta=-9\,\mathrm{id}.
\]
Multiplication by nine on the complex points of a
Jacobian is surjective. The asserted integral
cokernel surjection follows, and in particular the
ordinary rational restriction map is not surjective.

For a specified nontorsion class take
$o=(1:\eta:0)$ with $\eta^4=-1$ and
$p=(u:v:1)$ with $u$ transcendental over
$\overline{\mathbb Q}$ and $v^4=-1-u^4$.
The Abel closed immersion is defined over
$\overline{\mathbb Q}$ and every torsion point of
the Jacobian is algebraic over that field. Thus
$[p-o]$ is nontorsion, and $\theta([p-o])$ maps
to the nontorsion class $-9[p-o]$ in the cokernel.
All generic choices used for the surface and the
three sections have been specified separately; no
finite test of a Fermat surface is substituted for
the Hodge-generic surface in this proof.

The exact incidence and intersection arithmetic is retained in
\path{code/check_dimension_six_incidence.py}; its run took
$0.046509863$ seconds. The computation checks the displayed pushforward
polynomials, lattice pairings, coefficient counts and base degree.
The source-checked Hodge and genericity arguments above are separate
from those finite calculations.

\paragraph{C.3ah (Refuted, with explicit counterexamples): the regular ambient range $N=r+2$.}
For every integer $r\ge1$, there is a regular excellent strictly
henselian local complex ring of dimension $N=r+2$ and a principal
divisor with regular puncture such that
\[
 \mathrm{CH}^r(X)=0,
 \qquad
 \mathrm{CH}^r(Y)_{\mathrm{alg}}\otimes\mathbb Q\ne0.
\]
In particular ordinary $J^r$ restriction is not surjective,
even with rational coefficients. For $r=2$ this gives the
dimension-four case not covered by C.3ae or C.3ag.

\begin{proof}
Put $d=r+2$ and specify
\[
 \begin{aligned}
 R&=\mathbb C[x_0,\ldots,x_{r+1}]_{(x_0,\ldots,x_{r+1})}^{\mathrm{sh}},
 & f&=x_0^d+\cdots+x_{r+1}^d,\\
 X&=\operatorname{Spec}R\setminus\{\mathfrak m_R\},
 &Y&=X\cap V(f).
 \end{aligned}
\]
Strict henselization is taken at the evident complex geometric
point. The hypersurface cone has dimension $r+1\ge2$.
Its gradient is $(d x_i^{d-1})_i$, so its only singular point
is the vertex. It is integral and normal: its projective
base is smooth and connected, and the hypersurface is
Cohen--Macaulay and regular in codimension one. The quotient
$R/(f)$ is the strict henselization of its vertex local ring;
this is the usual compatibility of strict henselization with
quotients, with the same geometric residue point. Thus $R$
and $Y$ satisfy all the required local hypotheses.

We use the smooth projective base
\[
 T=V(x_0^d+\cdots+x_{r+1}^d)\subset\mathbb P^{r+1},
 \qquad H=c_1(\mathcal O_T(1)),
 \qquad D=T\cap V(x_{r+1}).
\]
Let $k_0=\overline{\mathbb Q}\subset\mathbb C$, choose
$u_1,\ldots,u_r$ algebraically independent over $k_0$, and choose
$w,\eta$ with
\[
 w^d=-1-\sum_{i=1}^r u_i^d,
 \qquad \eta^d=-1.
\]
We take the specified points and degree-zero cycle
\[
 P=(u_1:\cdots:u_r:w:1),\qquad
 a=(1:\eta:0:\cdots:0)\in D(k_0),\qquad
 \alpha=[P]-[a].
\]
The point $P$ is a geometric generic point of the
$k_0$-variety $T$. The class $\alpha$ is ordinarily
algebraically trivial: when $r=1$ the curve $T$ itself
provides the difference family; when $r\ge2$, a smooth
connected complete-intersection curve through $P,a$
provides it. Such a curve is obtained by Bertini with the
two prescribed smooth points, in a sufficiently ample
linear system. The diagonal on that curve followed by its
closed immersion in $T$ gives an actual fixed-product
cycle family.

We recall why this particular class survives the hyperplane
quotient, rather than merely appealing to nonvanishing of
some zero-cycle. Suppose a nonzero multiple of $\alpha$
belonged to $H\mathrm{CH}^{r-1}(T)_{\mathbb Q}$. Every
hyperplane product has a representative supported on $D$.
After clearing denominators, the finitely many cycles and
rational functions expressing the supposed relation can be
spread, specialized to $\overline{k_0(T)}$, and descended
to a finite extension of $k_0(T)$. Norm and closure then
give
\[
       b\Delta_T=\Theta_D+\Theta_Z,\qquad b\ne0,
\]
where $\Theta_D$ is supported on $T\times D$ and
$\Theta_Z$ on $Z\times T$ for a proper divisor $Z\subset T$.
These finite-data and localization operations are those
proved in C.3f and C.3s. Adjunction gives $K_T=\mathcal O_T$,
so $H^{r,0}(T)\ne0$. Both correspondences on the right
annihilate this space: restriction to a resolution of $Z$
kills an $r$-form, and pushforward from $D$ cannot produce
Hodge bidegree $(r,0)$. The diagonal acts as multiplication
by $b$, a contradiction. Consequently
\[
 \alpha\ne0\quad\hbox{in}\quad
 \frac{\mathrm{CH}^r(T)_{\mathbb Q}}
      {H\mathrm{CH}^{r-1}(T)_{\mathbb Q}}.
\]

The Fermat hypersurface is projectively normal with section
ring generated in degree one. The strict-local cone injection
C.3e therefore sends $\alpha$ to a non-torsion class
$\gamma\in\mathrm{CH}^r(Y)$. Its construction is the flat
pullback from $T$ to the punctured line bundle, followed by
the strict-local base change. Pulling back the fixed curve
family just described proves
$\gamma\in\mathrm{CH}^r(Y)_{\mathrm{alg}}$ integrally.
This is exactly the case $m=0$ of C.3s.

Finally the local map $\mathbb C\to R$ is geometrically
regular, Noetherian, and has henselian target. We apply
\citet[Corollary~6.4(i), author p.~27]{rootSkalit2017}
with zero-dimensional base. It gives
$\mathrm{CH}^r(\operatorname{Spec}R)=0$ for every $r>0$,
integrally. Chow localization surjects this group onto
$\mathrm{CH}^r(X)$, so the source is zero. The specified
non-torsion class $\gamma$ proves the failure of ordinary
restriction surjectivity.
\end{proof}

Combining $r=2$ here with C.3r at $N=5$, C.3ag at
$N=6$, and C.3ae at every $N\ge7$ yields the exact
universal dimension statement over $\mathbb C$:
\[
 \text{ordinary }J^2(X)\longrightarrow J^2(Y)
 \text{ is universally surjective}
 \quad\Longleftrightarrow\quad N\le3.
\]
For $N\le3$ the target has no codimension-two points and
its Chow group is zero. For every $N\ge4$ the cited
examples satisfy the regular-puncture hypotheses and give
failure after tensoring with $\mathbb Q$. The polynomial
and dimension checks for the present construction are
retained in \texttt{code/check\_regular\_ambient\_calabi\_yau.py}.
Its checks support the displayed equations; the generic-point
and Hodge-form proof above establishes non-torsion in all
dimensions.

\paragraph{C.3ai (Refuted, with explicit counterexamples): every nonempty low-dimensional ordinary restriction range.}
For every pair of integers
\[
       r\ge1,\qquad r+2\le N\le2r+1,
\]
there is a specified regular excellent strictly henselian local
complex ring $R$ of dimension $N$, a nonzero $f\in\mathfrak m_R$,
and regular punctures $X$ and $Y=X\cap V(f)$ for which
\[
    \mathrm{CH}^r(X)=0,
    \qquad \mathrm{CH}^r(Y)_{\mathrm{alg}}\otimes\mathbb Q\ne0.
\]
We prove this by explicit suspended Fermat hypersurfaces. In
particular we do not assume a generalized Hodge conjecture or
infer non-torsion from a finite computation.

\begin{proof}
Set
\[
 d=2r-N+2\ge1,\qquad k=N-r-2\ge0,
 \qquad e=d+2,
 \qquad r=d+k,
 \qquad N=d+2k+2.
\]
Over $k_0=\overline{\mathbb Q}\subset\mathbb C$ define
\[
 \begin{split}
 T&=V(y_0^e+\cdots+y_{d+1}^e)\subset\mathbb P^{d+1},\\
 V&=V\left(\sum_{j=0}^{k-1}(a_j^e+b_j^e)
                    +\sum_{i=0}^{d+1}y_i^e\right)
                         \subset\mathbb P^{d+2k+1}.
 \end{split}
\]
Thus $T$ is a Calabi--Yau $d$-fold, $V$ has dimension
$n=d+2k$, and the coordinate embedding $T\subset V$ sets
all $a_j,b_j$ equal to zero. Write $H=c_1(\mathcal O_V(1))$.
Choose $\eta,\zeta\in k_0$ such that $\eta^e=-1$ and
$\zeta$ is a primitive $e$th root of unity. The two distinct
scalars $\lambda_0=\eta$ and $\lambda_1=\eta\zeta$ both
have $e$th power $-1$.

For $k>0$, impose $a_j=\lambda_i b_j$ for every $j$.
The resulting subvariety $R_i\subset V$ is the join of $T$
with the linear vertex $\mathbb P^{k-1}$ in the paired
coordinates. Its incidence resolution $Z_i\subset T\times V$
is a $\mathbb P^k$-bundle over $T$: over a point
$[y_0:\cdots:y_{d+1}]$, its fibre consists of
\[
 [\lambda_i b_0:b_0:\cdots:\lambda_i b_{k-1}:b_{k-1}:
                         t y_0:\cdots:t y_{d+1}].
\]
These coordinates describe the projective bundle homogeneously.
The fibre lies on $V$ since $\lambda_i^e+1=0$.
Its dimension is $d+k=r$, so we obtain
\[
       \Gamma_i=[Z_i]\in\mathrm{CH}^{r}(T\times V).
\]
For $k=0$ we put $V=T$ and
$\Gamma_0=\Gamma_1=\Delta_T$.

We have an integral left inverse
\begin{equation}\label{eq:C3-suspended-left-inverse}
             {}^t\Gamma_1\circ\Gamma_0=\Delta_T
                         \quad\text{in }\mathrm{CH}^{d}(T\times T).
\end{equation}
For $k>0$, the two different linear equations in each pair
force $a_j=b_j=0$. Thus $R_0\cap R_1=T$ scheme-theoretically,
and the intersection avoids both vertices. At a point of
$T$, the tangent space of $V$ is the direct sum of $T_T$
and the $2k$ paired-coordinate directions. The tangent
spaces of the two joins use complementary $k$-dimensional
subspaces of those directions, since
$\lambda_0\ne\lambda_1$. Their intersection is therefore
transverse. The fibre product of the two incidence
resolutions over $V$ is exactly the diagonal copy of $T$,
with multiplicity one. The Chow composition formula proves
\eqref{eq:C3-suspended-left-inverse}. The assertion is immediate
for $k=0$.

Choose $u_1,\ldots,u_d$ algebraically independent over $k_0$,
choose $w$ with $w^e=-1-\sum u_i^e$, and put
\[
    P=(u_1:\cdots:u_d:w:1),\qquad
    o=(1:\eta:0:\cdots:0),\qquad
    \alpha=[P]-[o]\in\mathrm{CH}^{d}(T)_{\mathrm{alg}}.
\]
The point $P$ is a geometric generic point over $k_0$.
The ordinary algebraic triviality is supplied by an actual
smooth complete-intersection curve through $P,o$ when
$d>1$, and by $T$ itself when $d=1$. Thus
$(\Gamma_0)_*\alpha\in\mathrm{CH}^{r}(V)_{\mathrm{alg}}$.
We claim that its image in
\[
       \mathrm{CH}^{r}(V)_{\mathbb Q}/
                          H\mathrm{CH}^{r-1}(V)_{\mathbb Q}
\]
is nonzero.

Suppose otherwise. Descend the finitely many cycles and
rational functions in a relation
$(\Gamma_0)_*\alpha=H\beta$ to the algebraic closure of
$k_0(T)$, then to a finite extension of that field, and
take the norm. Here the descent from $\mathbb C$ uses
specialization of finite data over the algebraically closed
subfield, as in C.3s. Closing the resulting generic-fibre
cycles on $T\times V$ and applying localization gives
\begin{equation}\label{eq:C3-suspended-spread}
 b\Gamma_0
    =H B+b\bigl(T\times(\Gamma_0)_*[o]\bigr)+\Theta,
 \qquad b\ne0,
\end{equation}
where $B\in\mathrm{CH}^{r-1}(T\times V)_{\mathbb Q}$ and
$\Theta\in\mathrm{CH}^{r}(T\times V)_{\mathbb Q}$ is
supported on $Z\times V$ for a proper divisor $Z\subset T$.
We retain the constant correspondence term explicitly;
it is not absorbed into the support condition on $\Theta$.

Apply these correspondences to $H^{d,0}(T)\ne0$. The
constant correspondence on the right acts as zero, since
$d>0$. The term supported on $Z\times V$ also acts as
zero: restriction to a resolution of $Z$ kills a
holomorphic $d$-form. The codimension of $B$ is $r-1$,
so its action sends this space into
\[
       H^{d+k-1,k-1}(V)
                   \subset H^{d+2k-2}(V,\mathbb C).
\]
For $k=0$ this is zero because its second Hodge index is
negative. For $k\ge1$ the total degree is $n-2<n$.
Weak Lefschetz identifies it with the same cohomology of
projective space, which has only diagonal Hodge types.
For the precise range, the complement of $V$ in
$\mathbb P^{n+1}$ is affine. The vanishing in C.2b and
Poincare duality give $H_c^j(\mathbb P^{n+1}\setminus V)=0$
for $j<n+1$. The compact-support localization sequence
therefore makes restriction an isomorphism in every degree
less than $n$, compatibly with Hodge structures.
The two displayed Hodge indices differ by $d\ge1$, so
this space is again zero. It follows that $HB$ acts as
zero on $H^{d,0}(T)$.

On the other hand, \eqref{eq:C3-suspended-left-inverse}
makes $(\Gamma_0)_*$ injective on cohomology; its left
inverse is $({}^t\Gamma_1)_*$. Thus the left side of
\eqref{eq:C3-suspended-spread} cannot act as zero on the
nonzero space $H^{d,0}(T)$. This contradiction proves
the claim. The argument uses weak Lefschetz below the
middle degree and the explicit correspondence identity;
no assertion about algebraic classes in the middle is needed.

Finally specify the regular local ambient ring and its
equation by
\[
 \begin{split}
 R&=\mathbb C[a_0,b_0,\ldots,a_{k-1},b_{k-1},
                y_0,\ldots,y_{d+1}]_{(a,b,y)}^{\mathrm{sh}},\\
 f&=\sum_{j=0}^{k-1}(a_j^e+b_j^e)+\sum_{i=0}^{d+1}y_i^e.
 \end{split}
\]
The paired-coordinate list is omitted when $k=0$.
The ring has dimension $d+2k+2=N$. The hypersurface is normal
and regular off its vertex, by the diagonal Fermat
gradient and the hypersurface depth criterion. Thus
$Y=(\operatorname{Spec}R\cap V(f))\setminus\{\mathfrak m_R\}$
is the strict-local punctured cone on $V$ and is regular.
The projectively normal cone injection C.3e sends the
nonzero quotient class just proved to a non-torsion class
in $\mathrm{CH}^{r}(Y)_{\mathrm{alg}}$. Ordinary algebraic
triviality follows by pulling the actual curve family and
the displayed incidence correspondence to the fixed
puncture. Corollary~6.4(i) of \citet[author p.~27]{rootSkalit2017},
with base $\mathbb C$, gives
$\mathrm{CH}^r(\operatorname{Spec}R)=0$; localization then
gives $\mathrm{CH}^r(X)=0$. This proves the assertion for
every pair in the stated range.
\end{proof}

The endpoints $k=0$ and $d=1$ recover respectively the
Calabi--Yau construction C.3ah and the odd-dimensional
Fermat correspondence in C.3r. The full interval between
them is proved by the same cohomological degree calculation.

The retained program is
\begin{center}
\texttt{code/check\_suspended\_calabi\_yau.py}.
\end{center}
It checks the parameter conversions and equations, as well as the
tangent matrices, cycle codimensions and Hodge shifts.
The all-dimensional non-torsion argument is the proof above.

% C.3aj: full proof independently audited by root, Frobenius and cycles; source checks retained.
\paragraph{C.3aj (Refuted, with fully specified Hodge-generic counterexamples): ordinary higher-cycle restriction in every range $N\ge 2r+2$.}
For every pair of integers $r\ge2$ and $N\ge2r+2$, we specify
an excellent normal strictly henselian local complex ring $R$ of
dimension $N$, with regular puncture $X$, and a nonzero element
$f\in\mathfrak m_R$ whose punctured divisor $Y=X\cap V(f)$ is
regular, such that ordinary algebraically trivial cycles have a
surjection
\[
 \operatorname{coker}\bigl(
 \mathrm{CH}^r(X)_{\mathrm{alg}}\longrightarrow
 \mathrm{CH}^r(Y)_{\mathrm{alg}}\bigr)
 \longrightarrow
 Q_C:=\frac{\operatorname{Pic}^0(C)(\mathbb C)}
                  {\operatorname{Pic}^0(C)(\overline K)}\ne0.
\]
Here $C=V(z_0^4+z_1^4+z_2^4)\subset\mathbb P^2$ and
$\overline K\subset\mathbb C$ is a specified countable algebraically
closed field of definition of the construction. The group $Q_C$ is
uniquely divisible. Thus restriction fails to be surjective even
after tensoring with $\mathbb Q$. Every algebraic equivalence used
below is represented on an actual fixed product with a smooth
connected curve.

\emph{1. The surfaces, coefficient fields, and generic section data.}
Put
\[
 s=r-1,\qquad m=N-2r\ge2.
\]
Choose $s$ smooth quartic K3 surfaces $W_\nu\subset\mathbb P^3$
containing the respective fixed lines
$D_\nu=V(w_{\nu0},w_{\nu1})$, with hyperplane classes $H_\nu$.
The tuple is chosen Hodge-generic in the product of these families,
in the precise sense of the following auxiliary calculation:
\[
 \operatorname{NS}(W_\nu)=\mathbb ZH_\nu\oplus\mathbb ZD_\nu,
 \qquad
 \operatorname{Hg}\left(\bigoplus_{\nu=1}^s T_\nu\right)
       =\prod_{\nu=1}^s\operatorname{SO}(T_\nu),
 \quad
 T_\nu=\operatorname{NS}(W_\nu)_{\mathbb Q}^{\perp}.
\]
The Hodge-generic complex tuple itself is part of the data. We
do not claim that all embeddings of a coefficient field preserve
this Hodge-genericity.

% Auxiliary Hodge input for the higher-range local restriction construction.
% This fragment proves the Hodge assertion only, not the local geometry.
\paragraph{Auxiliary product Hodge calculation for quartics containing lines.}
Let $s\geq1$. We choose smooth complex quartics $W_1,\ldots,W_s$, each
containing a specified line $D_i$, at a jointly Hodge-generic point of
the product of the corresponding quartic families. Write
$H_i=c_1(\mathcal O_{W_i}(1))$ and
$T_i=(\mathbb QH_i+\mathbb QD_i)^\perp\subset H^2(W_i,\mathbb Q)$.
We can and do require
\begin{equation}\label{eq:product-K3-generic-data}
 \begin{gathered}
 \operatorname{NS}(W_i)=\mathbb ZH_i\oplus\mathbb ZD_i,
 \qquad
 \begin{pmatrix}H_i^2&H_iD_i\\H_iD_i&D_i^2\end{pmatrix}
       =\begin{pmatrix}4&1\\1&-2\end{pmatrix},\\
 \operatorname{End}_{\mathrm{Hdg}}(T_i)=\mathbb Q,
 \qquad
 \operatorname{Hom}_{\mathrm{Hdg}}(T_i,T_j)=0\quad(i\ne j).
 \end{gathered}
\end{equation}
Let $C$ be the Fermat plane quartic, set
$S=(\prod_{i=1}^sW_i)\times C$, and consider the indicated Kunneth
summand
\[
 K=\left(\bigotimes_{i=1}^sH^2(W_i,\mathbb Q)\right)
                       \otimes H^1(C,\mathbb Q)
       \subset H^{2s+1}(S,\mathbb Q).
\]
Then its largest rational sub-Hodge structure of level at most one is
\begin{equation}\label{eq:product-K3-level-one}
 K_{\leq1}=
 \left(\bigotimes_{i=1}^s\operatorname{NS}(W_i)_{\mathbb Q}\right)
                       \otimes H^1(C,\mathbb Q).
\end{equation}
In particular, for every smooth projective complex curve $P$ and every
$z\in\mathrm{CH}^{s+1}(S\times P)$, the component of its rational
cycle class in $K\otimes H^1(P,\mathbb Q)$ belongs to
$K_{\leq1}\otimes H^1(P,\mathbb Q)$. The parameter curve is arbitrary;
it is not required to be defined over the field of definition of the
quartics. We assert \eqref{eq:product-K3-level-one} for the specified
Kunneth summand $K$, not for all of $H^{2s+1}(S,\mathbb Q)$.

\emph{Existence of the chosen tuple.}\par
For a fixed line $D_0\subset\mathbb P^3$, let $B$ be the smooth open
subset of $\mathbb P H^0(\mathcal I_{D_0}(4))$. Theorem~1.1 and
Corollary~1.3(c) of \citet[p.~2 of the opened author version]
{frBrevikNollet2011} give $\operatorname{Pic}(W)=\mathbb ZH\oplus
\mathbb ZD_0$ for a very general point of $B$. Their hypotheses hold
because $\mathcal I_{D_0}(3)$ is globally generated and
$H^0(\mathcal I_{D_0}(2))\ne0$. Adjunction gives the intersection matrix
in \eqref{eq:product-K3-generic-data}.

We impose the additional Hodge conditions explicitly. By Local Torelli
and the line-bundle period description in
\citet[Proposition~6.2.8 and Section~6.2.4, pp.~107--108 of the
opened author version]{frHuybrechtsK3Book}, locally marked deformations
preserving $H,D_0$ fill an open subset of the $18$-dimensional period
quadric in $(\mathbb QH+\mathbb QD_0)^\perp_{\mathbb C}$.
They are locally realized in $B$: very ampleness of $H$ is open;
Riemann--Roch and $H\cdot D_0=1$ make the effective class $D_0$ a line;
and a projective coordinate change returns that line to the fixed line.
Thus the product family has independently varying local periods.

For a nonscalar rational endomorphism $a$ of a marked transcendental
space, the Hodge-endomorphism condition forces the period line to be
an eigenline of $a$. The finite union of its eigenspaces cannot contain
the irreducible nondegenerate period quadric, which spans the ambient
projective space. This condition therefore defines a proper locus.
After imposing $\operatorname{End}_{\mathrm{Hdg}}(T_i)=\mathbb Q$,
each $T_i$ is simple: a nontrivial direct summand would give a
nontrivial rational Hodge projector. A nonzero Hodge map between two
$T_i$ would consequently be an isomorphism. For each rational linear
isomorphism $a:T_i\to T_j$, its Hodge locus forces
$a\omega_i\in\mathbb C\omega_j$. The independent variation of the
two period lines makes this a proper locus in the product family.
There are only countably many such rational maps. Applied to the
weight-zero variations $\operatorname{End}(T_i)$ and
$\operatorname{Hom}(T_i,T_j)$, the Hodge-locus theorem
\citep[Theorem~1.1 and Corollary~1.3, p.~2]{frCDK1994}
makes the loci, with their monodromy translates, countable unions of
closed algebraic subsets. Avoiding these loci and the extra Picard
loci gives \eqref{eq:product-K3-generic-data}. This specifies a jointly
Hodge-generic complex point; it does not assert that every embedding
of a coefficientwise algebraically independent parameter field is
Hodge generic.

\emph{Independence of the orthogonal factors.}\par
By \citet[Theorem~2.2.1 and Example~2.2.3, p.~211]{frZarhin1983},
each $T_i$ has rank $20$ and $\operatorname{Hg}(T_i)=\operatorname{SO}(T_i)$.
Let $G=\operatorname{Hg}(\bigoplus_iT_i)$ and
$\mathfrak g=\operatorname{Lie}(G)$. Its projections onto the
absolutely simple Lie algebras
$\mathfrak g_i=\mathfrak{so}(T_i)$ are surjective.
The radical of $\mathfrak g$ maps to a solvable ideal in every
$\mathfrak g_i$, hence is zero. Decompose this semisimple rational
Lie algebra into its simple ideals. Each projection to
$\mathfrak g_i$ is an isomorphism on exactly one of these ideals and
zero on the others. Consequently, if
$\mathfrak g\ne\bigoplus_i\mathfrak g_i$, one simple ideal maps
isomorphically onto two different factors. The projection to that
pair is the graph of a rational Lie-algebra isomorphism
$\varphi:\mathfrak g_i\simeq\mathfrak g_j$.

We record why this would identify the two standard representations;
this is the point at which the rank matters. Over $\overline{\mathbb Q}$,
an automorphism of $\mathfrak{so}_{20}$ integrates to the simply
connected group, and, modulo an inner automorphism, acts by a
symmetry of its $D_{10}$ diagram
\citep[20.33, p.~108, and 21.4, p.~112]{frMilneReductive2018}.
The simple roots are
$\varepsilon_1-\varepsilon_2,\ldots,
\varepsilon_9-\varepsilon_{10},\varepsilon_9+\varepsilon_{10}$
\citep[18.26--18.27, pp.~91--92]{frMilneReductive2018}.
The only nontrivial diagram symmetry swaps its two short terminal
branches and fixes the first node. It therefore fixes the highest
weight $\varepsilon_1$ of the standard representation, whose weights
are $\{\pm\varepsilon_j\}_{j=1}^{10}$. Highest-weight uniqueness
\citep[20.3--20.4, pp.~102--103]{frMilneReductive2018}
shows that twisting the standard representation by any such
automorphism preserves its isomorphism class.

It follows that the intertwiner space
\[
 \{u\in\operatorname{Hom}_{\mathbb Q}(T_i,T_j):
          uX=\varphi(X)u\text{ for every }X\in\mathfrak g_i\}
\]
is nonzero after scalar extension. These are rational linear equations,
so the rational space itself is nonzero. A nonzero intertwiner is an
isomorphism, since both standard representations are absolutely
irreducible. It commutes with $\mathfrak g$ on the pair, hence with
the connected group $G$. Because $T_i,T_j$ have the same weight, it
is a rational Hodge isomorphism, contrary to
\eqref{eq:product-K3-generic-data}. We conclude that
\[
       \operatorname{Hg}\left(\bigoplus_{i=1}^sT_i\right)
                          =\prod_{i=1}^s\operatorname{SO}(T_i).
\]

\emph{The CM factor and the levels of the tensor constituents.}\par
By \citet[Lemma~2, p.~106, and Example~3, pp.~106--107]
{frBeauville2014}, $\operatorname{Jac}(C)$ is isogenous to $E^3$,
where $E=\mathbb C/\mathbb Z[i]$. Put $V=H^1(E,\mathbb Q)$ and
$U=\operatorname{Res}^1_{\mathbb Q(i)/\mathbb Q}\mathbb G_m$.
Then $H^1(C,\mathbb Q)\simeq V^{\oplus3}$ and
$\operatorname{Hg}(V)=U$. The joint Hodge group is a connected
subdirect product of $\prod_i\operatorname{SO}(T_i)$ and $U$.
Its derived group projects onto the full semisimple product and
trivially onto $U$, so it contains the semisimple product. Surjectivity
onto $U$ now gives
\[
 \operatorname{Hg}\left(\bigoplus_iT_i\oplus V\right)
             =\left(\prod_i\operatorname{SO}(T_i)\right)\times U.
\]

For a nonempty subset $I\subset\{1,\ldots,s\}$, put
$T_I=\bigotimes_{i\in I}T_i$ and $t=|I|$.
The representation $T_I$ of the product of the corresponding
orthogonal groups is absolutely irreducible. Over $\mathbb Q(i)$,
$T_I\otimes V$ is the sum of its twists by the two opposite
nontrivial characters of $U$. These are nonisomorphic absolutely
irreducible representations, exchanged by complex conjugation.
Therefore $T_I\otimes V$ is a simple rational Hodge structure.
After the Tate twist supplied by the remaining $s-t$ Neron--Severi
factors, it has a nonzero component of type
\[
                       (s+t+1,s-t),
\]
and hence has level $2t+1\geq3$. Each remaining Neron--Severi
factor has rank two, and $H^1(C)=V^{\oplus3}$; these give
multiplicities, not further simple constituents of smaller level.

Expand $K$ using $H^2(W_i)=\operatorname{NS}(W_i)_{\mathbb Q}\oplus T_i$.
The term with no transcendental factor is precisely the right-hand
side of \eqref{eq:product-K3-level-one}, and has types
$(s+1,s),(s,s+1)$. Every other term is a direct sum of the simple
constituents just described. Semisimplicity of polarizable rational
Hodge structures shows that none of these other terms has a nonzero
sub-Hodge structure of level at most one. This proves
\eqref{eq:product-K3-level-one}.

Finally, the selected Kunneth component of an algebraic cycle class
$\operatorname{cl}(z)$ is a rational Hodge class of type $(s+1,s+1)$.
Poincare duality on $P$ identifies it with a Hodge morphism
\[
                 H^1(P,\mathbb Q)(-s)\longrightarrow K.
\]
Its source has level at most one. Its image therefore lies in
$K_{\leq1}$, proving the stated consequence for every parameter $P$.


Let $K_0$ be generated over $\overline{\mathbb Q}$ by the finitely
many coefficients of these quartic equations. Set
\[
 \begin{gathered}
 A=\prod_{\nu=1}^s W_\nu,\qquad S=A\times C,\\
 M=\left(\boxtimes_{\nu=1}^s\mathcal O_{W_\nu}(1)\right)
                        \boxtimes K_C,\\
 a=c_1(M)=\sum_{\nu=1}^s H_\nu+k_C,
 \qquad k_C=c_1(K_C).
 \end{gathered}
\]
Choose the $(m+1)3\cdot4^s$ complex numbers
\[
 a_{i,\mathbf j,k},\qquad
 0\le i\le m,\quad\mathbf j=(j_1,\ldots,j_s)\in\{0,1,2,3\}^s,
 \quad 0\le k\le2,
\]
algebraically independent over $K_0$, and put
\[
 s_i=\sum_{\mathbf j,k}a_{i,\mathbf j,k}
                 \left(\prod_\nu w_{\nu j_\nu}\right)z_k.
\]
Let $K$ be generated by $K_0$ and
all the $a_{i,\mathbf j,k}$, and let $\overline K$ be its algebraic
closure inside $\mathbb C$. Both fields are countable.

The common zero scheme $B\subset S$ of the $m+1$ sections is
either empty or a smooth complete intersection of codimension
$m+1$. This is a nonempty Zariski-open condition on the tuples:
$M$ is very ample, and successive general hyperplanes give smooth
transverse intersections until their expected dimension is negative,
when their common intersection is empty. The condition is defined
over $K_0$, so algebraic independence ensures it for our tuple.
A second nonempty open condition, needed for the Chern obstruction,
is proved in step 5 below.

Set
\[
 V=\mathbb P^m_x\times S,\quad
 L=\mathcal O_{\mathbb P^m}(1)\boxtimes M,\quad
 \ell=\sum_{i=0}^m x_i s_i,\quad
 H=V(\ell)\subset V,
 \qquad p:H\longrightarrow S.
\]
Over $S\setminus B$, $p$ is the projective bundle of lines in
the kernel of $\mathcal O_S^{m+1}\to M$, with fibre
$\mathbb P^{m-1}$. At a point of $B$, the $ds_i$ are linearly
independent, so $\sum x_i ds_i$ is nonzero for every projective
point $x$. Thus $H$ is smooth everywhere. It is integral: its
projective-bundle open is integral and its complement
$\mathbb P^m\times B$ has codimension $m\ge2$. Its line bundle
of class $h+a$, where $h=c_1(\mathcal O_{\mathbb P^m}(1))|_H$,
is the restriction of the very ample $L$.

\emph{2. The actual local ring and its regular models.}
Let
\[
 B'=\mathbb P^m_x\times A\times\mathbb P^2_z,
 \qquad
 \mathcal N=\mathcal O_{\mathbb P^m}(-1)
       \boxtimes\left(\boxtimes_\nu\mathcal O_{W_\nu}(-1)\right)
       \boxtimes\mathcal O_{\mathbb P^2}(-1),
\]
and write $t$ for the tautological fibre vector. Put
\[
 F=z_0^4+z_1^4+z_2^4,\qquad
 G=z_0^2z_1^2+z_0^2z_2^2+z_1^2z_2^2,
 \qquad Z=V(F+t\ell G)\subset\operatorname{Tot}(\mathcal N).
\]
The expression is a section of the pullback of
$\mathcal O_{\mathbb P^2}(4)$ and is independent of
trivializations. The zero section on $Z$ is exactly $V$, with
normal bundle $L^{-1}$. The differential $dF$ is nonzero there,
so $Z$ is smooth near $V$. On $\ell G\ne0$ it is a graph over
$B'$. The locus $F=\ell G=0$ has codimension two in $B'$
and cannot supply an additional component of this Cartier
hypersurface. Since the hypersurface is Cohen--Macaulay and
generically reduced on its graph open, $Z$ is integral.

Contract the zero section by the Segre coordinates
\[
 y_{i,\mathbf j,k}=t x_i\left(\prod_\nu w_{\nu j_\nu}\right)z_k.
\]
Take the reduced closed affine image $Q$, its finite normalization
$Q^{\mathrm n}$, and the strict henselization of the local ring
at the unique point over the vertex. This defines $R$ without
an unspecified shrinking choice. The proper contraction is an
isomorphism off the zero section. Properness allows a neighbourhood
of the vertex to be chosen on which the model is regular. Its
normalization has only one point over the vertex: the connected
exceptional fibre $V$ maps surjectively onto that finite fibre.
After the strict henselization let $\mathcal Z$ denote this regular
projective model. As in C.3ae, the coordinates $y$ generate
$\mathcal I_V$ on every Segre chart, and every normalization
function in the maximal ideal vanishes on the reduced $V$.
Consequently
\[
 \mathfrak m_R\mathcal O_{\mathcal Z}=\mathcal I_V,
 \qquad X=\mathcal Z\setminus V.
\]
Pointed etale base change with the same residue field preserves
all these statements and the first infinitesimal neighbourhood.
The ring is excellent and normal, and
$\dim R=\dim V+1=m+2r=N$.

The function
\[
 f=\sum_{i,\mathbf j,k}a_{i,\mathbf j,k}y_{i,\mathbf j,k}
\]
pulls back to $t\ell$. Its strict transform is, before the
indicated local base changes, the total space of $L_H^{-1}$
over $H$. Denote its base change by $\mathcal D$. It is regular,
meets $V$ transversely in $H$, and
$Y=\mathcal D\setminus H=X\cap V(f)$ is regular. All the
schemes, line bundles and maps in this step are defined over
$\overline K$ before the final complex and strict-local base changes.

\emph{3. The incidence operator and its boundary on proper families.}
For each $\nu$, put $\tau_\nu=4D_\nu-H_\nu$ and set
$\tau=\prod_\nu\tau_\nu\in\mathrm{CH}^{s}(A)$. Adjunction gives
\[
 \deg(\tau_\nu H_\nu)=0,\qquad
 \deg(\tau_\nu D_\nu)=-9.
\]
For a smooth projective connected parameter $P$, including
a point, define
\[
 \Phi_P:\mathrm{CH}^r(H\times P)\longrightarrow
                  \operatorname{Pic}(C\times P),
 \qquad
 \Phi_P(\alpha)=\operatorname{pr}_{C\times P,*}
                       \bigl(\tau\,p_*(h^{m-1}\alpha)\bigr).
\]
All products with the displayed divisor classes are defined;
no flatness of $p$ is assumed. Since $[H]=h+a$ in
$\mathbb P^m\times S$, the projection formula in that ambient
projective bundle gives
\[
 p_*(h^{m-1})=1,\qquad p_*(h^m)=a,
 \qquad p_*(h^{m+j})=0\quad(j\ge1).
\]
The same identities hold after taking the product with $P$.

Write
\[
 \Psi_P(\gamma)=\Phi_P\bigl((h+a)\gamma\bigr),
 \qquad \gamma\in\mathrm{CH}^{r-1}(H\times P).
\]
We claim that when $P$ is a smooth projective curve the line
bundle $\Psi_P(\gamma)$ has constant fibre class in
$\operatorname{Pic}(C)$. We prove this on its mixed Hodge
component, without asserting that $\gamma$ lifts as a Chow
class to $V\times P$.

Indeed $\dim V=m+2r-1$ and $H$ is a smooth very ample
divisor. Weak Lefschetz gives an isomorphism of rational
Hodge structures
\[
 H^{2r-3}(V,\mathbb Q)\xrightarrow{\ \sim\ }
 H^{2r-3}(H,\mathbb Q).
\]
The exact range also follows directly from C.2b: the complement
is smooth affine, its cohomology vanishes above $\dim V$,
Poincare duality makes its compactly supported cohomology
vanish below $\dim V$, and the pair sequence gives the
displayed isomorphism since $2r-3<\dim V-1$.
Thus the $H^{2r-3}(H)\otimes H^1(P)$ component of the
cycle class of $\gamma$ has a unique ambient lift. Write the
first factor of that lift as
$\zeta=\sum_j h^j\zeta_j$, with
$\zeta_j\in H^{2r-3-2j}(S,\mathbb Q)$. The preceding
pushforwards give on this component exactly
\[
 \Psi_*\zeta
   =\operatorname{pr}_{C,*}
                 \bigl(\tau(2a\zeta_0+a^2\zeta_1)\bigr).
\]
Terms with a negative cohomological degree mean zero.

Every odd cohomology class of $S=A\times C$ has its odd
degree on $C$, because K3 surfaces have no odd cohomology.
The contribution relevant to $H^1(C)$ therefore uses
$H^1(C)$ and total $W_\nu$ degree $2(s-1-j)$ in
$\zeta_j$. A factor of $k_C=c_1(K_C)$ would make the degree
on $C$ exceed two and gives zero. Any $W_\nu$ degree four
already gives zero after multiplication by $\tau_\nu$.
For a remaining term to survive integration over every
$W_\nu$, exactly $j+1$ factors have initial degree zero,
and the required factors of $a^{j+1}$ must insert $H_\nu$
in those positions. Each such term contains
$\deg(\tau_\nu H_\nu)=0$. This proves
$\Psi_*\zeta=0$, for every $\zeta$, not just algebraic ones.
The line bundle $\Psi_P(\gamma)$ has zero mixed Hodge
component. Its classifying map $P\to\operatorname{Pic}(C)$,
after subtraction of one fibre, factors through the Jacobian
of $P$. Its differential is zero, hence the homomorphism of
complex abelian varieties is zero. This proves the claim
integrally on the fibre classes.

\emph{4. Descent of the entire boundary to a countable field.}
It follows that $\Psi=\Psi_{\operatorname{Spec}\mathbb C}$
kills ordinary algebraic-equivalence differences on the
projective variety $H$. Every cycle on $H_{\mathbb C}$ is
ordinarily algebraically equivalent to a cycle defined over
$\overline K$. Here is the exact finite-type justification.
An integral subvariety gives a complex point of a Hilbert
scheme with its fixed Hilbert polynomial. That Hilbert
scheme is projective of finite type over $\overline K$,
with a universal flat closed family
\citep[Theorems 3.1--3.2 and the Hilbert specialization,
pp.~259--260, 265--266]{D_GrothendieckHilbert1961}.
Because $\overline K$ is algebraically closed, the irreducible
component containing the given point is defined over
$\overline K$ and has a $\overline K$-point. The two points
can be joined by a chain of projective integral curves
in that component. Normalize those curves and pull back
the universal flat family. Their fundamental cycles give
actual families on $H$ times smooth projective connected
curves. They prove the asserted algebraic equivalence.
Applying this to the finitely many components of a cycle
and using that $\Psi$ is defined over $\overline K$ yields
\[
 \Psi\bigl(\mathrm{CH}^{r-1}(H)\bigr)
                  \subset\operatorname{Pic}(C)(\overline K).
\]
This statement requires no claim that cycles on a strictly
local puncture itself have a finite-type parameter space.
Only the smooth projective exceptional section $H$ is used.

Chow localization and refined pullback on
$H\subset\mathcal D$ now define
\[
 \rho:\mathrm{CH}^r(Y)\longrightarrow
                  \operatorname{Pic}(C)(\mathbb C)/
                  \operatorname{Pic}(C)(\overline K).
\]
Namely close a cycle on $\mathcal D$, restrict to $H$, and
apply $\Phi$. Changing the closure adds the pushforward of
a class $\gamma\in\mathrm{CH}^{r-1}(H)$; its refined
restriction changes by $-(h+a)\gamma$, whose $\Phi$-image
is zero in the indicated quotient. Localization also proves
independence of rational-equivalence witnesses. A closure
chosen on $\mathcal D\times P$ before taking a fibre gives
the same resulting $\rho$-class after taking that fibre,
because refined pullbacks and proper pushforwards commute.

Every degree class on $C$ is defined over $\overline K$
using a fixed $\overline K$-point. Thus the last quotient is
canonically $Q_C$. It is divisible because the group of
complex points of its Jacobian is divisible. It is
torsion-free: if $n\lambda$ is a $\overline K$-point of the
Jacobian, the finite morphism $[n]$, defined over the
algebraically closed $\overline K$, has all its preimages
over $\overline K$, so $\lambda$ is already such a point.
Consequently $Q_C$ is uniquely divisible.

\emph{5. The three-section obstruction in arbitrary codimension.}
All differentials in this step are relative to $\mathbb C$.
The first neighbourhood of $V$ in $\mathcal Z$ has ideal
$L$ and deformation class
\[
 \kappa=\sum_{i,\mathbf j,k}a_{i,\mathbf j,k}x_i
       \left(\prod_\nu w_{\nu j_\nu}\right)
                   \otimes\delta(z_kG)
       \in H^1(V,T_V\otimes L).
\]
It has only a $C$-tangent component. As calculated in
C.3m and C.3ae, the three classes $\delta(z_kG)$ form a
basis of $H^1(C,\mathcal O_C)$; their Serre-dual basis
is used for $H^0(C,K_C)$.

For each smooth quartic $W_\nu$, the Euler and normal
sequences and Serre duality give
\[
 H^1(W_\nu,\Omega^1_{W_\nu})/\mathbb C H_\nu
       \simeq\mathcal R_{\nu,4}^*,\qquad
 H^1(W_\nu,\Omega^1_{W_\nu}(1))
       \simeq\mathcal R_{\nu,3}^*,
\]
where $\mathcal R_\nu$ is the Jacobian ring of its quartic.
The maps $\mu_{\nu,w}$ of multiplication by a linear form
kill $H_\nu$ and have zero common kernel modulo $H_\nu$:
dually, the maps $w_{\nu j}\mathcal R_{\nu,3}$ span
$\mathcal R_{\nu,4}$. Hence we can choose a linear form
$w_\nu$ with $\mu_{\nu,w_\nu}(\tau_\nu)\ne0$.

Consider the obstruction map on
\[
 \left(\bigotimes_\nu\mathbb C\tau_\nu\right)
                       \otimes H^0(C,K_C).
\]
Its target is
\[
 H^0(\mathbb P^m,\mathcal O(1))\otimes
       \bigotimes_\nu H^1(W_\nu,\Omega^1_{W_\nu}(1)).
\]
For the special triple
$s_i=(\prod_\nu w_\nu)z_i$ for $i=0,1,2$, with all
other $s_i=0$, the three dual basis vectors map to
$x_i\otimes\bigotimes_\nu\mu_{\nu,w_\nu}(\tau_\nu)$.
They are independent. Thus injectivity is a nonempty
rank-three open condition on the full section tuple.
The cohomology maps, divisor classes and minors are
defined over $K_0$ by algebraic de Rham cohomology.
Our algebraically independent coefficients satisfy this
condition and the independent geometric open condition
from step 1. The special tuple is only a rank witness;
its base locus is not asserted to have the required geometry.

Let $P$ be a smooth projective connected curve and
$\Gamma\in\mathrm{CH}^r(X\times P)$. Close it on the
fixed $\mathcal Z\times P$ and let
\[
 b=\sum_{j=0}^{\min(m,r)}h^j b_j,
 \qquad b_j\in\mathrm{CH}^{r-j}(S\times P),
\]
be its refined restriction to $V\times P$.
The necessary Chern obstruction, in its degree-$r$ form, is
\[
 \kappa\mathbin{\lrcorner}\operatorname{cl}_{\mathrm{Hdg}}(b)=0
       \quad\hbox{in }H^{r+1}
       (V\times P,L\otimes\Omega^{r-1}_{V\times P}).
\]
For clarity this uses exactly the exterior-power sequence
\[
 0\longrightarrow L\otimes\Omega^{r-1}_{V\times P}
 \longrightarrow\Omega^r_{\mathcal Z\times P}|_{V\times P}
 \longrightarrow\Omega^r_{V\times P}\longrightarrow0.
\]
Its boundary is contraction with $\kappa$, as is seen by
comparing local splittings of the conormal sequence.
The Hodge Chern character of a vector bundle on
$\mathcal Z\times P$ lifts through the middle term:
the local Atiyah matrices exist before restriction, and
$\operatorname{tr}(\operatorname{At}^r)/r!$ gives the
degree-$r$ component. A rational virtual bundle with
Chow Chern character equal to the chosen closure exists
by the rational Chern-character isomorphism for this
finite-dimensional regular projective scheme over $R$
\citep[Lemma 42.58.1, Tag 0FEY]{cyStacksChernObstruction}.
Naturality under refined restriction identifies its
restricted degree-$r$ character with the class of $b$
\citep[Remark 42.59.9, Tag 0FF8]{cyStacksChernObstruction}.
Equivalently, the finite bundle and cycle data descend
to a smooth pointed etale finite-type complex model.
No sufficiency assertion about deforming cycles is used.

Project the obstruction to the summand
\[
 H^0(\mathbb P^m,\mathcal O(1))\otimes
 \bigotimes_\nu H^1(W_\nu,\Omega^1_{W_\nu}(1))\otimes
 H^1(C,K_C)\otimes H^1(P,\mathcal O_P).
\]
Only the component
\[
 \bigotimes_\nu H^1(W_\nu,\Omega^1_{W_\nu})\otimes
 H^0(C,K_C)\otimes H^1(P,\mathcal O_P)
\]
of $\operatorname{cl}(b_0)$ contributes. The auxiliary
product-Hodge calculation implies that its rational
antecedent has all its $W_\nu$ factors in
$\operatorname{NS}(W_\nu)_{\mathbb Q}$: the tensor
parts containing a transcendental factor have no
level-one Hodge constituent. Components with a
projective-space differential, including every $h^j b_j$
for $j>0$, cannot reach the displayed target because
$\kappa$ contracts only the curve differential.
Every tensor containing an $H_\nu$ factor is killed
by $\mu_{\nu,w}$. The remaining tensor line is
$\bigotimes_\nu\mathbb C\tau_\nu$, and the rank-three
map just proved is injective on its $H^0(C,K_C)$
factor, also after tensoring with $H^1(P,\mathcal O_P)$.
This component is therefore zero.
Its complex-conjugate component with
$H^1(C,\mathcal O_C)\otimes H^0(P,K_P)$ is also zero:
the cycle class, its Kunneth projection and the
Neron--Severi tensor factors are rational. Thus both
halves of the relevant mixed component vanish.

\emph{6. Source-family vanishing and a nonzero target.}
The exact incidence pushforwards give, for the fixed closure
of the preceding source family,
\[
 \Phi_P(b|_{H\times P})
   =\operatorname{pr}_{C\times P,*}\bigl(\tau(b_0+a b_1)\bigr).
\]
The first term has zero mixed component by step 5 and
$\deg(\tau_\nu H_\nu)=0$. The second also has zero mixed
component: the relevant part of $b_1$ has degree
$2r-3=2s-1$ on $S$, with its odd degree on $C$.
The same factor-by-factor degree argument in step 3
forces the single factor of $a$ into a missing $W_\nu$
position and gives $\deg(\tau_\nu H_\nu)=0$; its
$K_C$ component is zero for degree reasons.
Thus this line bundle has constant fibre class on $P$.
Fibre compatibility of $\rho$ proves that every actual
source-family difference maps to zero. A smooth open
parameter curve can be compactified and its cycle closed
on the fixed product, without changing the original
fibres. Consequently
\[
 \rho\bigl(\mathrm{CH}^r(X)_{\mathrm{alg}}|_Y\bigr)=0.
\]

Pull back the divisor classes $D_1,\ldots,D_s$ and a
line bundle from $C$ to $Y$, and multiply their first
Chern classes. This defines
\[
 \theta:\operatorname{Pic}^0(C)(\mathbb C)
          \longrightarrow\mathrm{CH}^r(Y)_{\mathrm{alg}}.
\]
The Poincare line bundle and smooth proper parameter
curves in its connected Jacobian give actual ordinary
families for these cycles. All the pullbacks are of
line bundles, so the possibly nonflat $p$ causes no
undefined operation. Taking the same divisor product
on $\mathcal D$ and using $p_*h^{m-1}=1$ gives
\[
 \rho\theta(\lambda)=(-9)^s\lambda\pmod{
                                      \operatorname{Pic}(C)(\overline K)}.
\]
Since multiplication by the nonzero integer $(-9)^s$
is surjective on $\operatorname{Pic}^0(C)(\mathbb C)$,
this proves the asserted integral cokernel surjection.

Finally choose $\eta\in\overline{\mathbb Q}$ with
$\eta^4=-1$, put $o=(1:\eta:0)$, choose $u$
transcendental over $\overline K$, choose
$v^4=-1-u^4$, and put $q=(u:v:1)\in C(\mathbb C)$.
No nonzero multiple of $[q-o]$ belongs to
$\operatorname{Pic}^0(C)(\overline K)$: finiteness of
$[n]$ would first make its Abel point a $\overline K$-point,
and the Abel closed immersion, defined over that field,
would then make $q$ a $\overline K$-point. This contradicts
the choice of $u$. Thus $Q_C\ne0$, and the target cycle
$\theta([q-o])$ cannot belong to the ordinary source image,
even rationally. The argument proves the stated result
for every $r\ge2$ and $N\ge2r+2$; a finite calculation
is not being used to infer any universal Hodge-genericity.

The retained program
\texttt{code/check\_high\_range\_incidence.py} checks 154
exact pushforward cases, 16 tensor-degree cancellation cases,
eight rank-witness minors and 88 dimension cases. Its runtime
was $0.057480419$ seconds. These checks verify the displayed
coefficient formulas and inequalities in their finite sample;
the arguments above prove the general identities, the actual
family statements and the Hodge-generic existence conditions.

\paragraph{C.3ak (Proved, with counterexamples in every complementary dimension): the complete ordinary local surjectivity range over $\mathbb C$.}
Fix integers $r\ge1$ and $N\ge1$. We use the manuscript's
ordinary fixed-product algebraic equivalence and put
\[
        J^r(U)=\mathrm{CH}^r(U)_{\mathrm{alg}}\otimes\mathbb Q.
\]
Consider all excellent normal strictly henselian local
complex rings $R$ of dimension $N$, with residue field
$\mathbb C$, regular puncture
$X=\operatorname{Spec}R\setminus\{\mathfrak m_R\}$, and
nonzero $f\in\mathfrak m_R$ for which
$Y=X\cap V(f)$ is regular. Then
\[
 \boxed{\quad
 J^r(X)\longrightarrow J^r(Y)
 \text{ is surjective for every such pair}
       \quad\Longleftrightarrow\quad N\le r+1.
 \quad}
\]
The same universal dimension classification holds before
tensoring the two ordinary groups with $\mathbb Q$.

\begin{proof}
If $N\le r+1$, the punctured divisor has no point of
codimension $r$: its dimension is at most $N-2<r$.
For $N=1$ it is empty. Consequently
$\mathrm{CH}^r(Y)=0$, proving the universal assertion
integrally in this entire range.

Suppose $N\ge r+2$. The counterexamples already proved
cover every possibility, as follows:
\begin{center}
\begin{tabular}{lll}
\toprule
Codimension & Local dimension & Verified construction \\
\midrule
$r\ge1$ & $r+2\le N\le2r+1$ & C.3ai: suspended Calabi--Yau \\
$r=1$ & $N\ge4$ & C.3m: deformed quartic cone \\
$r\ge2$ & $N\ge2r+2$ & C.3aj: product of quartics \\
\bottomrule
\end{tabular}
\end{center}
In the first row the whole source Chow group vanishes and
the specified ordinary target class is nontorsion. The
second row has zero ordinary source and a Jacobian subgroup
with nontorsion target points. The third row gives an
ordinary restriction cokernel surjecting onto a nonzero
uniquely divisible group. In particular every row supplies
failure after tensoring with $\mathbb Q$, as well as failure
of the integral ordinary map. All its rings, divisor
equations and target witnesses are specified in the cited
constructions, and all have regular punctures.
\end{proof}

This settles the universal ordinary local restriction
\emph{surjectivity} question in every positive codimension
over $\mathbb C$. Higher-codimension restriction
injectivity remains a separate question. The
characteristic-$p$ assertions retain the scopes and field
hypotheses stated in their own paragraphs.


\paragraph{Remaining work for C.3.}
The cokernel at the quintic, the divisor injectivity theorem, the exact
Griffiths-group obstruction sequence, the $J^1$ criterion, and the
punctured-cone counterexample are established. The universal
surjectivity classifications for $A_{\mathrm{alg}}$ and $J$ are now
given by C.1f and C.3d. C.3e--g settle the strict local cone
passage, refute every nonempty proposed high-degree finiteness
range, and give a nonzero local Gysin map and restriction kernel.
C.3h gives exact Beauville-graded tests for the two injectivity
questions on an abelian variety; C.3l proves that their joint
universal assertion is equivalent to the ordinary negative-index
Beauville conjecture. C.3i--k distinguish the full local Picard
group from its ordinary algebraically trivial part: the former
has explicit additive restriction kernels, while the latter is
injective for algebraizable characteristic-zero local rings of
dimension at least three. C.3k also proves prime-to-residue-
characteristic torsion injectivity under its stated excellent
strictly henselian hypotheses. C.3m gives completely specified
failures of ordinary local $J^1$ surjectivity in every dimension
at least four. C.3n extends its abelian-image injectivity argument
under the precise stated coefficient-field hypotheses. C.3o
settles all ordinary degree-one boundary dimensions over
$\mathbb C$, and C.3p proves the asserted prime-to-residue-
characteristic torsion surjectivity in dimensions at least four.
In its precise geometric scope, C.3q lifts target abelian
subgroups of the local Picard functor up to a power-of-$p$
isogeny. C.3w proves the required realization of every
point of a proper abelian subgroup by actual product
families in that algebraizable scope; hence ordinary
degree-one restriction is integrally surjective there
in characteristic $p$, with $p$-primary kernel.
C.3x completes the ordinary rational degree-one dimension
table in that scope, including explicitly specified boundary
counterexamples in characteristics two and three and the
vanishing over $\overline{\mathbb F}_p$. C.3z removes the
algebraizability restriction from the point-realization and
positive-characteristic conclusions: they hold for excellent
henselian normal local $k$-algebras with the specified
coefficient field $k$, residue field $k$, and regular puncture.
C.3r gives explicit ordinary $J^r$ restriction-surjectivity
failures at every boundary pair $(N,r)=(2r+1,r)$. C.3s
refutes all universal ordinary-$J^r$ and Chow local-finiteness
assertions in every nonzero range $1\le r<N$; the only
universal vanishing range is $r\ge N$.
C.3t supplies nonzero ordinary local Gysin maps and
open-restriction kernels for every $r+1\le N\le2r-1$,
including every odd critical dimension.
C.3u proves ordinary degree-one restriction surjectivity for
the stated smooth-cone sources in every characteristic and
identifies their ordinary source group with
$\operatorname{Pic}^0(V)_{\mathrm{red}}(k)$; its integral kernel
is zero in characteristic zero and finite $p$-primary in
characteristic $p$.
C.3v computes the ordinary divisor quotient on specified
strict local cones in every dimension at least two,
and refutes its proposed finiteness already in dimension two.
C.3y supplies a specified multiplicative divisor quotient
$k^\times\oplus(\mathbb Z/4\mathbb Z)^2$ in dimension two
in every characteristic, with ordinary algebraically
trivial subgroup zero.
C.3aa refutes the proposed finiteness of $B^r$ in every
non-top local degree $2\le r\le N-2$ over $\mathbb C$ by a
specified Segre--Nori cone. It proves nonzero rationalization,
without claiming infinite rational dimension. The top-degree
quotient $B^{N-1}$ for $N\ge3$ is not decided by that construction.
C.3ab and C.3ad prove the stated restricted top-degree
smooth-cone results; they do not settle that quotient for
every local source. C.3ac settles the two abelian
injectivity questions over $\overline{\mathbb F}_p$ by
classical rational Chow hard Lefschetz and separately
records its stronger consequences under the expressly
cited 2026 CM-Hodge theorem input.
C.3ae gives an integral ordinary $J^2$ restriction cokernel
surjecting onto a specified curve's Picard variety in every
dimension $N\ge7$. C.3ag supplies the dimension-six case,
and C.3ah supplies dimension four, as well as the uniform
regular-ambient range $N=r+2$. Together with C.3r at $N=5$,
these complete ordinary degree-two local restriction
surjectivity over $\mathbb C$: its universal assertion
holds exactly for $N\le3$, when the target has no
codimension-two cycles.
C.3ai extends the specified regular-source examples to every
$r\ge1$ and every $r+2\le N\le2r+1$, by suspended
Calabi--Yau correspondences and the exact Hodge shift below
the middle degree.
C.3af supplies completely specified ordinary $J^r$
restriction counterexamples for every $r\ge2$ and every
$N\ge4^{r-1}+2r-1$, with the same Picard-variety quotient
of the integral cokernel. C.3aj fills every higher range
$r\ge2$, $N\ge2r+2$, using jointly Hodge-generic quartics
and a nonzero uniquely divisible Picard quotient.
C.3ak combines these with C.3ai and the ordinary divisor
examples: over $\mathbb C$, universal ordinary $J^r$
restriction surjectivity holds exactly when $N\le r+1$,
for every $r\ge1$, integrally and rationally.
Universal higher injectivity, the exact individual-variety
higher $J^i$ kernels and cokernels, higher-codimension local
restriction injectivity, and higher-degree restriction in
positive characteristic remain undecided here, with the
precise proved ranges and coefficient-field hypotheses above.

\subsection{C.4: degree one, denominators, and Questions 1--2}

\paragraph{C.4a (Proved: degree-one Lefschetz, independent of $\ell$).}
Let $X/k$ be smooth projective connected of dimension $n\ge1$, with
$k$ algebraically closed. Let $L_1,\ldots,L_{n-1}$ be ample line bundles
and $\eta=\prod c_1(L_j)$. For every $\ell\ne\operatorname{char}k$,
\begin{equation}\label{eq:degree-one-lefschetz}
 \eta\smile-:H^1(X,\mathbb Q_\ell)\xrightarrow{\sim}
 H^{2n-1}(X,\mathbb Q_\ell)(n-1)
\end{equation}
has an inverse induced by one rational divisorial correspondence,
independent of $\ell$. It is unique as a normalized divisor
correspondence class, not as a divisor on $X\times X$.

\emph{Proof.}
The case $n=1$ is the diagonal. For $n>1$ we replace each $L_j$ by a
positive power so that general members meet in a smooth connected
curve $v:C\hookrightarrow X$. Bertini applies over the algebraically
closed, hence infinite, field. The replacement multiplies $\eta$ by
a positive integer $N$; a rational inverse for $N\eta$ gives one for
$\eta$ by multiplication by $N$, so this step loses no claimed
integrality.

Write $B_X=\operatorname{Pic}^0_{X/k,\mathrm{red}}$ and
$A_X=\operatorname{Alb}_{X/k}=B_X^\vee$. Restriction
$v^*:B_X\to\operatorname{Pic}^0(C)$ has finite kernel: on prime-to-
characteristic Tate modules this follows by repeated weak Lefschetz,
and a positive-dimensional kernel would have a nonzero such Tate
module. Pulling back an ample theta divisor under this finite-kernel
homomorphism gives an ample line bundle on $B_X$. Its polarization is
an isogeny. Up to the conventional sign of Jacobian autoduality, it is
the morphism of abelian varieties
\[
 u:B_X\xrightarrow{v^*}\operatorname{Pic}^0(C)
 \xrightarrow{d_C}\operatorname{Alb}(C)
 \xrightarrow{v_*}A_X.
\]
The correspondence $(v,v)_*[C]\in\mathrm{CH}_1(X\times X)$ acts on
$H^1$ as multiplication by $[C]=N\eta$ by the projection formula.
Its action identifies with $u^*$ under the canonical Picard and
Albanese cohomology identifications. Hence it is invertible. The inverse
$u^{-1}\in\operatorname{Hom}(A_X,B_X)\otimes\mathbb Q$ corresponds to
a normalized rational divisor on $X\times X$ by the Poincar\'e bundle.
This construction is algebraic before choosing $\ell$, which proves
independence. Faithfulness of the Tate-module functor on homomorphisms
of abelian varieties proves uniqueness in this normalized group.
The argument agrees with the opened
\citet[Appendix, Theorem 2A9(5)--(6)]{cyKleiman1968}.

The sign in ``polarization'' on 14:67 depends on whether $d_C$ denotes
the polarization or its negative under the chosen autoduality. The
isogeny statement and the inverse construction are independent of this
sign. The map marked $\varphi_X$ in the comparison square on 14:59 must
have the source indexed by $Y$; this is a typing correction, not a new
geometric assertion.

\paragraph{C.4b (Proved corrected integral form of 2.4--2.5).}
Let $X,Y$ be smooth projective connected varieties over $k$ and let
$Z\in\mathrm{CH}_1(X\times Y)$ be integral. It defines
$u_Z:B_X\to A_Y$ and, for $n=\dim X$,
\[
 a(Z):H^1(Y,\mathbb Q_\ell)\longrightarrow
 H^{2n-1}(X,\mathbb Q_\ell)(n-1).
\]
This is an isomorphism if and only if $u_Z$ is an isogeny. In that case
its inverse is induced by the unique class corresponding to
$u_Z^{-1}\in\operatorname{Hom}(A_Y,B_X)_{\mathbb Q}$. It is represented
by an integral normalized divisor correspondence if and only if
$u_Z$ is an isomorphism.

More precisely, for a positive integer $m$,
\begin{equation}\label{eq:denominator-exponent}
 m u_Z^{-1}\in\operatorname{Hom}(A_Y,B_X)
 \quad\Longleftrightarrow\quad [m]|_{\ker u_Z}=0
\end{equation}
as a morphism of finite group schemes. If the right side holds,
$[m]:B_X\to B_X$ factors through the fppf quotient
$B_X/\ker u_Z=A_Y$, yielding the left side. Conversely,
$(m u_Z^{-1})u_Z=[m]$ kills the kernel. Thus the least positive
denominator is the \emph{exponent} of $\ker u_Z$, not generally its
order. For $u=[2]$ on a $g$-dimensional abelian variety, it is $2$,
whereas the kernel order is $2^{2g}$, including in characteristic $2$
when its scheme structure must be retained. This refutes the strongest
literal ``best denominator'' reading of 14:64 while preserving its
assertion about the set of possible prime divisors.

The induced map
\[
 H^1(Y,\mathbb Z_\ell)\longrightarrow
 H^{2n-1}(X,\mathbb Z_\ell)(n-1)/\mathrm{tors}
\]
is an isomorphism precisely when $u_Z$ is an isogeny of degree prime
to $\ell$. Indeed it is the corresponding map of free Tate-module
lattices; its index has $\ell$-adic valuation
$v_\ell(\deg u_Z)$. An isomorphism before dividing out torsion also
requires the target to be torsion free. This is the exact distinction
made by the insertions in Corollary 2.5.

\paragraph{C.4c (Proved: Question 1 and the generic zero-cycle criterion).}
Fix $x_0\in X(k)$, let $A=A_X$, let $a_X:X\to A$ send $x_0$ to $0$,
and set $K=k(A)$. Let $\eta_A\in A(K)$ be the generic point. Then the
following are equivalent:
\begin{enumerate}
\item There is an integral cycle $\Gamma\in\mathrm{CH}^{n}(A\times X)$
whose action on $H^1$ is inverse to $a_X^*:H^1(A)\to H^1(X)$.
\item There is $z\in\mathrm{CH}_0(X_K)^0$ with
$\operatorname{alb}_{X_K}(z)=\eta_A$.
\end{enumerate}
These assertions use one cycle inducing the inverse for all
prime-to-characteristic realizations. They are equivalent to existence
of a universal zero-cycle, normalized by its fibre over $0$.

\emph{Proof.}
A cycle $\Gamma$ as in (1) induces a morphism
$\phi_\Gamma:A\to A$, $t\mapsto
\operatorname{alb}(\Gamma_t-\Gamma_0)$. Its action on $H^1$ is the
identity. Faithfulness of $T_\ell$ gives $\phi_\Gamma=\mathrm{id}_A$.
After subtracting the constant family and a suitable multiple of
$x_0$, its generic fibre has degree zero and Albanese class $\eta_A$.
Conversely a representative of $z$ spreads over a nonempty open
subset of $A$. We close its components in $A\times X$ to obtain an
integral correspondence. The normalized Albanese morphism associated
to it agrees with the identity at the generic point, hence on the
whole abelian variety. This gives the inverse on $H^1$.
Changing representatives, closing at the boundary, and subtracting
constant cycles do not alter that morphism; existence, rather than a
canonical choice of cycle, is what is asserted.

The generic rational point of a suitable symmetric power mentioned
on 14:71 is a sufficient condition. The degree-zero-cycle criterion
above is the exact equivalent statement; effectiveness of a chosen
generic cycle is an extra condition and is not inserted into it.
Question 1 holds for a curve by its normalized Poincar\'e divisor,
and for every abelian variety by its diagonal. It therefore holds for
all the abelian cases specifically listed in C.4. It is false for
arbitrary smooth projective varieties: the opened
\citet[Theorem 1.4]{cyVoisin2025} supplies smooth projective threefolds
with representable $\mathrm{CH}_0$ but without a universal zero-cycle.
We use that counterexample as a known theorem, without claiming to
reproduce its construction here.

\paragraph{C.4d (Proved characterization of Question 2 for abelian varieties).}
For an abelian variety $A/k$, Question 2 holds integrally if and only
if $A$ is a direct summand, as an unpolarized abelian variety, of a
finite product of Jacobians of smooth projective curves. ``Direct
summand'' means actual morphisms $s:A\to J$, $r:J\to A$ with
$rs=\mathrm{id}_A$, not merely an isogeny decomposition.

\emph{Proof.}
The canonical Poincar\'e divisor identifies
$H^{2g-1}(A)(g-1)$ with $H^1(\widehat A)$. Its integral inverse exists
exactly when an integral one-cycle
$Z\in\mathrm{CH}_1(A\times\widehat A)$ induces
$u_Z=\mathrm{id}_{\widehat A}$. Write $Z$ as an integral combination
of images of normalizations of curves $C_j$, with projection maps
$f_j:C_j\to A$ and $g_j:C_j\to\widehat A$. After choosing basepoints,
the induced homomorphism is
\[
 u_Z=\sum_j n_j(g_j)_*d_{C_j}f_j^*:
 \widehat A\longrightarrow\widehat A.
\]
Thus the identity factors through $\prod_j J(C_j)$. Dualizing gives
the claimed splitting of $A$.

Conversely, dualize a splitting to obtain
$\widehat A\xrightarrow{s}\prod J(C_j)\xrightarrow{r}\widehat A$.
Jacobian autoduality identifies each component of $s$ with $f_j^*$
for a based morphism $f_j:C_j\to A$; the universal property of the
Jacobian identifies each component of $r$ with a based morphism
$g_j:C_j\to\widehat A$. The sum of the one-cycles
$(f_j,g_j)_*[C_j]$, with the signs prescribed by the chosen autoduality,
induces $rs$. It therefore supplies the required inverse. Constant
components contribute zero and cause no problem.

\paragraph{C.4e (Proved equivalence for principal polarizations; known counterexamples).}
For a principally polarized complex abelian variety $(A,\theta)$ of
dimension $g$, Question 2 is equivalent to algebraicity, with
\emph{integral} coefficients, of
\[
 \gamma_\theta=\frac{\theta^{g-1}}{(g-1)!}
 \in H^{2g-2}(A,\mathbb Z).
\]
One direction can be read directly from a curve decomposition of
$\gamma_\theta$: the induced map $\widehat A\to A$ is the inverse
principal polarization, as is checked in a symplectic basis; composing
with that polarization gives the identity factorization in C.4d.
For the other direction, a direct summand of a product of Jacobians
satisfies the integral Hodge conjecture for curves. Indeed the theorem
for the product is \citet[Theorem 1.2]{cyBeckmannFortman2023}; pushing
an integral Hodge homology class along the inclusion and then pushing
an algebraic representative back along the retraction proves the
claim for the summand. Apply it to $\gamma_\theta$.
This agrees with \citet[Theorem 1.4, Proposition 2.1,
Corollary 2.5]{cyVoisin2024}; the published theorem number is 1.4,
although an earlier revision labels it 1.3.

Over arbitrary algebraically closed fields, the relevant extension is
\citet[Theorem 1.2]{cyScaviaSuzuki2025}: it retains the integral
$\ell$-adic conditions for every $\ell\ne p$ and the crystalline
condition at $p$. Thus a single prime-to-$p$ assertion must not be
silently promoted to an integral assertion including $p$. For a
principal polarization the polarization itself is already algebraic.

The modern universal integral assertion is false.
\citet[Theorem 1.1]{cyEngelFortmanSchreieder2026} proves that on a very
general complex principally polarized abelian variety of dimension
$g\ge4$, a curve class is an even multiple of $\gamma_\theta$.
Question 2 therefore fails there. Their Corollary 1.2 identifies the
cycle-class image in dimensions $4$ and $5$ as exactly
$2\mathbb Z\gamma_\theta$. We have checked the theorem and corollary
in version 3 of the primary source, and do not substitute a historical
claim of openness for this newer result.

\paragraph{C.4f (Proved cases and the precise outstanding small-dimensional case).}
Question 2 holds for a curve: the Abel map identifies its $H^1$ with
that of its Jacobian and its graph gives the needed one-cycle.
It holds for Jacobians and products of Jacobians, in every
characteristic, by the splitting in C.4d; this includes all products
of elliptic curves. It holds for principally polarized abelian
varieties of dimension at most three: over an algebraically closed
field these are Jacobians of compact-type curves, hence products of
Jacobians \citep[Theorem 4]{cyOortUeno1973}.

These statements do not decide Question 2 for every
\emph{non-principally-polarized} abelian surface or threefold.
Knowing the integral Hodge conjecture for curves on $A$ alone is
insufficient for that inference; the general criterion over
$\mathbb C$ uses $A\times\widehat A$, equivalently the actual splitting
in C.4d. Passing to an isogenous principally polarized variety is not
an integral splitting. This is a remaining gate in the precise scope
of C.4.

\paragraph{C.4g (Proved reduction of Question 2).}
Assume Question 1 for $X$ and Question 2 for $A_X$. Then Question 2
holds for $X$. To check the direction, take the integral
correspondence giving the inverse of
$H^1(A_X)\to H^1(X)$. Its transpose, under Poincar\'e duality, gives
the corresponding isomorphism in the dual degree. Compose it with
the integral one-cycle giving the inverse of the Poincar\'e map for
$A_X$. The identification
$\operatorname{Pic}^0(A_X)\simeq\operatorname{Pic}^0(X)$ identifies
the remaining degree-one target. Functoriality of the Poincar\'e
bundle makes the composite the inverse of
$H^{2n-1}(X)(n-1)\to H^1(B_X)$. Composition of integral
correspondences preserves integrality, so no denominator has been
introduced. The reduction is conditional on Question 1; it does not
prove that question for $X$.

% Independent contribution by /root/frobenius.
% Uses citation keys from cycles/references.bib and C4_references.bib.
\subsection{C.4 supplement: exact integral bounds without a principal polarization}
\label{sec:c4-nonprincipal-supplement}

\paragraph{C.4h (Proved: the ideal of scalar factorizations).}
Let $A$ be an abelian variety of positive dimension $g$ over an
algebraically closed field $k$, in arbitrary characteristic. We define
\[
 I_J(A)=\{n\in\mathbb Z:\ [n]_A=v\circ u
       \text{ for some }A\xrightarrow{u}\prod_jJ(C_j)
                           \xrightarrow{v}A\},
\]
where each $C_j$ is a smooth projective connected curve over $k$.
This is an ideal of $\mathbb Z$. Indeed, multiplication of $v$ by an
integer realizes integer multiples; the product of two intermediate
products of Jacobians realizes the sum of their two composites.
We will prove below that $I_J(A)\ne0$, and write
$I_J(A)=s_J(A)\mathbb Z$ with $s_J(A)>0$.
Then
\begin{equation}\label{eq:c4-exact-splitting}
 A\text{ is an actual direct summand of a product of Jacobians}
 \quad\Longleftrightarrow\quad s_J(A)=1.
\end{equation}
For completeness, if $v u=1_A$, the maps
\[
 A\times\ker v\longrightarrow J,\quad(a,b)\longmapsto u(a)+b,
 \qquad
 J\longrightarrow A\times\ker v,\quad
 x\longmapsto(v(x),x-u v(x))
\]
are mutually inverse as group schemes. The kernel is a connected smooth
proper group: smoothness follows from the splitting of the differential
of $v$, and connectedness from the displayed product decomposition.
It is consequently an abelian variety.

The definition also gives
\[
 s_J(\widehat A)=s_J(A),\qquad
 s_J(A\times B)=\operatorname{lcm}(s_J(A),s_J(B)).
\]
For the first assertion we dualize a factorization and use the canonical
principal polarizations of the Jacobians. For the second we restrict and
project a factorization to each factor, and conversely combine two
factorizations of the same integer.

Here is the exact isogeny calculation needed later. If
$f:A\to B$ is an isogeny and the \emph{group scheme} $\ker f$ is killed
by $e>0$, the quotient property gives a homomorphism $h:B\to A$ with
$h f=[e]_A$; surjectivity of $f$ also gives $f h=[e]_B$.
Thus
\begin{equation}\label{eq:c4-isogeny-divisibility}
 eI_J(B)\subset I_J(A),\qquad eI_J(A)\subset I_J(B).
\end{equation}
In particular an isogeny whose kernel is killed by an integer prime to
$\ell$ preserves the $\ell$-adic valuation of $s_J$.
We may always take $e=\deg f$; this remains true for inseparable
isogenies \citep[Section 8, paragraph after Proposition 8.1,
p.~12 of the corrected text]{frC4Milne1986}.
None of these statements replaces an integral retraction by a
rational inverse to an isogeny.

\paragraph{C.4i (Proved: a uniform factorial bound in every characteristic).}
For every $A$ as above,
\begin{equation}\label{eq:c4-factorial-bound}
             s_J(A)\mid(2g-1)!.
\end{equation}
We spell out the algebraic construction, including the integer in the
composite. This is a quantitative application of the proof of
\citet[Proposition 3.4, equations (3.1)--(3.3)]{cyScaviaSuzuki2025},
not an inference from rational cohomology alone.

We first record its useful form. Suppose $X$ has dimension $d$,
$\theta=c_1(L)\in\mathrm{NS}(X)$, and an integral curve cycle
$\alpha$ has numerical class
\[
                  \alpha=n\,\frac{\theta^{d-1}}{(d-1)!}.
\]
Write $\alpha=\sum_j m_j(f_j)_*[C_j]$ using normalizations of its
integral curve components. Such normalizations are smooth over the
algebraically closed field, including in characteristics $2$ and $3$.
Choose base points and let $\bar f_j:J(C_j)\to X$ be the induced
homomorphisms. Set $J=\prod_jJ(C_j)$, let
$\lambda_j:J(C_j)\xrightarrow{\sim}\widehat{J(C_j)}$ be its canonical
polarization, and define
\[
 U:X\longrightarrow J,\quad
 U=(\lambda_j^{-1}\widehat{\bar f_j}\,\varphi_\theta)_j,
 \qquad
 V:J\longrightarrow X,\quad V=\sum_jm_j\bar f_j.
\]
Matsusaka's identity gives the equality of homomorphisms
\begin{equation}\label{eq:c4-matsusaka-integer}
               VU=[\,n\chi(L)\,]_X,\qquad
               \chi(L)=\deg\frac{\theta^d}{d!}.
\end{equation}
The sign is the positive sign in equation (3.1) of the cited proof:
its two occurrences of Matsusaka's $\delta$ have the same negative
sign and cancel. Negative $m_j$ are carried by $V$, and cause no
positivity requirement on $\theta$. Translating $f_j$ to send its
base point to zero does not change its numerical curve class.

We can also check the coefficient in \eqref{eq:c4-matsusaka-integer}
directly. In an $\ell$-adic symplectic basis for the nondegenerate
alternating form of $\theta$, write the form as
$\sum_{a=1}^{d}t_a e_a\wedge f_a$.
The complementary exterior power
$\theta^{d-1}/(d-1)!$ has coefficients
$\prod_{b\ne a}t_b$. Multiplication by the original form therefore
gives the scalar $\prod_b t_b=\chi(L)$.
The same coefficient follows for every $\theta$ by a polynomial
identity. Faithfulness of the Tate module realization of
homomorphisms of abelian varieties makes this an equality of
homomorphisms, as in the cited Matsusaka proof.

Now take $X=A\times\widehat A$, $d=2g$, and
$\theta=c_1(\mathcal P_A)$. The associated symmetric homomorphism
$\varphi_\theta:X\to\widehat X=\widehat A\times A$ exchanges the two
factors, so it is an isomorphism. The Riemann--Roch formula
$\deg\varphi_\theta=\chi(\mathcal P_A)^2$ gives
$\chi(\mathcal P_A)=\pm1$ (in the usual orientation it is $(-1)^g$).
The integral intersection cycle $\alpha=\theta^{2g-1}$ satisfies the
preceding hypothesis with $n=(2g-1)!$.
Equation \eqref{eq:c4-matsusaka-integer} thus factors multiplication
by $\pm(2g-1)!$ on $X$ through $J$.
We adjust the sign of $V$, then compose with the inclusion of $A$
and the projection to $A$. This proves \eqref{eq:c4-factorial-bound}.
The Riemann--Roch and Poincar\'e degree facts are also explicitly
used in the proof of Corollary 1.3 of
\citet{cyScaviaSuzuki2025}.

\paragraph{C.4j (Proved: a simultaneous polarization criterion in dimensions at most three).}
Define the positive integer
\[
 \delta(A)=\gcd\{\chi(L):L\in\mathrm{Pic}(A)\}.
\]
Equivalently, we may take the gcd only over ample $L$.
To see this, fix an ample $H$ and any line bundle $L$.
The function $t\mapsto\chi(L\otimes H^t)$ is a polynomial of degree at
most $g$ with integer values. For all sufficiently large consecutive
integers $t$, the line bundles are ample. Integral finite differences,
or interpolation in the binomial basis shifted to these integers,
express its value at $0$ as an integral linear combination of these
values. Hence their common divisor divides $\chi(L)$.

If $1\le g\le3$, then
\begin{equation}\label{eq:c4-polarization-gcd}
 \boxed{\quad s_J(A)\mid
             \gcd\bigl((2g-1)!,\delta(A)\bigr).\quad}
\end{equation}
Indeed, let $L$ be ample. The group-scheme form of the polarization
descent theorem gives a maximal isotropic subgroup scheme
$H\subset K(L)$, a quotient isogeny $f:A\to B=A/H$, and a principal
polarization $M$ on $B$ with $f^*M\simeq L$.
This is \citet[Section 23, Theorem 4 and Corollary 1,
pp.~233--234]{frC4Mumford1970}; its proof explicitly includes
$\mu_p$ and $\alpha_p$, so it has no separability restriction.
Riemann--Roch gives
\[
                       \deg f=\chi(L).
\]
By \citet[Theorem 4, p.~2 of the opened preprint]{cyOortUeno1973}, the principally
polarized variety $B$, in these dimensions and over any algebraically
closed field, is the Jacobian of a compact-type curve, and therefore
the product of the Jacobians of its smooth irreducible components.
Thus $\chi(L)\in I_J(A)$ by \eqref{eq:c4-isogeny-divisibility}.
Taking the gcd and using \eqref{eq:c4-factorial-bound} proves the claim.

Consequently it suffices that, for \emph{each} prime
$\ell\mid(2g-1)!$, there exist an ample line bundle $L_\ell$ on $A$
with $\ell\nmid\chi(L_\ell)$.
The choices of $L_\ell$ may be different. We take the finitely many
resulting factorizations, together with the factorial factorization,
and integers $a_j$ with $\sum_j a_j n_j=1$. If
$v_j u_j=[n_j]_A$, then
\[
                 u=(u_j)_j,\qquad v=\sum_j a_jv_j
                 \quad\text{satisfy}\quad vu=1_A.
\]
This constructs a single actual retraction. In particular:
\[
\begin{array}{c|c|c}
 \dim A & \text{universal bound for }s_J(A)
        &\text{sufficient local polarization checks}\\ \hline
 2 & 6 & \ell=2,3\\
 3 & 120 & \ell=2,3,5.
\end{array}
\]
If $k$ has characteristic $2$ or $3$, the check at that very prime
is retained. The construction uses finite subgroup schemes rather
than their sets of geometric points.

The kernel exponent can improve the quantitative divisibility:
if a principalization has kernel killed by $e$, we obtain
$s_J(A)\mid\gcd((2g-1)!,e)$.
For a separable polarization of type
$(d_1,\ldots,d_g)$, its kernel has exponent $d_g$, so we may use
$e=d_g$. For an inseparable polarization we use its actual
group-scheme kernel exponent, or the always valid bound
$e=\chi(L)$; a prime-to-characteristic elementary-divisor
description must not be substituted for the full kernel.

\paragraph{C.4k (Proved: what the N\'eron--Severi calculation establishes).}
For an abelian surface, choose a basis $D_1,\ldots,D_\rho$ of
$\mathrm{NS}(A)$, and write
$q_i=D_i^2/2$ and $b_{ij}=D_iD_j$. Then
\begin{equation}\label{eq:c4-delta-surface}
 \delta(A)=\gcd\bigl(\{q_i\}_i,\{b_{ij}\}_{i<j}\bigr).
\end{equation}
Indeed,
$\chi(\sum n_iD_i)=\sum_iq_in_i^2+\sum_{i<j}b_{ij}n_in_j$;
the reverse divisibility follows by evaluating at $D_i$ and
$D_i+D_j$. Thus this particular sufficient test is a finite
integral lattice calculation.
In dimension three we may compute the analogous gcd from the
Euler polynomial on the finite grid $0\le n_i\le3$: its successive
finite differences are integral, and the binomial-basis
interpolation expresses every other integer value as an integral
combination of these grid values.

If $\mathrm{NS}(A)=\mathbb Z[L]$ for a primitive ample $L$, then
$\delta(A)=\chi(L)$, since $\chi(L^n)=n^g\chi(L)$.
In particular a complex surface with primitive polarization type
$(1,d)$ is settled whenever $\gcd(d,6)=1$; a complex threefold
with primitive polarization type $(1,d_2,d_3)$ is settled whenever
$\gcd(d_2d_3,30)=1$.
These are sufficient conditions and give no counterexample when
the gcd is larger than one.

There is a precise limitation on the proposed coprime-isogeny
method. An isogeny $A\to B$ of degree prime to $\ell$, with $B$
principally polarized, pulls a principal line bundle back to an
ample $L$ with $\ell\nmid\chi(L)$. The converse follows from
Mumford's principalization. Therefore, when $\ell\mid\delta(A)$,
there is \emph{no} such prime-to-$\ell$ principalization in the same
dimension, even if we allow all possible quotient polarizations.
This rules out that particular proof strategy in the residual
Picard-rank-one cases; it does not rule out a retraction through
Jacobians of higher dimension.

\paragraph{C.4l (Proved limitation of divisor products; no geometric counterexample asserted).}
We can make the residual surface calculation more explicit.
Let $A/\mathbb C$ have $\mathrm{End}(A)=\mathbb Z$ and
$\mathrm{NS}(A)=\mathbb Z[L]$, with $L$ of type $(1,d)$.
Let $\widehat L$ be the primitive dual polarization, normalized by
$\varphi_{\widehat L}\varphi_L=[d]$.
On $A\times\widehat A$ put
$D=\operatorname{pr}_1^*c_1(L)$,
$\widehat D=\operatorname{pr}_2^*c_1(\widehat L)$ and
$P=c_1(\mathcal P_A)$.
The N\'eron--Severi group is generated by
$D,\widehat D,P$: the cross summand is
$\operatorname{Hom}(A,\widehat{\widehat A})=\operatorname{End}(A)$.
The endomorphisms induced on $\widehat A$ by integral curve cycles
that are integer combinations of products of three divisors form
exactly
\begin{equation}\label{eq:c4-divisor-product-ideal}
                         \gcd(6,d)\,\mathbb Z\,1_{\widehat A}.
\end{equation}

Here is a direct calculation, avoiding an assumption that all
integral curve classes are generated by divisors. Choose an
integral symplectic basis $x_1,x_2,x_3,x_4$ for $H^1(A,\mathbb Z)$
and the dual basis $y_1,\ldots,y_4$ for $H^1(\widehat A,\mathbb Z)$
so that
\[
 D=x_1x_2+d\,x_3x_4,\quad
 \widehat D=d\,y_1y_2+y_3y_4,\quad
 P=\sum_{a=1}^{4}x_ay_a.
\]
Products here mean exterior products. The normalized class
$P^3/6$ has the cross component inducing $1_{\widehat A}$, and
expansion gives
\begin{equation}\label{eq:c4-surface-exterior}
                 D\widehat D P=-d\,P^3/6.
\end{equation}
The only degree-three monomials in these divisor classes with
K\"unneth bidegree $(3,3)$ are $P^3$ and $D\widehat D P$;
the other monomials have no component inducing a map
$\widehat A\to\widehat A$. They therefore give precisely the
integers generated by $6$ and $d$.
The identification of the cross component follows either from
the curve-factorization formula in C.4a or from
\eqref{eq:c4-matsusaka-integer} for $P$.
The script \texttt{verify\_c4\_exterior.py} checks the exterior
identity symbolically, together with its factorial normalization.

Thus for $d=2$, $3$, or $6$ the divisor-product calculation itself
cannot produce the desired identity. This is a proved statement
about that subring of algebraic cycles, not a statement about
all of $\mathrm{CH}_1(A\times\widehat A)$.
The rational Hodge group calculation does not justify replacing
the latter integral group by the former.

\paragraph{The scope of the available comparison theorems.}
The opened results impose the following precise limits.
The converse from the integral Hodge conjecture for curves on
$A$ to the direct-summand property requires a unimodular line
bundle in \citet[Theorem 1.1 and the paragraph immediately
following it]{cyBeckmannFortman2023} and
\citet[Theorem 1.1 and the paragraph following it]{cyScaviaSuzuki2025}.
Without that hypothesis the applicable complex criterion is on
$A\times\widehat A$
\citep[Theorem 1.4]{cyVoisin2024}.
Zarhin's construction makes $(A\times\widehat A)^4$ principally
polarizable, but its dimension is $8g$, beyond the
dimension-three Jacobian theorem; it alone supplies no retraction
through Jacobians \citep[Remarks 16.12--16.13]{frC4Milne1986}.
Finally the opened 2026 counterexample theorem concerns very
general principally polarized varieties of dimension at least
four \citep[Theorem 1.1 and Corollary 1.5]{cyEngelFortmanSchreieder2026}.
Its hypotheses do not include a product $A\times\widehat A$
with $A$ a nonprincipal surface. Applying its degeneration
obstruction to that restricted family would require a new argument.

% Independent addition by /root/frobenius, after C.4h--l.
% New keys are in C4_prym_references.bib; other keys are already in C4 references.
\subsection{A Prym improvement: only a possible factor of two for surfaces}
\label{sec:c4-prym-improvement}

\paragraph{C.4m (Proved: twice the identity for principally polarized varieties of dimension at most five).}
Let $(B,\Xi)$ be a principally polarized abelian variety of dimension
$1\le h\le5$ over an algebraically closed field of arbitrary characteristic.
Then
\begin{equation}\label{eq:c4-ppav-five-two}
                         2\in I_J(B).
\end{equation}
More precisely, there is an integral curve cycle $\alpha$ on $B$ with
\begin{equation}\label{eq:c4-twice-minimal-cycle}
       \operatorname{cl}(\alpha)
          =2\,\frac{\xi^{h-1}}{(h-1)!},\qquad \xi=c_1(\Xi),
\end{equation}
in rational $\ell$-adic cohomology for every $\ell$ different from the
characteristic. We only need this equality for one such $\ell$ to obtain
the equality of homomorphisms in \eqref{eq:c4-ppav-five-two}.

\emph{Proof in characteristic different from two.}
For $h\le3$ we can use the product-of-Jacobians theorem
\citep[Theorem 4, p.~2 of the opened preprint]{cyOortUeno1973}.
For $h=4,5$, \citet[Theorem 6.5(i), p.~178]{frC4Beauville1977}
realizes $(B,\Xi)$ as a generalized Prym. Let $\widetilde C$ be its
nodal covering curve, with smooth normalization components
$C_1,\ldots,C_t$. The composition
\[
 j:B\longrightarrow J(\widetilde C)
            \longrightarrow J:=\prod_{a=1}^{t}J(C_a)
\]
is a homomorphism of abelian varieties. The generalized Jacobian in
the middle is allowed to be semiabelian; the target $J$ is abelian.
Proposition 2.2, p.~155, and Theorem 5.4, pp.~175--176, of the same
paper give
\[
                         j^*\Theta_J\equiv_{\mathrm{alg}}2\Xi.
\]
Indeed, the first statement identifies the generalized theta class
with pullback of the product theta class of the normalization, and
the second identifies its restriction to the Prym.
Writing $\lambda_B$ and $\lambda_J$ for the principal polarizations,
we define the \emph{integral} polarized adjoint
\[
             q=j^\dagger:=\lambda_B^{-1}\widehat j\lambda_J:J\to B.
\]
The pullback identity gives
\begin{equation}\label{eq:c4-prym-adjoint}
 qj=\lambda_B^{-1}(\widehat j\lambda_Jj)
    =\lambda_B^{-1}(2\lambda_B)=[2]_B.
\end{equation}
We have inverted only principal polarizations, which are actual
isomorphisms. In particular, a possible finite kernel of $j$ has not
been divided out.

We also record the curve needed for specialization. Choose Abel maps
$\iota_a:C_a\to J(C_a)$, and let $q_a$ be the restriction of $q$ to
the $a$th factor. Put
\[
                 \alpha=\sum_a(q_a\iota_a)_*[C_a].
\]
Components mapped to a point contribute zero. The symmetric map
$\widehat B\to B$ associated with this curve class is
\[
 \sum_aq_a\lambda_a^{-1}\widehat{q_a}
   =q\lambda_J^{-1}\widehat q
   =\lambda_B^{-1}\widehat j\lambda_Jj\lambda_B^{-1}
   =2\lambda_B^{-1}.
\]
Under Poincar\'e duality, the class of a curve is the alternating
pairing obtained by integrating products of two degree-one classes
over that curve. The preceding equality says that this pairing is
twice the inverse principal-polarization pairing. In a symplectic
basis, the complementary exterior power
$\xi^{h-1}/(h-1)!$ gives precisely that inverse pairing.
Since degree-one products span $H^2(B,\mathbb Q_\ell)$, the pairing
determines the curve class. This proves
\eqref{eq:c4-twice-minimal-cycle}.

\emph{Characteristic two.}
By \citet[Theorem 2.4.1, p.~286, and Corollary 2.4.2 with its proof,
p.~287]{frC4Oort1971}, a principally polarized abelian variety lifts,
with its polarization, to an abelian scheme $\mathcal B$ over
$W(k)$. A principal polarization is separable even in characteristic
two, so the stated hypothesis is satisfied. On its geometric generic
fibre the preceding construction supplies an integral curve cycle
with class $2\xi^{h-1}/(h-1)!$.
We use only the separability-implies-smoothness direction of
Theorem 2.4.1: the converse in its printed statement is explicitly
retracted in \citet[p.~3 of the opened preprint]{cyOortUeno1973}.

The cycle and its finitely many components descend to a finite
extension of the fraction field. After the corresponding finite
extension of complete discrete valuation rings, we take the closures
of those components and their special-fibre cycles. This gives an
integral cycle $\alpha_0$ on $B$; the residue field remains $k$ because
it is algebraically closed. The specialization map commutes with
the $\ell$-adic cycle class map for $\ell\ne2$
\citep[Proposition 20.3, p.~399, and Example 20.3.5,
p.~400]{frC4Fulton1998}. Hence
\[
                 \operatorname{cl}(\alpha_0)
                    =2\,\frac{\xi^{h-1}}{(h-1)!}.
\]
In particular, the same equality holds numerically. Applying the
integral Matsusaka construction in
\eqref{eq:c4-matsusaka-integer}, with $n=2$ and $\chi(\Xi)=1$,
factors $[2]_B$ through the Jacobians of the normalizations of the
components of $\alpha_0$. This is an equality of algebraic
homomorphisms in characteristic two. We have specialized an actual
integral cycle; no inference from a crystalline realization, and
no inversion of two, enters the argument.\hfill$\square$

\paragraph{C.4n (Proved: the surface obstruction divides two, in every characteristic).}
For any abelian surface $A$ over an algebraically closed field,
\begin{equation}\label{eq:c4-surface-two}
 \boxed{\qquad s_J(A)\mid\gcd(2,\delta(A)).\qquad}
\end{equation}
Consequently every surface admitting an ample line bundle of odd
Euler characteristic is an actual direct summand of a product of
Jacobians. In particular, for a complex surface whose
N\'eron--Severi group is generated by a polarization of type $(1,d)$,
every odd $d$ satisfies Question 2 affirmatively.

\emph{Proof.}
We first make a useful consequence of C.4m explicit. If $X$ has
dimension four or five, principalizing any ample $L$ gives an
isogeny $f:X\to B$ of degree $\chi(L)$ to a principally polarized
variety of the same dimension. This is the all-characteristic
group-scheme construction used in C.4j
\citep[Section 23, Theorem 4 and Corollary 1,
pp.~233--234]{frC4Mumford1970}.
By \eqref{eq:c4-isogeny-divisibility} and C.4m,
$2\chi(L)\in I_J(X)$. Taking the gcd over ample $L$, and using
the equality of the ample and unrestricted gcds proved in C.4j,
we obtain
\begin{equation}\label{eq:c4-four-five-delta}
              s_J(X)\mid2\delta(X)\qquad(\dim X=4,5).
\end{equation}

Now take $X=A\times\widehat A$. The Poincar\'e line bundle is
unimodular and has Euler characteristic $1$ on this fourfold.
Therefore $\delta(X)=1$. Since $A$ is an actual direct factor of
$X$, restriction and projection of a factorization on $X$ show
$I_J(X)\subset I_J(A)$. Equation \eqref{eq:c4-four-five-delta}
therefore gives $2\in I_J(A)$.
C.4j also gives $\delta(A)\in I_J(A)$, proving
\eqref{eq:c4-surface-two}. If $\delta(A)$ is odd, an integral
B\'ezout combination of the two factorizations gives $1\in I_J(A)$;
the explicit splitting maps in C.4h then give the required actual
direct summand. For a Picard-rank-one surface of type $(1,d)$,
C.4k gives $\delta(A)=d$.\hfill$\square$

For a finite construction of the factorization on $X$, choose an
ample $H$, put $p(t)=\chi(\mathcal P_A\otimes H^t)$, and take
$T\ge1$ so large that the five line bundles with
$t=T,T+1,\ldots,T+4$ are ample. The integral interpolation identity
\[
 1=p(0)=\sum_{j=0}^{4}(-1)^j
       \binom{T+j-1}{j}\binom{T+4}{4-j}\,p(T+j)
\]
expresses $2$ as an integral combination of five available
factorizations of $2p(T+j)$. The identity follows by evaluating the
degree-at-most-four Lagrange interpolation formula at zero; its
displayed coefficients are integers. Thus the gcd step does not
require an infinite construction.

There is also a construction using just one principalization of
$A\times\widehat A$, besides the factorial factorization.
Fix an ample $H$ on $X$, and take $N$ sufficiently large and divisible
by $3\cdot4!=72$. Then
$L=\mathcal P_A\otimes H^N$ is ample and
\[
 \chi(L)
 =1+\sum_{r=1}^{4}\frac{N^r}{r!(4-r)!}
             \deg\bigl(c_1(\mathcal P_A)^{4-r}c_1(H)^r\bigr)
 \equiv1\pmod3.
\]
Every displayed coefficient after the constant is an integer
divisible by $3$. If $m=\chi(L)$, principalization and C.4m give
$2m\in I_J(A)$, whereas C.4i gives $6\in I_J(A)$.
The equality $\gcd(6,2m)=2$ provides the factorization of $[2]$.
This version makes the removal of the possible factor of three
particularly concrete.

% Independent continuation by /root/frobenius, after C.4m--n.
% The final scope paragraph supersedes the earlier intermediate scope paragraphs.
\subsection{The Prym--Tyurin improvement for arbitrary threefolds}
\label{sec:c4-pt-six}

\paragraph{C.4o (Proved: six times the identity on every principally polarized sixfold).}
Let $(B,\Xi)$ be a principally polarized abelian variety of dimension
six over an algebraically closed field of arbitrary characteristic.
There is an integral curve cycle $\alpha$ such that
\begin{equation}\label{eq:c4-six-minimal}
       \operatorname{cl}(\alpha)=6\,\frac{\xi^5}{5!},
       \qquad \xi=c_1(\Xi),
\end{equation}
and consequently $6\in I_J(B)$.

\emph{Proof.}
Over $\mathbb C$, the integral-cycle assertion is exactly
\citet[Corollary 2, p.~165]{frC4ADFIO2020}.
Their Theorem 1 on the same page proves dominance of the
Prym--Tyurin map of exponent six. The meaning of that exponent is
the identity $j^*\Theta_C\equiv6\Xi$ for the map
$j:B\hookrightarrow J(C)$ on a general member, as specified on
p.~164. Polarized adjunction gives $j^\dagger j=[6]$;
the corollary includes the specialization to \emph{every} complex
principally polarized sixfold. In particular, we do not replace the
word ``every'' by an unsupported assertion about the generic member.

We justify transfer from $\mathbb C$ to an arbitrary algebraically
closed field $F$ of characteristic zero in a form that also preserves
the integral coefficient. Define $(B,\Xi)$ over a finitely generated
subfield $K\subset F$, and choose an embedding $K\hookrightarrow
\mathbb C$. Write the curve furnished by the complex theorem as
\[
             \alpha=\sum_a m_a(f_a)_*[C_a],
\]
with smooth connected projective $C_a$, chosen base points, and
induced homomorphisms $\bar f_a:J(C_a)\to B_{\mathbb C}$.
Its class implies the identity
\begin{equation}\label{eq:c4-six-spread-identity}
 \sum_a m_a\bar f_a\lambda_a^{-1}
                  \widehat{\bar f_a}\lambda_B=[6]_B
\end{equation}
by the Matsusaka calculation in C.4i. The finitely many curves,
base points, maps, and identity spread to a nonempty scheme $S$ of
finite type over $K$, after shrinking so that the curves are smooth
and projective. Their relative Jacobians and principal polarizations
spread as well. A closed point of $S$ has residue field finite over
$K$, which embeds over $K$ into $F$. Its fibre therefore supplies
the same data over $F$, still satisfying
\eqref{eq:c4-six-spread-identity}.

As in the alternating-pairing argument in C.4m, the associated
curve on $B_F$ has class $6\xi^5/5!$: the symmetric map
$\widehat B\to B$ defined by the left side before composition
with $\lambda_B$ is $6\lambda_B^{-1}$. Thus this argument transfers
the full curve-class assertion, not merely an unpolarized
factorization whose adjoint relation might have been lost.

If $\operatorname{char}F=p>0$, lift $(B,\Xi)$ over $W(F)$ using
\citet[Corollary 2.4.2 and its proof, p.~287]{frC4Oort1971}.
The geometric generic fibre has characteristic zero, so the preceding
paragraph supplies its integral curve class. We descend the finite
cycle data to a finite extension of the fraction field, close the
cycle, and specialize it as in C.4m. The compatibility in
\citet[Example 20.3.5, p.~400]{frC4Fulton1998} yields
\eqref{eq:c4-six-minimal} in rational $\ell$-adic cohomology for
$\ell\ne p$, hence numerically. The integral Matsusaka formula,
with $\chi(\Xi)=1$, gives $6\in I_J(B)$.
This argument includes $p=2$ and $p=3$: the cycle remains integral
even when the coefficient six is divisible by the characteristic.
\hfill$\square$

\paragraph{C.4p (Proved: the threefold obstruction divides six).}
For every abelian threefold $A$ over an algebraically closed field,
\begin{equation}\label{eq:c4-threefold-six}
 \boxed{\qquad s_J(A)\mid\gcd(6,\delta(A)).\qquad}
\end{equation}
In particular, Question 2 holds for $A$ if, for each of
$\ell=2,3$, there is an ample line bundle $L_\ell$ with
$\ell\nmid\chi(L_\ell)$. The two line bundles may be different.
For a complex Picard-rank-one threefold with primitive polarization
of type $(1,d_2,d_3)$, the sufficient condition is
$\gcd(d_2d_3,6)=1$.

\emph{Proof.}
For a sixfold $X$, principalizing an ample line bundle $L$ and
applying C.4o gives
\[
                  6\chi(L)\in I_J(X).
\]
The same ample-gcd argument as before gives
\[
                         s_J(X)\mid6\delta(X).
\]
For $X=A\times\widehat A$, the Poincar\'e line bundle has Euler
characteristic $-1$, so $\delta(X)=1$. Restriction and projection
give $6\in I_J(A)$. C.4j gives
$\delta(A)\in I_J(A)$; taking the gcd proves
\eqref{eq:c4-threefold-six}. The local condition at two and three
means exactly that this gcd is one. Combining the finite
factorizations by B\'ezout and using C.4h produces an actual direct
summand. For Picard rank one, the Riemann--Roch calculation gives
$\delta(A)=\chi(L)=d_2d_3$.\hfill$\square$

The gcd construction is finite here too. Choose an ample $H$ on
$X=A\times\widehat A$ and $T$ so large that
$\mathcal P_A\otimes H^{T+j}$ is ample for $0\le j\le6$.
For $p(t)=\chi(\mathcal P_A\otimes H^t)$, the identity
\[
 -1=p(0)=\sum_{j=0}^{6}(-1)^j
          \binom{T+j-1}{j}\binom{T+6}{6-j}\,p(T+j)
\]
combines seven available factorizations of $6p(T+j)$ into one
of $-6$, and changing the sign of the return homomorphism gives
one of $6$.

\paragraph{The remaining low-dimensional cases.}
The strongest bounds established in this contribution are now
\[
\begin{array}{c|c|c}
 \dim A & \text{proved bound for }s_J(A)
        &\text{possible residual primes here}\\\hline
 2 & \gcd(2,\delta(A)) & 2\\
 3 & \gcd(6,\delta(A)) & 2,3.
\end{array}
\]
Thus the earlier possible prime three for surfaces and prime five
for threefolds have been removed. We have neither proved the
unrestricted direct-summand assertion for all remaining varieties
nor produced a counterexample among them. For example, a
Picard-rank-one surface of type $(1,2)$ remains untreated here;
type $(1,3)$ is affirmative. For threefolds, type $(1,1,5)$ is
affirmative, while the present bound alone does not decide
type $(1,1,3)$ or type $(1,2,2)$.

The bound for principally polarized sixfolds is sharp on a very
general complex member: \citet[Theorem 1.2]{frC4EGFSOptimality2025}
prove that the smallest algebraic multiple of its minimal curve
class is six. This does not give a counterexample for a threefold,
because the special sixfold $A\times\widehat A$ is not a very
general principally polarized sixfold. Likewise, the evenness
theorem for very general principally polarized fourfolds does not
decide the remaining nonprincipal surface cases. These distinctions
are essential to the exact integral question.

\paragraph{C.4q (Proved: the exact contribution of curves in proper abelian subvarieties).}
Let $A$ be a simple complex abelian variety of dimension
$g\in\{2,3\}$, with
$\operatorname{End}(A)=\mathbb Z$ and
$\operatorname{NS}(A)=\mathbb Z\theta$, where $\theta$ is an
ample primitive class of type $(1,d_2,\ldots,d_g)$.
Put $e=d_g$, and write
\[
 e=u v^2\quad(u\text{ square-free}),\qquad
 \kappa(e)=uv=\prod_{\ell\mid e}\ell^{\lceil v_\ell(e)/2\rceil}.
\]
Let $X=A\times\widehat A$, and let $\mathcal C_{\mathrm{prop}}(X)$
be the subgroup of $H^{4g-2}(X,\mathbb Z)$ generated by the
classes of integral curves contained in translates of proper
abelian subvarieties of $X$. We use the integral cohomological
Fourier transform, followed by the canonical interchange of the
factors of $\widehat X=\widehat A\times A$, to identify
$H^{4g-2}(X,\mathbb Z)$ with $H^2(X,\mathbb Z)$.
If $D$ and $\widehat D$ are the pullbacks of the primitive ample
classes on the two factors and $P=c_1(\mathcal P_A)$, then
\begin{equation}\label{eq:c4-proper-abelian-curve-lattice}
 \boxed{\quad
 \mathcal F_X\bigl(\mathcal C_{\mathrm{prop}}(X)\bigr)
 =\mathbb ZD\oplus\kappa(e)\mathbb ZP
                    \oplus\mathbb Z\widehat D.
 \quad}
\end{equation}
Here the factor interchange is understood in the notation
$\mathcal F_X$. Consequently the smallest positive multiple of
\[
 \mathcal R_A=\frac{P^{2g-1}}{(2g-1)!}
\]
that belongs to $\mathcal C_{\mathrm{prop}}(X)$ is exactly
$\kappa(e)$. In particular, this restricted source of curves
cannot remove any prime dividing $e$ from the direct-summand
problem. This is a statement about a specified subgroup of
curve classes, not a counterexample to algebraicity of
$\mathcal R_A$.

\emph{Proof.}
Write $\lambda=\varphi_\theta:A\to\widehat A$. The primitive
ample generator on $\widehat A$ has homomorphism
$\widehat\lambda=e\lambda^{-1}$; this is integral and primitive
because its elementary divisors are
$(e/d_g,e/d_{g-1},\ldots,e/d_1)$, where $d_1=1$.
The product decomposition of the N\'eron--Severi group and
$\operatorname{Hom}(A,\widehat{\widehat A})=\mathbb Z$
give
\[
 \operatorname{NS}(X)=\mathbb ZD\oplus\mathbb ZP
                         \oplus\mathbb Z\widehat D.
\]
The symmetric homomorphism associated to $aD+bP+c\widehat D$
is
\begin{equation}\label{eq:c4-polarization-block}
 \begin{pmatrix}
  a\lambda&b\,\operatorname{id}_{\widehat A}\\
  b\,\operatorname{id}_A&c\widehat\lambda
 \end{pmatrix}:A\times\widehat A\longrightarrow\widehat A\times A.
\end{equation}
After rationally identifying $\widehat A$ with $A$ by $\lambda$,
its two-by-two scalar matrix has determinant $eac-b^2$.
In particular, a nonzero such homomorphism with image dimension
$g$ satisfies
\begin{equation}\label{eq:c4-rank-one-discriminant}
                             eac=b^2.
\end{equation}

The integral Fourier isomorphism and its preservation of Hodge
type are \citet[Proposition 1, p.~241]{frC4Beauville1983}.
For an inclusion $i:B\hookrightarrow X$ of abelian varieties,
the projection formula and the identity of Poincar\'e bundles
\[
 (i\times\operatorname{id}_{\widehat X})^*\mathcal P_X
   =(\operatorname{id}_B\times\widehat i)^*\mathcal P_B
\]
give directly, on cohomology,
\[
               \mathcal F_X(i_*\alpha)
                   =\widehat i^{,*}\mathcal F_B(\alpha).
\]
This formula is proved by the displayed bundle identity, so we
do not invoke an isogeny-only statement for the inclusion $i$.
Translations act trivially on cohomology.

Every positive-dimensional proper abelian subvariety of $X$
has dimension $g$, since $A$ is simple and $X$ is isogenous to
$A^2$. The Fourier transform of an integral curve class on such
a subvariety $B$ is an integral Hodge class in
$H^2(\widehat B,\mathbb Z)$, hence a divisor class by the
Lefschetz $(1,1)$ theorem. Its pullback to $\widehat X$ has
homomorphism rank at most $g$.
Thus its nonzero block matrix satisfies
\eqref{eq:c4-rank-one-discriminant}. For every prime $\ell\mid e$,
\[
 2v_\ell(b)=v_\ell(e)+v_\ell(a)+v_\ell(c)
                  \ge v_\ell(e)
\]
when $abc\ne0$; if $ac=0$, then $b=0$.
It follows that $\kappa(e)$ divides $b$. This proves one inclusion
in \eqref{eq:c4-proper-abelian-curve-lattice}.

For the converse, we first recall that every integral Hodge
class of curves on any complex abelian variety of dimension
at most three is algebraic. In dimension two this is the
Lefschetz $(1,1)$ theorem. In dimension three it is
\citet[Corollary 3.1.9, pp.~29--30]{frC4Grabowski2004}; its
reduction to principally polarized varieties is Proposition
3.1.8, pp.~28--29. Therefore every integral divisor class on
the dual of such a variety is the Fourier transform of an
integral algebraic curve cycle. In particular, curves in the
two factors of $X$ supply the two diagonal generators $D$ and
$\widehat D$.

Consider the divisor class
\[
                 E=D+\kappa(e)P+u\widehat D.
\]
It is integral, and its associated Hermitian form is positive
semidefinite of rank $g$: both diagonal coefficients are
positive and $e\cdot1\cdot u=\kappa(e)^2$.
Let $q:X\to Y$ be the quotient by the connected kernel of that
form. Then $\dim Y=g$, and $E=q^*\eta$ for an integral ample
class $\eta\in\operatorname{NS}(Y)$. To see the integrality
without dividing by an isogeny degree, write $X=V/\Lambda$.
The alternating form of $E$ is integral on $\Lambda$ and has
rational radical. It descends to an integral positive Riemann
form on the quotient lattice by that radical. The
Appell--Humbert construction gives the asserted class $\eta$.

Choose an integral curve cycle $z$ on $\widehat Y$ whose
Fourier transform is $\eta$. The dual of $q$ is an embedding
$\widehat q:\widehat Y\hookrightarrow\widehat X$.
Fourier compatibility shows that its pushforward, followed by
the factor interchange, has Fourier class $E$ (an irrelevant
overall sign can be absorbed into $z$). Every component is
contained in a proper abelian subvariety. Subtracting the two
diagonal classes produces $\kappa(e)P$. Together with $D$ and
$\widehat D$, this proves the reverse inclusion.
Finally,
$\mathcal F_X(\mathcal R_A)=(-1)^{g+1}P$ under our factor
identification; this is also
\citet[Lemma 3.4, equation (6)]{cyBeckmannFortman2023}.
The exact assertion about multiples of $\mathcal R_A$ follows.
\hfill$\square$

\paragraph{The nearby integral Hodge and Tate theorems.}
The distinction in C.4q matters even though integral curve
classes on the original variety $A$ are understood in these
dimensions. In characteristic zero, Grabowski's theorem applies
to every abelian threefold, with no principal-polarization
hypothesis. Its proof transforms integral divisor classes,
expresses them as differences of ample classes, and principalizes
each ample class; this produces curves on the dual threefold.
It does not assert the algebraicity of the Poincar\'e curve
class on the sixfold $A\times\widehat A$.

Likewise, \citet[Theorem 7.1 and Proposition 7.2, pp.~17--19 of
the opened author version]{frC4Totaro2021} proves the integral
Tate conjecture for abelian threefolds over the separable
closure of a finitely generated field. The formulation uses
$\ell$-adic cohomology for $\ell\ne\operatorname{char}(k)$,
as specified in its Section 1; it does not itself supply the
crystalline statement at the characteristic prime.
Nor does it assert the required integral curve theorem on
$A\times\widehat A$. C.4q identifies an explicit limitation of
trying to obtain that cross class only by pushing curves
forward from its proper abelian subvarieties.

% Further consequences; no priority claim. Owned by /root/frobenius.
\paragraph{C.4r (Proved: a single integral numerical surface class detects the remaining cases).}
Let $A$ be an abelian variety of dimension $g\in\{2,3\}$ over an
algebraically closed field $k$ of arbitrary characteristic.
Put $X=A\times\widehat A$, write $P=c_1(\mathcal P_A)$, and set
\[
 N^r(X)=\mathrm{CH}^r(X)/{\equiv_{\mathrm{num}}},
 \qquad r=2g-2,
\]
with integral coefficients. We regard this torsion-free group as
a subgroup of $N^r(X)_{\mathbb Q}$.
The following are equivalent:
\begin{enumerate}
\item $A$ is an actual direct summand of a product of Jacobians;
\item the rational numerical class
\[
                 \gamma_A=\frac{P^{2g-2}}{(2g-2)!}
                 \in N^{2g-2}(X)_{\mathbb Q}
\]
belongs to the integral numerical lattice $N^{2g-2}(X)$.
\end{enumerate}
Thus we require an integral dimension-two cycle numerically
equivalent to $P^2/2$ for a surface, or to $P^4/24$ for a
threefold. Over $\mathbb C$ the same conditions are also
equivalent to integral algebraicity of the corresponding
integral Hodge class.

\emph{Proof.}
We first justify the integral lattice used in (2).
The rational numerical group is finite-dimensional: numerical
equivalence is coarser than homological equivalence for any
Weil cohomology. Choose finitely many integral complementary
cycle classes $b_1,\ldots,b_t$ whose intersection functionals
span its rational dual. Then
\[
 N^r(X)\longrightarrow\mathbb Z^t,\qquad
 z\longmapsto\bigl(\deg(zb_j)\bigr)_j
\]
is injective. Consequently $N^r(X)$ is a finite free abelian
group.

Suppose (1) holds. By \citet[Theorem 1.2 and the paragraph
immediately following it]{cyScaviaSuzuki2025}, every component
of $\operatorname{ch}(\mathcal P_A)$ is algebraic with
$\mathbb Z_\ell$ coefficients for every prime $\ell$; no
principal polarization on $A$ is required for this implication.
In characteristic $p>0$ this includes the crystalline cycle
map with source $\mathrm{CH}^r(X)\otimes\mathbb Z_p$, as
specified in their equation (1.2).
Map each such representative to $N^r(X)\otimes\mathbb Z_\ell$.
Compatibility of the cycle map with intersections shows that
its pairing with every $b_j$ agrees with that of $\gamma_A$.
The injectivity of the displayed pairing after extending
scalars to $\mathbb Q_\ell$ therefore gives
\[
 \gamma_A\in N^r(X)\otimes\mathbb Z_\ell
 \qquad\text{for every prime }\ell.
\]
For $\ell=p$ the same comparison takes place in the fraction
field of $W(k)$, with the intersection values lying in
$\mathbb Q_p$; the source already has $\mathbb Z_p$
coefficients. Intersecting these local lattices inside
$N^r(X)_{\mathbb Q}$ gives
\[
 N^r(X)_{\mathbb Q}\cap
       \bigcap_{\ell}\bigl(N^r(X)\otimes\mathbb Z_\ell\bigr)
       =N^r(X),
\]
which proves (2). This passage uses neither the Tate conjecture
nor equality of numerical and homological equivalence.

Conversely, choose an integral cycle representing the numerical
class in (2), and intersect it with the integral divisor class
$P$. The resulting integral curve cycle has numerical class
\[
 \frac{P^{2g-1}}{(2g-2)!}
       =(2g-1)\frac{P^{2g-1}}{(2g-1)!}.
\]
Apply the integral Matsusaka factorization
\eqref{eq:c4-matsusaka-integer} on $X$ with the unimodular
Poincar\'e line bundle. Since its Euler characteristic is
$(-1)^g$, we obtain $(2g-1)\in I_J(X)$, after adjusting the
sign of one homomorphism, and therefore $(2g-1)\in I_J(A)$.
For $g=2$, C.4n already gives $2\in I_J(A)$, so
$1=3-2\in I_J(A)$. For $g=3$, C.4p gives
$6\in I_J(A)$, so $1=6-5\in I_J(A)$.
The explicit splitting of C.4h proves (1).
The Matsusaka homomorphism used here depends only on the
numerical curve class, as in C.4i: its Fourier correspondence
is a divisor class, and numerical triviality for divisors
annihilates their homomorphism to the dual abelian variety.

Finally, over $\mathbb C$, (1) implies integral algebraicity
of every component of $\operatorname{ch}(\mathcal P_A)$ by
\citet[Theorem 1.4]{cyVoisin2024} and
\citet[Theorem 1.1]{cyBeckmannFortman2023}. Conversely an
integral cohomological representative supplies the numerical
representative in (2). These divided powers are integral in
cohomology; this does not by itself divide the indicated
intersection cycle in the integral Chow group.
\hfill$\square$

The numerical integrality assertion in (2) remains to be proved
or disproved in the residual cases. Rational algebraicity of
the cohomological Fourier isometry of an abelian surface alone
does not supply the integral numerical cycle required here.

\paragraph{C.4s (Known consequence: every supersingular surface is affirmative).}
Every supersingular abelian surface over an algebraically closed
field of positive characteristic is an actual direct summand of
a product of Jacobians, including in characteristic two.
Indeed, \citet[Lemma 6.2, p.~185]{frC4FuLi2021} prove that it
admits a principal polarization. Their Proposition 6.1,
p.~184, explicitly includes characteristic two in the
N\'eron--Severi lattice calculation used by the proof.
The principally polarized surface case C.4f then applies.
This does not settle arbitrary nonsupersingular surfaces whose
polarization Euler characteristics are all even.

% Supplemental exclusion of one attempted construction; no priority claim.
\paragraph{C.4t (Known exclusion, with the generic strengthening proved: the natural principally polarized product is not a genus-four Jacobian).}
Let $(A,L)$ be a complex abelian surface of polarization type $(1,2)$.
Write $\lambda=\lambda_L:A\to\widehat A$, and let
$\widehat\lambda:\widehat A\to A$ be its dual polarization, normalized by
$\widehat\lambda\lambda=[2]_A$ and
$\lambda\widehat\lambda=[2]_{\widehat A}$.
On $X=A\times\widehat A$, the symmetric homomorphism
\[
 \Lambda=
 \begin{pmatrix}\lambda&1\\1&\widehat\lambda\end{pmatrix}
 :X\longrightarrow\widehat X
\]
is a principal polarization. The principally polarized variety
$(X,\Lambda)$ is not the Jacobian of any smooth connected genus-four
curve. If, in addition,
\[
             \operatorname{End}(A)=\mathbb Z,\qquad
             \operatorname{NS}(A)=\mathbb Z[L],
\]
then every principal polarization on $X$ is carried to $\Lambda$
by an automorphism of $X$. In that case the unpolarized abelian
variety $X$ is not isomorphic to a genus-four Jacobian.

\emph{Proof of the principal polarization and the known exclusion.}
The inverse symmetric homomorphism is the integral matrix
\[
 \Lambda^{-1}=
 \begin{pmatrix}\widehat\lambda&-1\\-1&\lambda\end{pmatrix}.
\]
To check positivity, pull back along the isogeny
$j=1_A\times\lambda:A^2\to X$. The resulting polarization has
matrix
\[
 \lambda\begin{pmatrix}1&1\\1&2\end{pmatrix};
\]
it is represented by $m^*L\otimes\operatorname{pr}_2^*L$, where
$m(a,b)=a+b$. This is ample because the map
$(a,b)\mapsto(a+b,b)$ is an automorphism and the corresponding
external product of ample line bundles is ample. Ampleness
descends under the finite surjection $j$. Thus $\Lambda$ is
indeed a principal polarization, and its restriction to
$A\times0$ is $L$, of type $(1,2)$.

The exclusion is precisely
\citet[Corollary 2(1), p.~1540, with proof on p.~1541]{frC4Borowka2016}:
the locus $\operatorname{Is}^{2h}_D$ has empty intersection with
the smooth Jacobian locus when the $h$ entries of $D$ contain
both $1$ and $2$. We apply it with $h=2$ and $D=(1,2)$.
For completeness, in this case its mechanism can be seen
directly from the polarization-preserving involution
\[
 R=\begin{pmatrix}1&\widehat\lambda\\0&-1\end{pmatrix},
 \qquad R^2=1,\qquad R^\vee\Lambda R=\Lambda.
\]
Its connected plus and minus eigenspaces are respectively
$A\times0$ and the image of
$a\mapsto(-a,\lambda a)$; both restrictions of $\Lambda$
are $L$. If $(X,\Lambda)$ were a genus-four Jacobian,
the strong Torelli theorem would lift one of $\pm R$ to
an involution of the curve. Its quotient would have genus
two and the double cover would be ramified at two points.
The pullback of the quotient Jacobian is injective and has
the doubled principal polarization, of type $(2,2)$.
This contradicts the computed type $(1,2)$ of the
corresponding eigenspace. These precise Torelli and
ramified-pullback statements are recalled in
\citet[Theorems 5--6, p.~1539]{frC4Borowka2016}.

\emph{Proof of the generic strengthening.}
Under the additional hypotheses we have
\[
 \operatorname{Hom}(A,\widehat A)=\mathbb Z\lambda,
 \qquad
 \operatorname{Hom}(\widehat A,A)=\mathbb Z\widehat\lambda,
 \qquad
 \operatorname{End}(\widehat A)=\mathbb Z.
\]
Indeed all rational homomorphisms in the first group
are rational multiples of $\lambda$ and are symmetric;
over $\mathbb C$ the integral symmetric homomorphisms
are exactly the Neron--Severi group. The dual polarization
is primitive: if its ample class were $k$ times another
integral ample class, its Euler characteristic $2$ would
be divisible by $k^2$, forcing $k=1$.
Consequently every polarization on $X$ has matrix
\[
 \Lambda_{a,b,c}=
 \begin{pmatrix}a\lambda&b\\b&c\widehat\lambda\end{pmatrix},
 \qquad a,b,c\in\mathbb Z.
\]
Its pullback by $j$ is represented, relative to the
product polarization on $A^2$, by the positive definite
integral binary matrix
\[
                M=\begin{pmatrix}a&b\\b&2c\end{pmatrix}.
\]
If $\Lambda_{a,b,c}$ is principal, then
$\det M=2ac-b^2=1$. One can verify the degree formula
directly: $\deg j=4$, the product polarization on $A^2$
has degree $16$, and
$\deg\phi_{j^*\Theta}=16(\det M)^4
 =(\deg j)^2\deg\Lambda_{a,b,c}$.
Positivity selects the positive determinant.
It follows that $a,c,b$ are all odd and hence
\[
                 M\bmod2=
                 M_0\bmod2=
                 \begin{pmatrix}1&1\\1&0\end{pmatrix},
 \qquad M_0=\begin{pmatrix}1&1\\1&2\end{pmatrix}.
\]

Every positive definite integral unimodular binary
form is isometric over $\operatorname{GL}_2(\mathbb Z)$
to the identity form. Here is an elementary proof.
Choose a shortest primitive vector, of squared length
$k$, extend it to an integral basis, and subtract an
integer multiple of the first vector from the second
so that their inner product $d$ satisfies $|d|\le k/2$.
The other squared length $c'$ is at least $k$, whence
\[
             1=kc'-d^2\ge 3k^2/4.
\]
Thus $k=1$, then $d=0$ and $c'=1$.
We can therefore choose $G\in\operatorname{GL}_2(\mathbb Z)$
with $G^tMG=M_0$. Modulo two, the displayed common form
has the unique nonzero isotropic vector $(0,1)$.
The second column of $G$ must consequently be congruent
to $(0,1)$, so
\[
                  G=\begin{pmatrix}u&2v\\w&z\end{pmatrix}.
\]
This is exactly a matrix that lifts through $j$ to
an automorphism of $X$:
\[
 U=\begin{pmatrix}u&v\widehat\lambda\\w\lambda&z\end{pmatrix},
 \qquad Uj=jG.
\]
Since $uz-2vw=\det G=\pm1$, the usual adjugate is an
integral inverse for $U$. The equality of the pulled-back
polarizations gives $U^*\Lambda_{a,b,c}=\Lambda$:
pullback by an isogeny is injective on Neron--Severi groups.
This proves the asserted uniqueness up to automorphism.
An unpolarized Jacobian isomorphism would pull its canonical
principal polarization back to one of these polarizations,
contradicting the first part.

This excludes the proposed genus-four Jacobian construction.
It does not exclude an odd scalar factorization through
higher-dimensional Jacobians, and therefore does not
settle the remaining type-$(1,2)$ case of C.4.

% Independent continuation by /root/frobenius.
\paragraph{C.4u (Proved: all supersingular threefolds in characteristic at least five).}
Every supersingular abelian threefold over an algebraically closed
field of characteristic \(p\ge5\) is an actual direct summand of a
product of Jacobians. In characteristics two and three we obtain,
respectively,
\[
       s_J(A)\mid2\quad(p=2),\qquad
       s_J(A)\mid3\quad(p=3);
\]
these two divisibilities do not by themselves prove a retraction.
More generally, for a supersingular abelian variety \(A\) of
dimension \(g\ge1\),
\begin{equation}\label{eq:c4-supersingular-p-bound}
       s_J(A)\mid p^{\,v_p((2g-1)!)}.
\end{equation}
In particular every such \(A\) is a direct summand if \(p>2g-1\).

\begin{proof}
We first justify the prime support of the required isogeny,
which matters integrally. In
\citet[Lemma 3.2 and its proof, author-version p.~7,
using Lemma 2.4, p.~5]{frC4YuSupersingular2017},
the supersingular variety is mapped to a superspecial variety
by replacing its covariant Dieudonn\'e lattice \(M\) by a
finite-index superspecial overmodule \(M'\) in the same
isocrystal. Our field is already perfect, so the purely
inseparable field extension in the statement is unnecessary.
The quotient \(M'/M\) has finite length over \(W(k)\) and is
killed by \(p^a\) for some \(a\ge0\). The associated isogeny
\[
                         \varphi:A\longrightarrow B
\]
therefore has group-scheme kernel killed by \(p^a\).
The superspecial variety \(B\) is a product of supersingular
elliptic curves over the algebraically closed field, hence
itself a product of Jacobians. This is the definition recalled
in the cited paper on p.~2; its Theorem 1.1 on p.~1 gives the
stronger uniqueness statement, which we do not need.

The quotient property supplies an actual homomorphism
\(h:B\to A\) with
\[
                            h\varphi=[p^a]_A.
\]
Consequently \(p^a\in I_J(A)\). Together with C.4i this shows
that \(s_J(A)\) is a power of \(p\) dividing \((2g-1)!\),
which proves \eqref{eq:c4-supersingular-p-bound}.
This argument works over arbitrary algebraically closed
fields of characteristic \(p\), and does not require that
\(A\) descend to a finite field.

For \(g=3\) the improved integral construction C.4p gives
\(6\in I_J(A)\). Thus
\[
                         s_J(A)\mid\gcd(6,p^a).
\]
This is one for \(p\ge5\), and divides two or three in the
two remaining characteristics. We can display the actual
retraction when \(p\ge5\). Choose
\(A\xrightarrow{u_6}J_6\xrightarrow{v_6}A\)
with \(v_6u_6=[6]_A\). Since \(p^{2a}\equiv1\pmod6\),
the integer \(c=(1-p^{2a})/6\) is defined. The maps
\[
 A\xrightarrow{\,(\varphi,u_6)\,}B\times J_6
       \xrightarrow{\,p^a h+c\,v_6\,}A
\]
have composite
\([p^{2a}+6c]_A=1_A\).
The splitting argument in C.4h therefore gives the required
product decomposition as abelian varieties.
\end{proof}

This settles the supersingular threefold case in every
characteristic at least five, including varieties without a
chosen principal polarization. All supersingular surfaces,
including characteristic two, were already covered by C.4s.
No assertion of principal polarizability of every supersingular
threefold is used here. The residual supersingular cases in
characteristics two and three, and the generic nonprincipal
surface cases discussed above, remain separate questions in
this investigation.


% Independent continuation by /root/frobenius.
\paragraph{C.4v (Proved: full real multiplication removes every polarization gcd obstruction in characteristic zero).}
Let \(A\) be an abelian variety of dimension \(g\) over an
algebraically closed field of characteristic zero. Suppose that
there is an action
\[
                  \mathcal O_F\hookrightarrow\operatorname{End}(A)
\]
of the full ring of integers of a totally real field \(F\) of
degree \(g\), and a polarization whose Rosati involution fixes
\(F\) pointwise. Then
\[
                              \delta(A)=1.
\]
In particular, if \(g\le3\), \(A\) is an actual direct summand
of a product of Jacobians. This includes a simple abelian
surface or threefold with
\(\operatorname{End}(A)=\mathcal O_K\) for a CM field \(K\)
of degree \(2g\). No principal polarization on \(A\) is assumed.

There is a useful local refinement. Suppose an order
\(\mathcal O\subset F\), rather than the full integer ring,
acts with the same Rosati property, and put
\(c=[\mathcal O_F:\mathcal O]\). Then every prime divisor
of \(\delta(A)\) divides \(c\). Consequently an abelian
surface with \(c\) odd, or an abelian threefold with
\(\gcd(c,6)=1\), is still an actual direct summand of a
product of Jacobians. These conclusions require local
maximality only at two, or at two and three, respectively.

\begin{proof}
We first work over \(\mathbb C\). Put
\(\Lambda=H_1(A,\mathbb Z)\) and let \(J\) denote the complex
structure on \(\Lambda_{\mathbb R}\).
The \(\mathcal O_F\)-module \(\Lambda\) is torsion-free of
rank two, hence projective over this Dedekind domain.
Write \(\mathfrak D_F^{-1}\) for the trace-dual fractional
ideal of \(\mathcal O_F\). The module of integral alternating
forms on \(\Lambda\) satisfying
\[
                         E(ax,y)=E(x,ay)
                  \qquad(a\in\mathcal O_F)
\]
is naturally
\begin{equation}\label{eq:c4-rm-form-module}
 \mathcal P=
 \operatorname{Hom}_{\mathcal O_F}
       \bigl(\textstyle\bigwedge^2_{\mathcal O_F}\Lambda,
                         \mathfrak D_F^{-1}\bigr).
\end{equation}
It is an invertible \(\mathcal O_F\)-module.
Explicitly, an \(F\)-valued alternating form \(\psi\) in
\(\mathcal P\) gives \(E=\operatorname{Tr}_{F/\mathbb Q}\psi\).
Conversely, the nondegenerate trace pairing defines \(\psi\)
by
\[
          \operatorname{Tr}_{F/\mathbb Q}(a\psi(x,y))
                              =E(ax,y)\qquad(a\in F).
\]
The balancing condition makes it \(F\)-bilinear and alternating.
Integrality for every \(a\in\mathcal O_F\) is precisely the
codifferent condition in \eqref{eq:c4-rm-form-module}.

Every real form obtained from \(\mathcal P\) has Hodge type
\((1,1)\). Indeed,
\[
       \Lambda_{\mathbb R}=\bigoplus_{\sigma:F\hookrightarrow\mathbb R}
                                    W_\sigma,\qquad
                         \dim_{\mathbb R}W_\sigma=2.
\]
Balancing makes distinct \(W_\sigma\)'s orthogonal.
The action of \(F\) commutes with \(J\), so each real
two-plane is \(J\)-stable. Since \(J^2=-1\), its determinant
on that plane is one, and every alternating form on the
plane is \(J\)-invariant. This proves the assertion.
Let \(E_0\in\mathcal P\) be the integral positive Riemann
form of the given polarization; its balancing condition
is exactly the stated Rosati hypothesis.

Fix a rational prime \(\ell\), including \(\ell=2\).
Over the product of discrete valuation rings
\(\mathcal O_F\otimes\mathbb Z_\ell\), the module
\(\Lambda_\ell\) is free of rank two and
\(\mathcal P_\ell\) is free of rank one. Choose a generator
\(\psi_\ell\) of the latter.
In a basis of \(\Lambda_\ell\), its trace form is
\[
 ((x_1,x_2),(y_1,y_2))
       \longmapsto
 \operatorname{Tr}_{F\otimes\mathbb Q_\ell/\mathbb Q_\ell}
                    \bigl(b(x_1y_2-x_2y_1)\bigr),
\]
where \(b\) generates the local codifferent. The definition
of the codifferent makes this pairing perfect over
\(\mathbb Z_\ell\).
Choose \(E_1\in\mathcal P\) whose reduction modulo \(\ell\)
agrees with this form. This is possible because
\(\mathcal P/\ell\mathcal P
 \simeq\mathcal P_\ell/\ell\mathcal P_\ell\).
For a sufficiently large positive integer \(n\),
\[
                              E=E_1+\ell n E_0
\]
is positive definite in the Riemann-form sense, while its
reduction modulo \(\ell\) is still perfect. Positivity here
is an elementary fact about a fixed real symmetric form
plus a sufficiently large multiple of a positive definite
one; no local-to-global principal-polarization assertion
is being made.

Appell--Humbert and the positivity criterion produce an
ample algebraic line bundle \(L_\ell\) with this integral
Riemann form. The exact primary statement used is the
proof of \citet[Proposition 6.1, pp.~46--47]{frC4MilneCM2020RM},
where both the Appell--Humbert theorem and the Lefschetz
positivity theorem are stated. The induced homomorphism
on integral homology is the map given by \(E\), so
\(\deg\varphi_{L_\ell}=|\det E|\) is prime to \(\ell\);
see also the displayed polarization map on p.~26 of
the same source. Since
\(\deg\varphi_{L_\ell}=\chi(L_\ell)^2\), we have
\(\ell\nmid\chi(L_\ell)\). We construct this line bundle
separately for each prime. Hence no prime divides
\(\delta(A)\), proving \(\delta(A)=1\).

For an arbitrary algebraically closed characteristic-zero
field, define \(A\), the finite set of endomorphisms giving
the action, and the polarization over a finitely generated
subfield, and embed it into \(\mathbb C\).
For each \(\ell\), the resulting ample line bundle spreads
over a nonempty scheme of finite type over that subfield.
After shrinking, ampleness and its Euler characteristic
are preserved. A closed point has a finite residue extension,
which embeds into the given algebraically closed field.
Its line bundle therefore has the same Euler characteristic
prime to \(\ell\). This proves the same gcd assertion.
C.4j then constructs the actual retraction for \(g\le3\).

Finally, when \(\operatorname{End}(A)=\mathcal O_K\) for
the indicated CM field, take \(F=K^+\).
Any Rosati involution on \(K\) is complex conjugation:
it is positive, and
\citet[Proposition 1.39 and Corollary 1.40,
pp.~20--21]{frC4MilneCM2020RM}
identify the positive involution on a CM field uniquely.
Its restriction to \(F\) is the identity, and
\(\mathcal O_F\subset\mathcal O_K\). Thus all the stated
hypotheses hold.

For the local refinement, let \(\mathcal O\) be the given
order and replace \(\mathcal P\) by the integral lattice
\[
 \mathcal P_{\mathcal O}=
 \{E\in\operatorname{Hom}_{\mathbb Z}
       (\textstyle\bigwedge^2_{\mathbb Z}\Lambda,\mathbb Z):
             E(ax,y)=E(x,ay)\text{ for all }a\in\mathcal O\}.
\]
Using a finite integral basis of the order, these are
finitely many integral linear equations. Localization
is exact, so this lattice tensored with \(\mathbb Z_\ell\)
is exactly the lattice of locally balanced integral
forms. If \(\ell\nmid c\), then
\(\mathcal O\otimes\mathbb Z_\ell=
\mathcal O_F\otimes\mathbb Z_\ell\). The local homology
module is therefore free of rank two over each local
integer ring, and the codifferent calculation above
again supplies a perfect local form.
Lift its residue modulo \(\ell\) to
\(\mathcal P_{\mathcal O}\) and add \(\ell nE_0\).
The real-plane and positivity arguments are unchanged,
since \(\mathcal O\otimes\mathbb Q=F\).
We obtain an ample line bundle whose Euler
characteristic is prime to \(\ell\). Thus
\[
            \operatorname{rad}\delta(A)
                   \mid\operatorname{rad}c.
\]
For surfaces and threefolds we combine this with the
already proved universal bounds \(s_J(A)\mid2\) and
\(s_J(A)\mid6\), respectively, and with
\(s_J(A)\mid\delta(A)\). This proves the stated
local refinements; it does not assert that
\(\delta(A)=1\) at a prime where the order is nonmaximal.

In the simple CM case the same refinement applies if
the endomorphism order in \(K\) is locally maximal at
the indicated obstruction primes. Indeed its
intersection with \(K^+\) is then locally
\(\mathcal O_{K^+}\), and the positive Rosati
involution still restricts to the identity there.
\end{proof}

The same affirmative conclusion holds for any good
reduction of one of these varieties of dimension at most
three. Indeed, \(\delta(A)=1\) supplies finitely many
ample line bundles whose Euler characteristics have gcd
one. After a finite extension of a complete discrete
valuation ring of good reduction, all these line bundles
extend to the smooth, hence regular, abelian scheme by Cartier
closure. Its specialization need not stay ample, but
its Euler characteristic stays constant in this proper
flat family. The definition of \(\delta\) using all line
bundles therefore still gives \(\delta=1\) on the
geometric special fibre; the ample-gcd equivalence in
C.4j supplies the required retraction in every residue
characteristic. We retain the full-order and Rosati
hypotheses for the assertion \(\delta=1\).
The locally maximal refinement also specializes:
the finitely many line bundles needed at the
obstruction primes extend and preserve their Euler
characteristics by the same argument. We make no
assertion at a nonmaximal obstruction prime, nor
for the generic surface with endomorphism ring
\(\mathbb Z\).


\paragraph{C.4w (Proved: ordinary full-real-multiplication varieties in every characteristic).}
Let \(k\) be algebraically closed of characteristic \(p>0\),
including \(p=2\), and let \(A/k\) be ordinary of dimension
\(g\). Suppose that the full ring of integers of a totally
real field \(F\) of degree \(g\) acts on \(A\), and that a
polarization has Rosati involution fixing \(F\) pointwise.
Then \(\delta(A)=1\). In dimensions at most three,
\(A\) is therefore an actual direct summand of a product
of Jacobians.

More generally, if an order \(\mathcal O\subset F\)
acts with the same Rosati property, every prime
divisor of \(\delta(A)\) divides
\([\mathcal O_F:\mathcal O]\). Hence odd index
suffices for an ordinary surface, and index
prime to six suffices for an ordinary threefold,
including in characteristics two and three.

\begin{proof}
Let \(\mathcal A/W(k)\) be the canonical Serre--Tate lift.
The precise input is
\citet[Theorem 1.4 and the following paragraph,
p.~115]{frC4DworkOgus1986RM}: the canonical lift is an
abelian scheme, homomorphisms between ordinary abelian
varieties lift uniquely, duality is preserved, and every
polarization lifts. The source expressly explains that
only the converse from its additional crystalline
characterization requires odd characteristic. We do not
use that characterization, so this argument includes
characteristic two.

Lift the given \(\mathcal O_F\)-action and polarization.
Their ring relations and the identities
\[
                 \lambda\circ a=a^\vee\circ\lambda
                      \qquad(a\in\mathcal O_F)
\]
hold on the lift by uniqueness of lifted homomorphisms.
The action remains injective, since reducing a zero
lifted endomorphism would give a zero original
endomorphism. Thus the geometric generic fibre satisfies
all the hypotheses of C.4v and has \(\delta=1\).

Choose finitely many ample line bundles on that fibre
whose Euler characteristics have gcd one. They descend
to one finite extension of the fraction field of \(W(k)\).
After this extension, the abelian scheme is still smooth
over a complete discrete valuation ring, hence regular.
Cartier closure extends each generic line bundle to the
abelian scheme. Proper flatness keeps its Euler
characteristic unchanged on the special fibre.
Consequently the specialized line bundles, whether or
not ample, have Euler characteristics of gcd one.
By the definition of \(\delta\), and its ample-gcd
equivalence in C.4j, we obtain \(\delta(A)=1\).
The integral retraction for \(g\le3\) follows from C.4j.

For a nonmaximal order, the canonical lift preserves
that order and the Rosati identities just as above.
At each prime \(\ell\) where it is locally maximal,
the refinement in C.4v supplies a generic line bundle
whose Euler characteristic is prime to \(\ell\).
After a finite extension, Cartier closure and
constancy of the Euler characteristic give such a
line bundle on \(A\). This proves the asserted
prime-divisor constraint on \(\delta(A)\).
Combining it with \(s_J(A)\mid2\) for surfaces and
\(s_J(A)\mid6\) for threefolds gives the refined
actual retractions.
\end{proof}

Ordinarity is used for the canonical lift with all the
specified endomorphisms and the polarization. We make
no corresponding lifting assertion for an arbitrary
nonordinary variety. The stated local maximality
and Rosati hypotheses are retained.

% Separate continuation by /root/frobenius; C.4x.
\paragraph{C.4x (Proved: full real multiplication in arbitrary positive characteristic).}
Let \(k\) be algebraically closed of characteristic \(p>0\), and
let \(A/k\) have dimension \(g\). Suppose that the full integer
ring of a totally real field of degree \(g\) acts:
\[
                   \mathcal O_F\hookrightarrow\operatorname{End}(A),
                 \qquad [F:\mathbb Q]=g.
\]
We assume that \(A\) has a polarization whose Rosati involution
fixes \(F\) pointwise, equivalently an \(\mathcal O_F\)-linear
polarization. Then the integer \(\delta(A)\) of C.4j is a power
of \(p\), with exponent possibly zero. If, in addition,
\begin{equation}\label{eq:c4-hb-lie-condition}
 [\operatorname{Lie}(A)]
       =[\mathcal O_F\otimes_{\mathbb Z} k]
 \quad\text{in }G_0(\mathcal O_F\otimes_{\mathbb Z}k),
\end{equation}
where \(G_0\) is the Grothendieck group of finite modules, then
\[
                              \delta(A)=1.
\]
Condition \eqref{eq:c4-hb-lie-condition} is automatic if \(p\)
is totally ramified in \(F\).
There is no ordinarity or supersingularity hypothesis, and
\(p=2\) and \(p=3\) are included.

Consequently, in dimensions two and three respectively,
\begin{equation}\label{eq:c4-hb-splitting-bound}
 s_J(A)\mid p^{v_p(2)}
       \quad(g=2),\qquad
 s_J(A)\mid p^{v_p(6)}
       \quad(g=3).
\end{equation}
Every such surface in characteristic different from two, and
every such threefold in characteristic at least five, is an
actual direct summand of a product of Jacobians.
In every characteristic, including the remaining two small
primes, \eqref{eq:c4-hb-lie-condition} gives the same affirmative
conclusion for \(g\le3\). In particular this holds when the
characteristic is totally ramified in the acting real field.

\begin{proof}
We first record the exact input concerning integral
polarizations. In the original published paper,
\citet[Lemma 2.7, pp.~2117--2118]{frC4Yu2003HB}
proves that, for each rational prime \(\ell\ne p\), an abelian
\(\mathcal O_F\)-variety of dimension \([F:\mathbb Q]\) has an
\(\mathcal O_F\)-linear polarization of degree prime to \(\ell\).
His statement even allows an Artinian base; we use only the
algebraically closed field case. We explain the local
calculation, including the integrality step.

Let \(\lambda_0\) be the given compatible polarization. Write
\[
 \mathcal O_F\otimes\mathbb Z_\ell=\prod_{w\mid\ell}\mathcal O_w,
       \qquad T_\ell(A)=\bigoplus_{w\mid\ell}T_w .
\]
Each \(T_w\) is free of rank two over \(\mathcal O_w\).
The pairing of \(\lambda_0\) on \(T_w\) is the trace of an
alternating \(\mathcal O_w\)-bilinear pairing with values in
the local codifferent \(\mathfrak D_w^{-1}\).
After choosing a basis of \(T_w\), let \(c_w\ge0\) be the
valuation of its coefficient relative to a generator of
\(\mathfrak D_w^{-1}\). Thus this trace pairing is perfect
precisely when \(c_w=0\).

Choose \(a\in\mathcal O_F[1/\ell]\) with
\(\operatorname{ord}_w(a)=-c_w\) for every \(w\mid\ell\),
and with \(a\) totally positive. Such a choice follows
by prescribing the finitely many local valuations and
then adding a sufficiently large positive integer divisible
by a sufficiently high power of \(\ell\): the valuations
are unchanged, and all real embeddings become positive.
The resulting rational homomorphism
\(\lambda_0a:A\dashrightarrow A^\vee\) has no denominators
away from \(\ell\), and its \(\ell\)-adic pairing is perfect.
To check that it is an actual homomorphism, write
\(a=b/\ell^n\) with \(b\in\mathcal O_F\).
Integrality of its \(\ell\)-adic map says that \(\lambda_0b\)
annihilates the finite group scheme \(A[\ell^n]\), which is
étale because \(\ell\ne p\). Hence \(\lambda_0b\) factors
through \([\ell^n]\), as required. Total positivity makes
this symmetric homomorphism a polarization, exactly as
in the proof of Yu's lemma. Its perfect Tate pairing
makes its degree prime to \(\ell\).

Choose an ample line bundle \(L_\ell\) inducing this
polarization. The identity
\[
                   \deg\varphi_{L_\ell}=\chi(L_\ell)^2
\]
shows that \(\ell\nmid\chi(L_\ell)\).
We obtain this separately for every \(\ell\ne p\);
therefore the positive gcd \(\delta(A)\) has no prime
divisor other than \(p\).

At the characteristic prime we use
\citet[Proposition 2.8, p.~2118]{frC4Yu2003HB}.
It identifies \eqref{eq:c4-hb-lie-condition} with existence
of a compatible polarization of degree prime to \(p\)
(called a separable polarization there), and also with
the Deligne--Pappas and determinant conditions.
The notation in Sections 2--6 of that paper initially
assumes one prime of \(F\) over \(p\); the reduction to
the separate local components is explicitly stated on
p.~2109 and in Section 2.1, p.~2114.
We use its intrinsic condition (3), namely
\eqref{eq:c4-hb-lie-condition}, so the assertion applies
without that simplifying assumption.
The resulting ample \(L_p\) has \(p\nmid\chi(L_p)\).
Together with the preceding prime-by-prime construction,
this gives \(\delta(A)=1\).

For clarity, total ramification supplies the intrinsic
Lie condition without a smoothness assumption on a
Hilbert moduli space. If \(p\mathcal O_F=\mathfrak p^g\),
then \(R=\mathcal O_F\otimes_{\mathbb Z}k\) is an
Artinian local \(k\)-algebra of length \(g\), with residue
field \(k\). Its only simple finite module is \(k\).
Consequently \(G_0(R)\simeq\mathbb Z\), with the class
of a finite module equal to its \(k\)-dimension.
Both \(R\) and \(\operatorname{Lie}(A)\) have dimension
\(g\); hence their classes agree. This proves
\eqref{eq:c4-hb-lie-condition} also in wild
characteristic and for a nonfree Lie module.

Finally, C.4j gives \(s_J(A)\mid\delta(A)\) in dimensions
at most three. Combining this with the improved
universal bounds \(s_J(A)\mid2\) for surfaces and
\(s_J(A)\mid6\) for threefolds proves
\eqref{eq:c4-hb-splitting-bound}.
Whenever the resulting gcd is one, the finite
factorizations in C.4j and the universal factorization
combine by Bézout to give one actual retraction through
a product of Jacobians. Thus the conclusion is integral.
\end{proof}

The Lie condition is used only as a sufficient condition
for \(\delta(A)=1\); its failure is not asserted to
obstruct a Jacobian-product retraction. Nor does this
argument supply a real multiplication action on an
arbitrary supersingular threefold. The generic
\(\operatorname{End}(A)=\mathbb Z\), type-\((1,2)\)
surface and arbitrary supersingular threefolds in
characteristics two and three therefore retain the
previously isolated integral-cycle question.


\paragraph{Remaining work for C.4 and status of Lead C1.}
The degree-one theorem, its independence of $\ell$, the corrected
integral form and denominator, both proposed equivalences in their
stated settings, and the conditional reduction are proved above.
Lead C1 is correct with the stated coefficient and principal-
polarization qualifications. The full requested C.4 remains
incomplete in the residual nonprincipal cases isolated in the
supplements. C.4h--l supply an exact scalar-factorization ideal,
factorial bounds, simultaneous-polarization criteria, and the specific
limitation of the divisor-product method. The Prym argument improves
the surface bound to two and settles all surfaces with odd
$\delta(A)$, including type $(1,3)$; the remaining surface cases have
even $\delta(A)$. C.4o--p improve the threefold bound to
$\gcd(6,\delta(A))$, leaving only primes two and three in that
dimension. C.4r gives one exact all-characteristic residual test:
the specified rational numerical surface class on $A\times\widehat A$
must belong to the integral numerical cycle lattice. C.4s settles
every supersingular abelian surface, including characteristic two.
C.4u settles every supersingular abelian threefold in
characteristic at least five; only the indicated factors two
and three remain for supersingular threefolds in those respective
characteristics.
C.4v settles the full totally real multiplication cases in
characteristic zero, including its stated maximal-order CM
case, and their good reductions. C.4w settles ordinary
abelian varieties with the stated full real multiplication
and Rosati hypotheses in every positive characteristic,
including two. These families are removed from the
residual nonsplit cases; no nonmaximal-order extension
at an obstruction prime is asserted. The refined order-index
criterion covers real multiplication locally maximal at two
for surfaces, or at two and three for threefolds, including
the corresponding ordinary and good-reduction cases.
C.4x removes the ordinary hypothesis for full real multiplication
in positive characteristic: its polarization index is a power
of the characteristic, and is one under the stated determinant
condition, in particular under total ramification. Thus all such
surfaces in characteristic other than two and all such threefolds
in characteristic at least five split; the determinant condition
also settles their remaining small characteristics.
No general integral splitting in dimensions two and three is
asserted.

\subsection{C.5: rectangle relations and zero-cycles}

\paragraph{C.5a (Proved: products over $\overline{\mathbb F}_p$ vanish integrally).}
Let $k=\overline{\mathbb F}_p$. For every smooth projective $X/k$,
the group $\mathrm{CH}^r(X)_{\mathrm{alg}}$ is both torsion and
divisible. Algebraic equivalence is generated by pushforwards of
degree-zero divisor classes on smooth projective curves under
correspondences. Each such Jacobian point is defined over a finite
field and is therefore torsion. Multiplication by any positive
integer, including $p$, is surjective on the $k$-points of a Jacobian:
the corresponding isogeny is surjective as a morphism of varieties
over an algebraically closed field. Its images, and sums of these
images, are divisible. This proves both assertions.

If $\alpha,\beta$ are algebraically trivial cycle classes, choose
$m>0$ with $m\alpha=0$, and choose $\gamma$ with $\beta=m\gamma$.
Bilinearity gives
\[
 \alpha\cdot\beta=\alpha\cdot(m\gamma)=(m\alpha)\cdot\gamma=0.
\]
The identical argument applies to external products on two varieties.
There are no inverted primes in this proof. In particular the
assertions $1^\circ$ and $1^\circ\mathrm{bis}$ of 14:126 hold integrally
modulo rational equivalence over $\overline{\mathbb F}_p$.

\paragraph{C.5b (Known theorem and proved deductions: the Albanese kernel over
$\overline{\mathbb F}_p$).}
For a smooth projective connected $X/k$ as above,
\[
 \operatorname{alb}_X:\mathrm{CH}_0(X)^0\xrightarrow{\sim}A_X(k)
\]
is an integral isomorphism. The precise opened input is
\citet[Theorem 1.2]{cySpiessSzamuely2003}: it applies to a smooth
variety admitting a smooth compactification over an algebraic
closure of a finite field, and gives the full Albanese isomorphism
on degree-zero Suslin homology. For proper $X$ that homology is
$\mathrm{CH}_0(X)^0$, and $X$ is its own compactification. This includes
the $p$-primary part. The prime-to-$p$ assertion of their Theorem 1.1,
or prime-to-$p$ Rojtman alone, would not suffice.

For an abelian variety the inverse is $x\mapsto[x]-[0]$, so this map
is additive and the assertions $4^\circ$, $4^\circ\mathrm{bis}$ hold.
Assertion $5^\circ$ is precisely the displayed theorem. The canonical
identity $\operatorname{Alb}(X\times Y)=A_X\times A_Y$ gives
$6^\circ$. The kernel generators in $2^\circ$ and the rectangle
relations in $3^\circ$, $3^\circ\mathrm{bis}$ vanish by C.5a. Thus all
members of the chain hold over $\overline{\mathbb F}_p$, integrally.
We make no assertion that
$\mathrm{CH}_0(X/\mathbb F_q)^0=A_X(\mathbb F_q)$ for an arbitrary
finite-field model: the theorem just used is over its algebraic
closure.

\paragraph{C.5c (Proved explicit complex rectangle of infinite order).}
Let $E/\overline{\mathbb Q}$ be the elliptic curve
$y^2=x^3-x$ with origin $O$. Choose $t,s\in\mathbb C$ algebraically
independent over $\overline{\mathbb Q}$ and square roots in $\mathbb C$,
and put
\[
 P=(t,\sqrt{t^3-t}),\quad Q=(s,\sqrt{s^3-s}),\quad S=E\times E.
\]
Then
\begin{equation}\label{eq:explicit-rectangle}
 z=[P,Q]-[P,O]-[O,Q]+[O,O]\in\mathrm{CH}_0(S)
\end{equation}
has infinite order and Albanese class zero.

\emph{Proof.}
The coordinates define a geometric generic point of $S$ over
$\overline{\mathbb Q}$. If a positive multiple of $z$ were rationally
zero over $\mathbb C$, it would already vanish over an algebraic
closure of $K=\overline{\mathbb Q}(S)$. Here we use injectivity of
Chow groups under extensions of algebraically closed fields: a
rational equivalence spreads over a finite-type scheme over the
smaller algebraically closed field and can be specialized at a
rational point of a nonempty open of that scheme. The resulting
relation over $\overline K$ is defined over a finite extension of
$K$; norm then shows that some positive multiple of the generic
rectangle vanishes over $K$.

Write $r_1(x,y)=(x,O)$, $r_2(x,y)=(O,y)$, and let $c$ be constant
with value $(O,O)$. In $\mathrm{CH}^2(S\times S)$ set
\[
 \Gamma=\Delta_S-\Gamma_{r_1}-\Gamma_{r_2}+\Gamma_c.
\]
Its generic restriction over the first factor is the rectangle.
Localization, after spreading the rational equivalence, therefore
makes a positive multiple $N\Gamma$ rationally equivalent to a
cycle supported on $D\times S$ for a proper closed subset
$D\subset S$, which we may enlarge to a divisor. As a correspondence
acting by pullback, $\Gamma$ acts on $H^{2,0}(S)$ as the identity:
$r_1,r_2,c$ factor through varieties of dimension at most one and
their pullbacks annihilate holomorphic two-forms. A correspondence
supported on $D\times S$ has image in the coniveau-one part of
$H^2(S)$, spanned here by divisor classes; its $(2,0)$ part is zero.
This can also be checked by resolving each curve in $D$ and factoring
the pushforward through its degree-zero cohomology with a Tate twist.
But $H^{2,0}(S)$ contains the nonzero form
$\mathrm{pr}_1^*(dx/y)\wedge\mathrm{pr}_2^*(dx/y)$. Hence
$N\,\mathrm{id}_{H^{2,0}(S)}=0$ is impossible. Finally the Albanese
image of \eqref{eq:explicit-rectangle} is the sum of its two
coordinate sums, both zero. This proves the claim.

This is the generic-point and differential-form argument underlying
Mumford's obstruction; we have supplied its application to the
specified curve and points, rather than using ``Mumford'' as an
unexplained assertion that every chosen rectangle is nonzero.

\paragraph{C.5d (Refuted over $\mathbb C$: every universal member of the chain).}
The two divisor classes
$\mathrm{pr}_1^*([P]-[O])$ and $\mathrm{pr}_2^*([Q]-[O])$ on $S$ are
algebraically trivial and their product is $z$. This refutes
$1^\circ$. Their external product refutes $1^\circ\mathrm{bis}$;
the same $z$ lies in the kernel of both projections, refuting the
rectangle and product assertions $2^\circ,3^\circ,3^\circ\mathrm{bis}$
and $6^\circ$. For $A=S$, put $x=(P,O)$ and $y=(O,Q)$. Then
\[
 [x+y]-[x]-[y]+[0]=z\ne0,
\]
so $x\mapsto[x]-[0]$ is not a group homomorphism, refuting
$4^\circ$. The Albanese kernel contains $z$, refuting
$4^\circ\mathrm{bis}$ and $5^\circ$. All these are failures with
integral coefficients and survive tensoring with $\mathbb Q$.
The superscript $a$ on $A_0$ in 14:126 means the algebraically
trivial \emph{subgroup} of the rational Chow group. Reading it as
the quotient modulo algebraic equivalence would erase the problem.

\paragraph{C.5e (Proved universal implication $4^\circ\Rightarrow3^\circ\mathrm{bis}$).}
We retain the universal quantifier over abelian varieties in
$4^\circ$. For smooth projective curves $C,D$, apply it to
$J(C)\times J(D)$. Additivity of the point-to-cycle map kills every
rectangle made from points of its two factors. The normalized
Poincar\'e divisor on $J(C)\times C$ induces
$\mathrm{CH}_0(J(C))^0\to\mathrm{CH}_0(C)^0$ carrying
$[a]-[0]$ to the divisor class $a$; the analogous correspondence
exists for $D$. Applying their external product sends the killed
rectangle to any prescribed product of two degree-zero divisor
classes on $C,D$. Such products therefore vanish.

Any two points on a smooth projective connected variety over an
algebraically closed field can be joined by a connected chain of
curves; after normalizing the components, their difference is a
sum of pushforwards of differences on smooth projective curves.
Apply the curve result to each pair and push to $X\times Y$.
Every rectangle on $X\times Y$ vanishes, proving the implication.
Using $4^\circ$ only for one fixed abelian variety would not suffice.

\paragraph{C.5 scope: the requested fields are settled.}
The truth values of all universal assertions $1^\circ$--$6^\circ$
over the two requested fields are now determined: they all hold
over $\overline{\mathbb F}_p$ and all fail over $\mathbb C$.
In particular, $4^\circ\mathrm{bis}\Rightarrow5^\circ$ is true
in both requested settings: over the first field its conclusion
is independently proved in C.5b; over the second its universal
premise is false by C.5c--d. This decides the implication with
precisely the field quantifiers of C.5, and completes that
field-specific task. We have not supplied a uniform geometric
derivation of this implication over an arbitrary algebraically
closed field. That stronger structural question is distinct
from the two-field decision and remains explicitly open here.
